\chapter{Introduction}

This manual explains the correct equipment usage procedures for users of
the GL controller and is intended primarily for trained engineers and
technicians. Please carefully read this manual to ensure satisfactory
and safe usage of the test chamber.


\section{For restricted use}

The test chamber should be operated by experienced engineers or persons
who have received training in proper usage from an experienced engineer.

\begin{itemize}
\tightlist
\item
  Definition of an experienced engineer: A person who understands the
  purpose of use of the chamber; who has received training in the
  operation method, daily maintenance and inspections, etc.; who can
  foresee and prevent risks associated with common sources of hazards,
  such as electricity, etc.
\end{itemize}


\section{Safety Indications}

The following safety indications are used throughout this manual.

\begin{itemize}
\tightlist
\item
  The following signs represent the degrees of danger to users.
\end{itemize}

\begin{longtable}[]{@{}ll@{}}
\toprule
\endhead
\begin{minipage}[t]{0.28\columnwidth}\raggedright
DANGER\strut
\end{minipage} & \begin{minipage}[t]{0.66\columnwidth}\raggedright
Means that extremely dangerous consequences may arise, with the risk of
death or serious injury to the user, if the chamber is handled
improperly.\strut
\end{minipage}\tabularnewline
\begin{minipage}[t]{0.28\columnwidth}\raggedright
WARNING\strut
\end{minipage} & \begin{minipage}[t]{0.66\columnwidth}\raggedright
Means that extremely dangerous consequences may arise, with the risk of
death or serious injury to the user, if the chamber is handled
improperly.\strut
\end{minipage}\tabularnewline
\begin{minipage}[t]{0.28\columnwidth}\raggedright
CAUTION\strut
\end{minipage} & \begin{minipage}[t]{0.66\columnwidth}\raggedright
Means that extremely dangerous consequences may arise, with the risk of
death or serious injury to the user, if the chamber is handled
improperly.\strut
\end{minipage}\tabularnewline
\bottomrule
\end{longtable}

\begin{itemize}
\tightlist
\item
  Labels that instruct the user to avoid danger.
\end{itemize}

\begin{longtable}[]{@{}ll@{}}
\toprule
\endhead
\begin{minipage}[t]{0.28\columnwidth}\raggedright
PROHIBIT\strut
\end{minipage} & \begin{minipage}[t]{0.66\columnwidth}\raggedright
Means that extremely dangerous consequences may arise, with the risk of
death or serious injury to the user, if the chamber is handled
improperly.\strut
\end{minipage}\tabularnewline
\begin{minipage}[t]{0.28\columnwidth}\raggedright
Imperative Action Required\strut
\end{minipage} & \begin{minipage}[t]{0.66\columnwidth}\raggedright
Means that extremely dangerous consequences may arise, with the risk of
death or serious injury to the user, if the chamber is handled
improperly.\strut
\end{minipage}\tabularnewline
\bottomrule
\end{longtable}

\begin{itemize}
\tightlist
\item
  Labels that indicate information on physical damage and environment
  contamination.
\end{itemize}

\begin{longtable}[]{@{}ll@{}}
\toprule
\endhead
\begin{minipage}[t]{0.28\columnwidth}\raggedright
Notice\strut
\end{minipage} & \begin{minipage}[t]{0.66\columnwidth}\raggedright
This mark means dangerous consequences may arise, with the possibility
of damage to equipment and facilities, or environmental pollution, if
the equipment is handled incorrectly.\strut
\end{minipage}\tabularnewline
\bottomrule
\end{longtable}


\section{Keywords}

\begin{itemize}
\tightlist
\item
  The following keywords are used throughout this manual.
\end{itemize}

\begin{longtable}[]{@{}ll@{}}
\toprule
\endhead
\begin{minipage}[t]{0.28\columnwidth}\raggedright
Note\strut
\end{minipage} & \begin{minipage}[t]{0.66\columnwidth}\raggedright
Provides information necessary for gaining full performance from the GL
controller system or to prevent damage to the equipment.\strut
\end{minipage}\tabularnewline
\begin{minipage}[t]{0.28\columnwidth}\raggedright
Procedure\strut
\end{minipage} & \begin{minipage}[t]{0.66\columnwidth}\raggedright
Explains how to operate the GL controller system on a step-by-step
basis.\strut
\end{minipage}\tabularnewline
\begin{minipage}[t]{0.28\columnwidth}\raggedright
Reference\strut
\end{minipage} & \begin{minipage}[t]{0.66\columnwidth}\raggedright
Offers additional information\strut
\end{minipage}\tabularnewline
\bottomrule
\end{longtable}

\begin{longtable}[]{@{}l@{}}
\toprule
\begin{minipage}[b]{0.97\columnwidth}\raggedright
CAUTION\strut
\end{minipage}\tabularnewline
\midrule
\endhead
\begin{minipage}[t]{0.97\columnwidth}\raggedright
\textbf{This equipment is capable of being operated remotely.} When
operating this equipment in or around the test area, ensure that it is
not being remotely operated by LAN or transmitted communication. If
there is any chance it is being operated remotely, take necessary
precautions to coordinate awareness from all operators locally and
remotely. There is a risk of testing starting suddenly during operation
and causing injury to the operator.\strut
\end{minipage}\tabularnewline
\begin{minipage}[t]{0.97\columnwidth}\raggedright
Ensure the equipment clearly displays that it is being operated
remotely; and also properly notify the operator.\strut
\end{minipage}\tabularnewline
\bottomrule
\end{longtable}


\section{Display used in this
manual}

This manual was written based on the GL controller with temperature and
humidity. Display screens will differ slightly for a model without
humidity options. In particular, a model without humidity does not
display items related to humidity.


\subsection{GL Controller with Temperature and
Humidity}

The following figure depicts the standard display of the GL controller
with temperature and humidity. The display is provided by the
touchscreen monitor (HMI) mounted on the front of the chamber.

\begin{figure}
\centering

\includegraphics[width=0.95\textwidth]{figures/1866183695-S4_3-home-page-000.PNG}
\caption{GL Controller display with temperature and humidity}
\end{figure}\FloatBarrier


\section{Manuals}

The user's manual package that came with the chamber contains the
following reference materials. Refer to the appropriate manual according
to the information required.

\begin{longtable}[]{@{}ll@{}}
\toprule
\begin{minipage}[b]{0.28\columnwidth}\raggedright
Manual\strut
\end{minipage} & \begin{minipage}[b]{0.66\columnwidth}\raggedright
Content\strut
\end{minipage}\tabularnewline
\midrule
\endhead
\begin{minipage}[t]{0.28\columnwidth}\raggedright
GL Chamber Operation Manual\strut
\end{minipage} & \begin{minipage}[t]{0.66\columnwidth}\raggedright
Explains the operation and features of the GL test chamber\strut
\end{minipage}\tabularnewline
\begin{minipage}[t]{0.28\columnwidth}\raggedright
GL Controller Manual\strut
\end{minipage} & \begin{minipage}[t]{0.66\columnwidth}\raggedright
Explains the operation and features of the GL controller system and test
chamber (this manual you are reading)\strut
\end{minipage}\tabularnewline
\begin{minipage}[t]{0.28\columnwidth}\raggedright
GL Controller Communications Manual\strut
\end{minipage} & \begin{minipage}[t]{0.66\columnwidth}\raggedright
Explains the operation and features of the GL controller command
instructions and communication interface options\strut
\end{minipage}\tabularnewline
\begin{minipage}[t]{0.28\columnwidth}\raggedright
GL Options Manual\strut
\end{minipage} & \begin{minipage}[t]{0.66\columnwidth}\raggedright
Explains optional features of GL test chamber\strut
\end{minipage}\tabularnewline
\bottomrule
\end{longtable}


\chapter{Function Overview}

This chapter provides an overview of the functions of the GL controller.
It covers the specifications and methodologies of the GL controller
instrumentation.


\section{GL Controller Operation
Options}

The GL controller can operate as a standalone system or both as a
standalone and networked system. The following table outlines some
notable advantages between these two options. Refer to Chapter 8 for
further details.

\begin{longtable}[]{@{}ll@{}}
\toprule
\begin{minipage}[b]{0.48\columnwidth}\raggedright
Standalone System\strut
\end{minipage} & \begin{minipage}[b]{0.46\columnwidth}\raggedright
Networked System\strut
\end{minipage}\tabularnewline
\midrule
\endhead
\begin{minipage}[t]{0.48\columnwidth}\raggedright
No network connection is required.\strut
\end{minipage} & \begin{minipage}[t]{0.46\columnwidth}\raggedright
GL controller can join DHCP, static or private network.\strut
\end{minipage}\tabularnewline
\begin{minipage}[t]{0.48\columnwidth}\raggedright
System can operate in complete isolation.\strut
\end{minipage} & \begin{minipage}[t]{0.46\columnwidth}\raggedright
Only authorized users on the network can remotely access and operate the
GL controller.\strut
\end{minipage}\tabularnewline
\begin{minipage}[t]{0.48\columnwidth}\raggedright
The touchscreen provides complete control and operation of the chamber
and the GL controller\strut
\end{minipage} & \begin{minipage}[t]{0.46\columnwidth}\raggedright
Any web-based device on the network can access the GL controller and its
user interface for control and operation.\strut
\end{minipage}\tabularnewline
\begin{minipage}[t]{0.48\columnwidth}\raggedright
Update or upgrade is provided via USB external port.\strut
\end{minipage} & \begin{minipage}[t]{0.46\columnwidth}\raggedright
Update on the GL controller is provided via Mender cloud service.\strut
\end{minipage}\tabularnewline
\bottomrule
\end{longtable}


\subsection{GL Controller: Standalone
System}

By default, the GL controller operates as a standalone system. Complete
control and operation of the instrumentation are achieved via the
dedicated (touchscreen) HMI on the chamber, as depicted in the following
diagram.

\begin{figure}
\centering

\includegraphics[width=0.95\textwidth]{figures/921358877-S2_1_1N-standalone-004a.PNG}
\caption{GL controller as a standalone system}
\end{figure}\FloatBarrier


\subsection{GL Controller: Networked
System}

When the GL controller is connected to a network, you can use a Web
browser to access the GL controller user interface (UI) to control and
operate the chamber. The GL controller utilizes a standard Web browser
to provide (i.e., host) its UI for control and operation of the
instrumentation. Thanks to the ability of a Web browser to operate on
any computer (or device) on the network, GL controller operations can be
performed remotely by authorized users on the network.

Details of the network description, accessibility and configuration are
discussed in Chapter 8.

The following figure depicts the GL controller UI on Google Chrome Web
browser, accessed via its IP address.

\begin{figure}
\centering

\includegraphics[width=0.95\textwidth]{figures/3796224613-S2_1_2N-gl-homepage-network-003a.PNG}
\caption{GL controller as a networked system}
\end{figure}\FloatBarrier


\section{Organization of GL Controller User
Interface}

In the GL controller system, there are many ways to accomplish the same
task. This feature was purposely designed to provide the operator an
alternative or shortcut approach to get to the next task or a different
task without having to go through a procedural step. The organization of
the instrumentation screen and user interface is outlined as follows.

\begin{longtable}[]{@{}ll@{}}
\toprule
\begin{minipage}[b]{0.13\columnwidth}\raggedright
Menu\strut
\end{minipage} & \begin{minipage}[b]{0.81\columnwidth}\raggedright
Feature / Submenu\strut
\end{minipage}\tabularnewline
\midrule
\endhead
\begin{minipage}[t]{0.13\columnwidth}\raggedright
Home Page\strut
\end{minipage} & \begin{minipage}[t]{0.81\columnwidth}\raggedright
Display main overview page, menu buttons and operation status\strut
\end{minipage}\tabularnewline
\begin{minipage}[t]{0.13\columnwidth}\raggedright
Operation Mode\strut
\end{minipage} & \begin{minipage}[t]{0.81\columnwidth}\raggedright
Constant Run: Start, Stop Program Run: Start, Pause, Step, Stop\strut
\end{minipage}\tabularnewline
\begin{minipage}[t]{0.13\columnwidth}\raggedright
Monitor\strut
\end{minipage} & \begin{minipage}[t]{0.81\columnwidth}\raggedright
Temp/Humi Details Program External Output Trend Graph\strut
\end{minipage}\tabularnewline
\begin{minipage}[t]{0.13\columnwidth}\raggedright
Set Mode\strut
\end{minipage} & \begin{minipage}[t]{0.81\columnwidth}\raggedright
Constant Setup (No.~1, 2, 3): Equipment Control Options, Refrigeration,
Time Signals, Others Program Setup: New PRGM, Import, Export, Copy,
Move/Swap, Rename, Edit, Preview, Delete\strut
\end{minipage}\tabularnewline
\begin{minipage}[t]{0.13\columnwidth}\raggedright
Setup\strut
\end{minipage} & \begin{minipage}[t]{0.81\columnwidth}\raggedright
Set Timers, Set Sampling, Set Protection, Set Defrost, Set Time Meters,
Data Download, Firmware History, Alarm/Operation History, ROM/Firmware
Information, Configuration, Set Back Trace, Accessory\strut
\end{minipage}\tabularnewline
\bottomrule
\end{longtable}

Configuration under \textbf{Setup} has its own submenu with extensive
options as listed below. Details of these features are discussed in
Chapter 7 (\emph{Configuration Menu}).

\begin{longtable}[]{@{}ll@{}}
\toprule
\begin{minipage}[b]{0.13\columnwidth}\raggedright
Submenu\strut
\end{minipage} & \begin{minipage}[b]{0.81\columnwidth}\raggedright
Feature / Configuration Options\strut
\end{minipage}\tabularnewline
\midrule
\endhead
\begin{minipage}[t]{0.13\columnwidth}\raggedright
Configuration\strut
\end{minipage} & \begin{minipage}[t]{0.81\columnwidth}\raggedright
Set Communication, Operation Process, Control Attainment Range, Time
Signal Names, Display Setup, Set Language, Set Sound, Macros, Set
Date/Time, Users (Register User Password), Sensor Offset (Calibration),
Set Chamber Details, Set Option, Email, Service\strut
\end{minipage}\tabularnewline
\bottomrule
\end{longtable}

\textbf{\emph{Note:}} The \textbf{Macros} submenu (under Configuration)
is not operational on the HMI. It can only be operated remotely and is
grayed out on the HMI display. Refer to Section 7.15 for details.


\section{Operating Panel}

The operating panel is the 10-inch touchscreen HMI that provides the UI
for control and operation of the GL controller instrumentation.

\begin{figure}
\centering

\includegraphics[width=0.95\textwidth]{figures/3932870397-S1_3-main-monitor-001.jpg}
\caption{Operation panel}
\end{figure}\FloatBarrier

The UI consists of three major components: (1) Menu bar, (2) Status bar,
and (3) Standard Display area. The Operation Mode page (depicted in the
figure below) illustrates these components. The touchscreen is operated
by gently pressing the screen elements provided by these components.

\begin{figure}
\centering

\includegraphics[width=0.95\textwidth]{figures/2938447182-S4_3-Op-mode-001a.PNG}
\caption{Operation panel}
\end{figure}\FloatBarrier

\begin{longtable}[]{@{}lll@{}}
\toprule
\begin{minipage}[b]{0.04\columnwidth}\raggedright
No.\strut
\end{minipage} & \begin{minipage}[b]{0.23\columnwidth}\raggedright
Name\strut
\end{minipage} & \begin{minipage}[b]{0.65\columnwidth}\raggedright
Function/Use\strut
\end{minipage}\tabularnewline
\midrule
\endhead
\begin{minipage}[t]{0.04\columnwidth}\raggedright
1\strut
\end{minipage} & \begin{minipage}[t]{0.23\columnwidth}\raggedright
Menu Bar\strut
\end{minipage} & \begin{minipage}[t]{0.65\columnwidth}\raggedright
Displays common operation menus of the system\strut
\end{minipage}\tabularnewline
\begin{minipage}[t]{0.04\columnwidth}\raggedright
2\strut
\end{minipage} & \begin{minipage}[t]{0.23\columnwidth}\raggedright
Status Bar\strut
\end{minipage} & \begin{minipage}[t]{0.65\columnwidth}\raggedright
Displays common system statuses: information, alarm, operating mode,
date/time, user login, and system menus and accessories\strut
\end{minipage}\tabularnewline
\begin{minipage}[t]{0.04\columnwidth}\raggedright
3\strut
\end{minipage} & \begin{minipage}[t]{0.23\columnwidth}\raggedright
Standard Display\strut
\end{minipage} & \begin{minipage}[t]{0.65\columnwidth}\raggedright
Displays and sets detailed items and content for the selected menu\strut
\end{minipage}\tabularnewline
\bottomrule
\end{longtable}

The layout of the UI consists of the menu bar on the left and the status
bar on the top; both of which remain fixed throughout the operation. The
touch operation is interactive. Whenever text or numeral input is
required, the floating keyboard or keypad is automatically laid over the
bottom portion of the screen to receive the input.

\textbf{\emph{Note:}} Press on the touchscreen only with your fingers.
Hard or sharp-pointed objects should never be used on the touchscreen;
doing so will damage the screen.


\subsection{Home Page}

The home page of the GL controller is the first display to appear on the
HMI after the system has started and is ready for operation. By default,
the localhost user is automatically logged in to operate the chamber.

\textbf{\emph{Note:}} If the GL controller was accessed remotely on a
network via a Web browser, the user must log in to operate the chamber.
The localhost user is only available on the HMI operation.

\begin{figure}
\centering

\includegraphics[width=0.95\textwidth]{figures/3252919907-S4_3_1-001a.PNG}
\caption{Home page of GL Controller}
\end{figure}\FloatBarrier

\begin{longtable}[]{@{}lll@{}}
\toprule
\begin{minipage}[b]{0.04\columnwidth}\raggedright
No.\strut
\end{minipage} & \begin{minipage}[b]{0.23\columnwidth}\raggedright
Name\strut
\end{minipage} & \begin{minipage}[b]{0.65\columnwidth}\raggedright
Function/Use\strut
\end{minipage}\tabularnewline
\midrule
\endhead
\begin{minipage}[t]{0.04\columnwidth}\raggedright
1\strut
\end{minipage} & \begin{minipage}[t]{0.23\columnwidth}\raggedright
Menu Bar\strut
\end{minipage} & \begin{minipage}[t]{0.65\columnwidth}\raggedright
Displays common operation menus of the GL controller system\strut
\end{minipage}\tabularnewline
\begin{minipage}[t]{0.04\columnwidth}\raggedright
2\strut
\end{minipage} & \begin{minipage}[t]{0.23\columnwidth}\raggedright
Status Bar\strut
\end{minipage} & \begin{minipage}[t]{0.65\columnwidth}\raggedright
Displays common system statuses that include information, alarm,
operating mode, date/time, user login, system menu and accessories (such
as, screenshot, set timers and network monitor page)\strut
\end{minipage}\tabularnewline
\begin{minipage}[t]{0.04\columnwidth}\raggedright
3\strut
\end{minipage} & \begin{minipage}[t]{0.23\columnwidth}\raggedright
Menu Buttons\strut
\end{minipage} & \begin{minipage}[t]{0.65\columnwidth}\raggedright
Displays and provides access to common operation menus as buttons
(available in main home page only)\strut
\end{minipage}\tabularnewline
\begin{minipage}[t]{0.04\columnwidth}\raggedright
4\strut
\end{minipage} & \begin{minipage}[t]{0.23\columnwidth}\raggedright
Operation Buttons\strut
\end{minipage} & \begin{minipage}[t]{0.65\columnwidth}\raggedright
Displays control buttons for start, stop of selected constant mode or
start, stop, pause of select program mode\strut
\end{minipage}\tabularnewline
\begin{minipage}[t]{0.04\columnwidth}\raggedright
5\strut
\end{minipage} & \begin{minipage}[t]{0.23\columnwidth}\raggedright
Set Buttons\strut
\end{minipage} & \begin{minipage}[t]{0.65\columnwidth}\raggedright
Displays a set of quick-access setting buttons for users, date/time and
language\strut
\end{minipage}\tabularnewline
\begin{minipage}[t]{0.04\columnwidth}\raggedright
6\strut
\end{minipage} & \begin{minipage}[t]{0.23\columnwidth}\raggedright
Secondary Buttons\strut
\end{minipage} & \begin{minipage}[t]{0.65\columnwidth}\raggedright
Displays secondary buttons for the selected menu. The edit button (pen
icon) allows the user to edit the layout. For example, the four menu
buttons can be replaced with other menus or button selected from a
different menu for quick access.\strut
\end{minipage}\tabularnewline
\begin{minipage}[t]{0.04\columnwidth}\raggedright
7\strut
\end{minipage} & \begin{minipage}[t]{0.23\columnwidth}\raggedright
Standard Display\strut
\end{minipage} & \begin{minipage}[t]{0.65\columnwidth}\raggedright
Displays and sets the main menu screens\strut
\end{minipage}\tabularnewline
\bottomrule
\end{longtable}


\subsection{Operation Mode}

The Operation Mode page has four groups of components in the display
area. The function of these components are listed in the table that
follows.

\begin{figure}
\centering

\includegraphics[width=0.95\textwidth]{figures/715538957-S1_2_2-op-mode-000a.PNG}
\caption{Features of Operation Mode}
\end{figure}\FloatBarrier

\begin{longtable}[]{@{}lll@{}}
\toprule
\begin{minipage}[b]{0.04\columnwidth}\raggedright
No.\strut
\end{minipage} & \begin{minipage}[b]{0.23\columnwidth}\raggedright
Name\strut
\end{minipage} & \begin{minipage}[b]{0.65\columnwidth}\raggedright
Description: Function/Operations\strut
\end{minipage}\tabularnewline
\midrule
\endhead
\begin{minipage}[t]{0.04\columnwidth}\raggedright
1\strut
\end{minipage} & \begin{minipage}[t]{0.23\columnwidth}\raggedright
Operation Button\strut
\end{minipage} & \begin{minipage}[t]{0.65\columnwidth}\raggedright
Start, stop, pause buttons for operating Constant or Program mode\strut
\end{minipage}\tabularnewline
\begin{minipage}[t]{0.04\columnwidth}\raggedright
2\strut
\end{minipage} & \begin{minipage}[t]{0.23\columnwidth}\raggedright
Constant Run\strut
\end{minipage} & \begin{minipage}[t]{0.65\columnwidth}\raggedright
Constant 1, 2 or 3 can be selected for constant mode operation.\strut
\end{minipage}\tabularnewline
\begin{minipage}[t]{0.04\columnwidth}\raggedright
3\strut
\end{minipage} & \begin{minipage}[t]{0.23\columnwidth}\raggedright
Program Run\strut
\end{minipage} & \begin{minipage}[t]{0.65\columnwidth}\raggedright
A profile in the program list can be selected to run.\strut
\end{minipage}\tabularnewline
\begin{minipage}[t]{0.04\columnwidth}\raggedright
4\strut
\end{minipage} & \begin{minipage}[t]{0.23\columnwidth}\raggedright
Paging Button\strut
\end{minipage} & \begin{minipage}[t]{0.65\columnwidth}\raggedright
Pages through program list; the paging button can page through the
program list up or down the page per row, per page or to the last
page\strut
\end{minipage}\tabularnewline
\bottomrule
\end{longtable}


\subsection{Monitor}

The Monitor page consists of six separate submenus to allow separate
displays, controls and operations. They include three Monitor pages (1,
2 and 3), Details, Program and Trend Graph. Refer to Section 4.9 for
details on constant operation monitoring, for example.

\begin{figure}
\centering

\includegraphics[width=0.95\textwidth]{figures/2108324026-S4_3_3_display-page-01a.PNG}
\caption{GL system monitor page}
\end{figure}\FloatBarrier

\begin{longtable}[]{@{}lll@{}}
\toprule
\begin{minipage}[b]{0.04\columnwidth}\raggedright
No.\strut
\end{minipage} & \begin{minipage}[b]{0.23\columnwidth}\raggedright
Name\strut
\end{minipage} & \begin{minipage}[b]{0.65\columnwidth}\raggedright
Description: Function/Operations\strut
\end{minipage}\tabularnewline
\midrule
\endhead
\begin{minipage}[t]{0.04\columnwidth}\raggedright
1\strut
\end{minipage} & \begin{minipage}[t]{0.23\columnwidth}\raggedright
Monitor Submenu\strut
\end{minipage} & \begin{minipage}[t]{0.65\columnwidth}\raggedright
Displays separate submenus for display, control and operation.\strut
\end{minipage}\tabularnewline
\begin{minipage}[t]{0.04\columnwidth}\raggedright
2\strut
\end{minipage} & \begin{minipage}[t]{0.23\columnwidth}\raggedright
Display Area\strut
\end{minipage} & \begin{minipage}[t]{0.65\columnwidth}\raggedright
Displays set point and process values of temperature and humidity. When
the equipment control is stopped or set to OFF, ``OFF'' is displayed. If
the temperature (humidity) process value is abnormal or the humidity
control is set to OFF, ``---'' is displayed.\strut
\end{minipage}\tabularnewline
\bottomrule
\end{longtable}


\subsection{Set Mode}

The Set Mode page allows the operator to set and configure constant mode
options and to manage programming profiles. Refer to Chapter 4 and 5 for
further details.

\begin{figure}
\centering

\includegraphics[width=0.95\textwidth]{figures/3279578874-S4_3_1-set-mode-000a-fig.PNG}
\caption{Set Mode page}
\end{figure}\FloatBarrier

\begin{longtable}[]{@{}lll@{}}
\toprule
\begin{minipage}[b]{0.04\columnwidth}\raggedright
No.\strut
\end{minipage} & \begin{minipage}[b]{0.23\columnwidth}\raggedright
Name\strut
\end{minipage} & \begin{minipage}[b]{0.65\columnwidth}\raggedright
Description: Function/Operations\strut
\end{minipage}\tabularnewline
\midrule
\endhead
\begin{minipage}[t]{0.04\columnwidth}\raggedright
1\strut
\end{minipage} & \begin{minipage}[t]{0.23\columnwidth}\raggedright
Constant Set Mode\strut
\end{minipage} & \begin{minipage}[t]{0.65\columnwidth}\raggedright
Displays available constant mode options; Constant 1, 2 or 3 can be
selected for configuration.\strut
\end{minipage}\tabularnewline
\begin{minipage}[t]{0.04\columnwidth}\raggedright
2\strut
\end{minipage} & \begin{minipage}[t]{0.23\columnwidth}\raggedright
Program Set Mode\strut
\end{minipage} & \begin{minipage}[t]{0.65\columnwidth}\raggedright
Displays available profiles on the program list; import, export, edit,
copy, move/relocate, view, rename and delete are the available features
to manage profiles.\strut
\end{minipage}\tabularnewline
\begin{minipage}[t]{0.04\columnwidth}\raggedright
3\strut
\end{minipage} & \begin{minipage}[t]{0.23\columnwidth}\raggedright
Paging Buttons\strut
\end{minipage} & \begin{minipage}[t]{0.65\columnwidth}\raggedright
Pages through program list; the button can page through the program list
up or down the page per row, per page or to the last page\strut
\end{minipage}\tabularnewline
\bottomrule
\end{longtable}


\subsection{Setup}

The Setup page provides many options for the operator to perform a
system-wide setting or configuration. Twelve submenus are available to
provide further controls of the GL controller system and the chamber.
The owner (or administrator) who manages this chamber should take the
necessary precautions to limit or allow users with certain privileges to
access and control this \textbf{Setup} menu and its submenus. Details of
user policy and access privileges are discussed in Chapter 6.

\begin{figure}
\centering

\includegraphics[width=0.95\textwidth]{figures/160479785-S4_3_5-setup-page-001b.PNG}
\caption{The Setup page}
\end{figure}\FloatBarrier

\begin{longtable}[]{@{}lll@{}}
\toprule
\begin{minipage}[b]{0.04\columnwidth}\raggedright
No.\strut
\end{minipage} & \begin{minipage}[b]{0.23\columnwidth}\raggedright
Name\strut
\end{minipage} & \begin{minipage}[b]{0.65\columnwidth}\raggedright
Description: Function/Operations\strut
\end{minipage}\tabularnewline
\midrule
\endhead
\begin{minipage}[t]{0.04\columnwidth}\raggedright
1\strut
\end{minipage} & \begin{minipage}[t]{0.23\columnwidth}\raggedright
Set Timers\strut
\end{minipage} & \begin{minipage}[t]{0.65\columnwidth}\raggedright
Provides an option to schedule a start (or stop) of an operation\strut
\end{minipage}\tabularnewline
\begin{minipage}[t]{0.04\columnwidth}\raggedright
2\strut
\end{minipage} & \begin{minipage}[t]{0.23\columnwidth}\raggedright
Set Sampling\strut
\end{minipage} & \begin{minipage}[t]{0.65\columnwidth}\raggedright
Provides an option to configure sampling and data logging details\strut
\end{minipage}\tabularnewline
\begin{minipage}[t]{0.04\columnwidth}\raggedright
3\strut
\end{minipage} & \begin{minipage}[t]{0.23\columnwidth}\raggedright
Set Protection\strut
\end{minipage} & \begin{minipage}[t]{0.65\columnwidth}\raggedright
Provides an option for users to set HMI automatic login with On/Off
options (available for HMI operation only)\strut
\end{minipage}\tabularnewline
\begin{minipage}[t]{0.04\columnwidth}\raggedright
4\strut
\end{minipage} & \begin{minipage}[t]{0.23\columnwidth}\raggedright
Set Defrost\strut
\end{minipage} & \begin{minipage}[t]{0.65\columnwidth}\raggedright
Option currently unavailable with this model\strut
\end{minipage}\tabularnewline
\begin{minipage}[t]{0.04\columnwidth}\raggedright
5\strut
\end{minipage} & \begin{minipage}[t]{0.23\columnwidth}\raggedright
Set Time Meters\strut
\end{minipage} & \begin{minipage}[t]{0.65\columnwidth}\raggedright
Provides a list of the operating condition (odometer) of the hardware
components or devices installed inside the chamber\strut
\end{minipage}\tabularnewline
\begin{minipage}[t]{0.04\columnwidth}\raggedright
6\strut
\end{minipage} & \begin{minipage}[t]{0.23\columnwidth}\raggedright
Data Download\strut
\end{minipage} & \begin{minipage}[t]{0.65\columnwidth}\raggedright
Backup data logging using Samba for remote storage\strut
\end{minipage}\tabularnewline
\begin{minipage}[t]{0.04\columnwidth}\raggedright
7\strut
\end{minipage} & \begin{minipage}[t]{0.23\columnwidth}\raggedright
Firmware\strut
\end{minipage} & \begin{minipage}[t]{0.65\columnwidth}\raggedright
Provides a list of the firmware history of the system\strut
\end{minipage}\tabularnewline
\begin{minipage}[t]{0.04\columnwidth}\raggedright
8\strut
\end{minipage} & \begin{minipage}[t]{0.23\columnwidth}\raggedright
Alarm/Operation\strut
\end{minipage} & \begin{minipage}[t]{0.65\columnwidth}\raggedright
Provides a list of Alarm and Operation history of the system\strut
\end{minipage}\tabularnewline
\begin{minipage}[t]{0.04\columnwidth}\raggedright
9\strut
\end{minipage} & \begin{minipage}[t]{0.23\columnwidth}\raggedright
ROM/Firmware Info\strut
\end{minipage} & \begin{minipage}[t]{0.65\columnwidth}\raggedright
Provides a list of the current firmware information\strut
\end{minipage}\tabularnewline
\begin{minipage}[t]{0.04\columnwidth}\raggedright
10\strut
\end{minipage} & \begin{minipage}[t]{0.23\columnwidth}\raggedright
Configuration\strut
\end{minipage} & \begin{minipage}[t]{0.65\columnwidth}\raggedright
Provides a list of system-wide configuration\strut
\end{minipage}\tabularnewline
\begin{minipage}[t]{0.04\columnwidth}\raggedright
11\strut
\end{minipage} & \begin{minipage}[t]{0.23\columnwidth}\raggedright
Set Back Trace\strut
\end{minipage} & \begin{minipage}[t]{0.65\columnwidth}\raggedright
Provides a list of alarm conditions available for download for
analysis\strut
\end{minipage}\tabularnewline
\begin{minipage}[t]{0.04\columnwidth}\raggedright
12\strut
\end{minipage} & \begin{minipage}[t]{0.23\columnwidth}\raggedright
Accessory\strut
\end{minipage} & \begin{minipage}[t]{0.65\columnwidth}\raggedright
Provides access to system runtime, HMI screensaver, and controller power
options (with reboot and service restart options)\strut
\end{minipage}\tabularnewline
\bottomrule
\end{longtable}

An alternate way to directly access the \textbf{Setup} menu and its
submenus altogether is via the Hamburger menu in the status bar, as
depicted below.

\textbf{\emph{Procedure:}}

\begin{enumerate}
\def\labelenumi{\arabic{enumi}.}
\item
  Press the Hamburger button (1) in the status bar to access the system
  menu. Check to confirm the floating keyboard is enabled. A diagonal
  bar over the keyboard indicates it is disabled.
\item
  Press the keyboard (2) to enable it.
\item
  Press Setup (3) to access the Setup submenu list.

  \begin{figure}
  \centering
  
\includegraphics[width=0.95\textwidth]{figures/3886512048-S4_3_5-system-menu-001c.PNG}
  \caption{System menu via the Hamburger button}
  \end{figure}\FloatBarrier
\end{enumerate}

\textbf{\emph{Note:}} The system menu via the Hamburger button lists
\textbf{Configuration} and \textbf{User} as separate menus for their
quick access. This option provides a shortcut to administer the system,
such as adding new user accounts or performing a system-wide
configuration.


\section{Operating Panel: Elements and
Controls}


\subsection{Menu and Common Operation
Buttons}

The overall control features and operations of the GL controller system
are accessible via the menus in the menu bar, menu buttons in the home
page and the system menus in the status bar. Here are the detailed
descriptions of the operation buttons and components of the control
system in the menu bar, menu and setting buttons.

\begin{figure}
\centering

\includegraphics[width=0.95\textwidth]{figures/1373510695-menubar-buttons-002.PNG}
\caption{Menu bar, Menu Buttons and Secondary Buttons}
\end{figure}\FloatBarrier

\begin{enumerate}
\def\labelenumi{\arabic{enumi}.}
\tightlist
\item
  \textbf{Menu Bar}: Displays system menus and control settings.
  \textbf{Home Main window} is the home page of the GL controller UI; it
  is the overview page indicating the system is ready for operation.
  \textbf{Operation Mode} provides access to control the chamber with
  options to start or stop Constant mode or Program mode.
  \textbf{Monitor} provides access to display details of chamber
  operating condition. \textbf{Set Mode} provides access to modify
  constant mode, create or edit a profile. \textbf{Setup} provides
  access to perform a system-wide configuration or settings of the GL
  controller.\\
\item
  \textbf{Menu Buttons}: Available only in the home page as a quick
  access, these buttons are identical to the four menus in the menu
  bar.\\
\item
  \textbf{Secondary Buttons}: These buttons operate in conjunction with
  the primary buttons; they serve as negative actions (e.g., Cancel or
  Back) or additional actions (e.g.~Create, Import, Edit). The Edit
  button is useful for customizing the UI and its components in the Home
  page and Monitor, Page 1 and Page 2.
\item
  \textbf{Start/Stop Buttons}: These start, stop and pause buttons are
  the operation buttons in the \textbf{Operation Mode}. They are placed
  here as quick-access buttons to start/stop a constant operating mode,
  start/stop or pause a profile in program mode.\\
\item
  \textbf{Setting Buttons}: Three commonly used features (date/time,
  language and user settings) from the \textbf{Setup} menu are made
  available as setting buttons for quick access to their operation.
\end{enumerate}


\subsection{Status Bar}

The status bar consists of multiple elements and components to provide
status notifications (such as alarm, information, operating mode,
date/time and user login) and quick access to the GL controller
operation (such as operation mode, screenshot action, date/time
adjustment, set timers operation, network monitor display and system
menu). There is a subtle difference in the status bar between the HMI
display and the Web display (see figure below). In the figure, the top
status bar represents the HMI display, identifiable by date/time and
localhost log-in status. The bottom status bar represents the Web
display (via a remote network access), identifiable by the GL controller
hostname and username.

\begin{figure}
\centering

\includegraphics[width=0.95\textwidth]{figures/1559059139-status-bar-hmi-n-web-001a.PNG}
\caption{Task bar showing common operation buttons and components}
\end{figure}\FloatBarrier

\begin{longtable}[]{@{}lll@{}}
\toprule
\begin{minipage}[b]{0.04\columnwidth}\raggedright
No.\strut
\end{minipage} & \begin{minipage}[b]{0.23\columnwidth}\raggedright
Name\strut
\end{minipage} & \begin{minipage}[b]{0.65\columnwidth}\raggedright
Description: Function/Operations\strut
\end{minipage}\tabularnewline
\midrule
\endhead
\begin{minipage}[t]{0.04\columnwidth}\raggedright
1\strut
\end{minipage} & \begin{minipage}[t]{0.23\columnwidth}\raggedright
Home\strut
\end{minipage} & \begin{minipage}[t]{0.65\columnwidth}\raggedright
Displays home page of GL controller operating system\strut
\end{minipage}\tabularnewline
\begin{minipage}[t]{0.04\columnwidth}\raggedright
2\strut
\end{minipage} & \begin{minipage}[t]{0.23\columnwidth}\raggedright
Alarm\strut
\end{minipage} & \begin{minipage}[t]{0.65\columnwidth}\raggedright
Displays alarm notifications and a list of alarm issues\strut
\end{minipage}\tabularnewline
\begin{minipage}[t]{0.04\columnwidth}\raggedright
3\strut
\end{minipage} & \begin{minipage}[t]{0.23\columnwidth}\raggedright
Information\strut
\end{minipage} & \begin{minipage}[t]{0.65\columnwidth}\raggedright
Displays information of chamber operation history\strut
\end{minipage}\tabularnewline
\begin{minipage}[t]{0.04\columnwidth}\raggedright
4\strut
\end{minipage} & \begin{minipage}[t]{0.23\columnwidth}\raggedright
Operating Mode\strut
\end{minipage} & \begin{minipage}[t]{0.65\columnwidth}\raggedright
Displays current operating mode; \textbf{On HMI:} (1) press and hold for
2 seconds to take screenshot of the current display (USB device plug-in
required), (2) press to quickly access the Operation Mode; \textbf{On
Web display page:} click this area to quickly access the Operation
Mode.\strut
\end{minipage}\tabularnewline
\begin{minipage}[t]{0.04\columnwidth}\raggedright
5\strut
\end{minipage} & \begin{minipage}[t]{0.23\columnwidth}\raggedright
Date/Time\strut
\end{minipage} & \begin{minipage}[t]{0.65\columnwidth}\raggedright
Displays current date/time; press and hold to access \textbf{Set
Timers}\strut
\end{minipage}\tabularnewline
\begin{minipage}[t]{0.04\columnwidth}\raggedright
6\strut
\end{minipage} & \begin{minipage}[t]{0.23\columnwidth}\raggedright
Login Status\strut
\end{minipage} & \begin{minipage}[t]{0.65\columnwidth}\raggedright
Displays login status; the diagonal bar over the person icon indicates
no user has logged in; Section 4.5 explains how to log in to the
system.\strut
\end{minipage}\tabularnewline
\begin{minipage}[t]{0.04\columnwidth}\raggedright
7\strut
\end{minipage} & \begin{minipage}[t]{0.23\columnwidth}\raggedright
Hamburger Menu\strut
\end{minipage} & \begin{minipage}[t]{0.65\columnwidth}\raggedright
Displays a complete list of menus and submenus for system
operations\strut
\end{minipage}\tabularnewline
\begin{minipage}[t]{0.04\columnwidth}\raggedright
8\strut
\end{minipage} & \begin{minipage}[t]{0.23\columnwidth}\raggedright
Hostname\strut
\end{minipage} & \begin{minipage}[t]{0.65\columnwidth}\raggedright
Displays GL controller hostname; click to launch Network Monitor
page\strut
\end{minipage}\tabularnewline
\bottomrule
\end{longtable}


\subsection{Standard Display}

Contents associated with each menu in the menu bar are displayed in the
standard display area. The selected menu is highlighted and marked by
the left vertical bar to indicate its active status. Each menu has
different display contents associated with their features and control
options. Two menus (\textbf{Home} and \textbf{Monitor}) are selected for
illustration as follows.

\textbf{Home} page and Standard Display:

\begin{figure}
\centering

\includegraphics[width=0.95\textwidth]{figures/3129471487-S4_4_3-fig1B.PNG}
\caption{Home: Display area and activate menu}
\end{figure}\FloatBarrier

Refer to Section 2.3.1 for details on the separate elements in the
standard display area of the \textbf{Home} page.

\textbf{Monitor} page and Standard Display:

\begin{figure}
\centering

\includegraphics[width=0.95\textwidth]{figures/908849024-S4_4_3-fig2B.PNG}
\caption{Monitor: Display area and activate menu}
\end{figure}\FloatBarrier

Refer to Section 2.3.3 for details on the separate elements in the
standard display area of the \textbf{Monitor} page.


\subsection{Floating Keyboard}

The floating keyboard is a multifunctional input device for an operator
to log in to the system, create or edit a profile, set or perform a
system-wide configuration. It has an enable and disable option under the
Hamburger menu on the status bar. If it is disabled, a user cannot log
in to operate the chamber. When the system (chamber) is powered on, the
keyboard is disabled. To control the chamber, the user must enable the
keyboard (and its keypad).

\begin{figure}
\centering

\includegraphics[width=0.95\textwidth]{figures/904356026-keyboard-001b.PNG}
\caption{Floating keyboard}
\end{figure}\FloatBarrier

\begin{longtable}[]{@{}lll@{}}
\toprule
No. & Name & Description: Function/Operations\tabularnewline
\midrule
\endhead
1 & Input Field & Use this input field to enter alpha-numeric
values\tabularnewline
2 & Apply Button & Apply input values\tabularnewline
3 & Close Button & Close keyboard after submission or cancel
input\tabularnewline
4 & Backspace & Deletes one character to its left each time the key is
pressed\tabularnewline
5 & CLR & Clears all characters in the input field\tabularnewline
\bottomrule
\end{longtable}

The following procedure illustrates how to enable (or disable) the
floating keyboard, which in turn also enables (or disables) the keypad
for chamber operation.

\textbf{\emph{Procedure:}}

\begin{itemize}
\item
  Press the Hamburger menu (1) in the status bar. Check to confirm the
  floating keyboard is enabled under the system menu. A diagonal bar
  over the keyboard indicates it is disabled. Press the keyboard (2) to
  enable it. Press the right angle (3) to close the system menu.

  \begin{figure}
  \centering
  
\includegraphics[width=0.95\textwidth]{figures/927228904-S4_4_4-step1.PNG}
  \caption{Enable the floating keyboard and keypad}
  \end{figure}\FloatBarrier
\end{itemize}


\subsection{Floating Keypad}

The floating keypad is a numerical input device for input parameters. It
is enabled or disabled under the keyboard enable/disable options (see
Section 2.4.4). If the keyboard is disabled, the keypad is also
disabled. In order to use the keypad, the keyboard must be enabled.

\begin{figure}
\centering

\includegraphics[width=0.95\textwidth]{figures/3131468203-keypad-001b.PNG}
\caption{Floating keypad}
\end{figure}\FloatBarrier

\begin{longtable}[]{@{}lll@{}}
\toprule
\begin{minipage}[b]{0.04\columnwidth}\raggedright
No.\strut
\end{minipage} & \begin{minipage}[b]{0.23\columnwidth}\raggedright
Name\strut
\end{minipage} & \begin{minipage}[b]{0.65\columnwidth}\raggedright
Description: Function/Operations\strut
\end{minipage}\tabularnewline
\midrule
\endhead
\begin{minipage}[t]{0.04\columnwidth}\raggedright
1\strut
\end{minipage} & \begin{minipage}[t]{0.23\columnwidth}\raggedright
Input Field\strut
\end{minipage} & \begin{minipage}[t]{0.65\columnwidth}\raggedright
Use this input field to enter numerical values for the input
parameter\strut
\end{minipage}\tabularnewline
\begin{minipage}[t]{0.04\columnwidth}\raggedright
2\strut
\end{minipage} & \begin{minipage}[t]{0.23\columnwidth}\raggedright
Value Range\strut
\end{minipage} & \begin{minipage}[t]{0.65\columnwidth}\raggedright
Displays available and permissible range of input values\strut
\end{minipage}\tabularnewline
\begin{minipage}[t]{0.04\columnwidth}\raggedright
3\strut
\end{minipage} & \begin{minipage}[t]{0.23\columnwidth}\raggedright
Apply Button\strut
\end{minipage} & \begin{minipage}[t]{0.65\columnwidth}\raggedright
Apply input values\strut
\end{minipage}\tabularnewline
\begin{minipage}[t]{0.04\columnwidth}\raggedright
4\strut
\end{minipage} & \begin{minipage}[t]{0.23\columnwidth}\raggedright
Close Button\strut
\end{minipage} & \begin{minipage}[t]{0.65\columnwidth}\raggedright
Close keypad after submission or cancel input\strut
\end{minipage}\tabularnewline
\begin{minipage}[t]{0.04\columnwidth}\raggedright
5\strut
\end{minipage} & \begin{minipage}[t]{0.23\columnwidth}\raggedright
Backspace\strut
\end{minipage} & \begin{minipage}[t]{0.65\columnwidth}\raggedright
Deletes one character to its left each time the key is pressed\strut
\end{minipage}\tabularnewline
\begin{minipage}[t]{0.04\columnwidth}\raggedright
6\strut
\end{minipage} & \begin{minipage}[t]{0.23\columnwidth}\raggedright
CLR\strut
\end{minipage} & \begin{minipage}[t]{0.65\columnwidth}\raggedright
Clears all numerical values in the input field\strut
\end{minipage}\tabularnewline
\bottomrule
\end{longtable}


\section{Turning the Power
ON/OFF}

The GL controller software consists of the PLC and embedded computer
system powered by GNU/Linux that communicates with each other to provide
the user interface for the GL controller, collectively called the
\emph{\textbf{GL controller system}}. The GL controller system and the
test chamber are connected to the same main power switch.


\subsection{Turning On the Power}

Turn on the GL controller system and chamber by performing the following
procedure.

\textbf{\emph{Procedure:}}

\begin{enumerate}
\def\labelenumi{\arabic{enumi}.}
\item
  Turn on the circuit breaker of the instrumentation by setting the main
  power breaker in the ON position.
\item
  The GL controller system starts. Within in a minute, the HMI display
  comes on. The GL controller system automatically logs the default user
  in for operation.

  \begin{figure}
  \centering
  
\includegraphics[width=0.95\textwidth]{figures/2008961489-S4_3-fig1.PNG}
  \caption{GL controller system touch-screen display}
  \end{figure}\FloatBarrier
\item
  The system is ready for operation.
\end{enumerate}


\subsection{Turning Off the
Power}

Turning off the GL controller system and chamber requires the opposite
procedure of the previous section, namely by setting the main power
breaker (on the back of the chamber) in the OFF position.


\section{User Accounts, System Menu and Login
Status}

When the main power breaker is set in the ON position, within a minute,
the display on the HMI will come on. The system automatically logs the
default user, called \emph{\textbf{localhost}}, in to operate the
chamber. The localhost is a special user account that the system grants
access to control the GL controller on the HMI. The system status bar
(see figure below) indicates the GL controller system is ready for
operation. The operator must enable the floating keyboard/keypad,
accessible via the system menu (indicated by the arrow), to allow the
system to receive any input parameter during operation.

\begin{figure}
\centering

\includegraphics[width=0.95\textwidth]{figures/3142074145-S1_6-hmi-login-status-001a.PNG}
\caption{Login status on HMI; GL system is ready for operation.}
\end{figure}\FloatBarrier

If the GL controller is accessed remotely via a Web browser, the user
must log in to operate the chamber, as depicted in the figure. The
optional floating keyboard/keypad can be enabled for on-screen operation
to receive any input parameter during operation. However, the mouse and
keyboard of the PC that launched the Web browser can be used to control
and operate the chamber.

\begin{figure}
\centering

\includegraphics[width=0.95\textwidth]{figures/1615289703-S1_6-web-login-status-001a.PNG}
\caption{Login status on Web display; GL system is ready for operation.}
\end{figure}\FloatBarrier


\subsection{System Menu}

The system menu provides full access to the system operation and
configuration menus. The floating keyboard/keypad can be enabled under
this menu to provide login and input support on the system when
operating on the HMI.

\begin{figure}
\centering

\includegraphics[width=0.95\textwidth]{figures/1330325741-hamburger-menu-000c.PNG}
\caption{Hamburger Menu}
\end{figure}\FloatBarrier

\begin{longtable}[]{@{}lll@{}}
\toprule
\begin{minipage}[b]{0.04\columnwidth}\raggedright
No.\strut
\end{minipage} & \begin{minipage}[b]{0.23\columnwidth}\raggedright
Name\strut
\end{minipage} & \begin{minipage}[b]{0.65\columnwidth}\raggedright
Description: Function/Operations\strut
\end{minipage}\tabularnewline
\midrule
\endhead
\begin{minipage}[t]{0.04\columnwidth}\raggedright
1\strut
\end{minipage} & \begin{minipage}[t]{0.23\columnwidth}\raggedright
Show/Hide Button\strut
\end{minipage} & \begin{minipage}[t]{0.65\columnwidth}\raggedright
Press this right angle to hide (close) the Hamburger Menu\strut
\end{minipage}\tabularnewline
\begin{minipage}[t]{0.04\columnwidth}\raggedright
2\strut
\end{minipage} & \begin{minipage}[t]{0.23\columnwidth}\raggedright
Keyboard\strut
\end{minipage} & \begin{minipage}[t]{0.65\columnwidth}\raggedright
Keyboard is disabled when the diagonal bar is over it; the system cannot
receive input without the keyboard when operating on the HMI; press on
the keyboard to enable it.\strut
\end{minipage}\tabularnewline
\begin{minipage}[t]{0.04\columnwidth}\raggedright
3\strut
\end{minipage} & \begin{minipage}[t]{0.23\columnwidth}\raggedright
Home\strut
\end{minipage} & \begin{minipage}[t]{0.65\columnwidth}\raggedright
Displays home page of the GL operating system; the highlighted menu
indicates its active status in the standard display.\strut
\end{minipage}\tabularnewline
\begin{minipage}[t]{0.04\columnwidth}\raggedright
4\strut
\end{minipage} & \begin{minipage}[t]{0.23\columnwidth}\raggedright
Notification\strut
\end{minipage} & \begin{minipage}[t]{0.65\columnwidth}\raggedright
Displays system notification and information, e.g.~alarm, system error
and information\strut
\end{minipage}\tabularnewline
\begin{minipage}[t]{0.04\columnwidth}\raggedright
5\strut
\end{minipage} & \begin{minipage}[t]{0.23\columnwidth}\raggedright
Main menu\strut
\end{minipage} & \begin{minipage}[t]{0.65\columnwidth}\raggedright
Displays system main menu, identical to the menu bar\strut
\end{minipage}\tabularnewline
\begin{minipage}[t]{0.04\columnwidth}\raggedright
6\strut
\end{minipage} & \begin{minipage}[t]{0.23\columnwidth}\raggedright
Configuration\strut
\end{minipage} & \begin{minipage}[t]{0.65\columnwidth}\raggedright
Configuration is the \textbf{Setup} submenu that contains common and
most used configuration operations; provided here as a quick-access
button.\strut
\end{minipage}\tabularnewline
\begin{minipage}[t]{0.04\columnwidth}\raggedright
7\strut
\end{minipage} & \begin{minipage}[t]{0.23\columnwidth}\raggedright
User\strut
\end{minipage} & \begin{minipage}[t]{0.65\columnwidth}\raggedright
User login prompt; similar to the user icon in the status bar\strut
\end{minipage}\tabularnewline
\bottomrule
\end{longtable}

\textbf{\emph{Note:}} The system menu may appear differently depending
on when or where it was accessed from. The above figure shows the system
menu being accessed from within the main home page. The following
figures depict the system menu being accessed from within the
\textbf{Monitor} menu and the \textbf{Settings} menu, respectively.

\textbf{System menu} accessed from within the \textbf{Monitor} page

\begin{figure}
\centering

\includegraphics[width=0.95\textwidth]{figures/1621260357-S1_6_1-image2aa.PNG}
\caption{System menu accessed from under the Monitor menu}
\end{figure}\FloatBarrier

\textbf{System menu} accessed from within the \textbf{Settings} page

\begin{figure}
\centering

\includegraphics[width=0.95\textwidth]{figures/2534120869-S1_6_1-image3aa.PNG}
\caption{System menu accessed from under the Settings menu}
\end{figure}\FloatBarrier


\subsection{Login Status \& User
Accounts}

When an operator logs out of the GL controller system (via HMI or Web
browser), all vital operating controls in the UI are locked out, while
the chamber continues in its current operating mode. For example, if the
user initiates Constant 2 and then logs out, the chamber continues in
its Constant 2 mode, as depicted in the figure below. If the user
initiates Profile 1 and then logs out, the chamber remains in the
program mode until the program completes all the instructions, while all
operating controls in the UI are locked out. The status bar displays the
diagonal bar over the user icon indicating the user has logged out.

\begin{figure}
\centering

\includegraphics[width=0.95\textwidth]{figures/2961643740-S2_6_2N-login-002a.PNG}
\caption{GL controller UI is locked out}
\end{figure}\FloatBarrier

\begin{longtable}[]{@{}lll@{}}
\toprule
\begin{minipage}[b]{0.04\columnwidth}\raggedright
No.\strut
\end{minipage} & \begin{minipage}[b]{0.23\columnwidth}\raggedright
Name\strut
\end{minipage} & \begin{minipage}[b]{0.65\columnwidth}\raggedright
Description\strut
\end{minipage}\tabularnewline
\midrule
\endhead
\begin{minipage}[t]{0.04\columnwidth}\raggedright
1\strut
\end{minipage} & \begin{minipage}[t]{0.23\columnwidth}\raggedright
Login Status\strut
\end{minipage} & \begin{minipage}[t]{0.65\columnwidth}\raggedright
Diagonal bar over user icon indicates the GL controller system has been
logged out; no user is logged in on this UI.\strut
\end{minipage}\tabularnewline
\begin{minipage}[t]{0.04\columnwidth}\raggedright
2\strut
\end{minipage} & \begin{minipage}[t]{0.23\columnwidth}\raggedright
Home\strut
\end{minipage} & \begin{minipage}[t]{0.65\columnwidth}\raggedright
Home page of GL controller. If any menu button is pressed/clicked, the
login dialog appears for the user to log in.\strut
\end{minipage}\tabularnewline
\begin{minipage}[t]{0.04\columnwidth}\raggedright
3\strut
\end{minipage} & \begin{minipage}[t]{0.23\columnwidth}\raggedright
Operation Mode\strut
\end{minipage} & \begin{minipage}[t]{0.65\columnwidth}\raggedright
If a user is not logged in, operation buttons for Constant or Program
are locked and grayed out.\strut
\end{minipage}\tabularnewline
\begin{minipage}[t]{0.04\columnwidth}\raggedright
4\strut
\end{minipage} & \begin{minipage}[t]{0.23\columnwidth}\raggedright
Monitor\strut
\end{minipage} & \begin{minipage}[t]{0.65\columnwidth}\raggedright
This menu displays the current status of the GL controller.\strut
\end{minipage}\tabularnewline
\begin{minipage}[t]{0.04\columnwidth}\raggedright
5\strut
\end{minipage} & \begin{minipage}[t]{0.23\columnwidth}\raggedright
Set Mode\strut
\end{minipage} & \begin{minipage}[t]{0.65\columnwidth}\raggedright
If a user is not logged in, contents of Constant Setup or Program Setup
are viewable in read-only mode; they cannot be modified.\strut
\end{minipage}\tabularnewline
\begin{minipage}[t]{0.04\columnwidth}\raggedright
6\strut
\end{minipage} & \begin{minipage}[t]{0.23\columnwidth}\raggedright
Setup\strut
\end{minipage} & \begin{minipage}[t]{0.65\columnwidth}\raggedright
If a user is not logged in, only few submenus are accessible; the rest
are locked out.\strut
\end{minipage}\tabularnewline
\begin{minipage}[t]{0.04\columnwidth}\raggedright
7\strut
\end{minipage} & \begin{minipage}[t]{0.23\columnwidth}\raggedright
Operation/Set\strut
\end{minipage} & \begin{minipage}[t]{0.65\columnwidth}\raggedright
These buttons are grayed out when no user is logged in\strut
\end{minipage}\tabularnewline
\bottomrule
\end{longtable}

Only authorized users can log in to the system to operate the chamber.
Two users, localhost and admin, are available by default when the
chamber is shipped from the manufacturer. The credentials of these two
accounts are listed in the following table. Refer to Section 7.4 on how
to administer the GL controller system with different user accounts with
different privileges.

\begin{longtable}[]{@{}lll@{}}
\toprule
\begin{minipage}[b]{0.14\columnwidth}\raggedright
User Name\strut
\end{minipage} & \begin{minipage}[b]{0.14\columnwidth}\raggedright
Password\strut
\end{minipage} & \begin{minipage}[b]{0.63\columnwidth}\raggedright
Description\strut
\end{minipage}\tabularnewline
\midrule
\endhead
\begin{minipage}[t]{0.14\columnwidth}\raggedright
localhost\strut
\end{minipage} & \begin{minipage}[t]{0.14\columnwidth}\raggedright
\strut
\end{minipage} & \begin{minipage}[t]{0.63\columnwidth}\raggedright
GL system grants access to GL controller under localhost user\strut
\end{minipage}\tabularnewline
\begin{minipage}[t]{0.14\columnwidth}\raggedright
admin\strut
\end{minipage} & \begin{minipage}[t]{0.14\columnwidth}\raggedright
admin\strut
\end{minipage} & \begin{minipage}[t]{0.63\columnwidth}\raggedright
An account which has specific roles and privileges to manage and control
the GL controller system operation, settings and configurations\strut
\end{minipage}\tabularnewline
\bottomrule
\end{longtable}

\textbf{\emph{Note:}} To protect the GL controller and chamber from
unauthorized users and to keep the GL controller system secure, the
admin password must be changed to something much more secure.

The password for the localhost user is stored internally in the system,
and only the system can log in using this account; hence, there is no
password for this account listed in the table (above). If you logged out
of the localhost account by accident or on purpose, you have two options
to log back in.

\begin{enumerate}
\def\labelenumi{\arabic{enumi}.}
\tightlist
\item
  Press the \textbf{Auto login} button underneath the login prompt to
  log back in to the localhost account.\\
\item
  Log in to admin or your own user account. Refer to Section 2.6.3 for
  details.
\end{enumerate}


\subsection{How to Log in via
HMI}

The following procedure illustrates how to log in to the GL controller
system using the admin account via HMI after logging out of the
localhost. The admin login credentials are listed in the above table.

\textbf{\emph{Procedure:}}

\begin{enumerate}
\def\labelenumi{\arabic{enumi}.}
\item
  Press the hamburger menu in the status bar.

  \begin{figure}
  \centering
  
\includegraphics[width=0.95\textwidth]{figures/661183384-S1_6_3-system-menu-001a.PNG}
  \caption{Access the system menu}
  \end{figure}\FloatBarrier
\item
  Check to make sure the keyboard has been enabled. If a diagonal bar is
  over the keyboard icon, press the keyboard to enable it (see arrow),
  then press \textbf{User} to log in (see arrow).

  \begin{figure}
  \centering
  
\includegraphics[width=0.95\textwidth]{figures/895746391-S1_6_3-step2-001a.PNG}
  \caption{Enable the keyboard and log in via the User login prompt}
  \end{figure}\FloatBarrier
\item
  Press the \textbf{User Name} field under the Please Login prompt (see
  arrow), the keyboard will appear at the bottom; enter \textbf{admin}
  in the User Name field on the keyboard (see arrow) and press the check
  button (see arrow). Repeat the process for the password.
  \_\textbf{\emph{Note:}} To cancel login, press X anytime during user
  name or password input.

  \begin{figure}
  \centering
  
\includegraphics[width=0.95\textwidth]{figures/701013834-S1_6_3-step3-login-cridens-001a.PNG}
  \caption{Log in using the admin user account}
  \end{figure}\FloatBarrier
\item
  Press the \textbf{Submit} button to log in.

  \begin{figure}
  \centering
  
\includegraphics[width=0.95\textwidth]{figures/2383628226-S1_6_3-step4-001b.PNG}
  \caption{Press Submit to log in}
  \end{figure}\FloatBarrier
\item
  The system is now ready for operation under the admin user account.

  \begin{figure}
  \centering
  
\includegraphics[width=0.95\textwidth]{figures/3945223490-S1_6_3-step5-001a.PNG}
  \caption{Successful login; system is ready for operation}
  \end{figure}\FloatBarrier
\end{enumerate}

\textbf{\emph{Note:}} The above procedure can be used to log in to any
user account on the system (except the localhost), provided that account
exists. Refer to Section 7.4 for details on how to manage user accounts.


\subsection{How to Log in via Remote
Access}

The following procedure illustrates how to log in to the admin account
on the GL controller system via a Web browser. The admin login
credentials are listed in the table in Section 2.6.3.

\textbf{\emph{Note:}} Logging in via a Web browser (for remote access),
the keyboard and mouse of the PC on which the browser was launched can
be used to enter the admin credentials. No need to enable the on-screen
keyboard.

\textbf{\emph{Procedure:}}

\begin{enumerate}
\def\labelenumi{\arabic{enumi}.}
\item
  Click the user icon on the status bar (see arrow).

  \begin{figure}
  \centering
  
\includegraphics[width=0.95\textwidth]{figures/2289870150-S1_6_4-step1-001a.PNG}
  \caption{Access the user login dialog}
  \end{figure}\FloatBarrier
\item
  Click the User Name field and enter \textbf{admin} (see arrow). Click
  the Password field and enter the password for \textbf{admin} (see
  arrow), then press Enter or click the Submit button to log in (see
  arrow).

  \begin{figure}
  \centering
  
\includegraphics[width=0.95\textwidth]{figures/3778126336-S1_6_4-step2-002a.PNG}
  \caption{Entering the admin login credentials}
  \end{figure}\FloatBarrier
\item
  The system is now ready for operation under the admin user account.

  \begin{figure}
  \centering
  
\includegraphics[width=0.95\textwidth]{figures/3234756050-S1_6_4-step3-001.PNG}
  \caption{Successful login; system is ready for operation}
  \end{figure}\FloatBarrier
\end{enumerate}


\section{Three Different Ways to Access GL Controller System
Menu}

The GL software system was designed and implemented to be robust,
providing the UI with multiple ways to accomplish the same task. This
flexibility allows users--novice and expert alike--the ability to access
system menus to complete their task effectively and efficiently. As
depicted in the following figure, there are three known ways to access
the system menu to control and operate the chamber. We outline these
three approaches as follows.

\begin{figure}
\centering

\includegraphics[width=0.95\textwidth]{figures/2771776335-S1_6_1-system-menus.PNG}
\caption{GL system menu; viewed from within the main home page}
\end{figure}\FloatBarrier


\subsection{System Menu via Menu
Bar}

The following figure depicts the system menu via the menu bar. The menu
bar remains fixed throughout the operation in any operating mode. In
addition to the four GL controller menus, this menu bar allows the user
to return to the main home page using the home button.

\begin{figure}
\centering

\includegraphics[width=0.95\textwidth]{figures/3531631476-S1_6_1-menubar-000.PNG}
\caption{System menu via menu bar}
\end{figure}\FloatBarrier


\subsection{System Menu via Menu
Buttons}

The following figure depicts the GL controller main home page consisting
of four menu buttons for a quick access. These menu buttons are
available within the main home page only. Whenever the chamber is
powered on from a cold start (i.e., when the main power breaker is set
in the ON position), the GL controller system will display the main home
page when it is ready for operation.

\begin{figure}
\centering

\includegraphics[width=0.95\textwidth]{figures/3652124681-S1_6_1-system-menu-buttons-001a.PNG}
\caption{System menu via menu buttons}
\end{figure}\FloatBarrier


\subsection{System Menu via Hamburger
Button}

The following figure depicts the system menu via the hamburger button in
the status bar. As indicated by the arrows, each system menu displays
different contents depending on how (i.e., from within what menu page)
it was accessed. The system menu on the left was accessed from within
the main home page, while the second one on the right was accessed from
within the \textbf{Setup} page; hence, the highlight menu name.

\begin{figure}
\centering

\includegraphics[width=0.95\textwidth]{figures/1496763686-S1_6_3-system-menu-hamburger-butoon-001a.PNG}
\caption{System menu via hamburger button}
\end{figure}\FloatBarrier

The system menu also contains (and displays) additional information
about the GL controller system that includes end-user license and
agreement, user support page for online operation manual, specific
version of the GL software system, the embedded operation manual and
more.


\section{Setting the Alarm and Warning
Buzzer}

The following procedure outlines the steps to manage the beep volume in
the event of an alarm or a warning. Refer to Section 7.13 for detail.

\textbf{\emph{Procedure:}}

\begin{enumerate}
\def\labelenumi{\arabic{enumi}.}
\item
  From the main home page, press the hamburger button (1) in the status
  bar to access the system menu. Press \textbf{Configuration} (2),
  followed by \textbf{Set Sound} submenu (3).

  \begin{figure}
  \centering
  
\includegraphics[width=0.95\textwidth]{figures/3588522138-S1_7-step1.PNG}
  \caption{Accessing sound setting options}
  \end{figure}\FloatBarrier
\item
  By default, the Alarm volume and Warning volume are at maximum
  setting, as indicated by the volume slider bar.

  \begin{figure}
  \centering
  
\includegraphics[width=0.95\textwidth]{figures/1253536369-S1_7-step2-000.PNG}
  \caption{Alarm and Warning beep level}
  \end{figure}\FloatBarrier
\item
  To adjust the beep volume, change the volume level on the slider bar.
\item
  To turn off the beep volume, either press the speaker icon (1) or set
  the volume level on the slider bar to its minimum setting (2).

  \begin{figure}
  \centering
  
\includegraphics[width=0.95\textwidth]{figures/2697260505-S1_7-step4a.PNG}
  \caption{Beep volume adjustment}
  \end{figure}\FloatBarrier
\item
  Press the Save button (3) to save the setting. To cancel the setting,
  press the Back button (4) and press Yes to confirm the cancel.

  \begin{figure}
  \centering
  
\includegraphics[width=0.95\textwidth]{figures/4264269375-S1_7-step5a.PNG}
  \caption{Save the beep volume setting}
  \end{figure}\FloatBarrier
\end{enumerate}

\begin{longtable}[]{@{}l@{}}
\toprule
\begin{minipage}[b]{0.97\columnwidth}\raggedright
CAUTION\strut
\end{minipage}\tabularnewline
\midrule
\endhead
\begin{minipage}[t]{0.97\columnwidth}\raggedright
\textbf{It is strongly advised to keep the beep on whenever possible to
prevent any delay in the detection of an alarm or warning.} If the beep
is turned off, alarm and warning are not notified with sound and are
indicated only with red blinking of the operation lamp and on warning
occurrence screen. \textbf{Set the beep volume in accordance with the
ambient environment.}\strut
\end{minipage}\tabularnewline
\bottomrule
\end{longtable}


\section{GL Controller and Chamber
Date/Time}

The GL controller and chamber are configured with correct date/time
based on the local time zone of the manufacturer using Eastern Standard
Time (EST, New York) with UTC (Coordinated Universal Time) as a point of
reference to ensure the correct conversion to local time for any time
zone. If your local time zone is EST, then there is no need to perform
date/time configuration. If your local time zone is not EST, refer to
Section 7.2 for detailed configuration of the GL controller date/time
and its operation.


\chapter{Convenient Functions}

This chapter introduces convenient functions provided in this equipment.


\section{Automatic/Manual Switching of Refrigeration
Capacity}

The refrigeration capacity is controlled automatically so that it is
optimally controlled depending on the temperature (humidity) set point.
In ordinary operation, always set the refrigeration capacity to Auto.

\begin{longtable}[]{@{}l@{}}
\toprule
\begin{minipage}[b]{0.97\columnwidth}\raggedright
Notice\strut
\end{minipage}\tabularnewline
\midrule
\endhead
\begin{minipage}[t]{0.97\columnwidth}\raggedright
\textbf{When the refrigeration capacity is set manually to the maximum
value, it may not be possible to retain the set temperature. Check the
setting before conducting a test.} There is a possibility that the
temperature cannot be maintained and the specimen will be damaged.
\textbf{If the refrigeration capacity setting is changed during
operation, control may become unstable, and the temperature and humidity
fluctuation of the specification performance may not be guaranteed. (The
refrigeration capacity setting cannot be changed during program
operation.)}\strut
\end{minipage}\tabularnewline
\bottomrule
\end{longtable}

\textbf{\emph{Note:}} If an issue occurs during testing, set
refrigeration capacity to manual control.


\subsection{Automatic Control of Refrigeration
Capacity}

The following procedure outlines the steps to set refrigeration in Auto
mode. The procedure assumes that the chamber has already been turned On.
If not, set the main power breaker in the ON position. Within a minute,
the display on the HMI will come on. The system automatically logs in
using the default user account (localhost) to operate the chamber.

\begin{figure}
\centering

\includegraphics[width=0.95\textwidth]{figures/3074297625-S6_2_2-p1-step1.PNG}
\caption{Operating Panel at Startup}
\end{figure}\FloatBarrier

\textbf{\emph{Procedure:}}

\begin{enumerate}
\def\labelenumi{\arabic{enumi}.}
\item
  Press the \textbf{Hamburger} menu (1) in the status bar. Check to
  confirm the floating keyboard is enabled. A diagonal bar over the
  keyboard indicates it is disabled. Press the keyboard (2) to enable
  it. Press \textbf{Set Mode} (3) to access the \textbf{Set Mode} page.

  \begin{figure}
  \centering
  
\includegraphics[width=0.95\textwidth]{figures/764530790-S3_2-step2.PNG}
  \caption{Enable the floating keyboard in the system menu}
  \end{figure}\FloatBarrier
\item
  Press \textbf{Constant 1} in the title bar to access \textbf{Constant
  1} configuration page. Three constant modes (No.~1, 2, and 3) are
  available for configuration and operation. While Constant 1 is
  selected for configuration, the following procedure applies to any of
  them.

  \begin{figure}
  \centering
  
\includegraphics[width=0.95\textwidth]{figures/917521979-S6_2_2-p1-step4a.PNG}
  \caption{Select and configure Constant 1 Mode}
  \end{figure}\FloatBarrier
\item
  Configuration options are spread over in two pages, as depicted in the
  figure, to allow configuration for all options, such as Set Point,
  High \& Low Limits, Time Signals and Product Temperature Control (if
  available). To set Time Signals and Product Temperature Control
  options, select page 2.

  \begin{figure}
  \centering
  
\includegraphics[width=0.95\textwidth]{figures/2211636124-S6_2_2-p1-step5.PNG}
  \caption{Constant 1 configuration options}
  \end{figure}\FloatBarrier
\item
  Press the set point field of temperature (see arrow). Adjust the set
  point value of temperature in the keypad. Press the check button to
  apply the input value; or press anywhere outside of the keypad to
  close it and apply the setting. Set values take effect immediately
  after the check button is applied and the keypad closes.

  \begin{figure}
  \centering
  
\includegraphics[width=0.95\textwidth]{figures/41414174-S6_3_step6-001b.PNG}
  \caption{Adjust temperature set point value}
  \end{figure}\FloatBarrier
\item
  Humidity control can be turned ON/OFF as needed. When the Off button
  is unlit (black), humidity is OFF. To turn on humidity, press the On
  button. Press the set point field of humidity (see arrow) to adjust
  the set point value via the keypad. Press the check button in the
  keypad to apply the input value. Set values take effect immediately
  after the check button is applied and the keypad closes.

  \begin{figure}
  \centering
  
\includegraphics[width=0.95\textwidth]{figures/1582714919-S5_8_3-step7b.PNG}
  \caption{Setting humidity set point}
  \end{figure}\FloatBarrier
\item
  As depicted in the figure, refrigeration capacity is available in auto
  and manual (off, 25\%, 50\%, 75\% and 100\%). Default setting is in
  auto (see arrow). To select between the manual capacity, press the
  desired refrigeration capacity. The selected option will lit and take
  effect immediately.

  \begin{figure}
  \centering
  
\includegraphics[width=0.95\textwidth]{figures/2003532222-S6_5_1-step4a.PNG}
  \caption{Set refrigeration to Auto}
  \end{figure}\FloatBarrier
\item
  When finished setting temperature (and humidity), press the Operation
  Mode in the Menu bar (1) or press an area near or on the Standby in
  the status bar (2) to start \textbf{Constant 1} mode.

  \begin{figure}
  \centering
  
\includegraphics[width=0.95\textwidth]{figures/1669112367-S6_3-step8-000a.PNG}
  \caption{Returning to the Operation Mode page}
  \end{figure}\FloatBarrier
\end{enumerate}


\section{Manual Draining/Refill of Humidity
Pan}

For operation details, refer to the GL Chamber Operation Manual.


\section{Humidification Delay
Control}

The following function diagrams explain the behavior and operation of
humidity. Humidity control starts after the temperature reaches the set
point (within temperature attainment range or longer than temperature
attainment time). Humidity control starts in the following timing.


\subsection{Example 1}

The following diagram depicts the behavior of humidity control when the
temperature in the test area is higher than the temperature set point.

\begin{figure}
\centering

\includegraphics[width=0.95\textwidth]{figures/1910060217-S2_3-fig-ex1.PNG}
\caption{Example 1: Humidification delay control.}
\end{figure}\FloatBarrier


\subsection{Example 2}

The following diagram depicts the behavior of humidity control when the
temperature in the test area is lower than the temperature set point.

\begin{figure}
\centering

\includegraphics[width=0.95\textwidth]{figures/4227507477-S2_3-fig-ex2.PNG}
\caption{Example 2: Humidification delay control.}
\end{figure}\FloatBarrier

\begin{longtable}[]{@{}l@{}}
\toprule
\begin{minipage}[b]{0.97\columnwidth}\raggedright
\textbf{Reference}\strut
\end{minipage}\tabularnewline
\midrule
\endhead
\begin{minipage}[t]{0.97\columnwidth}\raggedright
The temperature attainment range of the humidification delay control is
set to 1 °C and the temperature attainment time is set to 60 seconds;
both settings are fixed and cannot be changed. Even if the
\textbf{attainment judgment criteria} on the maintenance setting screen
change, they have no bearing in the operation of this function.\strut
\end{minipage}\tabularnewline
\bottomrule
\end{longtable}

For humidification delay control, its ON/OFF option is under the Set
Point panel in Page 1 of Constant configuration page. Refer to Section
3.1.1.


\chapter{Constant Value Operation}

The chamber has four operating modes: \textbf{Constant},
\textbf{Program}, \textbf{Alarm} and \textbf{Standby}. The
\textbf{Alarm} mode requires an immediate attention to the chamber
operating condition with red alerts, while \textbf{Standby} simply
indicates the chamber is in a standby mode and is not under any
operating condition. Details of \textbf{Program} mode will be discussed
in Chapter 5. Some aspects of \textbf{Alarm} conditions will be
discussed in Chapter 10.

This chapter focuses on the constant mode operation. It discusses the
method of setting and performing constant value operation. It explains
how to start, stop, monitor and extract information of constant value
operation.


\section{What is a Constant Value
Operation?}

A constant value operation is a method of keeping temperature and
humidity in the test area at constant values throughout the operation.
Since the operation is performed in constant values for set points, it
is called the \emph{\textbf{constant value operation}}. The following
graph depicts an operation at constant temperature of 30 °C.

\begin{figure}
\centering

\includegraphics[width=0.95\textwidth]{figures/1695782616-S3_1-const-temp-op-01.PNG}
\caption{Graph of constant value operation}
\end{figure}\FloatBarrier


\section{Constant Temperature (\& Humidity)
Setup}

For constant mode operation, it is necessary to input the desired
temperature and humidity using one of the available constant options.
Three constant modes are available for Constant Value Operation. They
are Constant 1, Constant 2 and Constant 3.

The following procedure outlines the step for setting up Constant 1 with
temperature and humidity. While Constant 1 is selected for
configuration, the procedure applies to any of the three constant mode
options.

\textbf{\emph{Procedure:}}

\begin{enumerate}
\def\labelenumi{\arabic{enumi}.}
\item
  If the GL controller system is already on, proceed to Step 2. From a
  cold start, set the main power breaker in the ON position. Within a
  minute, the touch-screen display will come on. The system
  automatically logs in using the default localhost user account to
  operate the chamber.

  \begin{figure}
  \centering
  
\includegraphics[width=0.95\textwidth]{figures/3074297625-S6_2_2-p1-step1.PNG}
  \caption{Operating Panel at Startup}
  \end{figure}\FloatBarrier
\item
  Press the \textbf{Hamburger} menu (1) in the status bar. Check to
  confirm the floating keyboard is enabled. A diagonal bar over the
  keyboard indicates it is disabled. Press the keyboard (2) to enable
  it. Press \textbf{Set Mode} (3) to access the \textbf{Set Mode} menu.

  \begin{figure}
  \centering
  
\includegraphics[width=0.95\textwidth]{figures/764530790-S3_2-step2.PNG}
  \caption{Enable the floating keyboard in the system menu}
  \end{figure}\FloatBarrier
\item
  Press \textbf{Constant 1} in the title bar to access \textbf{Constant
  1} configuration page (see arrow). Three constant modes (No.~1, 2, and
  3) are available for configuration and operation.

  \begin{figure}
  \centering
  
\includegraphics[width=0.95\textwidth]{figures/2344501277-S3_2-step3-002a.PNG}
  \caption{Select and configure Constant 1 Mode}
  \end{figure}\FloatBarrier
\item
  Under the \textbf{Set Point} panel, configuration options are
  available for Temperature, Humidity and Refrigeration, followed by
  \textbf{Time Signals}, \textbf{Product Temperature Control} (if
  available) and \textbf{High \& Low Limits} panels through this
  scrollable page.

  \begin{figure}
  \centering
  
\includegraphics[width=0.95\textwidth]{figures/252542159-S3_2-step4-002.PNG}
  \caption{Constant 1 configuration options}
  \end{figure}\FloatBarrier
\item
  To set new temperature set point, press the set point field (see
  arrow). Adjust the set point value in the keypad. Press the check
  button (or anywhere outside the keypad) to apply the input value. Set
  temperature mode by selecting Product or Air (See arrow).

  \begin{figure}
  \centering
  
\includegraphics[width=0.95\textwidth]{figures/1731941424-S3_2-step5-002-1a.PNG}
  \caption{Adjust temperature set point value}
  \end{figure}\FloatBarrier
\item
  Humidity control can be turned On/Off as needed. The set point can
  only be adjusted when humidity is On. When the Off button is unlit
  (black), humidity is Off. Press the On button (see arrow) to turn it
  on. Press the set point field (see arrow) to adjust the set point
  value via the keypad.

  \begin{figure}
  \centering
  
\includegraphics[width=0.95\textwidth]{figures/2996209063-S3_2-step6-002a.PNG}
  \caption{Setting humidity set point}
  \end{figure}\FloatBarrier
\end{enumerate}


\section{Refrigeration Setting}

In default setting, refrigeration is controlled automatically. This
default setting is found in Page 1 of the Constant configuration page.
However, under normal circumstances refrigeration capacity can be
controlled manually. Refrigeration setting options are available in
Auto, Off or manual operation in 25\%, 50\%, 75\% and 100\%. The
following procedure illustrates these options.

\textbf{\emph{Procedure:}}

\begin{enumerate}
\def\labelenumi{\arabic{enumi}.}
\item
  The \textbf{Set Point} panel in the figure (below) shows the
  refrigeration set in Auto mode.

  \begin{figure}
  \centering
  
\includegraphics[width=0.95\textwidth]{figures/3555960265-S3_3-step1-001a.PNG}
  \caption{Refrigeration settings}
  \end{figure}\FloatBarrier
\item
  To control the refrigeration system manually, select the manual
  cooling output percentage 25\%, 50\%, 75\% or 100\%. Set value takes
  effect immediately after selection.
\end{enumerate}

\begin{longtable}[]{@{}l@{}}
\toprule
\begin{minipage}[b]{0.97\columnwidth}\raggedright
Reference\strut
\end{minipage}\tabularnewline
\midrule
\endhead
\begin{minipage}[t]{0.97\columnwidth}\raggedright
For details on refrigeration setup, see Section 3.1. When you want to
change the mode from Auto to Manual during operation, manually set the
capacity of the refrigerator using one of these options: 25\%, 50\%,
75\% or 100\%. When refrigeration is changed from Auto to Manual (with
specified capacity), refrigeration stops immediately. This may cause
control to become unstable.\strut
\end{minipage}\tabularnewline
\begin{minipage}[t]{0.97\columnwidth}\raggedright
When refrigeration is set to Auto, the manual setting is not reflected
in the operation of the refrigeration system.\strut
\end{minipage}\tabularnewline
\bottomrule
\end{longtable}

\textbf{\emph{Note:}} Proper testing is required to verify the above
reference.


\section{High/Low Limit Setting}

The values set with Absolute High and Absolute Low (Abs High and Abs
Low) form the range in which temperature and humidity set point can be
entered in Constant Setup or Program Setup. These two values set the
extreme high and extreme low set points in a chamber operation. Should a
temperature in test area rise above the Abs High set point, alarm will
be tripped. The same is true if it goes below the Abs Low set point. The
following diagram depicts the behavior of these two situations.

\begin{figure}
\centering

\includegraphics[width=0.95\textwidth]{figures/1025545634-S3_4-HL-graph.PNG}
\caption{High/Low limit setup condition}
\end{figure}\FloatBarrier

Two different scenarios are described as follows.

\textbf{Abs High and Abs Low}

Criteria used for performing protection in the test area. Set values
greater than the temperature set points in the test area by 10 °C or
more. If the equipment detects a fault condition, the equipment stops
operation after issuing an alarm.

\textbf{Upper Deviation}

Criteria used for performing for specimens. Change the setting depending
on the test specimen. If the equipment detects a deviation fault, the
equipment stops the header (humidifier) after issuing a warning. When
the temperature in the test area decreases to the temperature set point,
the equipment returns to the normal control.

The following procedure outlines the steps to adjust Abs High, Abs Low
and Upper Deviation for \textbf{Constant 1}. Chapter 5 outlines the
steps for Program setup.

\textbf{\emph{Procedure:}}

\begin{enumerate}
\def\labelenumi{\arabic{enumi}.}
\item
  Confirm that the floating keyboard has been enabled. Refer to Section
  2.4.4.
\item
  From within the main home page, press the \textbf{Set Mode} button and
  press \textbf{Constant 1} title bar to open its configuration page
  (see arrow).

  \begin{figure}
  \centering
  
\includegraphics[width=0.95\textwidth]{figures/1005458508-S3_4a-step2-000a.PNG}
  \caption{Constant 1 configuration page}
  \end{figure}\FloatBarrier
\item
  Press the \textbf{Abs. High} value field to bring up the numeric input
  keypad.

  \begin{figure}
  \centering
  
\includegraphics[width=0.95\textwidth]{figures/468327071-S3_4-step3-002a.PNG}
  \caption{Adjusting Abs. High, Abs. Low temperature/humidity}
  \end{figure}\FloatBarrier
\item
  Enter a new temperature value on the keypad, taking note of the high
  and low limits. Press the check button to apply the new setting. To
  cancel the procedure, press the X button. Repeat the same procedure
  for the upper/lower deviation or humidity setting.

  \begin{figure}
  \centering
  
\includegraphics[width=0.95\textwidth]{figures/1648987967-S3_4-step4-002a.PNG}
  \caption{Enter new temperature setting}
  \end{figure}\FloatBarrier
\end{enumerate}


\section{Time Signals}

A contact by which your external device can be turned on and off with
the keys on the instrumentation can be installed. Time signals can be
used to apply voltage to control external devices.

Two time signals are offered in the standard model. Additional time
signals (3 through 12) are offered in the model with such options.
Default names for these time signals are Time Signal 1, Time Signal 2,
etc. They can be changed to reflect the description or purpose of the
external device(s). The following two figures depict examples of time
signal options set as default names and custom names.

\textbf{Example 1:} Time signal options using default names.

\begin{figure}
\centering

\includegraphics[width=0.95\textwidth]{figures/1592840808-S3_5-TS-example-001.PNG}
\caption{Chamber with eight (8) time signal options}
\end{figure}\FloatBarrier

\textbf{Example 2:} Time signal options using default and custom names

\begin{figure}
\centering

\includegraphics[width=0.95\textwidth]{figures/3237501349-S3_5-TS-example-002.PNG}
\caption{Chamber with ten (10) time signal options}
\end{figure}\FloatBarrier

Time Signal naming is configured under the \textbf{Configuration} page
of the \textbf{Setup} menu. The following procedure outlines the steps
to set time signals and customize their names.

\textbf{\emph{Procedure:}}

\begin{enumerate}
\def\labelenumi{\arabic{enumi}.}
\item
  If the GL controller system is already on and its touchscreen is in
  sleep mode, tab your finger on the touchscreen and proceed to Step 2.
  From a cold start, set the main power breaker in the ON position.
  Within a minute, the touch-screen display will come on. The system
  automatically logs in using the default localhost user account to
  operate the chamber.

  \begin{figure}
  \centering
  
\includegraphics[width=0.95\textwidth]{figures/3074297625-S6_2_2-p1-step1.PNG}
  \caption{Operating Panel at Startup}
  \end{figure}\FloatBarrier
\item
  Press the \textbf{Hamburger} menu (1) in the status bar. Check to
  confirm the floating keyboard is enabled. A diagonal bar over the
  keyboard indicates it is disabled. Press the keyboard (2) to enable
  it. Press \textbf{Configuration} (3); then press \textbf{Time Signal
  Names} (4).

  \begin{figure}
  \centering
  
\includegraphics[width=0.95\textwidth]{figures/1569113829-S3_5-step3-000b.PNG}
  \caption{Accessing Time Signal Setting}
  \end{figure}\FloatBarrier
\item
  To customize the name of Time Signal 1, press the Full name field (1).
  Enter a new name using the pop-up input keyboard (2). \emph{Note: If
  the input keyboard does not pop up, it has not been enabled; go back
  to Step 2.} Press the Check button (3) to apply the setting. Press the
  Save button (4) to save the setting. Or, press the Back button (5) to
  cancel the setting (and press Yes to confirm the action).

  \begin{figure}
  \centering
  
\includegraphics[width=0.95\textwidth]{figures/3757672429-S3_5-step3-TS-000b.PNG}
  \caption{Customize Time Signal names}
  \end{figure}\FloatBarrier
\item
  Short name (second column) and On/Off buttons (third and fourth
  columns) can also take descriptive names.

  \begin{figure}
  \centering
  
\includegraphics[width=0.95\textwidth]{figures/2218355775-S3_5-step4-TS-rename.PNG}
  \caption{Renaming Time signal 1}
  \end{figure}\FloatBarrier
\item
  Time signal can be operated via the On/Off button. This operation can
  be set using the On/Off button in the Constant configuration page
  (Constant 1 setting is depicted in the figure below). In Program mode,
  time signal is identified by its number (discussed in Chapter 5). As
  shown below, Time Signal 1 (now called \emph{\textbf{Control Ext1}})
  is in the OFF position. To turn it on, press the \textbf{RUNNING}
  button.

  \begin{figure}
  \centering
  
\includegraphics[width=0.95\textwidth]{figures/1843822913-S3_5-step5-002a.PNG}
  \caption{Enabling Time Signal operation}
  \end{figure}\FloatBarrier
\item
  The Full name of a time signal will appear in the standard display
  area during a Constant or Program operating mode as depicted in the
  figure.

  \begin{figure}
  \centering
  
\includegraphics[width=0.95\textwidth]{figures/3939821181-S3_5-step6-001a.PNG}
  \caption{How time signal is displayed during an chamber operation}
  \end{figure}\FloatBarrier
\end{enumerate}


\section{Setting Constant Name}

The default name of Constant 1 is \textbf{constant\_1}. This name can be
changed to something descriptive.

\textbf{\emph{Procedure:}}

\begin{enumerate}
\def\labelenumi{\arabic{enumi}.}
\item
  Press the \textbf{Hamburger} menu (1) in the status bar. Check to
  confirm the floating keyboard is enabled. A diagonal bar over the
  keyboard indicates it is disabled. Press the keyboard (2) to enable
  it. Press \textbf{Set Mode} (3) to access the \textbf{Set Mode} page.

  \begin{figure}
  \centering
  
\includegraphics[width=0.95\textwidth]{figures/764530790-S3_2-step2.PNG}
  \caption{Enable the floating keyboard in the system menu}
  \end{figure}\FloatBarrier
\item
  Press \textbf{Constant 1} in the title bar (see arrow) to access
  \textbf{Constant 1} configuration page.

  \begin{figure}
  \centering
  
\includegraphics[width=0.95\textwidth]{figures/244552609-S3_6-step2-accessing-const1-000a.PNG}
  \caption{Select and configure Constant 1 Mode}
  \end{figure}\FloatBarrier
\item
  Press the constant name field (1), edit the name using the pop-up
  keyboard (2) and press the check button (3) to apply the setting. The
  new name takes effect immediately.

  \begin{figure}
  \centering
  
\includegraphics[width=0.95\textwidth]{figures/136012415-S3_5-step3-002a.PNG}
  \caption{Renaming Constant 1}
  \end{figure}\FloatBarrier
\item
  The same procedure applies to \textbf{Constant 2} and \textbf{Constant
  3}.
\end{enumerate}


\section{Product Temperature Control
Options}

Temperature control options are available as Air and/or Product
Temperature Control. Since these functions are all optional, the
function that is not installed is not displayed or available for
operation. Thus the following feature is available only in a chamber
with such optional feature(s) installed.

The deviation limits are the allowable air temperature deviation from
the product temperature set point upper product temperature control
mode. The ``plus'' value will be added to the product temperature set
point to create the ``upper'' air temperature limit. The ``minus'' value
will be subtracted from the product temperature set point to create the
``lower'' air temperature limit.

When the product temp. process value is far away from the product temp.
set point, the controller will ``drive'' the air temperature set point
to this maximum upper (or lower) deviation limit. As the product temp.
process value gets closer to the product temp. set point, the air temp
set point will change so that it is less than the maximum upper (or
lower) deviation limit.

\textbf{How to determine the upper and lower deviation limits}

If the operator is unsure of what value to set for the upper and lower
deviation limits, they can use the following procedure to determine the
proper temperature offset. \emph{Note: This method requires the operator
to have some way of measuring both product and air temperature process
values, such as a recorder or a temperature meter.}

\begin{enumerate}
\def\labelenumi{\arabic{enumi}.}
\item
  Run the chamber with the selected test product to the desired product
  temp. set point with the product temperature control turned OFF. Run
  this ``test'' in constant mode.
\item
  When the air temperature has reached the desired set point, record the
  temperature difference between the product and the air temperature.
\item
  Use this temperature difference as the deviation limit setting. Enter
  the value into the ``+'' setting for product heat up. Enter the value
  into the ``-'' setting for product cool down.
\item
  Execute this product for heat up and cooling cycles in order to
  determine the upper (``+'') and lower (``-'') deviation limits.
\end{enumerate}

During product temperature control, the response characteristic of
control differs depending on the heat capacity, heat-transfer
coefficient, and thermal conductivity of the control target specimen.
Since the required control parameters are also changed, set the
appropriate values.

The following procedure outlines the steps to adjust Product Temperature
Control for \textbf{Constant 1}. Chapter 5 will outline the steps for
adjusting Product Temperature Control Program setup.

\textbf{\emph{Procedure:}}

\begin{enumerate}
\def\labelenumi{\arabic{enumi}.}
\item
  Confirm that the floating keyboard has been enabled. Refer to Section
  2.4.4.
\item
  From within the main home page, press the \textbf{Set Mode} button and
  press \textbf{Constant 1} title bar to open its configuration page
  (see arrow).

  \begin{figure}
  \centering
  
\includegraphics[width=0.95\textwidth]{figures/382946587-S3_7-step2-002a.PNG}
  \caption{Constant 1 configuration page}
  \end{figure}\FloatBarrier
\item
  Press the \textbf{Product} button to turn on Product Temperature
  control (see arrow). With the \textbf{Product} option set to ON
  (button lit), the system performs product temperature control. With
  the \textbf{Product} option set to OFF (button unlit), the system does
  not perform product temperature control.

  \begin{figure}
  \centering
  
\includegraphics[width=0.95\textwidth]{figures/21606746-S3_7-step3-002a.PNG}
  \caption{Enabling product temperature control feature}
  \end{figure}\FloatBarrier
\item
  Scroll up the page and press the \textbf{Upper Dev.} value field to
  bring up the numeric input keypad.

  \begin{figure}
  \centering
  
\includegraphics[width=0.95\textwidth]{figures/33917969-S3_7-step4-002a.PNG}
  \caption{Accessing PTCON Upper/Lower Dev.}
  \end{figure}\FloatBarrier
\item
  Enter upper deviation value on the keypad, taking note of the high and
  low limits. \emph{Note: The system ignores any value higher or lower
  than the required range.} Press the check button (or anywhere outside
  of the keypad) to apply the new setting. New value takes effect
  immediately.

  \begin{figure}
  \centering
  
\includegraphics[width=0.95\textwidth]{figures/3799326284-S3_7-step5-002a.PNG}
  \caption{Enter new temperature setting}
  \end{figure}\FloatBarrier
\item
  Repeat the process for the \textbf{Lower Dev.} option. The X button on
  the keypad can be used to cancel the configuration.
\end{enumerate}


\section{Start/Stop Operation}

Before starting the operation, check to confirm that the settings of the
absolute high/low limits of temperature (and humidity) warning on the
Product Temp Protector device for protecting specimen, as specified in
the chamber User Manual, have been performed properly.

The start/stop operation is performed in the \textbf{Operation Mode}
page, accessible via the \textbf{Operation} menu or button (in the
\textbf{Home} page).

\begin{longtable}[]{@{}l@{}}
\toprule
\begin{minipage}[b]{0.97\columnwidth}\raggedright
\textbf{Reference}\strut
\end{minipage}\tabularnewline
\midrule
\endhead
\begin{minipage}[t]{0.97\columnwidth}\raggedright
When conducting a test of the electronic components make certain that
the specimen has not been affected by condensation. Condensation occurs
when the surface temperature of the specimen is lower than the dew-point
temperature of the air in the test area. To avoid condensation of the
specimen, it is necessary to perform temperature-only operation in
advance and start temperature/humidity control operation after the
surface temperature of the specimen reaches equilibrium with the
temperature in the test area. In addition, condensation can be avoided
by using humidification delay control. If the temperature and humidity
in the test area is 85 °C and 85\% RH, respectively, the dew-point
temperature of the air in the test area is 80.9 °C. Therefore,
condensation occurs if the surface temperature of the specimen is less
than 80.9 °C.\strut
\end{minipage}\tabularnewline
\bottomrule
\end{longtable}

The following table lists the dew-point temperature.

\begin{longtable}[]{@{}lll@{}}
\toprule
Dry-bulb Temp (C) & Relative Humi (\% RH) & Dew-point Temp
(C)\tabularnewline
\midrule
\endhead
60 & 85 & 56.5\tabularnewline
70 & 85, 90 & 66.3, 67.7\tabularnewline
85 & 85, 90 & 80.9, 82.3\tabularnewline
\bottomrule
\end{longtable}


\subsection{Starting Constant Value
Operation}

The following procedure outlines the steps to start a constant operation
mode for \textbf{Constant 1}. This procedure also applies to
\textbf{Constant 2} and \textbf{Constant 3}.

\textbf{\emph{Procedure:}}

\begin{enumerate}
\def\labelenumi{\arabic{enumi}.}
\item
  From within the main home page, press the \textbf{Operation Mode} menu
  (1) or button (2). Refer to Section 2.3.2 for details.

  \begin{figure}
  \centering
  
\includegraphics[width=0.95\textwidth]{figures/4158085756-S3_8_1-step1-main-home-page.PNG}
  \caption{Accessing Constant Operation Mode}
  \end{figure}\FloatBarrier
\item
  Check the set points (Temperature, Humidity) displayed in the panel of
  \textbf{Constant 1} (see arrow). \textbf{\emph{Note:}} To check other
  options, it may require accessing the \textbf{Constant 1} setup page.

  \begin{figure}
  \centering
  
\includegraphics[width=0.95\textwidth]{figures/1211506617-S3_8_1a-step1-000a.PNG}
  \caption{Checking Constant 1 set point values}
  \end{figure}\FloatBarrier
\item
  Press the \textbf{Constant 1} panel and press Yes in the Attention box
  to confirm the action to start the operation.
\end{enumerate}


\subsection{Stopping Constant Value
Operation}

The following procedure outlines the steps to stop a constant operation
mode for \textbf{Constant 1}. This procedure also applies to
\textbf{Constant 2} and \textbf{Constant 3}.

\textbf{\emph{Procedure:}}

\begin{enumerate}
\def\labelenumi{\arabic{enumi}.}
\item
  From within the main home page, press the \textbf{Operation Mode} menu
  (1) or button (2). Refer to Section 2.3.2 for details.

  \begin{figure}
  \centering
  
\includegraphics[width=0.95\textwidth]{figures/4273468635-S3_8_2-step1.PNG}
  \caption{Accessing the Operation Mode page}
  \end{figure}\FloatBarrier
\item
  Press the \textbf{Stop} button (1) or \textbf{Constant 1} panel (2),
  then press Yes in the Attention box (3) to confirm the action.

  \begin{figure}
  \centering
  
\includegraphics[width=0.95\textwidth]{figures/2158917054-S3_8_2a-step2-stop-0002a.PNG}
  \caption{Stopping Constant 1 Operation Mode}
  \end{figure}\FloatBarrier
\item
  \textbf{Constant 1} mode is terminated and the system returns to a
  standby mode.\\
\item
  To save energy, turn off the circuit breaker when the equipment will
  not be used for a long time. If the circuit breaker remains in the ON
  position, the heater that warms up the refrigerator also remains
  energized.
\end{enumerate}

\begin{longtable}[]{@{}l@{}}
\toprule
\begin{minipage}[b]{0.97\columnwidth}\raggedright
\textbf{Note}\strut
\end{minipage}\tabularnewline
\midrule
\endhead
\begin{minipage}[t]{0.97\columnwidth}\raggedright
If the chamber operation is stopped in low temperature operation,
depending on the ambient condition, condensation may occur on the
equipment surface. In some cases, water puddles may occur in the chamber
installation location. To avoid condensation effect, return the
temperature in the test area to ambient temperature before stopping
operation.\strut
\end{minipage}\tabularnewline
\bottomrule
\end{longtable}


\section{Monitoring Constant Value
Operation}

\textbf{Home} and \textbf{Monitor} pages display different level of
information about the current operating mode or condition of the
chamber. This section outlines and explains in detail how to read
information on each page.


\subsection{Home Page Standard Display:
Temperature/Humidity}

Chamber operating mode is displayed in the middle of the status bar at
all time. As depicted in the following figure, the chamber is in
Constant mode, performing \textbf{Constant 2} settings. Temperature and
humidity process values are also displayed.

Detailed information of \textbf{Constant 2} (i.e., current operating
mode) is displayed in the Standard Display area indicated by labels 1, 2
and 3 (in the figure). Their descriptions are listed in the following
table. Refrigeration and time signals can be controlled during operation
under the setup page of \textbf{Constant 2}.

\begin{figure}
\centering

\includegraphics[width=0.95\textwidth]{figures/1758750591-S3_9A-homepage-001a.PNG}
\caption{Information of Temp/Humi for Constant 2 in Home page}
\end{figure}\FloatBarrier

\begin{longtable}[]{@{}ll@{}}
\toprule
\begin{minipage}[b]{0.07\columnwidth}\raggedright
Number\strut
\end{minipage} & \begin{minipage}[b]{0.88\columnwidth}\raggedright
Description\strut
\end{minipage}\tabularnewline
\midrule
\endhead
\begin{minipage}[t]{0.07\columnwidth}\raggedright
1\strut
\end{minipage} & \begin{minipage}[t]{0.88\columnwidth}\raggedright
Displays set point(s) and process value(s) of temperature (Air/Product)
and humidity.\strut
\end{minipage}\tabularnewline
\begin{minipage}[t]{0.07\columnwidth}\raggedright
2\strut
\end{minipage} & \begin{minipage}[t]{0.88\columnwidth}\raggedright
Displays temperature/humidity set points and refrigeration capacity. Set
points can be used to reference operation progress in Item 1.\strut
\end{minipage}\tabularnewline
\begin{minipage}[t]{0.07\columnwidth}\raggedright
3\strut
\end{minipage} & \begin{minipage}[t]{0.88\columnwidth}\raggedright
Displays time signals and their statuses. \textbf{\emph{Note:}} On/Off
manipulation can be controlled in the \textbf{Set Mode} page.\strut
\end{minipage}\tabularnewline
\bottomrule
\end{longtable}


\subsection{Monitor: Page 1}

The most detailed display page for any operating mode is the
\textbf{Monitor} page. Six different submonitor pages provide details of
the chamber operating condition.

Starting with Page 1, the display focuses on the main information of the
operating condition, namely, the set points and process values of
temperature (Air or Product) and humidity. The following two figures
illustrate how information is displayed for Constant and Standby modes.
Display information for Program mode is deferred to Chapter 5.

\textbf{Monitor Page 1}: Standby Mode

{[}Monitor Page 1: Information page of Standby mode{]}

\includegraphics[width=0.95\textwidth]{figures/2518384625-S4_9_2-page1-stdby-001a.PNG}

\textbf{Monitor Page 1}: Constant Mode

\begin{figure}
\centering

\includegraphics[width=0.95\textwidth]{figures/267390292-S4_9_2-page1-const-001a.PNG}
\caption{Monitor Page 1: Information page of Constant mode}
\end{figure}\FloatBarrier

\begin{longtable}[]{@{}ll@{}}
\toprule
\begin{minipage}[b]{0.07\columnwidth}\raggedright
Number\strut
\end{minipage} & \begin{minipage}[b]{0.88\columnwidth}\raggedright
Description\strut
\end{minipage}\tabularnewline
\midrule
\endhead
\begin{minipage}[t]{0.07\columnwidth}\raggedright
1\strut
\end{minipage} & \begin{minipage}[t]{0.88\columnwidth}\raggedright
Displays set points and process values of temperature (Air/Product). In
Standby mode, set points are OFF and process values are indicated by ---
for an unavailable option.\strut
\end{minipage}\tabularnewline
\begin{minipage}[t]{0.07\columnwidth}\raggedright
2\strut
\end{minipage} & \begin{minipage}[t]{0.88\columnwidth}\raggedright
Displays set points and process values of humidity. In Standby mode, set
points are OFF and process values are indicated by ---. Refer to Item
1.\strut
\end{minipage}\tabularnewline
\bottomrule
\end{longtable}


\subsection{Monitor: Page 2}

Page 2 of \textbf{Monitor} displays details of Time Signal outputs and
operating status of the refrigeration system. The following figures
depict Page 2 of \textbf{Monitor} for Standby and Constant modes,
respectively.

\textbf{Monitor, Page 2:} Standby Mode

\begin{figure}
\centering

\includegraphics[width=0.95\textwidth]{figures/3369992751-S3_9_3-standby-page2aaa.PNG}
\caption{Monitor Page 2, displaying information of Standby mode}
\end{figure}\FloatBarrier

\textbf{Monitor, Page 2:} Constant Mode

\begin{figure}
\centering

\includegraphics[width=0.95\textwidth]{figures/1705672758-S3_9_3-const-page2-001a.PNG}
\caption{Monitor Page 2, displaying information of Constant mode}
\end{figure}\FloatBarrier

\begin{longtable}[]{@{}ll@{}}
\toprule
\begin{minipage}[b]{0.07\columnwidth}\raggedright
Number\strut
\end{minipage} & \begin{minipage}[b]{0.88\columnwidth}\raggedright
Description\strut
\end{minipage}\tabularnewline
\midrule
\endhead
\begin{minipage}[t]{0.07\columnwidth}\raggedright
1\strut
\end{minipage} & \begin{minipage}[t]{0.88\columnwidth}\raggedright
Displays set points and process values of temperature (Air/Product) and
humidity. In Standby mode, set points are OFF and process values are
indicated by ---.\strut
\end{minipage}\tabularnewline
\begin{minipage}[t]{0.07\columnwidth}\raggedright
2\strut
\end{minipage} & \begin{minipage}[t]{0.88\columnwidth}\raggedright
Displays operating status of the available time signals.\strut
\end{minipage}\tabularnewline
\begin{minipage}[t]{0.07\columnwidth}\raggedright
3\strut
\end{minipage} & \begin{minipage}[t]{0.88\columnwidth}\raggedright
Displays refrigeration system status controlled by the PLC. A two-stage
refrigeration system is shown in the display. Since the PLC controls the
refrigeration system and its stages, users have no control over it. The
display mainly provides information for the user how the refrigeration
is being controlled and operated, in terms of power, discharge or
evaporation temp., gas pressure or temperature.\strut
\end{minipage}\tabularnewline
\begin{minipage}[t]{0.07\columnwidth}\raggedright
4\strut
\end{minipage} & \begin{minipage}[t]{0.88\columnwidth}\raggedright
Edit button: Click to edit the layout of this page.\strut
\end{minipage}\tabularnewline
\bottomrule
\end{longtable}


\subsection{Monitor: Page 3}

Page 3 of \textbf{Monitor} displays details of chamber operating mode:
Standby, Constant and Program. Details of the display in Program mode
will be discussed in Chapter 4. The following figures depict display
details for Standby and Constant.

\textbf{Monitor (page3):} Standby Mode

\begin{figure}
\centering

\includegraphics[width=0.95\textwidth]{figures/2240131210-S3_9_4-fig-stdby-000aaa.PNG}
\caption{Detailed page of operating mode}
\end{figure}\FloatBarrier

\begin{longtable}[]{@{}ll@{}}
\toprule
\begin{minipage}[b]{0.07\columnwidth}\raggedright
Number\strut
\end{minipage} & \begin{minipage}[b]{0.88\columnwidth}\raggedright
Description\strut
\end{minipage}\tabularnewline
\midrule
\endhead
\begin{minipage}[t]{0.07\columnwidth}\raggedright
1\strut
\end{minipage} & \begin{minipage}[t]{0.88\columnwidth}\raggedright
Displays current condition or operating mode\strut
\end{minipage}\tabularnewline
\begin{minipage}[t]{0.07\columnwidth}\raggedright
2\strut
\end{minipage} & \begin{minipage}[t]{0.88\columnwidth}\raggedright
Displays trend graph of recent data and set mode prior to current
(operating) mode: Standby. Details of this display will be provided
later.\strut
\end{minipage}\tabularnewline
\begin{minipage}[t]{0.07\columnwidth}\raggedright
3\strut
\end{minipage} & \begin{minipage}[t]{0.88\columnwidth}\raggedright
Edit button: Click to edit the layout of this page.\strut
\end{minipage}\tabularnewline
\bottomrule
\end{longtable}

\textbf{Constant Mode:} The constant mode page.

\begin{figure}
\centering

\includegraphics[width=0.95\textwidth]{figures/906783355-S3_9_4-fig-const-000aa.PNG}
\caption{Details page of Constant Mode}
\end{figure}\FloatBarrier

\begin{longtable}[]{@{}ll@{}}
\toprule
\begin{minipage}[b]{0.07\columnwidth}\raggedright
Number\strut
\end{minipage} & \begin{minipage}[b]{0.88\columnwidth}\raggedright
Description\strut
\end{minipage}\tabularnewline
\midrule
\endhead
\begin{minipage}[t]{0.07\columnwidth}\raggedright
1\strut
\end{minipage} & \begin{minipage}[t]{0.88\columnwidth}\raggedright
Displays detailed parameters of Constant mode.\strut
\end{minipage}\tabularnewline
\begin{minipage}[t]{0.07\columnwidth}\raggedright
2\strut
\end{minipage} & \begin{minipage}[t]{0.88\columnwidth}\raggedright
Displays trend graph of recent data and set mode prior to current
(operating) mode: Constant. Details of this display will be provided
later.\strut
\end{minipage}\tabularnewline
\begin{minipage}[t]{0.07\columnwidth}\raggedright
3\strut
\end{minipage} & \begin{minipage}[t]{0.88\columnwidth}\raggedright
Edit button: Click to edit the layout of this page.\strut
\end{minipage}\tabularnewline
\bottomrule
\end{longtable}


\subsection{Monitor: Detailed
Parameters}

Page 4 of \textbf{Monitor} displays a summary of the current operating
mode. It lists all the parameters associated with the current operating
mode, including control types and their status, such as temperature
(On/Off), humidity (On/Off), time signals (On/Off), refrigeration system
(On/Off), etc.

\begin{figure}
\centering

\includegraphics[width=0.95\textwidth]{figures/1714391146-S3_9_5-fig-detail-const-000.PNG}
\caption{Detailed summary of parameters of current Constant Mode}
\end{figure}\FloatBarrier


\subsection{Monitor: Preview and Program
Output}

Page 5 of \textbf{Monitor} is reserved for displaying details of program
output and its operating status. In Standby or Constant mode, this page
is blank. Chapter 5 will discuss this page in detail.


\subsection{Monitor: Trend}

Page 6 of \textbf{Monitor}, called \emph{\textbf{Trend Graph}}, displays
detailed scatter plot of data collected from the chamber. This trend
graph allows the operator to view a scatter plot of data collected
during chamber operation. Data collection can occur in two distinct
settings: \textbf{Always} and \textbf{Running}. Default setting is
\textbf{Always}, which means data will be collected during any operating
mode. Section 6.3 (``Set Sampling'') provides details of these settings.

\begin{figure}
\centering

\includegraphics[width=0.95\textwidth]{figures/1180964563-S3_9_7-trend-graph-const-000.PNG}
\caption{Trend graph of collected data}
\end{figure}\FloatBarrier

By default, the trend graph provides a scatter plot of data points
collected during the last hour. Various options of plot range of data
accumulated more than one hour can be selected. Data points collected
during operation are stored in the internal memory of the GL controller
system; as a data log, it remains stored in internal memory regardless
of a power outage or a system shutdown. Data can be downloaded (in whole
or in portion) and stored on an external USB device (if connected) or on
the local computer if the GL system was accessed remotely.

Trend graph is made up of different components, as depicted in the
following figure. Trend graph in program mode displays similar
components (discussed in Chapter 5).

\begin{figure}
\centering

\includegraphics[width=0.95\textwidth]{figures/1729650781-S3_9_7-trend-graph-const-001a.PNG}
\caption{Trend graph showing plots of current data from the chamber}
\end{figure}\FloatBarrier

The nine (9) components called out in the trend graph (above) are
described as follows:

\begin{enumerate}
\def\labelenumi{\arabic{enumi}.}
\item
  \textbf{Trend Graph}: Data collected from the chamber are rendered as
  trend graph via scatter plot methodology; data points of temperature,
  air temperature, product temperature control and humidity are plotted
  as a function of time. The y-axis represents the scale of these
  values. The x-axis represents the data logging runtime.
\item
  \textbf{Trend Graph Manipulation Buttons}: Control and management of
  trend graph are provided by the manipulation buttons. Detailed
  functionality and operation of these buttons are discussed in a
  separate section (Section 4.9.8).
\item
  \textbf{Line Graph}: Different styles of graph--solid lines, dashed
  lines and color--are used to designate each type of data points for
  visualization. Solid lines represent process values, dashed lines set
  values, etc.
\item
  \textbf{Y-axis Label}: The y-axis displays the selected type of data
  points to be plotted (temperature, product temperature control,
  humidity). These values are populated based on the selection of the
  \textbf{Graph View} under item 2 (above).
\item
  \textbf{X-axis Label}: The x-axis displays log time of the data.
  Default plot range is one hour with grids; each grid sets a 5-minute
  time-scale.
\item
  \textbf{Y-axis Title}: The y-axis title provides clarity to the type
  of graph being represented.
\item
  \textbf{Legend of Trend Graph}: The legends provide identity of each
  item on the trend graph designated with color coding.
\item
  \textbf{Status Bar}: Displays progress bar of each operating mode.
\item
  \textbf{Status Mode}: Displays chamber operating mode with date and
  time. If a program is being executed, program name, slot number and
  step number will be displayed.
\end{enumerate}


\subsection{Trend Graph Manipulation
Buttons}

The following is a detailed list of the \textbf{Manipulation Buttons},
item 2 in the previous figure (Section 4.9.7).

\begin{figure}
\centering

\includegraphics[width=0.95\textwidth]{figures/3772339018-trend-graph-manip-buttons-001abc.PNG}
\caption{Trend Graph manipulation buttons}
\end{figure}\FloatBarrier

\begin{enumerate}
\def\labelenumi{\arabic{enumi}.}
\item
  \textbf{Auto Refresh}: This Auto Refresh button reconstructs the graph
  to include data points up to the current time.
\item
  \textbf{Plot Range Selectors}: Default plot range is one hour. With
  this button (and its options), plot range can be configured up to 12
  hours or 7 days. Custom set range is also available via the Custom
  Time Span option.
\item
  \textbf{Download Button}: Data can be downloaded for backup on an
  external USB (or local) device and saved in CSV format. Date/time
  format can be downloaded in UTC or local time, as indicated by the
  drop-down menu options. Data log is accumulated and stored in the GL
  controller system internally; it remains stored in internal memory
  regardless of a power outage or a system shutdown. It is strongly
  recommended to download and clear data log regularly. Data
  accumulation increases in size almost exponentially. As a result,
  storage space may be wasted; large data file may affect the operation
  and performance of the GL control system. Refer to Section 6.3 for
  details how to clear data file.
\item
  \textbf{Plot Configuration}: Elements of the trend graph can be
  configured via this button.
\item
  \textbf{Pan/Zoom Controls}: The collapsible Pan/Zoom Controls button
  allows the operator to control and adjust the viewable section in the
  trend graph. The up/down arrows can be used to collapse/expand the
  drop-down menu to manipulate the trend graph as follows.

  \begin{itemize}
  \item
    \textbf{Zoom In}: The \textbf{Zoom In} button allows the operator to
    zoom into a small section of the trend graph. Depending on the
    degree of zooming, the display area will be confined to a small set
    of data points ranging between minutes to hours. To reset the trend
    graph, click the \textbf{Zoom Extents} button, select \textbf{Last
    Hour} from the drop-down menu, then click the \textbf{Auto Refresh}
    button.
  \item
    \textbf{Zoom Out}: The \textbf{Zoom Out} button does the opposite by
    allowing the operator to zoom out on the trend graph, thereby giving
    the operator an expansive view of the trend graph. To reset the
    trend graph, click the \textbf{Zoom Extents} button, select
    \textbf{Last Hour} from the drop-down menu, then click the
    \textbf{Auto Refresh} button.
  \item
    \textbf{Move Up}: This button allows the operator to move up the
    graph along the vertical axis to adjust the viewable area of the
    scatter plot. To reset the trend graph, click the \textbf{Zoom
    Extents} button, select \textbf{Last Hour} from the drop-down menu,
    then click the \textbf{Auto Refresh} button.
  \item
    \textbf{Move Down}: This button allows the operator to move down the
    trend graph along the vertical axis with the purpose to adjust the
    viewable area of the scatter plot. To reset the trend graph, click
    the \textbf{Zoom Extents} button, select \textbf{Last Hour} from the
    drop-down menu, then click the \textbf{Auto Refresh} button.
  \item
    \textbf{Move Left}: This button allows the operator to pan left on
    the trend graph, offering a quick preview of a plot of data points
    tracing back the time in hours or days. With this feature, the
    operator can quickly gain a preview of past data points which the
    operator may have missed.
  \item
    \textbf{Move Right}: This button does the opposite to \textbf{Move
    Left} by allowing the operator to pan right on the trend graph to
    the current time. To reconstruct the trend graph to contain the most
    recent data points, the \textbf{Auto Refresh} button allows the
    quickest operation.
  \end{itemize}
\end{enumerate}


\section{How to Handle Downloaded Data
File}

If data from Section 4.9.7 were downloaded via the \textbf{Export CSV}
and not via the \textbf{Export CSV (Local Time)} option, the date/time
of the data must be properly converted prior to any attempt to analyze
the data. The recorded time of the data points is based on a UTC (or
Universal Time Coordinated, previously referred to as the GMT) instead
of the local time. To preserve the integrity of the CSV file and its UTC
time stamp, do not rename the data file with MS Excel extension or
attempt to open it directly with MS Excel. Any attempt to force open it
in MS Excel will result in the wrong time conversion as shown in the
following figures.

The following figure displays the CSV file which was directly opened
with MS Excel. The time stamp in the first column is wrong; it is still
based on UTC time stamp indicated by the arrows.

\begin{figure}
\centering

\includegraphics[width=0.95\textwidth]{figures/3599308467-raw-data-downlaoded-001a.PNG}
\caption{Data file improperly open in MS Excel}
\end{figure}\FloatBarrier

The following figure displays the same CSV file which was properly
imported into MS Excel with the correct time stamp conversion for local
time indicated by the arrows.

\begin{figure}
\centering

\includegraphics[width=0.95\textwidth]{figures/3580293698-imported-data-downlaoded-001a.PNG}
\caption{Data file properly imported into MS Excel}
\end{figure}\FloatBarrier

The proper way to handle this CSV file is to import it using the
following procedure to render the correct conversion of the UTC time to
the local time on all data points.

Complete the following steps:

\begin{enumerate}
\def\labelenumi{\arabic{enumi}.}
\item
  Select \textbf{Trend} in the main menu bar
\item
  Click the \textbf{Download} button (item 6 in the previous section) to
  download the data (current view or entire data) in CSV file
\item
  Launch MS Excel and start a new blank spreadsheet
\item
  Click the \textbf{Data} ribbon
\item
  Click \textbf{From Text/CSV} (that appears under the \textbf{Home}
  ribbon)
\item
  In the \textbf{Import Data} Window explorer, open the
  \textbf{Downloads} folder and select the downloaded CSV file (in step
  2) and click the \textbf{Import} button (in the bottom).
\item
  Click the \textbf{Load} button to load data into the current
  spreadsheet.
\item
  To save the current spreadsheet as Excel Workbook, click \textbf{Save
  As} and follow the standard steps to complete the process.
\end{enumerate}


\chapter{Program}

This chapter discusses and explains the setting methods for performing
program operation that include creating, editing and managing programs,
starting and stopping program operation, monitoring program display
screens.

The features and functionality of program operation are available in two
modes in separate panels: Program Setup and Program Run. With the
Program Setup mode, the operator can create a program; open, view or
edit the program; preview the output of the program prior to executing
it; edit and/or overwrite an existing program; delete program from the
list; rename program on the list; move a program to a new location;
download a program (in JSON file) as a backup or archive; upload a
program from a USB (or a local computer) to the GL controller system,
and much more. With the Program Run mode, the operator can start or stop
a program. These two features are outlined in Section 5.2 and 5.3,
respectively.


\section{Program Operation}

A program operation changes the settings of the temperature and humidity
in the test area in accordance with the instructions and steps contained
in the executed program. Instructions in the program control temperature
and humidity operation based on the set or specified parameters. The
following diagram depicts a sequence of control operation specified by
the time, temperature, humidity and ramp for each step.

\begin{figure}
\centering

\includegraphics[width=0.95\textwidth]{figures/714097117-S4_1-fig-001a.PNG}
\caption{Step-wise program operation}
\end{figure}\FloatBarrier


\section{Program Setup Panel}

The Program Setup panel is accessible via the \textbf{Set Mode} menu.
The following figure depicts an empty Program Setup panel with no
programs on the list. The panel sets a default program slot, with
program number and name as \textbf{Program 1}, to start a new program.
The plus symbol indicates the slot is available for creating (i.e.,
adding) a new program.

\begin{figure}
\centering

\includegraphics[width=0.95\textwidth]{figures/1416653409-S4_2-prog-setup-panel-001.PNG}
\caption{Program Setup panel containing empty list}
\end{figure}\FloatBarrier

The following figure depicts a Program Setup panel that contains
programs on the list as well as empty ones (slot 3 and 5).

\begin{figure}
\centering

\includegraphics[width=0.95\textwidth]{figures/3235188799-S4_1_2-exiting-progs-on-list-001.PNG}
\caption{Program Setup panel containing}
\end{figure}\FloatBarrier

Slots 1, 2 and 4 are being occupied by programs with names and number of
steps displayed. Slots 3 and 5 are empty; hence, the plus symbols. The
GL controller system automatically opens a new slot for the next program
as soon as the current one is occupied by a program. However, if
programs were initially created in sequential order from slot 1 to 4
(with slot 5 automatically open for the next program), and then a
program in slot 3 is deleted, that slot remains open. A new program can
be created to occupy slot 3 again. The reason slot 3 still remains open
is because program \textbf{PAUL4} occupies slot 4, forcing slot 3 to
stay on the list, for it was pre-allocated before slot 4. This design
was carefully considered to prevent randomly reshuffling programs into
different slots that might create confusion after deleting a program out
of order. However, if \textbf{PAUL4} is deleted, then slots 4 and 5 will
collapse, forcing only slot 3 to remain open for the next program.

The number of programs that the GL controller system can hold is limited
to its available storage space. In the current model, that storage space
is 16GB. Data log and configuration files are also stored here, with
data log files most likely occupying the most space. For this reason, it
is highly recommended to back up data log and clear data log on the GL
controller system frequently. Chapter 6 will discuss how to set sampling
and clear data logs.


\section{Creating a New Program}

There are only two options associated with creating a new program: (1)
create a new program in the selected slot, or (2) import a program file
from an external USB device (or a local computer) into the selected
slot. Option (1) opens the program editor, while option (2) searches for
an external device connected to the USB port (or a local directory) to
open a file browser to import the program file.

The Program Setup panel has three tabs (or subpages) to provide
programming features and options for the selected program, as depicted
in the following figure. Their features are listed and described in the
table that follows.

\begin{figure}
\centering

\includegraphics[width=0.95\textwidth]{figures/640037717-S4_2_2-prgm-setup-page-002c.PNG}
\caption{Three tabs of the Program Setup panel}
\end{figure}\FloatBarrier

\begin{longtable}[]{@{}lll@{}}
\toprule
\begin{minipage}[b]{0.04\columnwidth}\raggedright
Number\strut
\end{minipage} & \begin{minipage}[b]{0.23\columnwidth}\raggedright
Name\strut
\end{minipage} & \begin{minipage}[b]{0.65\columnwidth}\raggedright
Description/Function\strut
\end{minipage}\tabularnewline
\midrule
\endhead
\begin{minipage}[t]{0.04\columnwidth}\raggedright
1\strut
\end{minipage} & \begin{minipage}[t]{0.23\columnwidth}\raggedright
Step-wise Programming\strut
\end{minipage} & \begin{minipage}[t]{0.65\columnwidth}\raggedright
Program instructions can be created step by step. Refer to Section 5.3.1
for details.\strut
\end{minipage}\tabularnewline
\begin{minipage}[t]{0.04\columnwidth}\raggedright
2\strut
\end{minipage} & \begin{minipage}[t]{0.23\columnwidth}\raggedright
Table-wise Programming\strut
\end{minipage} & \begin{minipage}[t]{0.65\columnwidth}\raggedright
Program instructions can be created by filling in the table elements.
Refer to Section 5.3.2 for details.\strut
\end{minipage}\tabularnewline
\begin{minipage}[t]{0.04\columnwidth}\raggedright
3\strut
\end{minipage} & \begin{minipage}[t]{0.23\columnwidth}\raggedright
Details\strut
\end{minipage} & \begin{minipage}[t]{0.65\columnwidth}\raggedright
Details of each program operation mode can be configured. Refer to
Section 5.3.3 for details.\strut
\end{minipage}\tabularnewline
\bottomrule
\end{longtable}

\textbf{Step-Wise Programming}: A program can be put together by
assembling instructions step by step. Parameters such as time,
temperature/humidity, product temp control option, type of control mode
(i.e., ramp, ramp rate or soak) can all be configured on the fly using
the pop-up step panel that opens next to the selected step (see arrow).
\textbf{\emph{Note:}} The step panel (depicted in the following figure)
is actually open and placed between the menu bar and Step 1; it is shown
for illustration only. The next figure illustrates the actual layout of
this pop-up panel.

\begin{figure}
\centering

\includegraphics[width=0.95\textwidth]{figures/2110777582-S5_3-fig2-001a.PNG}
\caption{Creating program using step-wise option}
\end{figure}\FloatBarrier

The following figure depicts the layout and components (and parameters)
of the selected step for editing.

\begin{figure}
\centering

\includegraphics[width=0.95\textwidth]{figures/3576817940-S5_3-step-wise-programming-001a.PNG}
\caption{Components of program step}
\end{figure}\FloatBarrier

\begin{longtable}[]{@{}lll@{}}
\toprule
\begin{minipage}[b]{0.04\columnwidth}\raggedright
Number\strut
\end{minipage} & \begin{minipage}[b]{0.23\columnwidth}\raggedright
Name\strut
\end{minipage} & \begin{minipage}[b]{0.65\columnwidth}\raggedright
Description/Function\strut
\end{minipage}\tabularnewline
\midrule
\endhead
\begin{minipage}[t]{0.04\columnwidth}\raggedright
1\strut
\end{minipage} & \begin{minipage}[t]{0.23\columnwidth}\raggedright
Step\strut
\end{minipage} & \begin{minipage}[t]{0.65\columnwidth}\raggedright
Displays the step number being edited.\strut
\end{minipage}\tabularnewline
\begin{minipage}[t]{0.04\columnwidth}\raggedright
2\strut
\end{minipage} & \begin{minipage}[t]{0.23\columnwidth}\raggedright
Set Temp/Humi\strut
\end{minipage} & \begin{minipage}[t]{0.65\columnwidth}\raggedright
Set point fields for temperature/humidity. \textbf{\emph{Note:}}
Humidity must be enabled first (via Item 3 below) in order to enter the
set point.\strut
\end{minipage}\tabularnewline
\begin{minipage}[t]{0.04\columnwidth}\raggedright
3\strut
\end{minipage} & \begin{minipage}[t]{0.23\columnwidth}\raggedright
Enable\strut
\end{minipage} & \begin{minipage}[t]{0.65\columnwidth}\raggedright
Displays product temperature control button (i.e., \textbf{Product}) and
humidity control (\textbf{Enable}) button. When the \textbf{Product}
button is activated, product temperature control is turned On. If the
chamber in question does not have product temperature control feature,
this button is not available. When the \textbf{Enable} button is
activated, humidity control is set to ON; otherwise, humidity is set to
OFF. See Section 5.3.1 for details.\strut
\end{minipage}\tabularnewline
\begin{minipage}[t]{0.04\columnwidth}\raggedright
4\strut
\end{minipage} & \begin{minipage}[t]{0.23\columnwidth}\raggedright
Ramp Control\strut
\end{minipage} & \begin{minipage}[t]{0.65\columnwidth}\raggedright
Displays buttons to enable temperature and humidity ramp control.\strut
\end{minipage}\tabularnewline
\begin{minipage}[t]{0.04\columnwidth}\raggedright
5\strut
\end{minipage} & \begin{minipage}[t]{0.23\columnwidth}\raggedright
Soak Control\strut
\end{minipage} & \begin{minipage}[t]{0.65\columnwidth}\raggedright
Displays buttons to enable temperature and humidity soak control.\strut
\end{minipage}\tabularnewline
\begin{minipage}[t]{0.04\columnwidth}\raggedright
6\strut
\end{minipage} & \begin{minipage}[t]{0.23\columnwidth}\raggedright
Duration\strut
\end{minipage} & \begin{minipage}[t]{0.65\columnwidth}\raggedright
Set the time duration of a step in hours, minutes and seconds for
temperature and humidity.\strut
\end{minipage}\tabularnewline
\begin{minipage}[t]{0.04\columnwidth}\raggedright
7\strut
\end{minipage} & \begin{minipage}[t]{0.23\columnwidth}\raggedright
Time Signals\strut
\end{minipage} & \begin{minipage}[t]{0.65\columnwidth}\raggedright
A list of Time Signals can be set to ON/OFF.\strut
\end{minipage}\tabularnewline
\begin{minipage}[t]{0.04\columnwidth}\raggedright
8\strut
\end{minipage} & \begin{minipage}[t]{0.23\columnwidth}\raggedright
Details\strut
\end{minipage} & \begin{minipage}[t]{0.65\columnwidth}\raggedright
Time Signals, refrigeration, and counter settings can be controlled
here, in addition to the \textbf{Pause} button.\strut
\end{minipage}\tabularnewline
\begin{minipage}[t]{0.04\columnwidth}\raggedright
9\strut
\end{minipage} & \begin{minipage}[t]{0.23\columnwidth}\raggedright
X\strut
\end{minipage} & \begin{minipage}[t]{0.65\columnwidth}\raggedright
Displays cancel button.\strut
\end{minipage}\tabularnewline
\begin{minipage}[t]{0.04\columnwidth}\raggedright
10\strut
\end{minipage} & \begin{minipage}[t]{0.23\columnwidth}\raggedright
Previous Button\strut
\end{minipage} & \begin{minipage}[t]{0.65\columnwidth}\raggedright
Displays button to navigate to the previous step.\strut
\end{minipage}\tabularnewline
\begin{minipage}[t]{0.04\columnwidth}\raggedright
11\strut
\end{minipage} & \begin{minipage}[t]{0.23\columnwidth}\raggedright
Accept Button\strut
\end{minipage} & \begin{minipage}[t]{0.65\columnwidth}\raggedright
Displays accept button to enter changes to the current step.\strut
\end{minipage}\tabularnewline
\begin{minipage}[t]{0.04\columnwidth}\raggedright
12\strut
\end{minipage} & \begin{minipage}[t]{0.23\columnwidth}\raggedright
Next Button\strut
\end{minipage} & \begin{minipage}[t]{0.65\columnwidth}\raggedright
Displays button to navigate to the next step.\strut
\end{minipage}\tabularnewline
\bottomrule
\end{longtable}

\textbf{Table-Wise Programming}: A program can be put together by
filling in the value for each parameter (such as time, temperature,
humidity) and selecting control mode such as product temp, ramp, soak,
refrigeration, and time signals can all be selected and configured in
one step via this method.

\begin{figure}
\centering

\includegraphics[width=0.95\textwidth]{figures/128661161-S4_2_2-create-prgm-table-wise-001a.PNG}
\caption{Creating program using table option}
\end{figure}\FloatBarrier

\textbf{Details}: Program operation modes and other parameters (such as
program start conditions, end mode, etc.) can be configured in the
\textbf{Details} tab. The upper and lower limits of the input range are
determined by the temperature (and/or humidity) range of the equipment
associated with the Abs.High and Abs.Low limits specified in this tab.
Deviation High/Low values, as well as Deviation Soak Control (Dev.Soak
Ctrl.) are configured here. Section 5.3.3 discusses these features in
detail.

\begin{figure}
\centering

\includegraphics[width=0.95\textwidth]{figures/1432725829-S4_2_2-create-prgm-detail-wise-001a.PNG}
\caption{Detail page of program creation}
\end{figure}\FloatBarrier


\subsection{Step-Wise
Programming}

This section outlines the steps to create a program using the step-wise
editor page.

\textbf{\emph{Note:}} The following procedure outlines the steps based
on the HMI (touch-screen) operation. For a remote access, the procedure
is identical starting with Step 2 by replacing the word \emph{press}
with \emph{click}.

\textbf{\emph{Procedure:}}

\begin{enumerate}
\def\labelenumi{\arabic{enumi}.}
\item
  If the GL controller system is already on, proceed to Step 2.
  Otherwise, set the main power breaker in the ON position. Within a
  minute, the HMI will come on.
\item
  Press \textbf{Set Mode} in the menu bar or the \textbf{Set Mode}
  button in the home page.
\item
  Press Program 1 slot (plus symbol), followed by the \textbf{Editor}
  button (pen icon) in the pop-up window. The X button can be used to
  cancel the action.

  \begin{figure}
  \centering
  
\includegraphics[width=0.95\textwidth]{figures/527372288-S5_3_1-step3-003a.PNG}
  \caption{Open program editor}
  \end{figure}\FloatBarrier
\item
  Press the body of Step 1 (see arrow). To cancel the current setting,
  press the X button. We outline from top to bottom on how to apply each
  parameter as follows.

  \begin{figure}
  \centering
  
\includegraphics[width=0.95\textwidth]{figures/103026662-S5_3_1-step4-002a.PNG}
  \caption{Edit step elemenets}
  \end{figure}\FloatBarrier

  \begin{itemize}
  \tightlist
  \item
    \textbf{Temperature:} Enter the temperature set point.
  \item
    \textbf{Humidity:} Enable humidity using the \textbf{Enable} button
    (below the Humi set point field); then enter the set point value in
    the Humi field. \textbf{\emph{Note:}} If humidity is not enabled,
    the numeric value in the Humi field does not take effect. A new
    numeric value can only be entered if the \textbf{Enable} button is
    activated. If \textbf{Ramp Ctrl} or \textbf{Soak Ctrl} (to be
    discussed below) is selected, humidity will be enabled.
  \item
    \textbf{Product Temperature Ctrl:} The chamber system used for this
    illustration (see above figure) does not have the Product
    Temperature Control feature. It is therefore empty and the enable
    button is blank. However, if the feature is available, to turn it
    on, simply select the \textbf{Product} button. When selected (i.e.,
    set to ON), the system performs product temperature control. Refer
    to the special note at the end of this section for details on the
    Product Temperature Control operation.\\
  \item
    \textbf{Ramp Ctrl:} Enable the ramp control for temperature and/or
    humidity, respectively using their appropriate buttons. When enabled
    (i.e., set to ON), ramp control is performed between the set point
    of the previous step and that of the current step. If ramp control
    is OFF, temperature (or humidity) is maintained based on the set
    point for this step. If ramp control is performed in Step 1 (such as
    in this example), the previous set point will be based on the start
    condition in the \textbf{Details} setup tab (see Section 5.3.3 for
    details).
  \item
    \textbf{Ramp Rate:} Ramp rate measured in degree per minute is
    determined internally by the system based on the difference between
    the two set points and duration (hence, °C/m). When a new time
    duration is entered, the value of ramp rate will change accordingly.
    Ramp rate value can be displayed down to three significant figures
    (e.g., 0.002). If these three digits rendered roughly to 0.000, the
    display will flag an error. This error will occur when the duration
    time is large and the difference between the two set points is
    small.
  \item
    \textbf{Soak Ctrl:} When selected for temperature, soak time control
    will maintain operation at the temperature set point. When soak time
    control is set to ON, after the operation of the step is started,
    the equipment waits until the process value reaches the control
    attainment range, and it then starts counting the time. Thus, the
    soak time in the temperature set point is the same as the setting
    time. If disabled, the system does not perform soak time control,
    and the counting starts at the same time as the start of the step
    operation. \textbf{\emph{Note:}} Soak control and ramp control
    cannot be performed together simultaneously. One or the other can be
    enabled, not both. If \textbf{Soak Ctrl} is selected, \textbf{Ramp
    Ctrl} which was previously selected will be disabled (deselected).\\
  \item
    \textbf{Duration:} Enter duration in the format of hours, minutes
    and seconds. The system can accommodate hour value up to four
    digits, e.g., 9999. The range of values for minute and second is
    between 1 and 59.
  \item
    \textbf{Time Signal:} Set the time signal contacts; signals 3 to 12
    are optional. If an option is not installed, the time signal key
    (button) cannot be selected. Descriptive name can be assigned to
    each time signal to identify each contact; refer to Section 7.7 for
    details. For each available time signal, set the desired time signal
    to ON or OFF; when ON, the system outputs the time signal;
    immediately after the program moves to the next step, the time
    signal changes the contact output setting; when OFF, it does not
    output the time signal.\\
  \item
    \textbf{Details:} Further control of the program step can be
    configured in this subpage. Press the \textbf{Details} button (see
    arrow in the figure above) to configure refrigeration capacity, time
    signals, loop counter, and pause function. Refer to the special note
    at the end of this section for details on \emph{refrigeration} and
    \emph{counter} controls. When \textbf{Pause} is set to ON, program
    execution is paused when this step is ended. The \textbf{Pause} icon
    also appears next to the time duration on the Step panel. To
    continue to the next step after the program is paused, select the
    \textbf{Resume/Play} button (see Section 5.6). When \textbf{Pause}
    is set to OFF, program continues to the next step.
  \item
    \textbf{Apply Settings:} Press the checkmark button (see arrow,
    above) to apply the settings and close the current Step panel.
    \textbf{\emph{Note:}} The previous (\textless{}) and next
    (\textgreater{}) buttons can be used to navigate between the
    previous and next steps without closing the Step panel.
  \end{itemize}
\item
  To create Step 2, repeat the procedure in Step 4. Additional steps can
  be created using the same process.\\
\item
  To quickly create an empty step, press the body of the next available
  step (plus symbol), followed by the checkmark button.
\item
  To edit an existing step (say, Step 2), press the body of the step
  (see arrow). Modify the contents of Step 2 in the step panel (see
  arrow). Press the \textbf{Details} button to adjust any configuration.
  Press the checkmark button to apply the new settings. Press the Save
  button (which becomes visible after closing the step panel with the
  checkmark) to save the program.

  \begin{figure}
  \centering
  
\includegraphics[width=0.95\textwidth]{figures/634398287-S5_3_1-step7-003a.PNG}
  \caption{Editing existing steps}
  \end{figure}\FloatBarrier
\item
  To traverse between existing steps within the program, with Step 2
  selected as the starting point (see arrow), press the left-arrow
  button (see arrow) or the right-arrow button (see arrow). Press the
  checkmark followed by the save button to save the new setting.

  \begin{figure}
  \centering
  
\includegraphics[width=0.95\textwidth]{figures/3726944312-S5_3_1-step8-001a.PNG}
  \caption{Traversing between existing steps}
  \end{figure}\FloatBarrier
\item
  A selected step may be moved backward or forward using the
  \textbf{Move Before} (\textless{}) or \textbf{Move After}
  (\textgreater{}) button. Press the Step 2 header (see arrow), then
  press the \textbf{Move Before} button to move Step 2 to before Step 1.
  As a result, the locations of Step 1 and Step 2 are swapped. To move
  Step 2 to the position after Step 3, press the Step 2 header (see
  arrow), then press the \textbf{Move After} button. As a result, the
  locations of Step 2 and Step 3 are swapped.

  \begin{figure}
  \centering
  
\includegraphics[width=0.95\textwidth]{figures/2445926883-S4_3_1-step9-001a.PNG}
  \caption{Relocating a selected step}
  \end{figure}\FloatBarrier
\item
  To insert a step before or after Step 2, press the Step 2 header (see
  arrow), then press the \textbf{Insert Step Before} or \textbf{Insert
  Step After} button to insert a new step in the desired location. The
  \textbf{Delete Step(s)} button can be used to delete the selected step
  (Step 2, in this case).

  \begin{figure}
  \centering
  
\includegraphics[width=0.95\textwidth]{figures/2445926883-S4_3_1-step9-001a.PNG}
  \caption{Inserting a step between existing steps}
  \end{figure}\FloatBarrier
\item
  To copy contents of Step 2 into a new step between Step 2 and Step 3,
  press Step 2 header (see arrow) followed by \textbf{Copy Step(s)},
  then press Step 2 header (again) to confirm the selection; press Step
  3 header and press Yes to confirm the action. To cancel the operation,
  press \textbf{No} followed by \textbf{Cancel} in the lower corner.

  \begin{figure}
  \centering
  
\includegraphics[width=0.95\textwidth]{figures/2445926883-S4_3_1-step9-001a.PNG}
  \caption{Copying step contents}
  \end{figure}\FloatBarrier
\item
  To copy a range of steps (say, Step 1 through 2) and put them in
  between Step 2 and Step 3, press Step 1 header (1) followed by
  \textbf{Copy Step(s)}, then press Step 2 header (2) to confirm the
  range, press Step 3 header (3) and press Yes (4) to confirm the
  action. To cancel the operation, press \textbf{No} followed by
  \textbf{Cancel} in the lower corner. (Actual programming examples will
  be provided later.)

  \begin{figure}
  \centering
  
\includegraphics[width=0.95\textwidth]{figures/131982083-S4_3_1-step11b.PNG}
  \caption{Copying a range of steps}
  \end{figure}\FloatBarrier
\item
  To delete a range of steps (say, Step 2 through 3), press Step 2
  header (1) followed by \textbf{Delete Step(s)}, then press Step 3
  header (2) to confirm the range and press Yes (3) to confirm the
  action. To cancel the operation, press \textbf{No} followed by
  \textbf{Cancel} in the lower corner. \textbf{\emph{Note:}} Operation
  can be done in reserve order, namely, starting with Step 2 as the
  initial step followed by Step 1.

  \begin{figure}
  \centering
  
\includegraphics[width=0.95\textwidth]{figures/1018220975-S4_3_1-step13-000a.PNG}
  \caption{Deleting Step 2 and 3}
  \end{figure}\FloatBarrier
\item
  To save program (in the current slot), press the \textbf{Save} button
  in the secondary menu (see top arrow). To save program in the next
  slot, press the \textbf{Save} button in the secondary menu (see bottom
  arrow).

  \begin{figure}
  \centering
  
\includegraphics[width=0.95\textwidth]{figures/28056107-S4_2_2-step11-000a.PNG}
  \caption{Saving the program}
  \end{figure}\FloatBarrier
\end{enumerate}

\begin{longtable}[]{@{}l@{}}
\toprule
\begin{minipage}[b]{0.97\columnwidth}\raggedright
\textbf{Reference}: Product Temperature Control\strut
\end{minipage}\tabularnewline
\midrule
\endhead
\begin{minipage}[t]{0.97\columnwidth}\raggedright
As stated in Step 4, when product temperature control (if available) is
On, the system performs product temperature control. The deviation
limits are the allowable air temperature deviation from product
temperature (P.Temp) set point under P.Temp control mode. The ``plus''
(+) value will be added to the P.Temp set point to create the ``upper''
(or High) air temperature limit. The ``minus'' (-) value will be
subtracted from the P.Temp set point to create the ``lower'' (i.e., Low)
air temperature limit. When the P.Temp process value is further from the
P.Temp set point, the controller will ``drive'' the air temperature set
point to this maximum upper (or lower) deviation limit. As the P.Temp
process value gets closer to the P.Temp set point, the air temperature
set point will change such that the effect is less than the maximum
upper (or lower) deviation limit. The High and Low deviation can be
adjusted via the \textbf{Details} tab by setting the Dev.High and
Dev.Low values in the High \& Low Limits panel in the appropriate
fields. Refer to Section 5.3.3.\strut
\end{minipage}\tabularnewline
\bottomrule
\end{longtable}

\begin{longtable}[]{@{}l@{}}
\toprule
\begin{minipage}[b]{0.97\columnwidth}\raggedright
\textbf{Reference}: Dev.High and Dev.Low Limits\strut
\end{minipage}\tabularnewline
\midrule
\endhead
\begin{minipage}[t]{0.97\columnwidth}\raggedright
If the operator is unsure of what value to set for the Dev.High and
Dev.Low, the following procedure can help to determine the proper
temperature offset. \textbf{\emph{Note:}} This method requires that the
operator has some way of measuring both product and air temperature
process values, such as a recorder or a temperature meter. (1) Run the
chamber with the selected test product to the desired P.Temp set point
with product temperature control turned OFF. Run this ``test'' in
constant mode. (2) When the air temperature has reached the desired set
point, record the temperature difference between the product and the air
temperature. This temperature difference will be used as the deviation
limit setting. (3) Enter the recorded value (from Step 2) into the (+)
setting for product heat up. Enter the value into the (-) setting for
product cool down. (4) Execute this procedure for heat up and cooling
cycles in order to determine the upper (+) and lower (-) deviation
limits. (5) During product temperature control, the response
characteristic of control differs depending on the heat capacity,
heat-transfer coefficient, and thermal conductivity of the control
target specimen. Since the required control parameters are also changed,
set the appropriate values accordingly.\strut
\end{minipage}\tabularnewline
\bottomrule
\end{longtable}

\begin{longtable}[]{@{}l@{}}
\toprule
\begin{minipage}[b]{0.97\columnwidth}\raggedright
\textbf{Reference}: Setting Refrigeration\strut
\end{minipage}\tabularnewline
\midrule
\endhead
\begin{minipage}[t]{0.97\columnwidth}\raggedright
Occasionally, automatic refrigeration control during linear ramping,
test steps may cause non-linear ramping results. If this was observed to
be the case, please contact ESPEC North America Inc.~customer service
for assistance; you will be provided with some guidelines on how to
program your test profile to use manual refrigeration control for better
ramping. Additionally, there are some specific circumstances where you
may be required to use manual refrigeration control to run a specific
test profile. An example of this kind would be to use manual
refrigeration to execute a ``pre-chill'' function with the refrigeration
system of your test chamber. This ``pre-chill'' function can be used to
maximize the temperature change rate of the chamber. This setting should
be used when transitioning from a set point above ambient to a low
temperature. The pre-chill setting is accomplished by setting the manual
refrigeration control to 25\% output 3-5 minutes before the transition
taking place. This will start the refrigeration system and pre-cool the
cascade of the system to ensure that maximum capacity is available for
the transition. When the temperature transition is started, the
refrigeration should be returned to ``auto'' control. This method can be
used in constant mode or program mode operation. \textbf{\emph{Note:}}
When refrigeration is set Auto, the manual setting is not reflected in
the operation of the refrigeration system. You cannot change the
refrigeration capacity during operation. Set the refrigeration capacity
when setting the program prior to operation.\strut
\end{minipage}\tabularnewline
\bottomrule
\end{longtable}

\begin{longtable}[]{@{}l@{}}
\toprule
\begin{minipage}[b]{0.97\columnwidth}\raggedright
\textbf{Reference}: Setting Counter Loop (Step-Wise Programming)\strut
\end{minipage}\tabularnewline
\midrule
\endhead
\begin{minipage}[t]{0.97\columnwidth}\raggedright
In a counter loop, the number of repetitions of a step or group of steps
can be configured. The process consists of: (1) Repeating start step
(step that is started from the second time around and continues there
after), (2) Repeating end step, (3) Repeat cycle When counter setup is
enabled, steps from the start step (1) through the end step (2) are
repeated until the repeat cycle count (3) is completed. To set up
counter, select the start step and end step (with step number 2 and
greater) to mark the loop. The loop count sets the condition to repeat
these steps until the number of loops is met. The conditions for
selecting start step and end step are: (A) 1 ≤ start step \textless{}
number of registered steps (B) Repeating start step \textless{}
repeating end step; end step \textgreater{} 1 (C) There is no limit set
for the number of repeat cycle. \textbf{\emph{Note:}} The total number
of operations resulting from the repeating start step (1) to the
repeating end step (2) is the number of repeat cycle (3) plus 1. This
means the first loop must be completed before the repeat cycle
commences. Therefore, if step 1 and step 3 are to repeat five times, the
input number of cycle is four (i.e., 4 + 1). Non-continuous steps can be
marked to be executed repeatedly. Multiple loops (or counters) can exist
within a program. In the Step-Wise Program editor, each counter loop is
marked with a bracket that covers the start step and end step. To mark
the start step, select the desired step and invoke the \textbf{Details}
subpage to configure the counter. Refer to Section 5.3.2 for details on
counter setup in Table-Wise Programming.\strut
\end{minipage}\tabularnewline
\bottomrule
\end{longtable}

The following are few examples and guidelines on how to set up counters
in a program. \textbf{\emph{Note:}} Counter A and Counter B are simply
counter labels. The GL controller supports multiple counters in a single
program. Thus Counter A, Counter B, Counter C, etc., can be used to
identify each individual counter.

\textbf{Example 1:} Two counters in a program, each contains a group of
steps, can be set up such that they perform the loop independently in
the order that they appear (from left to right in the diagram).

\begin{figure}
\centering

\includegraphics[width=0.95\textwidth]{figures/1994035793-S5_3_1-counter-ex1-001.PNG}
\caption{Two separate counters in a program}
\end{figure}\FloatBarrier

The following figure depicts the actual markings of these two separate
counters within the program.

\begin{figure}
\centering

\includegraphics[width=0.95\textwidth]{figures/3647257293-S5_3_1-counter-ex1-002.PNG}
\caption{Actual program containing two separate coutners}
\end{figure}\FloatBarrier

To study the behavior and output of the above program, perform the
following steps.

\textbf{\emph{Procedure:}}

\begin{enumerate}
\def\labelenumi{\arabic{enumi}.}
\tightlist
\item
  Create program as shown above (via the Setup menu).
\item
  Save the program
\item
  Execute the program (via the Operate Mode menu)
\item
  Access the \textbf{Monitor} menu
\item
  Select Subpage 3
\end{enumerate}

\textbf{Example 2:} Two counters nested in a program; the inner counter
(Counter B) is performed first inside Counter A.

\begin{figure}
\centering

\includegraphics[width=0.95\textwidth]{figures/2591358377-S5_3_1-counter-ex2-001.PNG}
\caption{Two separate counters nested within the program}
\end{figure}\FloatBarrier

The following figure depicts the actual markings of these two separate
counters nested inside the program.

\begin{figure}
\centering

\includegraphics[width=0.95\textwidth]{figures/1348622529-S5_3_1-counter-ex2-002.PNG}
\caption{Actual program containing two nested counters}
\end{figure}\FloatBarrier

\textbf{Example 3:} On paper, two counters nested inside one another on
the same start and end steps can be considered as ``combining'' the
number of cycles.

\begin{figure}
\centering

\includegraphics[width=0.95\textwidth]{figures/3037704782-S5_3_1-counter-ex3-001.PNG}
\caption{Two separate counters combine the total number of cycles}
\end{figure}\FloatBarrier

In actual programming, these two counters are nothing more than the
combined number of cycles, thus, using a single loop to achieve the goal
by repeating the number cycles 9 (6+3) times.

\textbf{Example 4:} Here is an example of an unusual counter setup which
should be avoided. Logically, it is an impossible setup in which the
start steps and end steps of the two counters intersecting one another.
With the GL controller, it is not an impossible but an abnormal setup.
To avoid unpredictable program behavior and outcome, no two (or more)
counters should be configured with their start steps and end steps
intersecting each other, as shown below.

\begin{figure}
\centering

\includegraphics[width=0.95\textwidth]{figures/4175884769-S5_3_1-counter-ex4-001.PNG}
\caption{An impossible counter nesting}
\end{figure}\FloatBarrier

The following figure depicts the actual markings of these two separate
counters intersecting their start steps and end steps which form a
strange outcome during program execution.

\begin{figure}
\centering

\includegraphics[width=0.95\textwidth]{figures/2270811692-S5_3_1-counter-ex4-001a.PNG}
\caption{S5\_3\_1-counter-ex4-001a.PNG}
\end{figure}\FloatBarrier

The program starts from Step 1 to Step 5; these steps are marked by
counter 1. When it reaches Step 5, it returns to Step 1 to repeat the
cycle marked by that counter (loop 1). Once it completes the cycle at
Step 5, it proceeds to Step 6 in the second counter and continues to
Step 9. Then it returns to Step 3 to repeat its cycle marked by counter
2. Here is where weird things start to behave: when it reaches Step 5,
it returns to Step 1 for counter 1 again until the number cycle is
fulfilled (in this case one time). Then it moves to Step 6 and continues
to Step 9 until its cycle is complete. It then continues to the next
step (Step 10). This weird behavior is depicted in the following figure,
with all the steps carried out during execution mapped at the bottom of
the screenshot, where both counters repeat the loop once (a total of two
cycles in each counter).

\begin{figure}
\centering

\includegraphics[width=0.95\textwidth]{figures/2286653660-S5_3_1-counter-ex4-002.PNG}
\caption{Trace of steps executed in disordered behavior}
\end{figure}\FloatBarrier

\textbf{Example 5}: Two loops overlapping each other by the end step of
the first and the start step of the second loop. Since the repeat cycle
marked by each counter is controlled by the end step, a strange behavior
occurs when loop 2 returns to its start step that also marks the end
step of loop 1.

\begin{figure}
\centering

\includegraphics[width=0.95\textwidth]{figures/3968327399-S5_3_1-counter-ex5-001.PNG}
\caption{Two counters with end step and start step locked in together}
\end{figure}\FloatBarrier

Program starts with Step 1; when it reaches Step 3, it returns to Step 1
to complete the first loop (counter 1). When the repeat cycle is
complete at Step 3, program continues to Step 4 through Step 6, which it
then returns to Step 3 to complete the cycle required in counter 2.
However, Step 3 marks the repeat cycle for counter 1, which the program
must loop back to Step 1. This means every time counter 2 needs to
complete its cycle by returning to Step 3, the program is forced to
execute counter 1 repeat cycle.

\textbf{Bottom Line:} Program structures in Example 4 and 5 should be
avoided.


\subsection{Table-Wise
Programming}

This section outlines the steps to create a program using the table-wise
editor page.

\textbf{\emph{Note:}} The following procedure outlines the steps based
on the HMI (touch-screen) operation. For a remote access, the procedure
is identical starting with Step 2 by replacing the word \emph{press}
with \emph{click}.

\textbf{\emph{Procedure:}}

\begin{enumerate}
\def\labelenumi{\arabic{enumi}.}
\item
  If the GL controller system is already on, proceed to Step 2.
  Otherwise, set the main power breaker in the ON position. Within a
  minute, the touch-screen display will come on.
\item
  Press \textbf{Set Mode} in the menu bar or the \textbf{Set Mode}
  button in the home page.
\item
  Press Program 1 slot (plus symbol, first arrow), followed by the
  Editor button (pen icon, second arrow) in the pop-up window.

  \begin{figure}
  \centering
  
\includegraphics[width=0.95\textwidth]{figures/3607149254-S5_3_2-step3-002a.PNG}
  \caption{Selecting Table-wise program editor}
  \end{figure}\FloatBarrier
\item
  Press the \textbf{Table-Wise} subpage (1) to turn on the table-wise
  editor. The program begins with Step 1 with some preselected
  parameters, such as the High/Low Dev., preconfigured in the
  \textbf{Details} page.

  \begin{figure}
  \centering
  
\includegraphics[width=0.95\textwidth]{figures/765730781-S4_3_2-step4-002b.PNG}
  \caption{Creating program using table option}
  \end{figure}\FloatBarrier
\item
  Complete Step 1 by entering parameter values for temperature, humidity
  and operation time (hour, minute and second) in their designated
  fields. Click the radio button under Temperature Mode (see arrow) to
  select control mode from the list. Repeat the same procedure for
  humidity mode (see arrow) and refrigeration capacity mode (see arrow).
  Each time signal (TS) can be turned On by placing a checkmark in the
  box; press on the box under the desired TS. \textbf{\emph{Note:}} Time
  signal control can be turned On/Off after the parameter values are
  entered. Therefore, Step 1 instructions involve a 2-step process, with
  time signal control being the last step to complete.

  \begin{figure}
  \centering
  
\includegraphics[width=0.95\textwidth]{figures/2813460863-S4_3_2-step5-002a.PNG}
  \caption{Complete instructions in Step 1}
  \end{figure}\FloatBarrier
\item
  Create Step 2 by pressing the \textbf{Append} button and filling in
  all the required fields with parameter values. Additional program
  steps can be repeated using the same process. A program consisting of
  five steps is depicted below.

  \begin{figure}
  \centering
  
\includegraphics[width=0.95\textwidth]{figures/3810732043-S4_3_2-step6-002.PNG}
  \caption{Five program steps have been created.}
  \end{figure}\FloatBarrier
\item
  A new step can be added (inserted) in between existing steps via the
  \textbf{step manipulation} menu, which can be invoked by pressing on
  an input field (e.g., Hrs field or Set Point field) in the selected
  step (until the menu appears). In the figure, the Hrs input field of
  Step 3 (see arrow) was momentarily depressed to bring up the
  \textbf{step manipulation} menu (red rectangle), which consists of
  \textbf{add above}, \textbf{delete (x)} and \textbf{add below}
  buttons. A new step can be inserted above or below an existing step
  using the (+) above or (+) below button. Press the (+) button at Step
  2 to insert a new step in between Step 1 and Step 2. Press the (+)
  button at Step 4 to insert a new step in between Step 4 and Step 5.

  \begin{figure}
  \centering
  
\includegraphics[width=0.95\textwidth]{figures/3446697236-S4_3_2-step7-002a.PNG}
  \caption{Inserting steps in between existing steps}
  \end{figure}\FloatBarrier
\item
  The \textbf{step manipulation} menu can be shifted along existing
  steps with the \textbf{delete (x)} button on the selected step (see
  arrows). To delete Step 2, press the input field of Step 2 (see
  arrow), then press the \textbf{delete (x)} button. To delete Step 4,
  press the input field of Step 4 (see arrow), then press the
  \textbf{delete (x)} button.

  \begin{figure}
  \centering
  
\includegraphics[width=0.95\textwidth]{figures/2061477884-S4_3_2-step8-003ab.PNG}
  \caption{Deleting steps using the X button}
  \end{figure}\FloatBarrier
\item
  To save program (in the current slot), press the \textbf{Save} button
  in the secondary menu (1). To save program in the next slot, press the
  \textbf{Save} button in the secondary menu (2).

  \begin{figure}
  \centering
  
\includegraphics[width=0.95\textwidth]{figures/1955823295-S4_3_2-step9-003a.PNG}
  \caption{Saving current program}
  \end{figure}\FloatBarrier
\end{enumerate}

\begin{longtable}[]{@{}l@{}}
\toprule
\begin{minipage}[b]{0.97\columnwidth}\raggedright
\textbf{Reference}: Setting Counter Loop (Table-Wise Programming)\strut
\end{minipage}\tabularnewline
\midrule
\endhead
\begin{minipage}[t]{0.97\columnwidth}\raggedright
In the Table-Wise Program editor, each counter is created by the count
number placed in the desired step marked as the end step for that
counter (loop). The counter setup process is easier to manage than that
for the Step-Wise Program Editor: select a desired step to mark the end
step for the counter; input the number of repeat cycles, then set the
start step. The principles of counter setup remain the same. Refer to
Section 5.3.6 for examples on counter setup.\strut
\end{minipage}\tabularnewline
\bottomrule
\end{longtable}


\subsection{Program Details}

The \textbf{Details} tab contains four separate panels to provide
further control and operation of the program.

\begin{figure}
\centering

\includegraphics[width=0.95\textwidth]{figures/3522713967-S4_3_3-details-page-001ab.PNG}
\caption{Configuring options in program features and modes}
\end{figure}\FloatBarrier

\textbf{Program Name:} Enter or edit program name. The above figure
depicts an example how to edit or enter program name and save the
changes. \textbf{\emph{Note:}} If program name is not specified,
\textbf{Program \#?} will be used, where ? will be the numeral in
sequence of 1, 2, 3, etc.

\textbf{Program Start Conditions:} When you want to perform ramp
operation for temperature and/or humidity in Step 1, the Temp and Humi
must be enabled. Turn on \textbf{Process Value} when you want to start
process control from the current temperature and humidity in the test
area. Turn on \textbf{Set Point} and enter the desired value(s) when you
want to set arbitrary temperature and humidity for ramp control. The
following figure depicts Temp and Humi start conditions using set point
values.

\begin{figure}
\centering

\includegraphics[width=0.95\textwidth]{figures/1759555740-temp-humi-start-conditions-001.PNG}
\caption{Set start conditions for Temp and/or Humi}
\end{figure}\FloatBarrier

When the \textbf{Set Point} option is selected, the start temperature
and humidity can be changed. The ramp operation is to control the
temperatures from set point of the previous step to that of the present
step at a certain ramp. The following figure depicts graphically how
ramp control operation is performed from Step 1.

\begin{figure}
\centering

\includegraphics[width=0.95\textwidth]{figures/1123077062-S5_3_3-start-cond-001.PNG}
\caption{Graphs of ramp operation and conditions from Step 1}
\end{figure}\FloatBarrier

Ramp control operation (temp or humi) will not be performed based on the
following two conditions:

\begin{enumerate}
\def\labelenumi{\arabic{enumi}.}
\tightlist
\item
  When the start condition (temp or humi) is set to OFF in Step 1.
\item
  When the humidity control in the previous step is set to OFF in ramp
  control (humidity).
\end{enumerate}

\textbf{Program End Mode:} End mode condition can be set to control
(i.e., instruct) the equipment what to do next after program operation
is ended. Conditions include: Stop the chamber; hold the last step of
the current program; start Constant 1, 2, or 3; start next program (with
step number).

By default, \textbf{Stop} is selected for end mode. In this mode, the
chamber will stop the operation after the program is ended. A different
end mode can be selected simply by activating the appropriate button,
with the exception of program mode that involves several steps
(explained next).

To start a new program after the current program is ended, proceed as
follows.

\textbf{\emph{Procedure:}}

\begin{enumerate}
\def\labelenumi{\arabic{enumi}.}
\tightlist
\item
  Select the \textbf{Program} button.
\item
  Select the ``program name and step number'' button.
\item
  From the \textbf{Select a Program} pop-up window, select the desired
  program and press \textbf{Accept}
\item
  From the \textbf{Select Program Step} pop-up window, select the step
  number to start and press \textbf{Accept}. By default, Step 1 is
  selected.\\
\item
  Press \textbf{Save} to update the changes.
\end{enumerate}

\textbf{High \& Low Limits:} The values set with Abs.High and Abs.Low
are the range of values in which data can be entered during programming
in the Program Setup page. These values form the range in which
temperature and humidity set point can be entered in Program Setup
(similar to Constant Setup, see Section 4.4). These two values set the
extreme high and extreme low set points in a chamber operation. Should a
temperature in test area rise above the Abs High set point, alarm will
be tripped. The same is true if it goes below the Abs Low set point. The
following diagram depicts the behavior of these two situations.

\begin{figure}
\centering

\includegraphics[width=0.95\textwidth]{figures/1025545634-S3_4-HL-graph.PNG}
\caption{High/Low limit setup condition}
\end{figure}\FloatBarrier

Two different scenarios are described as follows.

\textbf{Abs High and Abs Low}

Criteria used for performing protection in the test area. Set values
greater than the temperature set points in the test area by 10 °C or
more. If the equipment detects a fault condition, the equipment stops
operation after issuing an alarm.

\textbf{Upper Deviation}

Criteria used for performing for specimens. Change the setting depending
on the test specimen. If the equipment detects a deviation fault, the
equipment stops the header (humidifier) after issuing a warning. When
the temperature in the test area decreases to the temperature set point,
the equipment returns to the normal control.


\subsection{Example 1: Creating a Program using Step-Wise
Method}

The following procedure outlines the steps to create a program using the
Step-Wise method. The program will consist of 15 steps. Operation time,
temperature and humidity set point values are listed in the following
table.

\begin{longtable}[]{@{}llll@{}}
\toprule
Step No. & Temp & Humi & Refrig\tabularnewline
\midrule
\endhead
1 & 85 & 10 & Auto\tabularnewline
2 & 60 & 10 & Auto\tabularnewline
3 & 40 & 20 & Auto\tabularnewline
4 & 20 & 45 & Auto\tabularnewline
5 & 10 & 80 & Auto\tabularnewline
6 & 10 & 90 & Auto\tabularnewline
7 & 15 & 98 & Auto\tabularnewline
8 & 30 & 50 & Auto\tabularnewline
9 & 30 & 65 & Auto\tabularnewline
10 & 40 & 98 & Auto\tabularnewline
11 & 70 & 50 & Auto\tabularnewline
12 & 85 & 85 & Auto\tabularnewline
13 & 85 & 98 & Auto\tabularnewline
14 & 83 & 10 & Auto\tabularnewline
15 & 23 & OFF & Auto\tabularnewline
\bottomrule
\end{longtable}

The procedure assumes the chamber has already been turned On, and the GL
controller system starts from its home page. If the chamber is not yet
turned On, set the main power breaker in the ON position. Within a
minute, the touch-screen display will come on. The system automatically
logs in using the default user account (localhost) to operate the
chamber.

\begin{figure}
\centering

\includegraphics[width=0.95\textwidth]{figures/2850490313-S6_2_3-step0.PNG}
\caption{Operating Panel at Startup}
\end{figure}\FloatBarrier

\textbf{\emph{Procedure:}}

\begin{enumerate}
\def\labelenumi{\arabic{enumi}.}
\item
  Press the \textbf{Hamburger} menu (1) in the status bar. Check to
  confirm the floating keyboard is enabled. A diagonal bar over the
  keyboard indicates it is disabled. Press the keyboard (2) to enable
  it. Press \textbf{Set Mode} (3) to access the \textbf{Set Mode} page.

  \begin{figure}
  \centering
  
\includegraphics[width=0.95\textwidth]{figures/4051585419-S6_2_3-step1.PNG}
  \caption{Enable the floating keyboard in the system menu}
  \end{figure}\FloatBarrier
\item
  Press \textbf{Program 1} panel (see top arrow), then press the Edit
  button (see bottom arrow) in the pop-up window to access Program 1
  edit page for programming.

  \begin{figure}
  \centering
  
\includegraphics[width=0.95\textwidth]{figures/3607149254-S5_3_2-step3-002a.PNG}
  \caption{Select Program 1 for profile creation}
  \end{figure}\FloatBarrier
\item
  The step-wise editing page (intended for this example) is displayed by
  default with Step 1, unedited with 0 operation time and empty trend
  graph at the bottom. Press the Step 1 body (see arrow) to access its
  editing template.

  \begin{figure}
  \centering
  
\includegraphics[width=0.95\textwidth]{figures/3115392718-S6_2_3-step4.PNG}
  \caption{Accessing Step 1 of program}
  \end{figure}\FloatBarrier
\item
  Press the \textbf{Temp} box (see arrow) and enter 85. Press
  \textbf{Enable} (see arrow) under Humi to turn it on and adjust set
  value (from 5) to 10 RH. Set operation time to 1 hour: \textbf{Hrs} to
  1, \textbf{Mins} and \textbf{Secs} to 0. Step 1 summary page is shown
  in the right for reference. The first nine time signals are set to Off
  (unselected).

  \begin{figure}
  \centering
  
\includegraphics[width=0.95\textwidth]{figures/2365826148-S4_3_4-step4-000c.PNG}
  \caption{Creating Step 1; setting Temp value and duration}
  \end{figure}\FloatBarrier
\item
  By default, refrigeration is set to Auto and time signals are set to
  Off. Press the \textbf{Details} button at the bottom of Step 1 summary
  page (refer to the above figure and/or see arrow in the following
  figure) to confirm refrigeration and time signals settings. Press the
  checkmark (see arrow) to complete Step 1 programming.
  \textbf{\emph{Note}}: The right/left buttons are for navigating
  through program steps forward or backward. Since only one step is
  being created thus far, these buttons are not yet fully operational.

  \begin{figure}
  \centering
  
\includegraphics[width=0.95\textwidth]{figures/680538036-S4_3_4-step5-003ab.PNG}
  \caption{Confirm refrigeration and time signals configuration}
  \end{figure}\FloatBarrier
\item
  Press the \textbf{Details Page} (see arrow) to access program details
  page for configuring program start/stop mode, Absolute High/Low limits
  for temperature and humidity. \textbf{\emph{Note}}: The slider bar and
  left/right paging buttons are for traversing through program steps.

  \begin{figure}
  \centering
  
\includegraphics[width=0.95\textwidth]{figures/1996146587-S4_3_4-step5-002a.PNG}
  \caption{Accessing detail page of Step 1}
  \end{figure}\FloatBarrier
\item
  The default program name is \textbf{Program 1}. Press the program name
  field (see arrow) and enter \textbf{Humidity Points} for program name.
  By default, program start mode consists of both temperature and
  humidity without initial set point values, which means the chamber
  will raise ambient temperature (inside the chamber) to the set point
  specified in Step 1. Program will stop when it completes all the
  steps. Adjust the Absolute High and Low Limits (see arrows) for
  temperature values as desired. \textbf{\emph{Note}}: The High/Low
  limits only affect the current program being edited; they will not
  affect any other programs or constant modes. Press the \textbf{Save}
  button (see arrow) to save the settings.

  \begin{figure}
  \centering
  
\includegraphics[width=0.95\textwidth]{figures/87363343-S6_2_3-step8a.PNG}
  \caption{Set Program name and temp/humi high/low limits}
  \end{figure}\FloatBarrier

  \textbf{\emph{Note:}} The high/low limit values cannot be set to a
  value greater than the default high/low limit values set at the
  factory. Refer to the following figure for detail.

  \begin{figure}
  \centering
  
\includegraphics[width=0.95\textwidth]{figures/1698428879-S6_2_3-figure1.PNG}
  \caption{Absolute high/low Limit and UPPER Deviation Limit of
  Temperature}
  \end{figure}\FloatBarrier

  \textbf{Absolute high/low temperature limits}: These settings are
  designed to protect the chamber against thermal damage. Set the values
  to at least 10 °C higher/lower than the target temperature. When
  tripped, an alarm is generated and the chamber stops running.

  \textbf{Upper deviation temperate limit}: This setting is designed to
  protect specimens against damage by heat. It also triggers a safety
  device inside the chamber, which causes the heater output to set at
  0\% output. When temperature returns within range, normal control is
  restored automatically.

  \textbf{Absolute high/low temperature limits}: These settings are
  designed to protect the specimens against humidity damage. When
  tripped, it triggers a safety device inside the chamber. Absolute high
  will cause the humidity heater to go to 0\% output. When humidity
  returns within range, normal control is restored automatically.
\item
  The following table outlines the detailed elements for the remaining
  14 steps. Repeat Steps 3 through 5 to complete these steps using
  values in the table.
\item
  The completed program is illustrated in the following figure. The
  slider bar and paging buttons can now to be used to traverse between
  program steps. The trend graph illustrates the output of the program,
  showing what the program will do when it is executed. This trend graph
  is extremely useful for previewing the program output to study its
  results. Press the \textbf{Save} button (top arrow) to save the
  program in the current slot (i.e., Program 1). \textbf{\emph{Note}}:
  The \textbf{Save As} button (middle arrow) can be used to save the
  program in a different slot. The reverse button (last arrow) can be
  used to cancel or revert the action.

  \begin{figure}
  \centering
  
\includegraphics[width=0.95\textwidth]{figures/2752159224-S4_3_4-step9-000a.PNG}
  \caption{Program 1 with 15 steps}
  \end{figure}\FloatBarrier
\item
  Press Yes to confirm the action (saving). Program 1 is now saved in
  slot 1.

  \begin{figure}
  \centering
  
\includegraphics[width=0.95\textwidth]{figures/2410566821-S4_3_4-step10-000.PNG}
  \caption{Saving Program 1}
  \end{figure}\FloatBarrier
\end{enumerate}


\subsection{Example 2: Creating a Program with Table-Wise
Method}

The following procedure outlines the steps to create a program using the
Table-Wise method. The program will consist of 8 steps, with operation
time, temperature and humidity set point values listed in the following
table. Program 2 (in slot 2) will be used to create the program.

\begin{longtable}[]{@{}llll@{}}
\toprule
Step No. & Temp & Humi & Refrig\tabularnewline
\midrule
\endhead
1 & 85 & 10 & Auto\tabularnewline
2 & 60 & 10 & Auto\tabularnewline
3 & 40 & 20 & Auto\tabularnewline
4 & 20 & 45 & Auto\tabularnewline
5 & 10 & 80 & Auto\tabularnewline
6 & 10 & 90 & Auto\tabularnewline
7 & 15 & 98 & Auto\tabularnewline
8 & 23 & OFF & Auto\tabularnewline
\bottomrule
\end{longtable}

The procedure assumes the chamber has already been turned On, and the GL
controller system starts from its home page. If the chamber is not yet
turned On, set the main power breaker in the ON position. Within a
minute, the touch-screen display will come on. The system automatically
logs in using the default user account (localhost) to operate the
chamber.

\begin{figure}
\centering

\includegraphics[width=0.95\textwidth]{figures/2850490313-S6_2_3-step0.PNG}
\caption{Operating Panel at Startup}
\end{figure}\FloatBarrier

\textbf{\emph{Procedure:}}

\begin{enumerate}
\def\labelenumi{\arabic{enumi}.}
\item
  Press the \textbf{Hamburger} menu (1) in the status bar. Check to
  confirm the floating keyboard is enabled. A diagonal bar over the
  keyboard indicates it is disabled. Press the keyboard (2) to enable
  it. Press \textbf{Set Mode} (3) to access the \textbf{Set Mode} page.

  \begin{figure}
  \centering
  
\includegraphics[width=0.95\textwidth]{figures/4051585419-S6_2_3-step1.PNG}
  \caption{Enable the floating keyboard in the system menu}
  \end{figure}\FloatBarrier
\item
  Press \textbf{Program 2} panel (see top arrow), then press the Edit
  button (see bottom arrow) in the pop-up window to access Program 2
  edit page for programming.

  \begin{figure}
  \centering
  
\includegraphics[width=0.95\textwidth]{figures/103476112-S5_3_5-step2-002a.PNG}
  \caption{Select Program 2 for profile creation}
  \end{figure}\FloatBarrier
\item
  Press the Table-wise editing page (see arrow) to display the program
  edit table.

  \begin{figure}
  \centering
  
\includegraphics[width=0.95\textwidth]{figures/2698228481-S4_3_5-step3-000a.PNG}
  \caption{Selecting Table-Wise programming method}
  \end{figure}\FloatBarrier
\item
  Edit Step 1 by entering the values as indicated by the four arrows. By
  default, humidity mode is set to Off. Press the down arrow icon (see
  arrow) to select On from the menu and enter 10 for humidity set point.
  The rest of the parameters are set by default. Press the
  \textbf{Append} button (see arrow) to add Step 2 and proceed to the
  next step.

  \begin{figure}
  \centering
  
\includegraphics[width=0.95\textwidth]{figures/482900626-S4_3_5-step4-001a.PNG}
  \caption{Edit Step 1 of Program 2}
  \end{figure}\FloatBarrier
\item
  Complete Step 2 of the program by editing the temperature set point to
  60 (refer to the table above); the rest of the values and parameters
  remain the same (as in Step 1).

  \begin{figure}
  \centering
  
\includegraphics[width=0.95\textwidth]{figures/2936436928-S4_3_5-step5-000.PNG}
  \caption{Complete Step 2 of the program}
  \end{figure}\FloatBarrier
\item
  Use the values in the above table to complete the remaining 6 steps in
  the program, as depicted below. The default settings for program
  start/stop modes are identical to those in Section 5.3.4; and they
  will be used here also and therefore no need to access the
  \textbf{Details} page (with the exception if program name needs
  editing). The vertical slider bar and paging buttons (on the right)
  can be used to traverse through program steps for editing. The
  horizontal slider bar has two small arrows. The right arrow slides the
  table leftwards to display the remaining time signals; the left arrow
  slides back the table. Press the \textbf{Save} button to save the
  program.

  \begin{figure}
  \centering
  
\includegraphics[width=0.95\textwidth]{figures/2812727144-S4_3_5-step6-000a.PNG}
  \caption{Program 2 completion}
  \end{figure}\FloatBarrier
\item
  Program 2 is now saved.
\end{enumerate}


\subsection{Example 3: Setting Multiple Counters in a
Program}

In this example, we illustrate how to set up multiple counters in a
program, using Step-Wise and Table-Wise method. The main focus is how to
set up counters, giving less attention to temperature (humidity) or
other options; we specify the duration for each step with 2 minutes. The
following table lists two counters: (1) Counter 1 (Step 1 through Step
3) repeats the cycle one time; (2) Counter 2, with start step at Step 5
and end step at Step 9, repeats the cycle 3 times.

\begin{longtable}[]{@{}llllllll@{}}
\toprule
Step No. & Time & Temp & Humi & Refrig & Start Step & End Step &
Counter\tabularnewline
\midrule
\endhead
1 & 2 min & 21 & Off & Auto & yes & &\tabularnewline
2 & 2 min & 30 & Off & Auto & & &\tabularnewline
2 & 2 min & 30 & Off & Auto & & &\tabularnewline
3 & 2 min & 35 & Off & Auto & & yes & 1\tabularnewline
4 & 2 min & 30 & Off & Auto & & &\tabularnewline
5 & 2 min & 35 & Off & Auto & yes & &\tabularnewline
6 & 2 min & 21 & Off & Auto & & &\tabularnewline
7 & 2 min & 30 & Off & Auto & & &\tabularnewline
8 & 2 min & 35 & Off & Auto & & &\tabularnewline
9 & 2 min & 30 & Off & Auto & & yes & 3\tabularnewline
10 & 2 min & 21 & Off & Auto & & &\tabularnewline
\bottomrule
\end{longtable}

Refer to Section 5.3.4 on how to create a program that consists of 10
steps with the specified temperature and refrigeration, with humidity
feature disabled. The table above (without counters) produces the
following:

\begin{figure}
\centering

\includegraphics[width=0.95\textwidth]{figures/2317935606-S5_3_6-image-001.PNG}
\caption{Program without counters}
\end{figure}\FloatBarrier

\textbf{\emph{Procedure:}} Step-Wise Method

\begin{enumerate}
\def\labelenumi{\arabic{enumi}.}
\item
  Press/click the body of Step 3 (see arrow).

  \begin{figure}
  \centering
  
\includegraphics[width=0.95\textwidth]{figures/1447191847-S5_3_6-step2-table-wise-001a.PNG}
  \caption{Setting counter with Step 3 as the end step}
  \end{figure}\FloatBarrier
\item
  Press/click the Details button (see arrow in figure below).
\item
  Enter 1 in the Count field (see arrow). \textbf{\emph{Note:}} By
  default Step 1 is selected for Start Step.

  \begin{figure}
  \centering
  
\includegraphics[width=0.95\textwidth]{figures/3059144716-S5_3_6-step3-step-wise-prog-001a.PNG}
  \caption{Setting Counter 1 with repeat cycle at Step 3}
  \end{figure}\FloatBarrier
\item
  Press/click the checkmark button in the step panel (below the Details
  button).
\item
  Press/click the body of Step 9, followed by the Details button.
\item
  Enter 3 in the Count field (see arrow).

  \begin{figure}
  \centering
  
\includegraphics[width=0.95\textwidth]{figures/2869678099-S5_3_6-step6-step-wise-prog-001a.PNG}
  \caption{Setting Counter 2 with repeat cycle at Step 9}
  \end{figure}\FloatBarrier
\item
  Enter 5 in the Start Step field (see arrow in above figure).
\item
  Press/click the checkmark button.
\item
  Press/click \textbf{Save}, followed by Yes, to save the changes.
\end{enumerate}

\textbf{\emph{Note:}} To remove a counter, open a program in the
Step-Wise editor. Select the end step of the counter; access the Details
page in the step panel, reset the number in the Count field to 0, select
the checkmark button to apply the changes and save the program. For the
above procedure, press/click the body of Step 3, followed by the Details
button in the step panel, then reset the number in the Count field to 0
and press/click the checkmark button. Repeat the procedure for the next
counter. Save the changes via the \textbf{Save} button.

\textbf{\emph{Procedure:}} Table-Wise Method

\begin{enumerate}
\def\labelenumi{\arabic{enumi}.}
\item
  Switch the editor to Table-Wise method (see arrow).

  \begin{figure}
  \centering
  
\includegraphics[width=0.95\textwidth]{figures/2540780506-S5_3_6-example-step1-001a.PNG}
  \caption{Switching to table-wise editing}
  \end{figure}\FloatBarrier
\item
  Press/click the Count column (under Counter) on Step 3 and enter 1 in
  the cell (see arrow). \textbf{\emph{Note:}} By default, 1 is selected
  for the start step (in the Step cell, see arrow).

  \begin{figure}
  \centering
  
\includegraphics[width=0.95\textwidth]{figures/2074568092-S5_3_6-example-step2-001a.PNG}
  \caption{Selecting end step of first counter}
  \end{figure}\FloatBarrier
\item
  Press/click the Count column (under Counter) on Step 9 and enter 2 in
  the cell (see arrow). Enter 5 in the Step cell for start step (see
  arrow).

  \begin{figure}
  \centering
  
\includegraphics[width=0.95\textwidth]{figures/3200941779-S5_3_6-example-step3-001a.PNG}
  \caption{selecting end step of second counter}
  \end{figure}\FloatBarrier
\item
  Press/click \textbf{Save} (in the secondary menu), followed by Yes to
  save the changes.
\end{enumerate}

\textbf{\emph{Note:}} To remove the counters, open the program in the
Table-Wise editor and remove the count numbers 1 and 3, respectively,
from under the Count fields and click \textbf{Save}.


\section{Managing a Program: Export, Import and All
That}

Several operation options are available for managing programs in the GL
controller. They include Import, Export, Copy, Move or Swap, Rename,
Edit, Preview and Delete. These operations can be performed during any
chamber operating mode.

The following figure depicts a chamber operating in program mode,
currently executing Program 1 (see Section 5.3.4). Technically, when a
program is executed, the GL controller system loads that program into
its memory for execution. This means Program 1 (saved in the secondary
storage) can be manipulated as shown by the editing options in the
pop-up menu. Notice how slot 3 provides only two options (Editor and
Import) compared to slots 1 and 2.

\begin{figure}
\centering

\includegraphics[width=0.95\textwidth]{figures/828421102-S5_4-import-export-opt-001a.PNG}
\caption{Editing, Exporting, Importing options}
\end{figure}\FloatBarrier

The following table lists the available options for managing and
handling programs. Program step(s) manipulation will be discussed in
Section 5.5.

\begin{longtable}[]{@{}ll@{}}
\toprule
\begin{minipage}[b]{0.20\columnwidth}\raggedright
Name\strut
\end{minipage} & \begin{minipage}[b]{0.74\columnwidth}\raggedright
Description/Function\strut
\end{minipage}\tabularnewline
\midrule
\endhead
\begin{minipage}[t]{0.20\columnwidth}\raggedright
Import\strut
\end{minipage} & \begin{minipage}[t]{0.74\columnwidth}\raggedright
Copy program from external device (USB) into the GL controller
storage\strut
\end{minipage}\tabularnewline
\begin{minipage}[t]{0.20\columnwidth}\raggedright
Export\strut
\end{minipage} & \begin{minipage}[t]{0.74\columnwidth}\raggedright
Export selected program to an external USB device (or stored on a local
device that accesses the GL controller remotely). Program can be
exported in JSON or MS Excel via the appropriate icon.\strut
\end{minipage}\tabularnewline
\begin{minipage}[t]{0.20\columnwidth}\raggedright
Copy\strut
\end{minipage} & \begin{minipage}[t]{0.74\columnwidth}\raggedright
Program and its contents can be copied into a new location\strut
\end{minipage}\tabularnewline
\begin{minipage}[t]{0.20\columnwidth}\raggedright
Move/Swap\strut
\end{minipage} & \begin{minipage}[t]{0.74\columnwidth}\raggedright
Programs can be moved or swapped in different slots\strut
\end{minipage}\tabularnewline
\begin{minipage}[t]{0.20\columnwidth}\raggedright
Rename\strut
\end{minipage} & \begin{minipage}[t]{0.74\columnwidth}\raggedright
Program name can be quickly edited via this button\strut
\end{minipage}\tabularnewline
\begin{minipage}[t]{0.20\columnwidth}\raggedright
Editor\strut
\end{minipage} & \begin{minipage}[t]{0.74\columnwidth}\raggedright
Open selected program in editor mode\strut
\end{minipage}\tabularnewline
\begin{minipage}[t]{0.20\columnwidth}\raggedright
Preview\strut
\end{minipage} & \begin{minipage}[t]{0.74\columnwidth}\raggedright
Program can be quickly previewed in trend graph via a step-by-step
temp/humi output\strut
\end{minipage}\tabularnewline
\begin{minipage}[t]{0.20\columnwidth}\raggedright
Delete\strut
\end{minipage} & \begin{minipage}[t]{0.74\columnwidth}\raggedright
Program can be purged from the GL controller storage\strut
\end{minipage}\tabularnewline
\bottomrule
\end{longtable}


\subsection{Exporting a Program}

Programs can be exported to a USB device for backup or stored on a local
device that accesses the GL controller remotely (via a Web browser). As
depicted in the following figure, program can be exported in JSON, MS
Excel or ESPEC Corp.~GL compatibility mode via the appropriate icon.

\begin{figure}
\centering

\includegraphics[width=0.95\textwidth]{figures/1888251991-S5_4-import-export-003a.PNG}
\caption{Three export options are available}
\end{figure}\FloatBarrier

\begin{longtable}[]{@{}ll@{}}
\toprule
No. & Description\tabularnewline
\midrule
\endhead
1 & Export in JSON file for ESPEC North America, Inc.\tabularnewline
2 & Export in MS Excel file\tabularnewline
3 & Export in ESPEC CORP. GL compatibility mode\tabularnewline
\bottomrule
\end{longtable}

Program filename begins with chamber model and program slot number;
e.g., \textbf{GL-EPX-2-RM\_program\_1}. The following procedure outlines
the steps to export Program 1 (created in Section 5.3.4) to a USB device
(via the HMI).

\textbf{\emph{Procedure}}:

\begin{enumerate}
\def\labelenumi{\arabic{enumi}.}
\item
  Plug in a USB device into the USB port (see arrow) at the front of the
  chamber.

  \begin{figure}
  \centering
  
\includegraphics[width=0.95\textwidth]{figures/1703896270-S5_4_1N-external-ports-001a.PNG}
  \caption{External ports}
  \end{figure}\FloatBarrier
\item
  Access (or press) \textbf{Set Mode}
\item
  Press \textbf{Program 1}

  \begin{figure}
  \centering
  
\includegraphics[width=0.95\textwidth]{figures/692073832-S5_4_1-step4-001a.PNG}
  \caption{Exporting Program 1}
  \end{figure}\FloatBarrier
\item
  Export Options:

  \begin{itemize}
  \item
    \textbf{JSON Format (ESPEC North America, Inc)}: Press (click) the
    export button (see arrow 1 in the previous figure). Program will be
    automatically stored in the external USB device. Program filename
    will be: \textbf{GL-model\_program\_1.json}
  \item
    \textbf{MS Excel}: Press (click) the Excel (X) icon to download the
    program in MS Excel file. Program will be saved with \emph{.xlsx}
    extension.
  \item
    \textbf{ESPEC CORP. GL Compatibility mode}: Press (click) the JSON
    (\textless{}\textgreater{}) file icon to download the program.
    Program filename will be: \textbf{GL-model\_program\_1.json}
  \end{itemize}
\end{enumerate}

The following figures depict Program-1 in JSON and MS Excel file. It
requires some programming skills to read JSON file. However, the MS
Excel file lends itself readily to a wider audience.

\textbf{JSON Format:} Program 1 displayed as standard ASCII format

\begin{figure}
\centering

\includegraphics[width=0.95\textwidth]{figures/3523103571-S5_4_1-prog1-json.PNG}
\caption{Program 1 in JSON format}
\end{figure}\FloatBarrier

\textbf{MS Excel:} Program 1 displayed as a spreadsheet using MS Excel
program

\begin{figure}
\centering

\includegraphics[width=0.95\textwidth]{figures/3129773071-S5_4-excel-003b.PNG}
\caption{Program 1 in MS Excel format}
\end{figure}\FloatBarrier


\subsection{Importing a Program}

Programs can be imported into the GL controller, as long as the format
of the program structure is in the JSON format accepted by the GL
controller system. To avoid loading issues, imported programs should be
those that were originally exported out of the GL controller.

Imported programs can be created or edited outside of the GL controller
in JSON format, as long as they bear the structure required by the GL
controller programming system. For this reason, program(s) created by
the GL controller system should be used as a template or a model to
produce new programs for importing into the GL controller system.

\textbf{\emph{Procedure}}:

\begin{enumerate}
\def\labelenumi{\arabic{enumi}.}
\item
  Plug in a USB device into the USB port (see arrow) at the front of the
  chamber.

  \begin{figure}
  \centering
  
\includegraphics[width=0.95\textwidth]{figures/1703896270-S5_4_1N-external-ports-001a.PNG}
  \caption{External ports}
  \end{figure}\FloatBarrier
\item
  Access (or press) \textbf{Set Mode}
\item
  Press \textbf{Program 3}

  \begin{figure}
  \centering
  
\includegraphics[width=0.95\textwidth]{figures/3508985049-S5_4_2-image-001a.PNG}
  \caption{Importing Program 3}
  \end{figure}\FloatBarrier
\item
  A USB folder will be open. Browse the folder to select and import the
  program.
\item
  The imported program is now stored in slot 3.

  \begin{figure}
  \centering
  
\includegraphics[width=0.95\textwidth]{figures/1493040258-S4_4_2-step6-00-.PNG}
  \caption{Imported program in Slot 3}
  \end{figure}\FloatBarrier
\end{enumerate}


\subsection{Moving or Swapping a
Program}

A program can be relocated (moved) into a new slot. For demonstration,
Program 3 will be moved into slot 4.

\textbf{\emph{Procedure}}:

\begin{enumerate}
\def\labelenumi{\arabic{enumi}.}
\item
  Press Program 3 (see top arrow). Press \textbf{Move/Swap} button in
  the Program 3 drop-down panel (see arrow).

  \begin{figure}
  \centering
  
\includegraphics[width=0.95\textwidth]{figures/2943978485-S5_4_3-step1-003a.PNG}
  \caption{Selecting Program 3 for Move/Swap operation}
  \end{figure}\FloatBarrier
\item
  The blue pop-up notification (in lower right) indicates Program 3 has
  been selected for Move/Swap operation. To cancel the action, press
  \textbf{Cancel} in the notification. \textbf{\emph{Note:}} Moving over
  to a different menu will also cancel the Move/Swap action. Click
  Program 4 (see arrow) to move Program 3 in 4.

  \begin{figure}
  \centering
  
\includegraphics[width=0.95\textwidth]{figures/511542172-S5_4_3-step3-001a.PNG}
  \caption{Select Program 4 to move Program 3 in}
  \end{figure}\FloatBarrier
\item
  Confirm the action with Yes in the dialog box.

  \begin{figure}
  \centering
  
\includegraphics[width=0.95\textwidth]{figures/4087364197-S5_4_3-step4-001.PNG}
  \caption{Confirm the Move action}
  \end{figure}\FloatBarrier
\item
  Program 3 has been move into Program 4. Slot 3 is now empty, as shown
  in the figure.

  \begin{figure}
  \centering
  
\includegraphics[width=0.95\textwidth]{figures/534048827-S5_4_3-step5-001.PNG}
  \caption{Program 3 has been move and slot3 is now empty}
  \end{figure}\FloatBarrier
\end{enumerate}


\subsection{Renaming a Program}

A program can be renamed to something more descriptive. For
demonstration, Program 2 will be renamed to ``TempTest'' that actually
describes its program purpose.

\textbf{\emph{Procedure}}:

\begin{enumerate}
\def\labelenumi{\arabic{enumi}.}
\item
  Press Program 2 (see arrow). Press \textbf{Rename} button in the
  Program 2 drop-down panel (see arrow).

  \begin{figure}
  \centering
  
\includegraphics[width=0.95\textwidth]{figures/2459717929-S5_4_4-step1-002a.PNG}
  \caption{Rename Program 2}
  \end{figure}\FloatBarrier
\item
  Clear the ``Program 2'' and enter ``TempTest'', then press Ok. To
  cancel the action, press \textbf{Cancel}. \textbf{\emph{Note:}} Moving
  over to a different menu will also cancel the Rename action.

  \begin{figure}
  \centering
  
\includegraphics[width=0.95\textwidth]{figures/4022731286-S5_4_4-step2-002.PNG}
  \caption{Edit program name}
  \end{figure}\FloatBarrier
\item
  Program 2 has been renamed to ``TempTest'', as shown in the figure.

  \begin{figure}
  \centering
  
\includegraphics[width=0.95\textwidth]{figures/395812864-S5_4_4-step3-001.PNG}
  \caption{Program 2 has been renamed.}
  \end{figure}\FloatBarrier
\end{enumerate}


\subsection{Previewing a Program}

A program can be previewed using the \textbf{Preview} button in the
program drop-down panel. For demonstration, Program 4 (Humidity
Fluctuation) will be selected for preview.

\textbf{\emph{Procedure}}:

\begin{enumerate}
\def\labelenumi{\arabic{enumi}.}
\item
  Press Program 4 (see arrow). Press \textbf{Preview} button in the
  Program 4 drop-down panel (see arrow).

  \begin{figure}
  \centering
  
\includegraphics[width=0.95\textwidth]{figures/1150784192-S5_4_5-step1-002a.PNG}
  \caption{Selecting Program 4 for preview}
  \end{figure}\FloatBarrier
\item
  By default, contents of Program 4 are displayed for preview in
  Step-Wise mode. Click Table-Wise option to display Program 4 in that
  mode. \textbf{\emph{Note:}} In preview mode, program and its contents
  cannot be edited or changed.

  \begin{figure}
  \centering
  
\includegraphics[width=0.95\textwidth]{figures/3001426492-S5_4_5-step2-001.PNG}
  \caption{Program 4 in preview mode}
  \end{figure}\FloatBarrier
\item
  Press the \textbf{Back} button (in the secondary menu) to return to
  Program Setup.
\end{enumerate}


\subsection{Deleting a Program}

A program can be deleted from the program list. For demonstration,
Program 4 (Humidity Fluctuation) will be selected for removal.

\textbf{\emph{Procedure}}:

\begin{enumerate}
\def\labelenumi{\arabic{enumi}.}
\item
  Press Program 4 (see arrow). Press \textbf{Delete} button in the
  Program 4 drop-down panel (see arrow).

  \begin{figure}
  \centering
  
\includegraphics[width=0.95\textwidth]{figures/2644696631-S5_4_6-step1-002a.PNG}
  \caption{Selecting Program 4 for removal from the program list}
  \end{figure}\FloatBarrier
\item
  Program 4 (and its slot) is immediately removed from the list, along
  with empty slot 5, leaving just slot 3 ready to be populated by a new
  program.

  \begin{figure}
  \centering
  
\includegraphics[width=0.95\textwidth]{figures/3848750984-S5_4_6-step2-001a.PNG}
  \caption{Slots prior to slot 4 are available}
  \end{figure}\FloatBarrier
\end{enumerate}


\section{Editing a Program}

This section describes how to open an existing program for viewing and
editing. Changes made in the program can be saved back into the slot
(using Save); they can also be saved as a new program to occupy a new
slot (with Save As). The editing features include delete, move, insert,
copy, cut and paste range of steps. By default, the program editor
operates in Step-Wise editing mode.


\subsection{Delete Steps in a
Program}

Program steps can be removed as individual step or in a range of steps.
For demonstration, Program 2 (``TempTest'') will be selected to remove a
single step, followed by removing a range of steps. Program 2 consists
of 13 steps. Step 2 will be removed as a single step; step 5 through
step 7 will be removed as a range of steps.

\textbf{\emph{Note:}} To prepare Program 2 to contain 13 steps to follow
along with this demonstration, refer to Section 5.3.4. It is not
necessary to use slot 2 (Program 2); any available slot will work.

\textbf{\emph{Note:}} The following procedure outlines the steps based
on the HMI operation. For a remote access, the procedure is identical
starting with Step 2 by replacing the word \emph{press} with
\emph{click}.

\textbf{\emph{Procedure}}: Delete Step 2 in the Program

\begin{enumerate}
\def\labelenumi{\arabic{enumi}.}
\item
  If the GL controller system is already on, proceed to Step 2.
  Otherwise, set the main power breaker in the ON position. Within a
  minute, the HMI will come on.
\item
  Press \textbf{Set Mode} in the menu bar or the \textbf{Set Mode}
  button in the home page.
\item
  Press Program 2 (see top arrow), followed by the \textbf{Editor}
  button in the Program 2 drop-down panel.

  \begin{figure}
  \centering
  
\includegraphics[width=0.95\textwidth]{figures/1494949465-S5_5_1-step3-002ab.PNG}
  \caption{Selecting Program 2 for editing}
  \end{figure}\FloatBarrier
\item
  Press Step 2 header (see top arrow), followed by the \textbf{Delete}
  button in the drop-down panel (bottom arrow).

  \begin{figure}
  \centering
  
\includegraphics[width=0.95\textwidth]{figures/442520040-S5_5_1-step4-001a.PNG}
  \caption{Selecting Step 2 in the program}
  \end{figure}\FloatBarrier
\item
  Press Step 2 header again (see top arrow) to mark the range of steps
  for deletion, as instructed in the blue pop-up notification (bottom
  right). Press Yes to confirm the action in the pop-up dialog box. To
  cancel the current action, press No in the dialog box, followed by the
  \textbf{Cancel} button in the blue pop-up notification (lower right).
  \textbf{\emph{Note:}} Moving over to a different menu (e.g.,
  Table-Wise editing page) will also cancel the current action.

  \begin{figure}
  \centering
  
\includegraphics[width=0.95\textwidth]{figures/4207147614-S5_5_1-step5-001a.PNG}
  \caption{Marking Step 2 for deletion}
  \end{figure}\FloatBarrier
\item
  Step 2 is removed immediately; and Step 3 now becomes Step 2 (as can
  be verified by their parameters). Press the \textbf{Save} button,
  followed by Yes in the dialog box, to write the changes back into slot
  2.

  \begin{figure}
  \centering
  
\includegraphics[width=0.95\textwidth]{figures/1660515446-S5_5_1-step6-001a.PNG}
  \caption{Saving the changes in Program 2}
  \end{figure}\FloatBarrier
\end{enumerate}

\textbf{\emph{Procedure}}: Delete Step 5 through 7 in the Program

\begin{enumerate}
\def\labelenumi{\arabic{enumi}.}
\item
  Continuing from the previous procedure, press Step 5 header (see
  arrow), followed by the \textbf{Delete} button in the drop-down panel.

  \begin{figure}
  \centering
  
\includegraphics[width=0.95\textwidth]{figures/2020061316-S5_5_1-step1a-002a.PNG}
  \caption{Selecting Step 5 to set range of step removal}
  \end{figure}\FloatBarrier
\item
  Press Step 7 (see arrow) to mark the range of steps for deletion.

  \begin{figure}
  \centering
  
\includegraphics[width=0.95\textwidth]{figures/2633614511-S5_5_1-step2a-002a.PNG}
  \caption{Selecting Step 7 to mark selection}
  \end{figure}\FloatBarrier
\item
  Press Yes to confirm the action.

  \begin{figure}
  \centering
  
\includegraphics[width=0.95\textwidth]{figures/2810292390-S5_5_1-step3a-002a.PNG}
  \caption{Confirm step 5-7 removal}
  \end{figure}\FloatBarrier
\item
  Step 5 through 7 are removed immediately, as shown in the following
  figure. Press Save (see arrow) to save the changes in the program.

  \begin{figure}
  \centering
  
\includegraphics[width=0.95\textwidth]{figures/3179553772-S5_5_1-step4a-003a.PNG}
  \caption{Save changes to the program}
  \end{figure}\FloatBarrier
\end{enumerate}


\subsection{Moving Steps within the
Program}

Program step(s) can be moved to a different location, such as before or
after a certain step in the program with the \textbf{Move Before} or
\textbf{Move After} button. These buttons allow a selected step to be
swapped with its neighboring steps. For example, if Step 3 is selected
and the \textbf{Move After} button is pressed, Step 3 and Step 4 swap
their positions (with Step 3 now becoming Step 4 and vice versa). If the
\textbf{Move After} is pressed again, Step 4 and Step 5 are swapped.
This process can be repeated until a desired location for Step 3 is met.
The procedure is straight-forward.

\textbf{\emph{Procedure:}}

\begin{enumerate}
\def\labelenumi{\arabic{enumi}.}
\tightlist
\item
  Access the \textbf{Set Mode} menu.
\item
  Press the desired program on the program list, followed by the
  \textbf{Editor} button in the drop-down panel.
\item
  Press the desired step to mark selection.
\item
  Press the \textbf{Move Before} or \textbf{Move After} button in the
  drop-down panel to relocate the selected step.
\item
  Close the drop-down panel.
\item
  Save the changes.
\end{enumerate}

Program 2 from Section 5.5.1 is used for demonstration. The original
steps and their locations are depicted in the following figure.

\begin{figure}
\centering

\includegraphics[width=0.95\textwidth]{figures/3572524222-S5_5_2-orig-image.PNG}
\caption{Contents of Program 2}
\end{figure}\FloatBarrier

\textbf{\emph{Procedure:}} Move Step 3 to after Step 6.

\begin{enumerate}
\def\labelenumi{\arabic{enumi}.}
\item
  Assuming this activity is continued from Section 5.5.1, proceed to the
  next step (Step 2). Otherwise, if the GL system is already On, press
  \textbf{Set Mode} in the menu bar or the \textbf{Set Mode} button in
  the home page. Press \textbf{Program 2}, followed by the
  \textbf{Editor} button in the drop-down panel.
\item
  Press Step 3 (see top arrow). The drop-down panel appears.
  \textbf{\emph{Note:}} To cancel current operation, press X to close
  out the drop-down panel.

  \begin{figure}
  \centering
  
\includegraphics[width=0.95\textwidth]{figures/1702048269-S5_5_2-step2-001a.PNG}
  \caption{Selecting Step 3 to be moved}
  \end{figure}\FloatBarrier
\item
  Press \textbf{Move After} four times until Step 7 is highlighted as
  shown below; that is, Step 3 has moved into Step 7 (with Step 7 now
  becoming Step 6).

  \begin{figure}
  \centering
  
\includegraphics[width=0.95\textwidth]{figures/722400340-S5_5_2-step3-001a.PNG}
  \caption{Relocate Step 3 to after Step 7 with \textbf{Move After}}
  \end{figure}\FloatBarrier
\item
  Press X close the drop-down panel.

  \begin{figure}
  \centering
  
\includegraphics[width=0.95\textwidth]{figures/3886719421-S5_5_2-step4-001a.PNG}
  \caption{Close the drop-down panel}
  \end{figure}\FloatBarrier
\item
  Press the \textbf{Save} button (see arrow) to save the changes.

  \begin{figure}
  \centering
  
\includegraphics[width=0.95\textwidth]{figures/840120071-S5_5_2-step5-001a.PNG}
  \caption{Save the changes back in the original program}
  \end{figure}\FloatBarrier
\end{enumerate}


\subsection{Inserting Steps in the
Program}

An existing program can be edited by inserting new steps. Similar to the
preceding sections, Step-Wise or Table-Wise method can be used for the
task (refer to Section 5.3.1 and 5.3.2).

With the Step-Wise method, the \textbf{Insert Before} or \textbf{Insert
After} button can be used to insert a new step in desired location. If
Step 3 is selected, for instance, pressing the \textbf{Insert Before}
button will insert a new step to the left of Step 3; with \textbf{Insert
After} a new step is inserted to the right of Step 3. Step panel for the
new step is open with contents of Step 3 prefilled as the parameters.
These parameters can be modified to meet the requirement. Use the
\textbf{Check} button in the panel to apply settings to the inserted
step.

With the Table-Wise method, the add above or add below button (see
Section 5.3.2) can be used to insert the step. This approach is quite
subtle since all the parameters can be entered once the new step has
been created.

The following procedure presents both methods using the following
parameters for the new step.

\begin{longtable}[]{@{}lllll@{}}
\toprule
Step No. & Temp & Humi & Duration & Time Signal\tabularnewline
\midrule
\endhead
New Step & 72 & Off & 1HR 0Min 0Sec & 1 = On\tabularnewline
\bottomrule
\end{longtable}

\textbf{\emph{Note:}} The same procedure can be applied using the
\textbf{Insert After} button.

\textbf{\emph{Procedure:}} Step-Wise Method

\begin{enumerate}
\def\labelenumi{\arabic{enumi}.}
\item
  Assuming this activity is continued from Section 5.5.2, press
  \textbf{Program 2}, followed by the \textbf{Editor} button in the
  drop-down panel.
\item
  Press Step 3 (see arrow), followed by the \textbf{Insert Before}
  button in the drop-down panel (see arrow). \textbf{\emph{Note:}} To
  cancel the current action before committing to insert, press X.

  \begin{figure}
  \centering
  
\includegraphics[width=0.95\textwidth]{figures/26673975-S5_5_2-step2-001a.PNG}
  \caption{Inserting new step before Step 3}
  \end{figure}\FloatBarrier
\item
  Step panel is open with parameters prefilled. Modify the parameters
  using the contents in the above table. Press the \textbf{Check} button
  to apply the settings. \textbf{\emph{Note:}} The left or right button
  in the step panel has the same effect as the \textbf{Check} button.
  However, the X button does not apply the changes; it instead uses
  (i.e., writes) the default parameters in the inserted step.

  \begin{figure}
  \centering
  
\includegraphics[width=0.95\textwidth]{figures/2059538172-S5_5_2-step3-002a.PNG}
  \caption{Editing contents for the new step}
  \end{figure}\FloatBarrier
\item
  Press \textbf{Save} to apply the changes in Program 2.
\end{enumerate}

\textbf{\emph{Procedure:}} Table-Wise Method

\begin{enumerate}
\def\labelenumi{\arabic{enumi}.}
\item
  Assuming this activity is continued from Section 5.5.2, press
  \textbf{Program 2}, followed by the \textbf{Editor} button in the
  drop-down panel.
\item
  Press the Table-Wise subpage (see arrow).

  \begin{figure}
  \centering
  
\includegraphics[width=0.95\textwidth]{figures/1154104584-S5_5_2-step2a-001a.PNG}
  \caption{Setting Table-wise editing mode}
  \end{figure}\FloatBarrier
\item
  Press Step 4 (see arrow), then press the add above icon on Step 3 (see
  arrow).

  \begin{figure}
  \centering
  
\includegraphics[width=0.95\textwidth]{figures/1635463010-S5_5_3-step3aa-001a.PNG}
  \caption{Inserting new step above Step 3}
  \end{figure}\FloatBarrier
\item
  Edit the parameters in Step 3 using the contents in the above table.
  Press \textbf{Save}, followed by Yes in the dialog (pop-up) box to
  apply the changes in Program 2.

  \begin{figure}
  \centering
  
\includegraphics[width=0.95\textwidth]{figures/2418808331-S5_5_3-step4a-001a.PNG}
  \caption{Saving changes in Program 2}
  \end{figure}\FloatBarrier
\end{enumerate}


\subsection{Cutting/Pasting Steps within the
Program}

The cut-and-paste feature, available in the Step-Wise mode, allows
certain steps to be relocated within the program. One or more steps can
be relocated. The procedure is straight-forward:

\textbf{\emph{Procedure:}}

\begin{enumerate}
\def\labelenumi{\arabic{enumi}.}
\tightlist
\item
  Select a program for editing, followed by the \textbf{Editor} button
\item
  Select the initial step by its header (i.e., step number); press
  \textbf{Cut/Paste Step(s)}; select the last step to mark the range.
\item
  Select the step to mark the location for pasting.
  \textbf{\emph{Note:}} The selected step(s) will be pasted in front of
  the chosen step.\\
\item
  Press Yes in the dialog box to confirm the action. The following
  procedure outlines the steps to cut Step 3-5 and paste them in front
  of Step 1.
\end{enumerate}

\textbf{\emph{Procedure:}}

\begin{enumerate}
\def\labelenumi{\arabic{enumi}.}
\item
  Assuming this activity is continued from Section 5.5.3, press
  \textbf{Program 2}, followed by the \textbf{Editor} button in the
  drop-down panel.
\item
  Press Step 3 header, followed by the \textbf{Cut/Paste Step(s)} button
  in the drop-down panel.

  \begin{figure}
  \centering
  
\includegraphics[width=0.95\textwidth]{figures/2928284305-S5_5_4-step2-001a.PNG}
  \caption{Selecting initial step}
  \end{figure}\FloatBarrier
\item
  Press Step 5 to mark the range.

  \begin{figure}
  \centering
  
\includegraphics[width=0.95\textwidth]{figures/2651171131-S5_5_4-step3-002a.PNG}
  \caption{Selecting final step to set the range}
  \end{figure}\FloatBarrier
\item
  Press Step 1, followed by Yes in the dialog box to confirm the action.

  \begin{figure}
  \centering
  
\includegraphics[width=0.95\textwidth]{figures/2300914669-S5_5_4-step4-001a.PNG}
  \caption{Selecting Step 1 to paste the selected steps}
  \end{figure}\FloatBarrier
\item
  Step 3-5 should now be pasted in front of Step 1. Press \textbf{Save}
  to save changes to the program, followed by Yes to confirm the action.

  \begin{figure}
  \centering
  
\includegraphics[width=0.95\textwidth]{figures/2767006148-S5_5_4-step5-001a.PNG}
  \caption{Save the changes to the program}
  \end{figure}\FloatBarrier
\end{enumerate}


\section{Executing Programs}

Programs are executed in the Program Run panel under the
\textbf{Operation Mode} menu, which is accessible via the menu bar (1),
the button in the home page (2), or the system menu via the hamburger
button (3), as depicted in the following figure. Depending on how the
home page is configured (i.e., customized), a program may be executed
directly from within the home page via the \textbf{Start Prog. \#}
button (4). Any program available in the Program Run panel can be pinned
to the home page using the edit button (pen icon) in the secondary menu.
Section 5.6.3 explains in detail how to pin a program.

\begin{figure}
\centering

\includegraphics[width=0.95\textwidth]{figures/3645024435-S5_6-start-stop-program-001a.PNG}
\caption{How to execute a program}
\end{figure}\FloatBarrier

\begin{longtable}[]{@{}lll@{}}
\toprule
\begin{minipage}[b]{0.04\columnwidth}\raggedright
No.\strut
\end{minipage} & \begin{minipage}[b]{0.23\columnwidth}\raggedright
Name\strut
\end{minipage} & \begin{minipage}[b]{0.65\columnwidth}\raggedright
Description / Operation\strut
\end{minipage}\tabularnewline
\midrule
\endhead
\begin{minipage}[t]{0.04\columnwidth}\raggedright
1\strut
\end{minipage} & \begin{minipage}[t]{0.23\columnwidth}\raggedright
Operation Mode\strut
\end{minipage} & \begin{minipage}[t]{0.65\columnwidth}\raggedright
Direct access to the \textbf{Operation Mode} in the menu bar\strut
\end{minipage}\tabularnewline
\begin{minipage}[t]{0.04\columnwidth}\raggedright
2\strut
\end{minipage} & \begin{minipage}[t]{0.23\columnwidth}\raggedright
Operation Mode button\strut
\end{minipage} & \begin{minipage}[t]{0.65\columnwidth}\raggedright
Direct access to the \textbf{Operation Mode} via a button in the home
page\strut
\end{minipage}\tabularnewline
\begin{minipage}[t]{0.04\columnwidth}\raggedright
3\strut
\end{minipage} & \begin{minipage}[t]{0.23\columnwidth}\raggedright
Operation Mode under System Menu\strut
\end{minipage} & \begin{minipage}[t]{0.65\columnwidth}\raggedright
Alternative method to accessing the \textbf{Operation Mode} via the
system menu\strut
\end{minipage}\tabularnewline
\begin{minipage}[t]{0.04\columnwidth}\raggedright
4\strut
\end{minipage} & \begin{minipage}[t]{0.23\columnwidth}\raggedright
Start Prog\strut
\end{minipage} & \begin{minipage}[t]{0.65\columnwidth}\raggedright
Short cut to executing a program pinned on the home page; any program
can be pinned for a direct and quick execution.\strut
\end{minipage}\tabularnewline
\bottomrule
\end{longtable}

Any program available in the Program Run panel can be executed. The
following figure depicts the layout and operation buttons of the Program
Run panel.

\begin{figure}
\centering

\includegraphics[width=0.95\textwidth]{figures/3657323916-S5_6_1-prog-list-001b.PNG}
\caption{Program Run panel layout and operation buttons}
\end{figure}\FloatBarrier

\begin{longtable}[]{@{}lll@{}}
\toprule
\begin{minipage}[b]{0.04\columnwidth}\raggedright
No.\strut
\end{minipage} & \begin{minipage}[b]{0.23\columnwidth}\raggedright
Name\strut
\end{minipage} & \begin{minipage}[b]{0.65\columnwidth}\raggedright
Description and Operation\strut
\end{minipage}\tabularnewline
\midrule
\endhead
\begin{minipage}[t]{0.04\columnwidth}\raggedright
1\strut
\end{minipage} & \begin{minipage}[t]{0.23\columnwidth}\raggedright
Program List\strut
\end{minipage} & \begin{minipage}[t]{0.65\columnwidth}\raggedright
Displays all programs available in the panel.\strut
\end{minipage}\tabularnewline
\begin{minipage}[t]{0.04\columnwidth}\raggedright
2\strut
\end{minipage} & \begin{minipage}[t]{0.23\columnwidth}\raggedright
Scroll bar\strut
\end{minipage} & \begin{minipage}[t]{0.65\columnwidth}\raggedright
This scroll bar can be used to view all programs in the panel by sliding
up or down.\strut
\end{minipage}\tabularnewline
\begin{minipage}[t]{0.04\columnwidth}\raggedright
3\strut
\end{minipage} & \begin{minipage}[t]{0.23\columnwidth}\raggedright
Paging buttons\strut
\end{minipage} & \begin{minipage}[t]{0.65\columnwidth}\raggedright
Three up/down paging buttons can be used to page through the long list
of programs; button 3a or 3f takes to the top or bottom page; button 3b
or 3e views the list one page at a time; button 3c or 3d pages through
the list one row at a time.\strut
\end{minipage}\tabularnewline
\begin{minipage}[t]{0.04\columnwidth}\raggedright
4\strut
\end{minipage} & \begin{minipage}[t]{0.23\columnwidth}\raggedright
Operation buttons\strut
\end{minipage} & \begin{minipage}[t]{0.65\columnwidth}\raggedright
Three buttons are available to manage the program during execution; the
\textbf{Stop} button (top button) terminates the program immediately;
the \textbf{Pause/Resume} button (middle button) pauses/resumes the
program; the \textbf{Skip} button (bottom button) skips to the next step
in the program.\strut
\end{minipage}\tabularnewline
\bottomrule
\end{longtable}


\subsection{Executing a Program in the Program Run
Panel}

To execute a program, press (or click on) the desired program in the
panel. The following figure depicts the action to start Program 2. By
default, execution begins with Step 1. However, any step can be selected
to begin the execution, using the step selection field. The step pager
can be used to page through the steps (back and forth). To start the
execution, press (or click) \textbf{Start Program}. Press
\textbf{Cancel} to cancel the action.

\begin{figure}
\centering

\includegraphics[width=0.95\textwidth]{figures/1159962881-S5_6_2-start-001a.PNG}
\caption{Program execution options and features}
\end{figure}\FloatBarrier

\begin{longtable}[]{@{}lll@{}}
\toprule
\begin{minipage}[b]{0.04\columnwidth}\raggedright
No.\strut
\end{minipage} & \begin{minipage}[b]{0.23\columnwidth}\raggedright
Name\strut
\end{minipage} & \begin{minipage}[b]{0.65\columnwidth}\raggedright
Description\strut
\end{minipage}\tabularnewline
\midrule
\endhead
\begin{minipage}[t]{0.04\columnwidth}\raggedright
1\strut
\end{minipage} & \begin{minipage}[t]{0.23\columnwidth}\raggedright
Program No.\strut
\end{minipage} & \begin{minipage}[t]{0.65\columnwidth}\raggedright
Lists Program Number\strut
\end{minipage}\tabularnewline
\begin{minipage}[t]{0.04\columnwidth}\raggedright
2\strut
\end{minipage} & \begin{minipage}[t]{0.23\columnwidth}\raggedright
Program Name\strut
\end{minipage} & \begin{minipage}[t]{0.65\columnwidth}\raggedright
Lists program name; program number is displayed if program has no
name.\strut
\end{minipage}\tabularnewline
\begin{minipage}[t]{0.04\columnwidth}\raggedright
3\strut
\end{minipage} & \begin{minipage}[t]{0.23\columnwidth}\raggedright
Step No.\strut
\end{minipage} & \begin{minipage}[t]{0.65\columnwidth}\raggedright
Step 1 is selected by default to start execution.\strut
\end{minipage}\tabularnewline
\begin{minipage}[t]{0.04\columnwidth}\raggedright
4\strut
\end{minipage} & \begin{minipage}[t]{0.23\columnwidth}\raggedright
Step Selection\strut
\end{minipage} & \begin{minipage}[t]{0.65\columnwidth}\raggedright
A desired step number can be selected to begin the execution.\strut
\end{minipage}\tabularnewline
\begin{minipage}[t]{0.04\columnwidth}\raggedright
5\strut
\end{minipage} & \begin{minipage}[t]{0.23\columnwidth}\raggedright
Step Pager\strut
\end{minipage} & \begin{minipage}[t]{0.65\columnwidth}\raggedright
Use these buttons to page through the program to view or select a step
to start the execution.\strut
\end{minipage}\tabularnewline
\begin{minipage}[t]{0.04\columnwidth}\raggedright
6\strut
\end{minipage} & \begin{minipage}[t]{0.23\columnwidth}\raggedright
Start Program\strut
\end{minipage} & \begin{minipage}[t]{0.65\columnwidth}\raggedright
Use this button to start the program execution.\strut
\end{minipage}\tabularnewline
\begin{minipage}[t]{0.04\columnwidth}\raggedright
7\strut
\end{minipage} & \begin{minipage}[t]{0.23\columnwidth}\raggedright
Cancel\strut
\end{minipage} & \begin{minipage}[t]{0.65\columnwidth}\raggedright
Program execution can still be canceled with the \textbf{Cancel}
button.\strut
\end{minipage}\tabularnewline
\bottomrule
\end{longtable}

\textbf{\emph{Procedure:}}

\begin{enumerate}
\def\labelenumi{\arabic{enumi}.}
\tightlist
\item
  Press (or click) \textbf{Operation Mode}
\item
  Press (or click) the desired program in the Program Run panel.
\item
  Press (or click) \textbf{Start Program} to start the execution with
  Step 1.
\item
  If a desired step is required to start the execution, enter step
  number in the Step Selection field or select the step in the program,
  then press \textbf{Start Program}.
\end{enumerate}

\textbf{\emph{Note:}} If the GL controller and chamber are operated
remotely (via a Web browser), the following error message will appear if
the operator has insufficient privilege to execute the program or if the
user did not log in to their account.

\begin{figure}
\centering

\includegraphics[width=0.95\textwidth]{figures/960069550-S5_6_2-exec-error-001.PNG}
\caption{Program execution error messages and queue}
\end{figure}\FloatBarrier


\subsection{Executing a Program from within the Home
Page}

To execute a program on the home page, press (or click) \textbf{Start
Prog \#} (see arrow), followed by Yes in the dialog box to commence the
action. \textbf{\emph{Note:}} If this option is not available, proceed
to Section 5.6.3 for details how to pin a program. By default, the
\textbf{Start Constant 1} is configured to provide the quick execute
button.

\begin{figure}
\centering

\includegraphics[width=0.95\textwidth]{figures/3602099353-S5_6_3-execute-in-home-001a.PNG}
\caption{Executing a Program from within the Home Page}
\end{figure}\FloatBarrier

\textbf{\emph{Note:}} If the GL controller and chamber are operated
remotely (via a Web browser), the following error message identical to
that in Section 5.6.1 will appear if the operator has insufficient
privilege to execute the program or if the user did not log in to their
account.


\subsection{How to Pin a Program to the Home
Page}

The following procedure outlines the steps to pin a specific program to
the home page. By default, \textbf{Constant 1} is pinned to the home
page (see figure below). The entire home page can be customized using
the pin/unpin action via the \textbf{Edit} button in the secondary menu
(see above figure).

\begin{figure}
\centering

\includegraphics[width=0.95\textwidth]{figures/3611225745-S5_6_4-home-page-001a.PNG}
\caption{Default home page pin setting}
\end{figure}\FloatBarrier

The following figure depicts the home page in edit mode. In general,
home page editing can be performed during any operating mode. However,
it is recommended to perform this action during the standby mode.

\begin{figure}
\centering

\includegraphics[width=0.95\textwidth]{figures/3605228041-S5_6_4-home-page-002a.PNG}
\caption{Home page in edit mode}
\end{figure}\FloatBarrier

\begin{longtable}[]{@{}lll@{}}
\toprule
\begin{minipage}[b]{0.04\columnwidth}\raggedright
No.\strut
\end{minipage} & \begin{minipage}[b]{0.23\columnwidth}\raggedright
Name\strut
\end{minipage} & \begin{minipage}[b]{0.65\columnwidth}\raggedright
Description\strut
\end{minipage}\tabularnewline
\midrule
\endhead
\begin{minipage}[t]{0.04\columnwidth}\raggedright
1\strut
\end{minipage} & \begin{minipage}[t]{0.23\columnwidth}\raggedright
Edit\strut
\end{minipage} & \begin{minipage}[t]{0.65\columnwidth}\raggedright
Use this button to edit the selected widget or link.\strut
\end{minipage}\tabularnewline
\begin{minipage}[t]{0.04\columnwidth}\raggedright
2\strut
\end{minipage} & \begin{minipage}[t]{0.23\columnwidth}\raggedright
Add\strut
\end{minipage} & \begin{minipage}[t]{0.65\columnwidth}\raggedright
Use this \textbf{(+)} button to add a new component and pin it to the
home page. A new widget is added to the bottom of the standard display
area; it can be moved to a desired location. Refer to Section 6.12 for
details.\strut
\end{minipage}\tabularnewline
\begin{minipage}[t]{0.04\columnwidth}\raggedright
3\strut
\end{minipage} & \begin{minipage}[t]{0.23\columnwidth}\raggedright
Apply\strut
\end{minipage} & \begin{minipage}[t]{0.65\columnwidth}\raggedright
Use this \textbf{Check} mark to confirm and apply new settings to the
home page.\strut
\end{minipage}\tabularnewline
\begin{minipage}[t]{0.04\columnwidth}\raggedright
4\strut
\end{minipage} & \begin{minipage}[t]{0.23\columnwidth}\raggedright
Cancel\strut
\end{minipage} & \begin{minipage}[t]{0.65\columnwidth}\raggedright
Use this \textbf{Back} button to cancel the operation.\strut
\end{minipage}\tabularnewline
\bottomrule
\end{longtable}

The following procedure outlines the steps to pin Program 1 to the home
page. Step 1 of the program will be the default step to start the
execution.

\textbf{\emph{Procedure:}}

\begin{enumerate}
\def\labelenumi{\arabic{enumi}.}
\item
  Press \textbf{Home} (see arrow), followed by the \textbf{Edit} button
  in the secondary menu (see arrow).

  \begin{figure}
  \centering
  
\includegraphics[width=0.95\textwidth]{figures/3241351464-S5_6_3-execute-in-home-002a.PNG}
  \caption{Activating home page edit mode}
  \end{figure}\FloatBarrier
\item
  Press the \textbf{Edit} button of \textbf{Start Constant 1} (see
  arrow).

  \begin{figure}
  \centering
  
\includegraphics[width=0.95\textwidth]{figures/665574343-S5_6_3-step2-001a.PNG}
  \caption{Editing the Start/Stop button}
  \end{figure}\FloatBarrier
\item
  Press \textbf{Operation Mode} (see arrow). \textbf{\emph{Note:}} To
  cancel, press \textbf{Close}; to delete the current button, press
  \textbf{Delete}. Use the \textbf{Delete} button with caution as it
  requires adding the widget back into the spot. If you accidentally
  pressed \textbf{Delete}, press \textbf{Back} in the secondary menu to
  cancel this task. Refer to Section 6.12 for details.

  \begin{figure}
  \centering
  
\includegraphics[width=0.95\textwidth]{figures/1329716152-S5_6_3-step3-001a.PNG}
  \caption{Accessing the Operation Mode}
  \end{figure}\FloatBarrier
\item
  Press \textbf{Program 1} (see arrow) to select Program 1 (Range
  Checker).

  \begin{figure}
  \centering
  
\includegraphics[width=0.95\textwidth]{figures/2509727079-S5_6_3-step4-001a.PNG}
  \caption{Select Program 1}
  \end{figure}\FloatBarrier
\item
  Press \textbf{Start Program} (see arrow) to set execution with step 1
  by default. A desired step may be selected before pressing
  \textbf{Start Program}. However, in this example, there is only one
  step; and therefore, step selection is defaulted to step 1.

  \begin{figure}
  \centering
  
\includegraphics[width=0.95\textwidth]{figures/1600152636-S5_6_3-step5-001a.PNG}
  \caption{Set default step to start execution}
  \end{figure}\FloatBarrier
\item
  Press the check mark (see arrow), followed by Yes in the Attention
  dialog box to save the current settings.

  \begin{figure}
  \centering
  
\includegraphics[width=0.95\textwidth]{figures/4290535785-S5_6_3-step6-001a.PNG}
  \caption{Saving the changes to the home page}
  \end{figure}\FloatBarrier
\item
  Program 1 is now pinned to the home page as shown.

  \begin{figure}
  \centering
  
\includegraphics[width=0.95\textwidth]{figures/2745408708-S5_6_3-step7-001a.PNG}
  \caption{Program 1 pinned to the home page}
  \end{figure}\FloatBarrier
\end{enumerate}


\subsection{Manipulating a Program during its
Execution}

As stated in Section 5.4, when a program is selected for execution, the
GL controller loads that program from its secondary memory (i.e.,
storage) into its primary memory (i.e., RAM) to start the execution.
Thus, during execution, that program can still be loaded (from storage)
into the program editor for \emph{editing}, but several editing features
are forbidden. If the program has been edited or modified, such as with
steps being edited, added or removed, it cannot be saved (in the current
program slot), as depicted in the following figure.

\begin{figure}
\centering

\includegraphics[width=0.95\textwidth]{figures/1693827026-S5_6_4-exec-during-op-001.PNG}
\caption{Program being executed cannot be edited and saved}
\end{figure}\FloatBarrier

However, the edited version may be saved in a different slot under a
different name. In addition, that same program cannot be deleted during
its execution, as depicted in the following figure. In other words, if
you select Program \#1 (see arrow), then select \textbf{Delete} (see
arrow) in the pop-up dialog, an error flag will occur, indicating
\emph{cannot delete running program}.

\begin{figure}
\centering

\includegraphics[width=0.95\textwidth]{figures/1336002044-S5_6_4-delete-during-exec-001b.PNG}
\caption{Deleting a program during its execution}
\end{figure}\FloatBarrier


\section{Monitoring Program}

Chamber operating mode is displayed in the middle of the status bar at
all time. As depicted in the following figure, the chamber is in Program
mode. Detailed information of the program being executed is displayed in
the Standard Display area (see callout 2 and 3). Their descriptions are
listed in the following table.

\begin{figure}
\centering

\includegraphics[width=0.95\textwidth]{figures/822631156-S5_7_0-monitor-002a.PNG}
\caption{Operating status of program mode}
\end{figure}\FloatBarrier

\begin{longtable}[]{@{}ll@{}}
\toprule
\begin{minipage}[b]{0.07\columnwidth}\raggedright
Number\strut
\end{minipage} & \begin{minipage}[b]{0.88\columnwidth}\raggedright
Description\strut
\end{minipage}\tabularnewline
\midrule
\endhead
\begin{minipage}[t]{0.07\columnwidth}\raggedright
1\strut
\end{minipage} & \begin{minipage}[t]{0.88\columnwidth}\raggedright
Displays operating mode and status that include program name, process
values of temperature and/or humidity.\strut
\end{minipage}\tabularnewline
\begin{minipage}[t]{0.07\columnwidth}\raggedright
2\strut
\end{minipage} & \begin{minipage}[t]{0.88\columnwidth}\raggedright
Displays set point(s) and process value(s) of temperature (Air/Product)
and/or humidity.\strut
\end{minipage}\tabularnewline
\begin{minipage}[t]{0.07\columnwidth}\raggedright
3\strut
\end{minipage} & \begin{minipage}[t]{0.88\columnwidth}\raggedright
Displays program execution status, such as program number and name, step
being executed, total execution time, progress bar, remaining time and
date/time when program execution completes.\strut
\end{minipage}\tabularnewline
\bottomrule
\end{longtable}


\subsection{Monitor: Page 1}

The most detailed display page for any operating mode is the
\textbf{Monitor} page. Six different submonitor pages provide details of
the chamber and program operating condition.

Starting with Page 1, the display focuses on the main information of the
operating condition, namely, the set points and process values of
temperature (Air or Product) and humidity.

\textbf{Submonitor, Page 1}: Program Mode

\begin{figure}
\centering

\includegraphics[width=0.95\textwidth]{figures/555507061-S5_7_0-monitor-001a.PNG}
\caption{Monitor Page 1, displaying information of Program mode}
\end{figure}\FloatBarrier

\begin{longtable}[]{@{}ll@{}}
\toprule
\begin{minipage}[b]{0.07\columnwidth}\raggedright
Number\strut
\end{minipage} & \begin{minipage}[b]{0.88\columnwidth}\raggedright
Description\strut
\end{minipage}\tabularnewline
\midrule
\endhead
\begin{minipage}[t]{0.07\columnwidth}\raggedright
1\strut
\end{minipage} & \begin{minipage}[t]{0.88\columnwidth}\raggedright
Displays set points and process values of temperature (Air/Product). In
Standby mode, set points are OFF and process values are indicated by ---
for any unavailable option.\strut
\end{minipage}\tabularnewline
\begin{minipage}[t]{0.07\columnwidth}\raggedright
2\strut
\end{minipage} & \begin{minipage}[t]{0.88\columnwidth}\raggedright
Displays set points and process values of humidity. In Standby mode, set
points are OFF and process values are indicated by ---. Refer to Item
1.\strut
\end{minipage}\tabularnewline
\bottomrule
\end{longtable}


\subsection{Monitor: Page 2}

Submonitor Page 2 displays details of Time Signal outputs and operating
status of the refrigeration system during a program execution.

\textbf{Submonitor, Page 2:} Program Mode

\begin{figure}
\centering

\includegraphics[width=0.95\textwidth]{figures/3574379521-S5_7_2-monitor-p2-001a.PNG}
\caption{Monitor Page 2, displaying information of Program mode}
\end{figure}\FloatBarrier

\begin{longtable}[]{@{}ll@{}}
\toprule
\begin{minipage}[b]{0.07\columnwidth}\raggedright
Number\strut
\end{minipage} & \begin{minipage}[b]{0.88\columnwidth}\raggedright
Description\strut
\end{minipage}\tabularnewline
\midrule
\endhead
\begin{minipage}[t]{0.07\columnwidth}\raggedright
1\strut
\end{minipage} & \begin{minipage}[t]{0.88\columnwidth}\raggedright
Displays set points and process values of temperature (Air/Product)
and/or humidity.\strut
\end{minipage}\tabularnewline
\begin{minipage}[t]{0.07\columnwidth}\raggedright
2\strut
\end{minipage} & \begin{minipage}[t]{0.88\columnwidth}\raggedright
Displays operating status of the available time signals.\strut
\end{minipage}\tabularnewline
\begin{minipage}[t]{0.07\columnwidth}\raggedright
3\strut
\end{minipage} & \begin{minipage}[t]{0.88\columnwidth}\raggedright
Displays refrigeration system status controlled by the PLC. A two-stage
refrigeration system is shown in the display. Since the PLC controls the
refrigeration system and its stages, users have no control over it. The
display mainly provides information for the user how the refrigeration
is being controlled and operated, in terms of power, discharge or
evaporation temp., gas pressure or temperature.\strut
\end{minipage}\tabularnewline
\begin{minipage}[t]{0.07\columnwidth}\raggedright
4\strut
\end{minipage} & \begin{minipage}[t]{0.88\columnwidth}\raggedright
Edit button: Click to edit the layout of this page. Refer to Section
6.12 for details.\strut
\end{minipage}\tabularnewline
\bottomrule
\end{longtable}


\subsection{Monitor: Page 3}

Submonitor Page 3 displays details of chamber operating mode in Standby,
Constant or Program.

\textbf{Submonitor, Page 3:} Program Mode

\begin{figure}
\centering

\includegraphics[width=0.95\textwidth]{figures/883617183-S5_7_3-monitor-p3-001aa.PNG}
\caption{Detailed page of operating mode}
\end{figure}\FloatBarrier

\begin{longtable}[]{@{}ll@{}}
\toprule
\begin{minipage}[b]{0.07\columnwidth}\raggedright
Number\strut
\end{minipage} & \begin{minipage}[b]{0.88\columnwidth}\raggedright
Description\strut
\end{minipage}\tabularnewline
\midrule
\endhead
\begin{minipage}[t]{0.07\columnwidth}\raggedright
1\strut
\end{minipage} & \begin{minipage}[t]{0.88\columnwidth}\raggedright
Displays current condition or operating mode\strut
\end{minipage}\tabularnewline
\begin{minipage}[t]{0.07\columnwidth}\raggedright
2\strut
\end{minipage} & \begin{minipage}[t]{0.88\columnwidth}\raggedright
Displays trend graph of recent data and set mode prior to current
(operating) mode: Standby. Details of this display will be provided
later.\strut
\end{minipage}\tabularnewline
\begin{minipage}[t]{0.07\columnwidth}\raggedright
3\strut
\end{minipage} & \begin{minipage}[t]{0.88\columnwidth}\raggedright
Edit button: Click to edit the layout of this page. Refer to Section
6.12 for details.\strut
\end{minipage}\tabularnewline
\bottomrule
\end{longtable}


\subsection{Monitor: Detailed
Parameters}

Submonitor Page 4 displays a summary of the current operating mode. It
lists all the parameters associated with the current operating mode that
include control types and their status, such as temperature (On/Off),
humidity (On/Off), time signals (On/Off), refrigeration system (On/Off),
etc.

\begin{figure}
\centering

\includegraphics[width=0.95\textwidth]{figures/2942269458-S5_7_4-monitor-details-002a.PNG}
\caption{Detailed summary of parameters of current Program Mode}
\end{figure}\FloatBarrier


\subsection{Monitor: Preview and Program
Output}

Submonitor Page 5 displays details of program output and status. In
Standby or Constant mode, this submonitor page is blank. The following
figure depicts details of program status, showing set points and process
values of temperature and/or humidity, program runtime (both remaining
time and elapsed time), and date/time when program will be completed. It
also displays program steps and the current step (highlighted) being
executed. The scroll bar or paging buttons can be used to preview these
steps while the program is being executed.

\begin{figure}
\centering

\includegraphics[width=0.95\textwidth]{figures/923638922-S5_7_5-monitor-preview-002a.PNG}
\caption{Details outputs of program during execution}
\end{figure}\FloatBarrier

\begin{longtable}[]{@{}lll@{}}
\toprule
\begin{minipage}[b]{0.04\columnwidth}\raggedright
No.\strut
\end{minipage} & \begin{minipage}[b]{0.23\columnwidth}\raggedright
Name\strut
\end{minipage} & \begin{minipage}[b]{0.65\columnwidth}\raggedright
Description\strut
\end{minipage}\tabularnewline
\midrule
\endhead
\begin{minipage}[t]{0.04\columnwidth}\raggedright
1\strut
\end{minipage} & \begin{minipage}[t]{0.23\columnwidth}\raggedright
SP/PV\strut
\end{minipage} & \begin{minipage}[t]{0.65\columnwidth}\raggedright
Displays set points and process values of temperature and/or
humidity.\strut
\end{minipage}\tabularnewline
\begin{minipage}[t]{0.04\columnwidth}\raggedright
2\strut
\end{minipage} & \begin{minipage}[t]{0.23\columnwidth}\raggedright
Runtime\strut
\end{minipage} & \begin{minipage}[t]{0.65\columnwidth}\raggedright
Displays details of program execution time, its elapsed time and
remaining time.\strut
\end{minipage}\tabularnewline
\begin{minipage}[t]{0.04\columnwidth}\raggedright
3\strut
\end{minipage} & \begin{minipage}[t]{0.23\columnwidth}\raggedright
Step\strut
\end{minipage} & \begin{minipage}[t]{0.65\columnwidth}\raggedright
Display program steps; paging buttons or scroll bar can used to browse
the steps; current step being executed is highlighted.\strut
\end{minipage}\tabularnewline
\bottomrule
\end{longtable}

The following figure illustrates another program containing four steps
with Step 4 being carried out.

\begin{figure}
\centering

\includegraphics[width=0.95\textwidth]{figures/3856978094-S5_7_5-monitor-preview-003.PNG}
\caption{Execution status of a program, step by step}
\end{figure}\FloatBarrier


\subsection{Monitor: Trend}

Submonitor Page 6, called \emph{\textbf{Trend Graph}}, displays detailed
scatter plot of data points from the data log (past and present). This
graph allows the operator to view a scatter plot of data collected
during chamber operation. Data collection can occur in two distinct
settings: \textbf{Always} and \textbf{Running}. Default setting is
\textbf{Always}, which means data will be collected during any operating
mode. Section 6.3 (``Set Sampling'') provides details of these settings.

\begin{figure}
\centering

\includegraphics[width=0.95\textwidth]{figures/2746091548-S5_7_6-trend-003.PNG}
\caption{Trend graph of collected data}
\end{figure}\FloatBarrier

By default, the trend graph provides a scatter plot of data points
collected during the last hour. Various options of plot range of data
accumulated more than one hour can be selected. Data can be downloaded
(in whole or in portion) and stored on an external USB device (if
connected) or on the local computer if the GL system was accessed
remotely.

\begin{figure}
\centering

\includegraphics[width=0.95\textwidth]{figures/2268754256-S5_7_6-trend-003a.PNG}
\caption{Trend graph showing plots of current data from the chamber}
\end{figure}\FloatBarrier

The trend graph is made up of different components are described as
follows:

\begin{enumerate}
\def\labelenumi{\arabic{enumi}.}
\tightlist
\item
  \textbf{Trend Graph}: Data collected from the chamber are rendered as
  trend graph via scatter plot methodology; data points of temperature,
  air temperature, product temperature control and humidity are plotted
  as a function of time. The y-axis represents the scale of these
  values. The x-axis represents the data logging runtime.
\item
  \textbf{Trend Graph Manipulation Buttons}: Control and management of
  trend graph are provided by the manipulation buttons. Detailed
  functionality and operation of these buttons are discussed in a
  separate section (Section 5.7.7).
\item
  \textbf{Line Graph}: Different styles of graph--solid lines, dashed
  lines and color--are used to designate each type of data points for
  visualization. Solid lines represent process values, dashed lines set
  values, etc.\\
\item
  \textbf{Y-axis Label}: The y-axis displays the selected type of data
  points to be plotted (temperature, product temperature control,
  humidity). These values are populated based on the selection of the
  \textbf{Graph View} under item 2 (above).
\item
  \textbf{X-axis Label}: The x-axis displays log time of the data.
  Default plot range is one hour with grids; each grid sets a 5-minute
  time-scale.
\item
  \textbf{Y-axis Title}: The y-axis title provides clarity to the type
  of graph being represented.
\item
  \textbf{Legend of Trend Graph}: The legends provide identity of each
  item on the trend graph designated with color coding.
\item
  \textbf{Status Bar}: Displays progress bar of each operating mode.
\item
  \textbf{Status Mode}: Displays chamber operating mode with date and
  time. If a program is being executed, program name, slot number and
  step number will be displayed.
\end{enumerate}


\subsection{Trend Graph Manipulation
Buttons}

The following is a detailed list of the \textbf{Manipulation Buttons},
item 2 in the previous figure (Section 5.7.6).

\begin{figure}
\centering

\includegraphics[width=0.95\textwidth]{figures/3772339018-trend-graph-manip-buttons-001abc.PNG}
\caption{Trend Graph manipulation buttons}
\end{figure}\FloatBarrier

\begin{enumerate}
\def\labelenumi{\arabic{enumi}.}
\item
  \textbf{Auto Refresh}: This Auto Refresh button reconstructs the graph
  to include data points up to the current time.
\item
  \textbf{Plot Range Selectors}: Default plot range is one hour. With
  this button (and its options), plot range can be configured up to 12
  hours or 7 days. Custom set range is also available via the Custom
  Time Span option.
\item
  \textbf{Download Button}: Data can be downloaded for backup on an
  external USB (or local) device and saved in CSV format. Date/time
  format can be downloaded in UTC or local time, as indicated by the
  drop-down menu options. It is strongly recommended to download and
  clear data regularly. Data accumulation increases in size almost
  exponentially. As a result, storage space may be wasted; large data
  file may affect the operation and performance of the GL control
  system. Refer to Section 6.3 for details how to clear data file.
\item
  \textbf{Plot Configuration}: Elements of the trend graph can be
  configured via this button.
\item
  \textbf{Pan/Zoom Controls}: The collapsible Pan/Zoom Controls button
  allows the operator to control and adjust the viewable section in the
  trend graph. The up/down arrows can be used to collapse/expand the
  drop-down menu to manipulate the trend graph as follows.

  \begin{itemize}
  \item
    \textbf{Zoom In}: The \textbf{Zoom In} button allows the operator to
    zoom into a small section of the trend graph. Depending on the
    degree of zooming, the display area will be confined to a small set
    of data points ranging between minutes to hours. To reset the trend
    graph, click the \textbf{Zoom Extents} button, select \textbf{Last
    Hour} from the drop-down menu, then click the \textbf{Auto Refresh}
    button.
  \item
    \textbf{Zoom Out}: The \textbf{Zoom Out} button does the opposite by
    allowing the operator to zoom out on the trend graph, thereby giving
    the operator an expansive view of the trend graph. To reset the
    trend graph, click the \textbf{Zoom Extents} button, select
    \textbf{Last Hour} from the drop-down menu, then click the
    \textbf{Auto Refresh} button.
  \item
    \textbf{Move Up}: This button allows the operator to move up the
    graph along the vertical axis to adjust the viewable area of the
    scatter plot. To reset the trend graph, click the \textbf{Zoom
    Extents} button, select \textbf{Last Hour} from the drop-down menu,
    then click the \textbf{Auto Refresh} button.
  \item
    \textbf{Move Down}: This button allows the operator to move down the
    trend graph along the vertical axis with the purpose to adjust the
    viewable area of the scatter plot. To reset the trend graph, click
    the \textbf{Zoom Extents} button, select \textbf{Last Hour} from the
    drop-down menu, then click the \textbf{Auto Refresh} button.
  \item
    \textbf{Move Left}: This button allows the operator to pan left on
    the trend graph, offering a quick preview of a plot of data points
    tracing back the time in hours or days. With this feature, the
    operator can quickly gain a preview of past data points which the
    operator may have missed.
  \item
    \textbf{Move Right}: This button does the opposite to \textbf{Move
    Left} by allowing the operator to pan right on the trend graph to
    the current time. To reconstruct the trend graph to contain the most
    recent data points, the \textbf{Auto Refresh} button allows the
    quickest operation.
  \end{itemize}
\end{enumerate}


\chapter{GL Setup Menu}

The \textbf{Setup} menu provides a wide range of configuration options
for control and operation of the GL controller system and the chamber.
As such, this menu should be off limit to regular users. To protect the
GL controller system and chamber from accidental damages, the owner
(supervisor or administrator) who manages this GL controller system must
grant access to only the selected users to this menu. Access privileges
can be enforced on all users under the \textbf{Users (Register User
Password)} setup page. Refer to Section 7.4 for details.

According to Section 2.3.5, \textbf{Setup} can be accessed in two ways:
(1) Menu bar, (2) Hamburger button.

\textbf{Setup via Menu Bar}: This is the primary feature of the menu,
where each submenu is accessible under the \textbf{Setup} page (depicted
in the figure).

\begin{figure}
\centering

\includegraphics[width=0.95\textwidth]{figures/1028195047-S5-main-setup-menu-000a.PNG}
\caption{Setup menu in the menu bar}
\end{figure}\FloatBarrier

\textbf{System Menu via Hamburger button}: \textbf{Setup} can be
accessed via the hamburger button in the status bar (see arrow). It
provides a direct access via the system menu without affecting the
standard display.

\begin{figure}
\centering

\includegraphics[width=0.95\textwidth]{figures/1370477211-S5-main-setup-sys-menu-000ab.PNG}
\caption{Setup menu via the system menu of the hamburger button}
\end{figure}\FloatBarrier

\textbf{\emph{Note}}: The \textbf{Setup} menu can also be accessed from
within the Home page if it is custom configured to contain that menu.

The following sections describe these submenus in the order they appear
according to the list in the system menu (via the hamburger button),
with the except of \textbf{Configuration}, to be discussed separately as
it deserves a chapter of its own.


\section{Set Timers}

\textbf{\emph{Note}}: \textbf{Set Timers} option cannot be operated via
the HMI by the local user at the chamber. This feature must be operated
remotely, where the GL controller system is be part of a network and the
user accesses it from a PC via a Web browser. Refer to Chapter 8 for
details on how to access the GL controller system remotely on a network.

The \textbf{Timers} tab provides a quick way to apply the Macro Editor
to automate a simple macro action to extend the operating mode of the
chamber. There are two timer types: One-time timer and repeated timer.
The one-time timer or one-shot timer allows an operator to set up a
single (hence, one-time) operation on the chamber, while the repeated
timer permits an operator to set up a repeated operation on the chamber.
The repeated timer will operate indefinitely until it is disabled. The
following figure depicts these two timer types.

\begin{figure}
\centering

\includegraphics[width=0.95\textwidth]{figures/4083199896-illustrationtimer--001a.PNG}
\caption{Example of one-time and repeated timer operation}
\end{figure}\FloatBarrier

\begin{enumerate}
\def\labelenumi{\arabic{enumi}.}
\item
  \textbf{One-Time}: This timer fires the constant mode operation at
  3:06:58 PM on 9/6/2023. The timer instruction and 3:06:58 PM must be
  configured ahead of time, later than the current time in order for the
  operation to take effect. For this reason, the GL controller date/time
  must synchronize with the current or local date/time (refer to
  \textbf{Date/Time Settings} under the \textbf{Settings} menu for
  detail). The one-shot timer is date/time dependent.
\item
  \textbf{Repeated Daily/Weekly}: This line is added after the first
  line got executed. Since it is set to occur after the first line, it
  will trip at 3:15 PM on Wednesday and Thursday of the week (in
  reference to the current date/time of the GL controller). The action
  executes a program from slot 2 (also referenced by the program name).
  When program execution is complete, the chamber will be set in the
  operating mode specified by the instruction in the program. This timer
  type is date/time independent; the operation will recur until the
  timer itself is removed from the \textbf{Timers} list or disabled with
  the \textbf{Enabled} box unchecked.
\item
  \textbf{Add Timer}: A new timer can be added using the \textbf{Add
  Timer} button, which will be added at the bottom of the list.
\item
  \textbf{Remove/Delete}: Each timer can be removed using the \textbf{X}
  button to its right.
\item
  \textbf{Save/Apply}: To start the timer, the \textbf{Save} or
  \textbf{Apply} button must be used. The two timer instructions
  (depicted above) are saved under the Macro Editor called
  \textbf{Operation Timer 1} and \textbf{Operation Timer 2},
  respectively. If the timers are edited with new instructions and
  saved, the system generates a new name \textbf{Operation Timer 3} for
  the third timer, with a numeral incremented by one per each update. If
  the timer instructions are all removed from the \textbf{Timers} list,
  and a new timer is created and saved, the system starts over with
  \textbf{Operation Timer 1}, with its numeral incremented by one on
  each new update. This information can be viewed under the
  \textbf{Macros} submenu (under the \textbf{Settings} menu).
\end{enumerate}


\subsection{Program Mode: One-Time
Timer}

A one-shot timer sets the \textbf{Timers} to fire a certain action once
based on the set date and time. With this feature, a selected program
can start at a certain day and time once. The running program can be
paused, resumed or stopped by another (subsequent) one-shot timer. An
operator can plan to stop a certain program and switch the chamber to a
different mode with the one-shot timer.

The following example illustrates how to set a one-shot timer to switch
the chamber from \textbf{Standby} to \textbf{Program} and execute the
said program starting with step 1.

\begin{enumerate}
\def\labelenumi{\arabic{enumi}.}
\item
  Log in using the \textbf{admin} account
\item
  Access the \textbf{Start Stop} menu
\item
  Touch or click the \textbf{ADD TIMER} bar
\item
  Select/confirm that \textbf{One Time} is selected from the drop-down
  list of the \textbf{Timer} type; by default, \textbf{One Time} is
  selected.

  \begin{figure}
  \centering
  
\includegraphics[width=0.95\textwidth]{figures/689094357-one-shot-timer-000a.PNG}
  \caption{Setting a one-time automation}
  \end{figure}\FloatBarrier
\item
  Touch or click the date/time field under the \textbf{Timer Properties}
  and apply the calendar and time to set the desired time.
  \textbf{Note}: Date and time must be ahead of the current date and
  time on the GL controller system.

  \begin{figure}
  \centering
  
\includegraphics[width=0.95\textwidth]{figures/3455739059-date-time-001aaa.PNG}
  \caption{Set date/time for a one-shot action}
  \end{figure}\FloatBarrier

  \textbf{Note}: For a one-time operation, it is easier to use the
  Offset button (on the right of the date/time field) to adjust the
  future time, as depicted by the arrows in the following figure. Here,
  the GL controller will add the time to the current time to calculate
  when this action will fire.

  \begin{figure}
  \centering
  
\includegraphics[width=0.95\textwidth]{figures/2503503837-calendar-offset-001aa.PNG}
  \caption{Using the time-offset calculator}
  \end{figure}\FloatBarrier
\item
  Under \textbf{Operation}, select \textbf{Program: Run Program}
\item
  Under \textbf{Program}, select from name (or slot number)
\item
  Under \textbf{Step}, select what step to start
\item
  Apply the setting with the \textbf{APPLY} button

  \begin{figure}
  \centering
  
\includegraphics[width=0.95\textwidth]{figures/1332814146-set-program-mode-1.PNG}
  \caption{Program mode for a one-time start}
  \end{figure}\FloatBarrier

  \textbf{Note}: Both the \textbf{Save} and \textbf{APPLY} buttons
  provide the same effect. Either button can be used to save the
  settings and apply the action.
\item
  The program in slot 21 (as depicted in the example) will commence at
  5:59:28PM (according to the GL controller local time).
\end{enumerate}


\subsection{Removing or Adding a
Timer}

The one-shot or repeated timer can be removed from or added to the
\textbf{Timers} list with the \textbf{X} or \textbf{ADD TIMER} button.
Each new timer is added to the bottom of the list and takes effect after
each update.

To add a new timer and use it, simply apply the \textbf{ADD TIMER}
button. The following figure illustrates how a second timer is added to
switch the chamber from a program mode (activated by the first timer) to
\textbf{Constant}. Here, program from slot 21 gets executed (by the
timer in the first line) at 1:20:15 PM, with execution time to last
several hours. The timer in line 2 will switch the chamber to
\textbf{Constant} mode at 2:00:00 PM.

\begin{figure}
\centering

\includegraphics[width=0.95\textwidth]{figures/2397509582-set-new-timmer-003.PNG}
\caption{Setting additional one-shot timer after the first}
\end{figure}\FloatBarrier

Different timers on the list can be enabled or disabled using the
\textbf{Enabled} box. Additional timer can be added to the
\textbf{Timers} list and updated to take effect, provided the time is
set for future time.


\subsection{Constant Mode: One-Time
Timer}

The following example illustrates how to configure a one-shot timer to
stop the current program and switch the chamber to \textbf{Constant} at
a specific time. The timer is configured as a new instance with one line
by itself, as depicted in the following figure. The action is instant
and the program being executed will be stopped immediately; therefore,
the action must be applied with care.

\textbf{\emph{Procedure:}}

\begin{enumerate}
\def\labelenumi{\arabic{enumi}.}
\item
  Log in using the \textbf{admin} account
\item
  Access the \textbf{Start Stop} menu
\item
  Clear all or current timer settings with the \textbf{X} button (if
  necessary).
\item
  Touch or click the \textbf{ADD TIMER} bar
\item
  Confirm that \textbf{One Time} is selected (or make selection from the
  drop-down list).
\item
  Touch or click the date/time field under the \textbf{Timer Properties}
  and apply the calendar and time to set the desired time.
\item
  Under \textbf{Operation}, select \textbf{Standby: Stop Operation}
\item
  Apply the setting with the \textbf{APPLY} button. The timer is now
  saved in the Macros list.

  \begin{figure}
  \centering
  
\includegraphics[width=0.95\textwidth]{figures/122092262-p300-timer-prog-to-const-001.PNG}
  \caption{Scheduled a one-time timer to set chamber in Constant mode}
  \end{figure}\FloatBarrier
\item
  The timer is activated at the set time, as depicted in the Macros list
  under the Operation Timer 1 script. To view the Operation Timer 1
  action, access the Settings menu and select Macros from the submenu.

  \begin{figure}
  \centering
  
\includegraphics[width=0.95\textwidth]{figures/4142389927-p300-timer-prog-to-const-001a.PNG}
  \caption{Macros action recorded on the one-time Timer}
  \end{figure}\FloatBarrier
\end{enumerate}


\subsection{Program Mode: Weekly
Repetition}

With a weekly repeated timer, a program can start on a certain day, at a
certain time. The following example illustrates how to execute a program
on a weekly basis (one day during each week) at a certain time. Here, we
set a program execute every Tuesday at 4:00 PM.

\textbf{\emph{Procedure:}}

\begin{enumerate}
\def\labelenumi{\arabic{enumi}.}
\item
  Log in using the \textbf{admin} account
\item
  Access the \textbf{Start Stop} menu
\item
  Clear all or current timer settings with the \textbf{X} button (if
  necessary).
\item
  Touch or click the \textbf{ADD TIMER} bar
\item
  Touch or click \textbf{Timer Type} and select \textbf{Repeated
  Daily/Weekly}
\item
  Edit the time to 4:00 PM
\item
  Check the box under T. The calendar format is Sunday(S) to Saturday
  (S).
\item
  Under \textbf{Operation}, select \textbf{Program: Run Program Mode}
\item
  Under \textbf{Program}, select program in the desired slot number
  (slot 8, PROGTEST8, was used as an example).
\item
  Set and select program step 1 to begin execution.
\item
  Apply the setting using the \textbf{APPLY} button. The complete
  configuration is depicted in the following example.

  \begin{figure}
  \centering
  
\includegraphics[width=0.95\textwidth]{figures/2240313602-p300-timer-prog-repeatedly-001.PNG}
  \caption{Single-day weekly timer}
  \end{figure}\FloatBarrier
\end{enumerate}

This timer can be disabled or enabled by manipulating the
\textbf{Enabled} box. If you stop the program or switch the chamber to a
different mode (such as, \textbf{Constant}), while this timer is still
enabled, the selected program will recur on Tuesday at 4:00 PM the
following week. To disable the timer permanently, it must be removed
from the list (with the X and APPLY button).


\subsection{Program Mode: Daily
Repetition}

The following example illustrates how to execute a program on a daily
basis, from Monday through Friday, at 5:00 PM during each week.

\textbf{\emph{Procedure:}}

\begin{enumerate}
\def\labelenumi{\arabic{enumi}.}
\item
  Log in using the \textbf{admin} account
\item
  Access the \textbf{Start Stop} menu
\item
  Clear all or current timer settings with the \textbf{X} button (if
  necessary).
\item
  Touch or click the \textbf{ADD TIMER} bar
\item
  Touch or click \textbf{Timer Type} and select \textbf{Repeated
  Daily/Weekly}
\item
  Edit the time to 5:00 PM
\item
  Check the boxes under M through F. The calendar format is Sunday(S) to
  Saturday (S).
\item
  Under \textbf{Operation}, select \textbf{Program: Run Program Mode}
\item
  Under \textbf{Program}, select program from the desired slot number
  slot number (PROTEST7 from slot 8 was used as an example). Set and
  select program to begin at step 1.
\item
  Apply the settings with the \textbf{APPLY} button. The complete
  configuration is depicted in the following figure.

  \begin{figure}
  \centering
  
\includegraphics[width=0.95\textwidth]{figures/777446160-es102-prog-run-weekly-001b.PNG}
  \caption{Timer setting for everyday of the week}
  \end{figure}\FloatBarrier
\end{enumerate}

The following trend graph is a result of the timer from a different
chamber (used for illustration here) with its configuration set to run
over several days. The trend graph was produced by extrapolating the
data from the last three days (Friday through Sunday). It shows how
program was repeated each day at 6:05 PM and ended shortly before 9:00
AM, as indicated by the status bar at the bottom of the graph. The
chamber was put in a \textbf{Standby} mode until 6:05 PM, at which time
the program recurred.

\begin{figure}
\centering

\includegraphics[width=0.95\textwidth]{figures/2656486419-ex10_7_4ab.PNG}
\caption{Trend graph of chamber controlled by a timer}
\end{figure}\FloatBarrier

This timer can be disabled or enabled by manipulating the
\textbf{Enabled} box. If you stop the program or switch the chamber to a
different mode (such as, \textbf{Constant}), while this timer is still
enabled, the selected program will recur on the next day. To disable it
permanently, the timer can be removed from the list (with the X and
APPLY button).


\subsection{Multiple Timers: One-Time or Repeated
Daily/Weekly}

The \textbf{Timers} list may contain multiple timers with different
instructions for a specific operation. Multiple timers allow for
sequential activations, order of script execution or specific execution
via the \textbf{enabled} box selection, etc.

The following figure depicts a list of multiple timers, each with a
specific function and operation, which has been prepared to execute in
the top-to-bottom paradigm.

\begin{figure}
\centering

\includegraphics[width=0.95\textwidth]{figures/3002991730-multiple-tiemrs-s1076.PNG}
\caption{Multiple sequential timer oeprations}
\end{figure}\FloatBarrier

Each entry in the above figure is added after the previous one (above
it) has been executed. All of which are enabled by default. Only timer 1
(top line) and timer 4 (bottom line) remain in effect, since the middle
two are one-shot timers. In addition, the last timer is also a one-shot
timer and it will end after the program it has executed has ended. Each
timer causes action to take over the previous one.


\section{Firmware History}

The \textbf{Firmware History} page can be used to check against the
current release of GL controller system by ESPEC to ensure your system
is current and up to date. This firmware can also be used to check
against the original version shipped with the chamber.


\section{Set Sampling}

Contents in the trend graph of the \textbf{Monitor} menu are dictated by
the settings on this page. Data points collected from the chamber will
bear a date/time stamp of the GL controller system. For practical
purposes, date/time of the GL controller should reflect the current date
and local time zone. Refer to Section 7.2 for details on how to adjust
date/time on the GL controller.

The \textbf{Set Sampling} page consists of three main components: (1)
data logging frequency, (2) data logging mode and (3) logging items. As
depicted in the following figure, item 4 is the button to clear and
delete data in the log file. It is highly recommended to clear (i.e.,
delete) the data log regularly to avoid mass data accumulation both to
protect and maintain data storage space and the overall responsiveness
of the GL controller system. Data log can be downloaded for archive.
Refer to Section 4.9.8 or 5.7.7 for details.

\begin{figure}
\centering

\includegraphics[width=0.95\textwidth]{figures/3644375677-S5_3-set-sampling-001a.PNG}
\caption{Options of the Set Sampling page}
\end{figure}\FloatBarrier

\begin{longtable}[]{@{}lll@{}}
\toprule
\begin{minipage}[b]{0.04\columnwidth}\raggedright
Number\strut
\end{minipage} & \begin{minipage}[b]{0.23\columnwidth}\raggedright
Name\strut
\end{minipage} & \begin{minipage}[b]{0.65\columnwidth}\raggedright
Description/Operation\strut
\end{minipage}\tabularnewline
\midrule
\endhead
\begin{minipage}[t]{0.04\columnwidth}\raggedright
1\strut
\end{minipage} & \begin{minipage}[t]{0.23\columnwidth}\raggedright
Sampling Frequency\strut
\end{minipage} & \begin{minipage}[t]{0.65\columnwidth}\raggedright
Data logging interval selection\strut
\end{minipage}\tabularnewline
\begin{minipage}[t]{0.04\columnwidth}\raggedright
2\strut
\end{minipage} & \begin{minipage}[t]{0.23\columnwidth}\raggedright
Logging Mode\strut
\end{minipage} & \begin{minipage}[t]{0.65\columnwidth}\raggedright
Data logging mode selection; available in \textbf{Always} and
\textbf{Running}\strut
\end{minipage}\tabularnewline
\begin{minipage}[t]{0.04\columnwidth}\raggedright
3\strut
\end{minipage} & \begin{minipage}[t]{0.23\columnwidth}\raggedright
Logging Item\strut
\end{minipage} & \begin{minipage}[t]{0.65\columnwidth}\raggedright
Data logging items; items to be included in data log\strut
\end{minipage}\tabularnewline
\begin{minipage}[t]{0.04\columnwidth}\raggedright
4\strut
\end{minipage} & \begin{minipage}[t]{0.23\columnwidth}\raggedright
Clear Data\strut
\end{minipage} & \begin{minipage}[t]{0.65\columnwidth}\raggedright
Clears data file for new data log\strut
\end{minipage}\tabularnewline
\bottomrule
\end{longtable}

Data logging is controlled by the sampling frequency (Item 1) and
chamber operating mode (Item 2). The default sampling frequency occurs
at a 5-second interval. An arbitrary value can be used to set the
sampling frequency in Item 1. Two operating modes are available to
control data logging: \textbf{Always} and \textbf{Running}. When
\textbf{Always} is selected, data will be logged at all time. When
\textbf{Running} is selected, data will be logged only when the chamber
is running, both in Constant and Program mode.

Data points and types (such as, temperature, humidity, etc.) can be
manually selected to appear in the trend graph. Selection can be made
based on the check mark placed in their respective boxes (shown in the
following figure). This page is scrollable to allow all available items
be selected and included in the data log. Color decoration can be
selected for plot legends. Different logs for a trend graph can be
created using the (+) button. The y-axis (vertical) scaling can be
adjusted via the Y Axis Max and Y Axis Min buttons. Detailed steps for
these settings are discussed in the following sections.

\begin{figure}
\centering

\includegraphics[width=0.95\textwidth]{figures/3968747127-S6_3-set-sampling-003a.PNG}
\caption{Item selection, data type and color options}
\end{figure}\FloatBarrier

\begin{longtable}[]{@{}ll@{}}
\toprule
\begin{minipage}[b]{0.07\columnwidth}\raggedright
No.\strut
\end{minipage} & \begin{minipage}[b]{0.88\columnwidth}\raggedright
Description\strut
\end{minipage}\tabularnewline
\midrule
\endhead
\begin{minipage}[t]{0.07\columnwidth}\raggedright
1\strut
\end{minipage} & \begin{minipage}[t]{0.88\columnwidth}\raggedright
Default log name for trend graph preview. The first log is called
\emph{default}.\strut
\end{minipage}\tabularnewline
\begin{minipage}[t]{0.07\columnwidth}\raggedright
2\strut
\end{minipage} & \begin{minipage}[t]{0.88\columnwidth}\raggedright
A new log can be created using the \textbf{+} button; default name is
\textbf{\emph{new}}. This name can be changed via Item No.~4
(below).\strut
\end{minipage}\tabularnewline
\begin{minipage}[t]{0.07\columnwidth}\raggedright
3\strut
\end{minipage} & \begin{minipage}[t]{0.88\columnwidth}\raggedright
Use this button to delete a log. The selected log is deleted when the
trash bin is pressed (or clicked).\strut
\end{minipage}\tabularnewline
\begin{minipage}[t]{0.07\columnwidth}\raggedright
4\strut
\end{minipage} & \begin{minipage}[t]{0.88\columnwidth}\raggedright
Log name can be created or edited using this field. This option operates
in concert with No.~2 (above).\strut
\end{minipage}\tabularnewline
\begin{minipage}[t]{0.07\columnwidth}\raggedright
5\strut
\end{minipage} & \begin{minipage}[t]{0.88\columnwidth}\raggedright
Lists available items selectable to include in log preview.\strut
\end{minipage}\tabularnewline
\begin{minipage}[t]{0.07\columnwidth}\raggedright
6\strut
\end{minipage} & \begin{minipage}[t]{0.88\columnwidth}\raggedright
Minimum and maximum vertical scaling can be adjusted using these two
fields.\strut
\end{minipage}\tabularnewline
\begin{minipage}[t]{0.07\columnwidth}\raggedright
7\strut
\end{minipage} & \begin{minipage}[t]{0.88\columnwidth}\raggedright
Legends of different data points can be characterized with unique
properties for easy identification.\strut
\end{minipage}\tabularnewline
\begin{minipage}[t]{0.07\columnwidth}\raggedright
8\strut
\end{minipage} & \begin{minipage}[t]{0.88\columnwidth}\raggedright
The \textbf{Save} button can be used to apply the new settings; use the
\textbf{Back} button in the secondary menu to cancel the current
action.\strut
\end{minipage}\tabularnewline
\bottomrule
\end{longtable}


\subsection{Set Data Log in any
Mode}

The following procedure outlines the steps to configure data log in any
operating mode with logging frequency at a 5-second interval; data
points will be logged every five seconds.

\textbf{\emph{Procedure:}}

\begin{enumerate}
\def\labelenumi{\arabic{enumi}.}
\item
  If the GL controller system is already on and its HMI is in sleep
  mode, tab your finger on the HMI and proceed to Step 2. From a cold
  start, set the main power breaker in the ON position. Within a minute,
  the HMI will come on. The system automatically logs in using the
  default localhost user account to operate the chamber.
\item
  Press the \textbf{Hamburger} menu (1) in the status bar. Check to
  confirm the floating keyboard is enabled. A diagonal bar over the
  keyboard indicates it is disabled. Press the keyboard (2) to enable
  it. Press \textbf{Setup} (3) followed by \textbf{Set Sampling} (4).

  \begin{figure}
  \centering
  
\includegraphics[width=0.95\textwidth]{figures/2960242964-S5_3_1-step1-001a.PNG}
  \caption{Accessing the Set Sampling page}
  \end{figure}\FloatBarrier
\item
  Press the \textbf{Delete all Logged Data} button (see callout 1) to
  reset data log file. Press the sampling frequency field (see callout
  2) and set the value to 5 via the pop-up keypad (not shown). Press the
  \textbf{Always} button (see callout 2). Press the \textbf{Save} button
  to save the current settings.

  \begin{figure}
  \centering
  
\includegraphics[width=0.95\textwidth]{figures/3028734154-S5_3_1-step3-001aa.PNG}
  \caption{Setting up data log file}
  \end{figure}\FloatBarrier
\item
  Data points will consist of the selected items in the default setting.
  These items are depicted in the above figure with the check mark in
  their boxes.
\end{enumerate}


\subsection{Set Data Log in Program
Mode}

The following procedure outlines the steps to configure data log in
Program mode with logging frequency at a 5-second interval; data points
will be logged every five seconds.

\textbf{\emph{Procedure:}}

\begin{enumerate}
\def\labelenumi{\arabic{enumi}.}
\item
  Repeat Step 1 and 2 in Section 6.3.1.
\item
  Press the \textbf{Delete all Logged Data} button (see callout 1) to
  reset data log file. Press the sampling frequency field (see callout
  2) and set the value to 5 via the pop-up keypad (not shown). Press the
  \textbf{Always} button (see callout 2). Press the \textbf{Save} button
  to save the current settings.

  \begin{figure}
  \centering
  
\includegraphics[width=0.95\textwidth]{figures/3028734154-S5_3_1-step3-001aa.PNG}
  \caption{Setting up data log file}
  \end{figure}\FloatBarrier
\item
  Data points will consist of the selected items in the default setting.
  These items are depicted in the above figure with the check mark in
  their boxes.
\end{enumerate}


\subsection{Example: Creating Custom
Log}

This section outlines the steps to create a custom log file for trend
graph preview. In this example, we set the chamber to operate in program
mode (using Program 2 discussed in Section 5.5). Our custom log will
include program ID and name, program step, refrigeration system and
alarm, as well as temp/humi SP and PV; these parameters will be checked
against the \emph{default} log.

\textbf{\emph{Procedure:}}

\begin{enumerate}
\def\labelenumi{\arabic{enumi}.}
\item
  Press \textbf{Setup}, followed by \textbf{Set Sampling}.
\item
  Apply \emph{default} setting and select \emph{status}, \emph{Temp Air
  SP}, \emph{Temp Prod. SP}, and their counter PV, then save the
  settings with the \textbf{Save} button in the secondary menu. Access
  \textbf{Monitor} and select \textbf{Trend} subpage as shown below.

  \begin{figure}
  \centering
  
\includegraphics[width=0.95\textwidth]{figures/2342323574-S6_3_3-step2-001.PNG}
  \caption{Default log setting}
  \end{figure}\FloatBarrier
\item
  Press \textbf{Setup}, followed by \textbf{Set Sampling}.
\item
  Create a new log called \emph{custom1} as follows: Press \textbf{+},
  enter \emph{custom1} in the name field.
\item
  Select the following items:

  \begin{itemize}
  \tightlist
  \item
    \emph{Status}
  \item
    \emph{PGM Step}
  \item
    \emph{PGM ID}
  \item
    \emph{PGM}
  \item
    \emph{Temp Air SP, Temp Prod. SP, Temp Air PV, Temp Prod. PV, and
    Temp EN}
  \item
    \emph{Refrig}
  \item
    \emph{Alarms}
  \end{itemize}
\item
  Press \textbf{Save}, followed by Yes in the dialog box.
\item
  Press \textbf{Operation Mode}, then execute Program 2.
\item
  Press \textbf{Monitor}, followed by the \textbf{Trend} submonitor
  page. The log being displayed is the \emph{default} log.

  \begin{figure}
  \centering
  
\includegraphics[width=0.95\textwidth]{figures/3813115736-S6_3_3-step8-002.PNG}
  \caption{Trend of default log setting}
  \end{figure}\FloatBarrier
\item
  Press the \textbf{Trend} icon in the upper-right (see arrow 1) and
  check the \emph{custom1} box (arrow 2) to display \emph{custom1} log.

  \begin{figure}
  \centering
  
\includegraphics[width=0.95\textwidth]{figures/1616083375-S6_3_3-step9-0022a.PNG}
  \caption{Selecting custom log}
  \end{figure}\FloatBarrier
\item
  Trend graph based on \emph{custom1} is now displayed, which can be
  compared to the \emph{default} log.

  \begin{figure}
  \centering
  
\includegraphics[width=0.95\textwidth]{figures/186329913-S6_3_3-step10-002.PNG}
  \caption{Trend graph of custom setting}
  \end{figure}\FloatBarrier
\end{enumerate}


\section{Alarm/Operation History}

The \textbf{Alarm/Operation History} page displays operation history of
the chamber that includes alarms/warnings, operating modes and
statistics. Any alarms or alerts tripped during the chamber operation
(and logged in the data file) are displayed here. Information of the
history log is linked to the data log. Therefore, history logs can go
back to as far as the available data in the data log. This page
therefore displays the chamber operation history based on the
information extracted from the data log.

\textbf{\emph{Note:}} Since the data log is recommended to be cleared
regularly in order to maintain storage space as well as the overall
performance of the GL controller, it is recommended that information of
the \textbf{Alarm/Operation History} be reviewed or analyzed regularly
prior to deleting and resetting the log file.

The quickest way to access and view \textbf{Alarm/Operation History} is
as follows.

\textbf{\emph{Procedure:}}

\begin{enumerate}
\def\labelenumi{\arabic{enumi}.}
\item
  If the GL controller system is already on and its touchscreen is in
  sleep mode, tab your finger on the touchscreen and proceed to Step 2
  (below). From a cold start, set the main power breaker in the ON
  position. Within a minute, the touch-screen display will come on. The
  system automatically logs in using the default localhost user account
  to operate the chamber.
\item
  Press \textbf{Setup} in the menu bar (see arrow 1). Press the
  \textbf{Alarm/Operation History} tab (see arrow 2).

  \begin{figure}
  \centering
  
\includegraphics[width=0.95\textwidth]{figures/4247892763-S5_4-accessing-hist-page-001a.PNG}
  \caption{Accessing the Alarm/Operation History page}
  \end{figure}\FloatBarrier
\end{enumerate}

The \textbf{Alarm/Operation History} consists of three separate
subpages, each providing different information of the operation history.

\begin{figure}
\centering

\includegraphics[width=0.95\textwidth]{figures/3791077674-S5_4-alarm-main-page-001aa.PNG}
\caption{Operation history of the chamber}
\end{figure}\FloatBarrier

\begin{longtable}[]{@{}lll@{}}
\toprule
\begin{minipage}[b]{0.04\columnwidth}\raggedright
No.\strut
\end{minipage} & \begin{minipage}[b]{0.23\columnwidth}\raggedright
Name\strut
\end{minipage} & \begin{minipage}[b]{0.65\columnwidth}\raggedright
Description/Function\strut
\end{minipage}\tabularnewline
\midrule
\endhead
\begin{minipage}[t]{0.04\columnwidth}\raggedright
1\strut
\end{minipage} & \begin{minipage}[t]{0.23\columnwidth}\raggedright
Alarm/Warning\strut
\end{minipage} & \begin{minipage}[t]{0.65\columnwidth}\raggedright
List of alarm and warning occurred during chamber operation\strut
\end{minipage}\tabularnewline
\begin{minipage}[t]{0.04\columnwidth}\raggedright
2\strut
\end{minipage} & \begin{minipage}[t]{0.23\columnwidth}\raggedright
Operating Mode\strut
\end{minipage} & \begin{minipage}[t]{0.65\columnwidth}\raggedright
List of operating modes taking place during chamber operation\strut
\end{minipage}\tabularnewline
\begin{minipage}[t]{0.04\columnwidth}\raggedright
3\strut
\end{minipage} & \begin{minipage}[t]{0.23\columnwidth}\raggedright
Operation Statistics\strut
\end{minipage} & \begin{minipage}[t]{0.65\columnwidth}\raggedright
List of operation statistics during chamber operation\strut
\end{minipage}\tabularnewline
\bottomrule
\end{longtable}


\subsection{Alarm/Warning}

Details of alarms and warnings that occurred during chamber operation
are tabulated in five separate columns to provide specific information
of each event. Events are listed in a top-down order, with recent ones
at the top, earlier ones at the bottom. The table can scroll up or down
by swiping the touchscreen in the down/up direction, namely, swipe the
touchscreen up to scroll down the table.

\begin{figure}
\centering

\includegraphics[width=0.95\textwidth]{figures/2664764610-S5_4_1-alarm-list-002a.PNG}
\caption{List of Alarm/Warning}
\end{figure}\FloatBarrier

\begin{longtable}[]{@{}lll@{}}
\toprule
\begin{minipage}[b]{0.04\columnwidth}\raggedright
No.\strut
\end{minipage} & \begin{minipage}[b]{0.23\columnwidth}\raggedright
Name\strut
\end{minipage} & \begin{minipage}[b]{0.65\columnwidth}\raggedright
Description\strut
\end{minipage}\tabularnewline
\midrule
\endhead
\begin{minipage}[t]{0.04\columnwidth}\raggedright
1\strut
\end{minipage} & \begin{minipage}[t]{0.23\columnwidth}\raggedright
Type\strut
\end{minipage} & \begin{minipage}[t]{0.65\columnwidth}\raggedright
Lists the type of alert as alarm (ALM) or warning (WAR); alarm or
warning ID's are included in the list\strut
\end{minipage}\tabularnewline
\begin{minipage}[t]{0.04\columnwidth}\raggedright
2\strut
\end{minipage} & \begin{minipage}[t]{0.23\columnwidth}\raggedright
Alarm\strut
\end{minipage} & \begin{minipage}[t]{0.65\columnwidth}\raggedright
List alarm or warning message to provide specific information\strut
\end{minipage}\tabularnewline
\begin{minipage}[t]{0.04\columnwidth}\raggedright
3\strut
\end{minipage} & \begin{minipage}[t]{0.23\columnwidth}\raggedright
Trip Time\strut
\end{minipage} & \begin{minipage}[t]{0.65\columnwidth}\raggedright
List date and time when alarm or warning occurred\strut
\end{minipage}\tabularnewline
\begin{minipage}[t]{0.04\columnwidth}\raggedright
4\strut
\end{minipage} & \begin{minipage}[t]{0.23\columnwidth}\raggedright
Clear Time\strut
\end{minipage} & \begin{minipage}[t]{0.65\columnwidth}\raggedright
List date and time when the trip was resolved (cleared)\strut
\end{minipage}\tabularnewline
\begin{minipage}[t]{0.04\columnwidth}\raggedright
5\strut
\end{minipage} & \begin{minipage}[t]{0.23\columnwidth}\raggedright
Trend Graph\strut
\end{minipage} & \begin{minipage}[t]{0.65\columnwidth}\raggedright
Provides an access link to view the trend graph when the trip
occurred\strut
\end{minipage}\tabularnewline
\bottomrule
\end{longtable}

Each entry in the table is attached to it a trend graph (in the last
column) to provide a graphical overview in the timeline when an event
occurred.

The following procedure outlines the steps to open a trend graph to view
details of an overheating alarm that occurred during chamber operation,
then return to the alarm/warning table. \textbf{Overheating Alarm} is
selected for this example, since it may reflect real event during a
production run.

\textbf{\emph{Procedure:}} Example

\begin{enumerate}
\def\labelenumi{\arabic{enumi}.}
\item
  Locate the \textbf{Overheating Alarm}, as depicted in the figure
  below. Press the trend graph icon (see arrow).

  \begin{figure}
  \centering
  
\includegraphics[width=0.95\textwidth]{figures/2310942176-S5_4_1-step1-002a.PNG}
  \caption{Viewing trend graph and details of the selected alarm}
  \end{figure}\FloatBarrier
\item
  The timeline in the trend graph indicated that \textbf{Overheating
  Alarm} occurred between 3:41 PM and 3:48 PM (see arrow 1) and was
  tripped by and during the constant operation. Overheating occurred
  when the Temp Air process value (PV) reached the SP value. For further
  analysis, data points of the trend graph may be downloaded via the
  trend graph manipulation buttons (see arrow 2). Refer to Section 4.9.7
  for further details. This trend graph is essentially the trend graph
  invoked in the \textbf{Monitor} page.

  \begin{figure}
  \centering
  
\includegraphics[width=0.95\textwidth]{figures/3000926666-S5_4_1-step2-002a.PNG}
  \caption{Details of the selected alarm}
  \end{figure}\FloatBarrier
\item
  Press the back button (see arrow 3 in above figure) to return to the
  Alarm/Warning table.
\end{enumerate}


\subsection{List of Operating
Modes}

The second subpage of \textbf{Alarm/Operation History} provides a list
of chamber operating modes, also consisting of five columns. Operating
modes are listed in a top-down order, with recent ones at the top,
earlier ones at the bottom. The table can scroll up or down by swiping
the touchscreen in the down/up direction, namely, swipe the touchscreen
up to scroll down the table. The table in the figure only shows few
items on the list; scroll down the table to view earlier operating modes
recorded on the table.

\begin{figure}
\centering

\includegraphics[width=0.95\textwidth]{figures/4127612676-S5_4_2-mode-display-003a.PNG}
\caption{A table listing operating modes}
\end{figure}\FloatBarrier

\begin{longtable}[]{@{}lll@{}}
\toprule
\begin{minipage}[b]{0.04\columnwidth}\raggedright
No.\strut
\end{minipage} & \begin{minipage}[b]{0.23\columnwidth}\raggedright
Name\strut
\end{minipage} & \begin{minipage}[b]{0.65\columnwidth}\raggedright
Description\strut
\end{minipage}\tabularnewline
\midrule
\endhead
\begin{minipage}[t]{0.04\columnwidth}\raggedright
1\strut
\end{minipage} & \begin{minipage}[t]{0.23\columnwidth}\raggedright
Entry\strut
\end{minipage} & \begin{minipage}[t]{0.65\columnwidth}\raggedright
Lists the entry number of each operating mode; entry number 1 is
assigned to the most recent operating mode (at the top of the
list)\strut
\end{minipage}\tabularnewline
\begin{minipage}[t]{0.04\columnwidth}\raggedright
2\strut
\end{minipage} & \begin{minipage}[t]{0.23\columnwidth}\raggedright
Operation (mode)\strut
\end{minipage} & \begin{minipage}[t]{0.65\columnwidth}\raggedright
Lists the type of operating mode: standby, constant, program, remote
access, etc.\strut
\end{minipage}\tabularnewline
\begin{minipage}[t]{0.04\columnwidth}\raggedright
3\strut
\end{minipage} & \begin{minipage}[t]{0.23\columnwidth}\raggedright
Details\strut
\end{minipage} & \begin{minipage}[t]{0.65\columnwidth}\raggedright
Lists a brief description of each operating mode; detail is listed only
for alarm and program mode\strut
\end{minipage}\tabularnewline
\begin{minipage}[t]{0.04\columnwidth}\raggedright
4\strut
\end{minipage} & \begin{minipage}[t]{0.23\columnwidth}\raggedright
Date/Time\strut
\end{minipage} & \begin{minipage}[t]{0.65\columnwidth}\raggedright
Lists date and time when the operating mode or event took place\strut
\end{minipage}\tabularnewline
\begin{minipage}[t]{0.04\columnwidth}\raggedright
5\strut
\end{minipage} & \begin{minipage}[t]{0.23\columnwidth}\raggedright
Trend Graph\strut
\end{minipage} & \begin{minipage}[t]{0.65\columnwidth}\raggedright
Provides an access link to view the trend graph when the operating mode
(or event) occurred\strut
\end{minipage}\tabularnewline
\bottomrule
\end{longtable}

The following procedure outlines the steps to open a trend graph to view
details of a tripped alarm, entry number 135 in the table, that occurred
during the chamber constant operation. The alarm (entry 135) is selected
for this example to coincide with that in Section 6.4.1, since this
incident may reflect real event during a production run.

\textbf{\emph{Procedure:}} Example

\begin{enumerate}
\def\labelenumi{\arabic{enumi}.}
\item
  Locate entry 135 (first column) for the \textbf{Alarm}, as depicted in
  the figure below. Press the trend graph icon (see arrow).

  \begin{figure}
  \centering
  
\includegraphics[width=0.95\textwidth]{figures/2676937678-S5_4_2-step1-002a.PNG}
  \caption{Viewing trend graph and details of the selected alarm}
  \end{figure}\FloatBarrier
\item
  The timeline in the trend graph specifically covers the alarm mode
  that occurred at 3:41 PM. According to the plot, alarm was tripped due
  to overheating that occurred when the Temp Air process value (PV)
  reached the SP value. Data points of this graph may be downloaded (see
  arrow 2) for analysis. Refer to Section 4.9.7 for further details.

  \begin{figure}
  \centering
  
\includegraphics[width=0.95\textwidth]{figures/3151710646-S5_4_2-step2-002ab.PNG}
  \caption{Details of the selected alarm}
  \end{figure}\FloatBarrier
\item
  Press the back button (see arrow 3 in above figure) to return to the
  operating mode table.
\end{enumerate}


\subsection{Operation Statistics: Donut
Chart}

Percentage of each operation mode according to the logged data is
presented as a donut chart to provide a quick overview of the overall
operating statistics of the chamber, as depicted in the following
figure.

\begin{figure}
\centering

\includegraphics[width=0.95\textwidth]{figures/3995959319-S5_4-op-stats-001.PNG}
\caption{Donut chart display of the overall history of the chamber}
\end{figure}\FloatBarrier


\section{Set Protection (User
Settings)}

The \textbf{Set Protection} page allows the operator to protect the
log-in process on the HMI. It has the On/Off option for the user to set
HMI automatic login. By default, the HMI automatic login is On. To
protect the user HMI login, set it to Off.


\section{ROM/Firmware
Information}

The GL controller system is quite robust. It is designed with a
dual-partition operating system (i.e., firmware) to complement one
another. Should one firmware fail, the other one can be booted to
operate in place of the failed one. Technically, the system keeps two
versions of the firmware: (1) Current version of the operating firmware,
(2) Previous stable version. Initially, the GL controller system is
shipped with identical firmware on both partitions. These two partitions
eventually hold different versions of the firmware as updates have been
applied. In this process, one partition is active and the other is
passive.

Firmware update with bug fixes or upgrade with new features can be
performed in one of two approaches:

\begin{enumerate}
\def\labelenumi{\arabic{enumi}.}
\tightlist
\item
  Automatic firmware update,
\item
  Offline (manual) firmware update.
\end{enumerate}

It incorporates a mechanism to protect its stable operation during a
firmware update.

After a successful update, the GL controller system switches to run on
the new firmware. All configuration and data log files are brought over
to operate on the new firmware. If firmware update fails, the GL
controller system returns to run on its previously known stable state
version until the update issue is resolved. This process ensures that
the GL controller system (and chamber) is always up and running without
disruptions.

The \textbf{ROM/Firmware Information} page consists of three subpages:
(1) Installed Version, (2) Over-the-Air Update, and (3) Manual Update.
They are described in the following sections.


\subsection{Installed Version}

The \textbf{Installed Version} page displays information of the GL
controller system that includes name of the package, operating system ID
and control version, front-end and back-end versions.

\begin{figure}
\centering

\includegraphics[width=0.95\textwidth]{figures/2028707115-S5_6_1-subpage1a.PNG}
\caption{Installed version page}
\end{figure}\FloatBarrier

Information on this page can be used to check against the version
released by ESPEC.


\subsection{Over-the-Air Update}

The \textbf{Over-the-Air Update} page displays information about
automatic update and remote management service. The process provides
remote service support with risk tolerant and efficient over-the-air
update via a cloud service provider (called Mender IO) through a unique
registration number linked to each GL controller system (and chamber).

\begin{figure}
\centering

\includegraphics[width=0.95\textwidth]{figures/1079910449-S5_6_2-subpage2a.PNG}
\caption{Over-the-Air Update setup page}
\end{figure}\FloatBarrier

Automatic firmware update is enabled by default with the check mark in
the box (depicted in the figure). Over-the-Air update can only happen if
the target GL controller system has access to the Internet. This update
option can be disabled by the customer by unchecking the box. If a new
update package is available, update process will commence when the
target chamber is in standby mode.

The Remote Management Service is a feature that allows ESPEC software
engineering team to access and troubleshoot the GL controller remotely
through Mender. It is disabled by default indicated by the red triangle
status icon, depicted in the figure.

Customer and their IT department have complete control and freedom to
dictate the use of Mender cloud service and Remote Management Service.
Should they decide to use the service, they must proceed as follows:

\begin{enumerate}
\def\labelenumi{\arabic{enumi}.}
\tightlist
\item
  Set GL controller system with access to the Internet
\item
  Start the GL controller system
\item
  Access the \textbf{Firmware} submenu under \textbf{Settings}
\item
  Enable the Remote Management Service (check the ``Enabled'' box)
\item
  Save the setting via the \textbf{Save} button (in the upper-right
  corner)
\item
  Reboot the GL controller system via the \textbf{Accessory} page (refer
  to 6.11 for details)
\end{enumerate}


\subsection{Manual Update via
HMI}

Manual firmware update can be performed right on the chamber via its
HMI. This option is made possible for the GL controller system that
operates as a stand-alone system, isolated and not part of a network,
where automatic Over-the-Air update cannot be completed over the
Internet.

\textbf{\emph{Procedure:}}

\begin{enumerate}
\def\labelenumi{\arabic{enumi}.}
\item
  Obtain Firmware Update Package from ESPEC and put it on a USB device
\item
  Set the chamber in Standby mode
\item
  Plug the USB device (from step 1) into the USB port as shown:

  \begin{figure}
  \centering
  
\includegraphics[width=0.95\textwidth]{figures/1703896270-S5_4_1N-external-ports-001a.PNG}
  \caption{External ports}
  \end{figure}\FloatBarrier
\item
  Press \textbf{Setup} (see arrow 1) followed by \textbf{ROM/Firmware
  Information}; press \textbf{Manual Update} subpage (2), then press
  \textbf{Browse} (3) to select the update package from the list on the
  USB device and press \textbf{Accept}.

  \begin{figure}
  \centering
  
\includegraphics[width=0.95\textwidth]{figures/2491955005-S5_6_3-ROM-info-p3-002ab.PNG}
  \caption{Apply firmware update}
  \end{figure}\FloatBarrier
\item
  The system begins to upload the update file which will take several
  minutes.
\item
  When update is complete, reboot the system with the \textbf{REBOOT}
  button.
\item
  The GL controller system has been updated. Its new firmware version
  can be found under the \textbf{Installed Version} page.
\end{enumerate}


\subsection{Manual Update via Local
Network}

Manual firmware update via a local network is another alternative
option. Instead of updating the firmware on the HMI, the process is done
remotely from a PC on the network.

\textbf{\emph{Procedure:}}

\begin{enumerate}
\def\labelenumi{\arabic{enumi}.}
\item
  Obtain Firmware Update Package from ESPEC and put it on the PC (or
  USB).
\item
  Launch a Web browser on the PC. Access the GL controller system via
  its hostname or its IP address. Refer to Chapter 8 for details on the
  URL format.
\item
  Set the chamber in standby mode.
\item
  Access the \textbf{ROM/Firmware Information} page.

  \begin{figure}
  \centering
  
\includegraphics[width=0.95\textwidth]{figures/2349967590-S5_6_4-ROM-info-p3-003a.PNG}
  \caption{Drag and drop the update package}
  \end{figure}\FloatBarrier
\item
  Click \textbf{Browse} (see arrow). Navigate to the folder that holds
  the firmware update package. Drag-and-drop the package in the white
  box (indicated by the arrow), or click the white box and browse
  through a folder for the package on the local computer.

  \begin{figure}
  \centering
  
\includegraphics[width=0.95\textwidth]{figures/3439596792-offline-local-network-firmware-update-001a.PNG}
  \caption{Manual update configuration}
  \end{figure}\FloatBarrier
\item
  Once an update package has been loaded, the system begins to install
  the new firmware, as shown in the following figure.

  \begin{figure}
  \centering
  
\includegraphics[width=0.95\textwidth]{figures/560545122-firmware-update-manual-003.PNG}
  \caption{System moves to perform firmware update}
  \end{figure}\FloatBarrier
\item
  Click \textbf{REBOOT}, as indicated in the figure, to restart the
  system.

  \begin{figure}
  \centering
  
\includegraphics[width=0.95\textwidth]{figures/1919695065-firmware-update-manual-004.PNG}
  \caption{Firmware update is complete}
  \end{figure}\FloatBarrier
\end{enumerate}


\subsection{Rollback Firmware}

The GL controller can rollback its firmware to the previous stable state
via the \textbf{Rollback} button (depicted in the figure). Rollback can
be performed at any time.

\begin{figure}
\centering

\includegraphics[width=0.95\textwidth]{figures/4148802009-S5_6_5-rollback-001a.PNG}
\caption{Rollback a firmware update}
\end{figure}\FloatBarrier


\section{Set Defrost}

Frost may form on the cooler in temperature operations below 40 °C.
Defrost the chamber in the following cases.

\begin{itemize}
\tightlist
\item
  If temperature inside the chamber is uncontrollable or rises slowly.
\item
  If air blowing from the chamber is weak (when the door is opened).
\item
  If frost or ice forms on test area wall.
\end{itemize}

\begin{figure}
\centering

\includegraphics[width=0.95\textwidth]{figures/3827918121-S8_4_3-defrost-0001b.PNG}
\caption{Operate Defrost mode}
\end{figure}\FloatBarrier

There is no automatic defrost operation; hence, this option is greyed
out (see above figure). Defrost must be operated manually. The following
procedure outlines the general steps to conduct defrost on GL chambers.
Refer to your specific chamber manual for instructions on manual defrost
operation.

\textbf{\emph{Procedure:}}

\begin{enumerate}
\def\labelenumi{\arabic{enumi}.}
\item
  Check the main power switch is in the ON position. If the GL
  controller system is already on and its touchscreen is in sleep mode,
  tab your finger on the touchscreen and proceed to Step 2. Otherwise,
  set the main power breaker in the ON position. Within a minute, the
  touch-screen display will come on. The system automatically logs in
  using the default user account (called localhost) to operate the
  chamber.

  \begin{figure}
  \centering
  
\includegraphics[width=0.95\textwidth]{figures/3959919147-S8_3_8-step1-001.PNG}
  \caption{Operating Panel at Startup}
  \end{figure}\FloatBarrier
\item
  Press the \textbf{Hamburger} menu (1) in the status bar. Check to
  confirm the floating keyboard is enabled. A diagonal bar over the
  keyboard indicates it is disabled. Press the keyboard (2) to enable
  it. Press \textbf{Configuration} (3); then press \textbf{Operation
  Process} (4).

  \begin{figure}
  \centering
  
\includegraphics[width=0.95\textwidth]{figures/3010205281-S8_3_7-step2-0001a.PNG}
  \caption{Accessing Operation Process configuration page}
  \end{figure}\FloatBarrier
\item
  Confirm that the following settings (see arrows) are configured such
  that the operation is not interrupted or an alarm generated if the
  chamber is run with the door cracked slightly.

  \begin{figure}
  \centering
  
\includegraphics[width=0.95\textwidth]{figures/3269330728-S8_3_7_step2-002a.PNG}
  \caption{Disable warning or alarm on door conditions}
  \end{figure}\FloatBarrier
\item
  Select one Constant mode (say, Constant 1); make notes of its original
  settings. Set target temperature to a minimum 70 degrees and set
  humidity and refrigeration control OFF. (It is unnecessary with
  temperature-only chambers.)
  
\includegraphics[width=0.95\textwidth]{figures/307414286-S8_4_3-Step4-TH-off-Const1a.PNG}
\item
  Start Constant 1 for about 60 minutes with the door closed.
\item
  Start Constant 1 again for 15 minutes with the door slightly cracked.
\item
  Return settings to their original configuration if modification was
  made in Step 3 (using the recorded notes).
\end{enumerate}

\textbf{Power OFF}: Set the main power breaker in the OFF position.


\section{Set Time Meters}

Running time meters (i.e., accumulation of operation time) of the vital
hardware devices inside the chamber indicate how long each device has
been in operation. The high-stage and low-stage compressor, as well as
the humidity system, are part of the list. These components are operated
by motors; and motors with moving parts are subject to wear. Running
time meters are recorded as Total Elapsed time and Time Since Service;
they provide information for recommended maintenance when each device
has reached its critical service and total runtime.

\begin{figure}
\centering

\includegraphics[width=0.95\textwidth]{figures/854185372-S5_8-time-meter-page-001a.PNG}
\caption{Device total run time}
\end{figure}\FloatBarrier

\begin{longtable}[]{@{}lll@{}}
\toprule
\begin{minipage}[b]{0.04\columnwidth}\raggedright
No.\strut
\end{minipage} & \begin{minipage}[b]{0.23\columnwidth}\raggedright
Type/Name\strut
\end{minipage} & \begin{minipage}[b]{0.65\columnwidth}\raggedright
Description\strut
\end{minipage}\tabularnewline
\midrule
\endhead
\begin{minipage}[t]{0.04\columnwidth}\raggedright
1\strut
\end{minipage} & \begin{minipage}[t]{0.23\columnwidth}\raggedright
Device Index\strut
\end{minipage} & \begin{minipage}[t]{0.65\columnwidth}\raggedright
Lists all vital devices in operation by index\strut
\end{minipage}\tabularnewline
\begin{minipage}[t]{0.04\columnwidth}\raggedright
2\strut
\end{minipage} & \begin{minipage}[t]{0.23\columnwidth}\raggedright
Device Name\strut
\end{minipage} & \begin{minipage}[t]{0.65\columnwidth}\raggedright
Lists all vital devices in operation by name\strut
\end{minipage}\tabularnewline
\begin{minipage}[t]{0.04\columnwidth}\raggedright
3\strut
\end{minipage} & \begin{minipage}[t]{0.23\columnwidth}\raggedright
Total Elapsed Time\strut
\end{minipage} & \begin{minipage}[t]{0.65\columnwidth}\raggedright
Lists total runtime of each device in operation\strut
\end{minipage}\tabularnewline
\begin{minipage}[t]{0.04\columnwidth}\raggedright
4\strut
\end{minipage} & \begin{minipage}[t]{0.23\columnwidth}\raggedright
Total Since Service\strut
\end{minipage} & \begin{minipage}[t]{0.65\columnwidth}\raggedright
Lists total runtime since previous service maintenance\strut
\end{minipage}\tabularnewline
\begin{minipage}[t]{0.04\columnwidth}\raggedright
5\strut
\end{minipage} & \begin{minipage}[t]{0.23\columnwidth}\raggedright
Reset Button\strut
\end{minipage} & \begin{minipage}[t]{0.65\columnwidth}\raggedright
Lists reset button for each listed device\strut
\end{minipage}\tabularnewline
\bottomrule
\end{longtable}

If a component has been serviced (or replaced), its record of operation
time should be reset. The following procedure outlines the steps to
reset the meter on the Humidity System (item 4 in the previous figure).

\textbf{\emph{Procedure:}}

\begin{enumerate}
\def\labelenumi{\arabic{enumi}.}
\item
  Press \textbf{Setup} in the (left) menu bar (see arrow 1). Press
  \textbf{Set Time Meters} (arrow 2).

  \begin{figure}
  \centering
  
\includegraphics[width=0.95\textwidth]{figures/2310940074-S5_8-step1-001bba.PNG}
  \caption{Steps to accessing Set Timer Meters from the Menu bar}
  \end{figure}\FloatBarrier
\item
  Press the reset button of humidity system (arrow 1). Press Yes (arrow
  2).

  \begin{figure}
  \centering
  
\includegraphics[width=0.95\textwidth]{figures/1216486535-S5_8-step2-000a.PNG}
  \caption{Resetting time meters of humidity system}
  \end{figure}\FloatBarrier
\item
  The humidity system time meter has been reset.

  \begin{figure}
  \centering
  
\includegraphics[width=0.95\textwidth]{figures/335941189-S5_8-step3-001a.PNG}
  \caption{Time meter of humidity system has bee nreset}
  \end{figure}\FloatBarrier
\end{enumerate}


\section{Set Back Trace}

The GL controller automatically records back trace data during
operation. The data contain set points and process values of
temperature/humidity and control value information of other components
required to control the GL controller instrumentation. If an alarm
occurs during operation, the GL controller completes the recording of
the back trace data. The data can then be downloaded and sent over to
Espec diagnostic service for analysis to determine the cause of
equipment failure (if any) which lead to an alarm. The customer will be
provided with the diagnostic results of the equipment before and after
the occurrence of an alarm.

Up to ten back trace data are kept on the list at all time. The first
back trace data drops from the list to allow the next back trace data to
take the spot; and this process repeats indefinitely in a first-in
first-out queue. Back trace data can be manually generated via the
\textbf{Manual Trigger} button (item 6, below).

\begin{figure}
\centering

\includegraphics[width=0.95\textwidth]{figures/2707886627-S5_9-back-trace-data-000a.PNG}
\caption{Back trace page listing back trace data for each alarm
occurrence}
\end{figure}\FloatBarrier

\begin{longtable}[]{@{}lll@{}}
\toprule
\begin{minipage}[b]{0.04\columnwidth}\raggedright
No.\strut
\end{minipage} & \begin{minipage}[b]{0.23\columnwidth}\raggedright
Type/Name\strut
\end{minipage} & \begin{minipage}[b]{0.65\columnwidth}\raggedright
Description\strut
\end{minipage}\tabularnewline
\midrule
\endhead
\begin{minipage}[t]{0.04\columnwidth}\raggedright
1\strut
\end{minipage} & \begin{minipage}[t]{0.23\columnwidth}\raggedright
Alarm\strut
\end{minipage} & \begin{minipage}[t]{0.65\columnwidth}\raggedright
Alarm notification\strut
\end{minipage}\tabularnewline
\begin{minipage}[t]{0.04\columnwidth}\raggedright
2\strut
\end{minipage} & \begin{minipage}[t]{0.23\columnwidth}\raggedright
Alarm ID\strut
\end{minipage} & \begin{minipage}[t]{0.65\columnwidth}\raggedright
Lists ID number of alarm\strut
\end{minipage}\tabularnewline
\begin{minipage}[t]{0.04\columnwidth}\raggedright
3\strut
\end{minipage} & \begin{minipage}[t]{0.23\columnwidth}\raggedright
Device\strut
\end{minipage} & \begin{minipage}[t]{0.65\columnwidth}\raggedright
Lists device or component that tripped an alarm\strut
\end{minipage}\tabularnewline
\begin{minipage}[t]{0.04\columnwidth}\raggedright
4\strut
\end{minipage} & \begin{minipage}[t]{0.23\columnwidth}\raggedright
Date/Time\strut
\end{minipage} & \begin{minipage}[t]{0.65\columnwidth}\raggedright
Lists date and time alarm occurred\strut
\end{minipage}\tabularnewline
\begin{minipage}[t]{0.04\columnwidth}\raggedright
5\strut
\end{minipage} & \begin{minipage}[t]{0.23\columnwidth}\raggedright
Download Button\strut
\end{minipage} & \begin{minipage}[t]{0.65\columnwidth}\raggedright
Lists button to download back trace data\strut
\end{minipage}\tabularnewline
\begin{minipage}[t]{0.04\columnwidth}\raggedright
6\strut
\end{minipage} & \begin{minipage}[t]{0.23\columnwidth}\raggedright
Manual Trigger\strut
\end{minipage} & \begin{minipage}[t]{0.65\columnwidth}\raggedright
This button allows a manual generation of the back trace data\strut
\end{minipage}\tabularnewline
\bottomrule
\end{longtable}


\subsection{Download Back Trace
Data}

The following procedure outlines the steps to access and download the
back trace data associated with an overheating alarm ID 003-3000
(depicted in the previous figure) as an example. The data will be
exported out to a USB device connected to the GL controller system.

\textbf{\emph{Procedure:}}

\begin{enumerate}
\def\labelenumi{\arabic{enumi}.}
\item
  If the GL controller system is already on and its touchscreen is in
  sleep mode, tab your finger on the touchscreen and proceed to Step 2
  (below). From a cold start, set the main power breaker in the ON
  position. Within a minute, the touch-screen display will come on. The
  system automatically logs in using the default localhost user account
  to operate the chamber.
\item
  Plug in a USB device into the USB port (see arrow) at the front of the
  chamber.

  \begin{figure}
  \centering
  
\includegraphics[width=0.95\textwidth]{figures/1703896270-S5_4_1N-external-ports-001a.PNG}
  \caption{External ports}
  \end{figure}\FloatBarrier
\item
  Press \textbf{Setup} in the menu bar (see arrow 1). Press the
  \textbf{Set Back Trace} tab (see arrow 2).

  \begin{figure}
  \centering
  
\includegraphics[width=0.95\textwidth]{figures/3438037417-S5_9-step1-001a.PNG}
  \caption{Accessing the Back Trace page}
  \end{figure}\FloatBarrier
\item
  Locate alarm ID 003-3000 (see arrow 1). Press the export button (see
  arrow 2) to export the data file out to the USB device.

  \begin{figure}
  \centering
  
\includegraphics[width=0.95\textwidth]{figures/1226316684-S5_9_1-step3-000ab.PNG}
  \caption{Exporting back trace data for ALM 003-3000}
  \end{figure}\FloatBarrier
\item
  Back trace data will be stored (automatically) in the external USB
  device as MS Excel format with extension \textbf{.xlsx} with with
  date/time stamp associated with the alarm.
\end{enumerate}


\subsection{Create Back Trace
Data}

The back trace data can be created manually in case a diagnostic
analysis is needed.

\textbf{\emph{Procedure:}}

\begin{enumerate}
\def\labelenumi{\arabic{enumi}.}
\item
  Press the \textbf{Manual Trigger} button (see arrow 6).

  \begin{figure}
  \centering
  
\includegraphics[width=0.95\textwidth]{figures/873379190-S5_9_2-step1-000a.PNG}
  \caption{Manually creating back trace data}
  \end{figure}\FloatBarrier
\item
  The GL controller reads the most recent back trace log and generates
  the data which may take between 5 to 8 minutes to complete, as
  indicated by the notification message and the rendering arrow.

  \begin{figure}
  \centering
  
\includegraphics[width=0.95\textwidth]{figures/694942249-S5_9_2-step2-002a.PNG}
  \caption{Rendering back trace data}
  \end{figure}\FloatBarrier
\item
  Once data generation is complete, the data file is populated at the
  bottom of the list. To fetch the data, press the export button (see
  arrow).

  \begin{figure}
  \centering
  
\includegraphics[width=0.95\textwidth]{figures/2413965673-S5_9_2-step3-001a.PNG}
  \caption{Back trace data from manual rendering}
  \end{figure}\FloatBarrier
\item
  Data file will be exported out to the attached USB device. If this
  procedure is performed on a PC accessing the GL controller system
  remotely on a network, the data file will be downloaded and stored in
  the Downloads folder on the PC. Filename of the data is associated
  with the current date and the key word ``manual'' with extension
  \textbf{.xlsx}. The data can be open and viewed on MS Excel.
\end{enumerate}


\section{Data Download (Log Data
Backup)}

Contents in the trend graph discussed in the \textbf{Trend} menu are
dictated by the settings on this page. Data points collected from the
chamber will bear a date/time stamp of EWC, shown in the upper-left in
the system bar, regardless of the date/time on the controller (see the
extension bar of the status tab in \textbf{Overview}). For practical
purposes, date/time of EWC should reflect the current date and local
time zone.

Two panes for data logging are available to support the configuration
(as depicted in the following figure). They are described as follows:

\begin{figure}
\centering

\includegraphics[width=0.95\textwidth]{figures/3542610259-new-data-log-001a.PNG}
\caption{Options for data logging}
\end{figure}\FloatBarrier

\begin{enumerate}
\def\labelenumi{\arabic{enumi}.}
\item
  \textbf{Data Logging Configuration/Settings}: Data logging is
  controlled by the logging frequency and chamber operating mode. By
  default, logging frequency occurs at a 10-second interval. An
  arbitrary value can be set but the setting is limited to the response
  of the chamber for a frequency less than 10 seconds. Two operating
  modes control data logging: \textbf{Always} and \textbf{When Chamber
  is on}. The following section discusses in detail how to set logging
  frequency and operating mode. Data logs can be archived via a Samba
  technology that allows file sharing between two different platforms
  (such as, Linux and MS Windows or Linux and Mac). Our EWC is already
  configured to use Samba to transfer and store logged data. Samba
  configuration on the target platform is the responsibility of the
  user.
\item
  \textbf{Trend Graph Views}: Data points and types (such as,
  temperature, humidity, etc.) can be manually selected to appear in the
  trend graph. Selection can be made based on the check mark placed in
  their respective boxes (shown in the following figure). Color
  decoration can be selected for plot legends. Different logs for a
  trend graph can be created. Detailed steps for these settings are
  discussed in the following sections.

  \begin{figure}
  \centering
  
\includegraphics[width=0.95\textwidth]{figures/769489772-trend-graph-views-001.PNG}
  \caption{Data type and color options}
  \end{figure}\FloatBarrier
\end{enumerate}


\subsection{Setting Logging Interval and Operating
Mode}

To set a new data logging frequency and operating mode, complete the
following steps.

\textbf{\emph{Procedure:}}

\begin{enumerate}
\def\labelenumi{\arabic{enumi}.}
\item
  Log in as admin (or into an account with access privilege to the
  \textbf{Settings} menu).
\item
  Access \textbf{Data Logging Settings} page from the \textbf{Settings}
  menu
\item
  Enter a new frequency value in the \textbf{Logging Period} field
  (indicated by the arrow) or click the spin button next to the
  \textbf{Seconds} unit to set a new value.
\item
  Click and select \textbf{Always} or \textbf{When Chamber is on} as
  shown by the arrow. Default setting is \textbf{Always}.

  \begin{figure}
  \centering
  
\includegraphics[width=0.95\textwidth]{figures/142093444-data-logging-interval-mode-001a.PNG}
  \caption{Setting options for data logging}
  \end{figure}\FloatBarrier

  \textbf{Note}: If the operating mode \textbf{When Chamber is on} is
  selected, data logging only occurs when the chamber is on. This
  setting can dramatically change the behavior of the trend graph in the
  \textbf{Trend} menu.
\item
  Click the \textbf{Save} button (upper-right corner) to save the
  settings.
\item
  To delete current setting, click the trash bin icon.
\end{enumerate}


\subsection{Data Backup via
Samba}

To set data logging backup using Samba, you must first configure or
create a Samba share on your computer or server. Then complete the
following steps:

\begin{enumerate}
\def\labelenumi{\arabic{enumi}.}
\item
  Turn on Samba with the check mark in the box to enable the feature.
\item
  Set the backup frequency via the drop-down. Three options are
  available: (1) Once per day, (2) Once per hour, and (3) every 5
  minutes.
\item
  Enter the correct path to the shared folder on your server with the
  complete credentials (username and password).

  \begin{figure}
  \centering
  
\includegraphics[width=0.95\textwidth]{figures/631357192-data-archiving-samba-001.PNG}
  \caption{Setting up Samba for data backup}
  \end{figure}\FloatBarrier
\item
  Click the \textbf{Save} button (upper-right corner) to save the
  settings.
\end{enumerate}


\subsection{Trend Graph View
Options}

Trend graph contents can be controlled via several options that include
the type of data, color for legend and plot range. These settings can be
stored in different view options with unique view names to represent
different trend graphs.

The following steps outline the procedure to create a new log view
(called \textbf{Log1}) for a trend graph to include data type (Temp and
Vibration) with set and process values, chamber status, program name,
step number, refrigeration and set the appropriate y-range for display.
Only five items to be enabled and included in the trend graph.

\textbf{\emph{Procedure:}}

\begin{enumerate}
\def\labelenumi{\arabic{enumi}.}
\item
  Click \textbf{Views} and select \textbf{Create New View} (shown by the
  arrow)
\item
  Enter \textbf{Log1} in the \textbf{New view name} field
\item
  Click the check mark on the right to accept the name.

  \begin{figure}
  \centering
  
\includegraphics[width=0.95\textwidth]{figures/724171531-data-logging-new-view-name-001a.PNG}
  \caption{Creating a new log view for Trend graph}
  \end{figure}\FloatBarrier
\item
  Set the min and max values for the y-range to -25 and 65,
  respectively.
\item
  Enable Status, Program Step, Program Name, Temperature set and process
  values, vibration set and process values, and refrigeration with the
  check marks in their boxes. Scroll down the page to enable all of them
  (as shown in the figure). Color picker, line width and line type for
  the selected data will be based on the default setting.

  \begin{figure}
  \centering
  
\includegraphics[width=0.95\textwidth]{figures/2858589674-new-view-select-items-001b.PNG}
  \caption{Selection of data types for a trend graph}
  \end{figure}\FloatBarrier

  \textbf{Note}: View name can be edited or removed using the trash bin
  or pen icons to the right of \textbf{Log1}.
\item
  Click the color picker to select a desired color for each data type.
\item
  Click the \textbf{Save} button at the upper-right corner.
\item
  Access \textbf{Trend} on the menu bar.
\item
  Click the \textbf{Graph View} icon (see figure)

  \begin{figure}
  \centering
  
\includegraphics[width=0.95\textwidth]{figures/2480407825-pick-trend-view-option-001.PNG}
  \caption{Selecting trend view option}
  \end{figure}\FloatBarrier
\item
  Click \textbf{Customize} then select \textbf{Log1}.

  \begin{figure}
  \centering
  
\includegraphics[width=0.95\textwidth]{figures/777525324-pick-trend-view-option-002.PNG}
  \caption{Select custom view Log1}
  \end{figure}\FloatBarrier
\item
  Trend graph now displays the custom view \textbf{Log1}.

  \begin{figure}
  \centering
  
\includegraphics[width=0.95\textwidth]{figures/369609872-pick-trend-view-option-003.PNG}
  \caption{Trend graph view of Log1}
  \end{figure}\FloatBarrier
\end{enumerate}


\subsection{Clear Data Log}

It may be necessary to start a new data log. Click the \textbf{Trash
bin} to clear the current log. A warning dialog box appears to reconfirm
the action with a Yes/No option to proceed.

\begin{figure}
\centering

\includegraphics[width=0.95\textwidth]{figures/456497475-clear-data-log-001.PNG}
\caption{Confirm to delete current data log}
\end{figure}\FloatBarrier

The structure of data log will be deleted from the data file immediately
after confirmation. EWC moves to prepare a new data log in the file with
another pop-up warning (depicted in the following figure).

\begin{figure}
\centering

\includegraphics[width=0.95\textwidth]{figures/1525553755-trash-bin-of-data-001b.PNG}
\caption{Deleting data from log file}
\end{figure}\FloatBarrier

Once data begins to accumulate, a trend graph can be produced. Refer to
the \textbf{Trend} menu in the menu bar for detail on the trend graph.


\section{Accessory \& GL Controller Power
Management}

The \textbf{Accessory} page allows the operator with two control
options: (1) HMI screensaver and (2) GL controller power management. The
HMI screensaver option is available (and visible) only when operated on
the HMI. The following figure depicts both the HMI screensaver and power
management options. The hour meter information is posted at the top of
this page to indicate how long the instrumentation has been in
operation.

\begin{figure}
\centering

\includegraphics[width=0.95\textwidth]{figures/1616749110-S5_11-accessory-001ab.PNG}
\caption{Accessory page with HMI screensaver and power management
options}
\end{figure}\FloatBarrier

\begin{longtable}[]{@{}lll@{}}
\toprule
No. & Name & Description / Function\tabularnewline
\midrule
\endhead
1 & Hour Meter & Displays hour meter information of the
chamber\tabularnewline
2 & HMI Screensaver & Timeout screensaver; only available on HMI
operation\tabularnewline
3 & Power Management & GL controller power management
options\tabularnewline
\bottomrule
\end{longtable}


\subsection{HMI Screensaver}

The GL controller system has a screensaver feature to control the HMI
display with timeout options in 10 minutes, 30 minutes and 1 hour, as
depicted in the above figure. With a 10-min timeout setting, the HMI
display goes to sleep if there are no activities on the HMI screen for
10 minutes. The GL controller system continues in its operating mode,
while the HMI is in sleep mode. To turn the HMI display back on, simply
tab the finger anywhere on the HMI screen once and wait for about four
seconds for it to illuminate.


\subsection{Restart Services or Reboot GL
Controller}

Sometimes, certain services of the GL controller system may stop working
due to a momentary disruption, such as a network was unplugged; and that
service refuses to restart automatically. The \textbf{Restart Services}
button can be used to manually restart those services.

Sometimes, a certain setting may require rebooting the GL controller
system. For example, if a firmware update is performed manually via the
HMI screen, the GL controller requires rebooting via the \textbf{Reboot
Controller} button. This action will reboot the operating system of the
GL controller, while the chamber and its components remain powered on.

During a reboot process, the operating system starts all processes to
reconnect communication and verify the configuration on the GL
controller and the chamber operation. Within in a minute, the HMI
display comes on. The GL controller system automatically logs the
localhost user in for operation.


\section{Display Customization: Home and Submonitor
Pages}

The display and layout of Home and Submonitor Page 2 and Page 3 are
customizable by manipulating the elements of the UI, called
\emph{widgets}, via the edit button in the secondary menu. With this
feature, the user can personalize the widgets on these pages to
streamline the operation and functionality, such as creating a shortcut.
The entire layout of the Home and Submonitor Page 2 is customizable;
only a few elements or widgets on Submonitor Page 3 are customizable.

\textbf{Home:} The following figure depicts the home page in edit mode.
All customizable widgets are indicated by the edit button (see arrow).
\textbf{\emph{Note:}} Submonitor Page 2 and 3 are similar.

\begin{figure}
\centering

\includegraphics[width=0.95\textwidth]{figures/251006935-S6_12-home-003a.PNG}
\caption{Home page in edit mode}
\end{figure}\FloatBarrier

\begin{longtable}[]{@{}lll@{}}
\toprule
\begin{minipage}[b]{0.04\columnwidth}\raggedright
No.\strut
\end{minipage} & \begin{minipage}[b]{0.23\columnwidth}\raggedright
Name\strut
\end{minipage} & \begin{minipage}[b]{0.65\columnwidth}\raggedright
Operation\strut
\end{minipage}\tabularnewline
\midrule
\endhead
\begin{minipage}[t]{0.04\columnwidth}\raggedright
1\strut
\end{minipage} & \begin{minipage}[t]{0.23\columnwidth}\raggedright
Add/Create\strut
\end{minipage} & \begin{minipage}[t]{0.65\columnwidth}\raggedright
Create a new widget; a new widget is placed at the bottom of the current
page.\strut
\end{minipage}\tabularnewline
\begin{minipage}[t]{0.04\columnwidth}\raggedright
2\strut
\end{minipage} & \begin{minipage}[t]{0.23\columnwidth}\raggedright
Save\strut
\end{minipage} & \begin{minipage}[t]{0.65\columnwidth}\raggedright
Apply the changes to the current page\strut
\end{minipage}\tabularnewline
\begin{minipage}[t]{0.04\columnwidth}\raggedright
3\strut
\end{minipage} & \begin{minipage}[t]{0.23\columnwidth}\raggedright
Back\strut
\end{minipage} & \begin{minipage}[t]{0.65\columnwidth}\raggedright
Cancel; press this \textbf{Back} button, followed by Yes, to cancel the
current settings and return to default setting.\strut
\end{minipage}\tabularnewline
\bottomrule
\end{longtable}

A widget selection window appears when any widget button is activated
(see arrows), as depicted in the following figure.

\begin{figure}
\centering

\includegraphics[width=0.95\textwidth]{figures/2901153180-S6_12-home-003bba.PNG}
\caption{S6\_12-home-003bba.PNG}
\end{figure}\FloatBarrier


\subsection{Widget Page and Selection
Options}

Four separate menus (\textbf{Links}, \textbf{Monitor},
\textbf{Start/Stop} and \textbf{Color}) of the widget page provide
different options to customize the functionality of Home, Submonitor
Page 2 and Page 3.

\begin{figure}
\centering

\includegraphics[width=0.95\textwidth]{figures/3176247319-S6_12-widget-selection-page-001a.PNG}
\caption{Widget selection and its menus}
\end{figure}\FloatBarrier

\begin{longtable}[]{@{}lll@{}}
\toprule
\begin{minipage}[b]{0.04\columnwidth}\raggedright
No.\strut
\end{minipage} & \begin{minipage}[b]{0.23\columnwidth}\raggedright
Name\strut
\end{minipage} & \begin{minipage}[b]{0.65\columnwidth}\raggedright
Description/Operation\strut
\end{minipage}\tabularnewline
\midrule
\endhead
\begin{minipage}[t]{0.04\columnwidth}\raggedright
1\strut
\end{minipage} & \begin{minipage}[t]{0.23\columnwidth}\raggedright
Links\strut
\end{minipage} & \begin{minipage}[t]{0.65\columnwidth}\raggedright
This menu provides four different categories to select and link a widget
for the shortcut: \emph{General Links}, \emph{Monitor Links},
\emph{Setup Links} and \emph{Configuration Links}.\strut
\end{minipage}\tabularnewline
\begin{minipage}[t]{0.04\columnwidth}\raggedright
2\strut
\end{minipage} & \begin{minipage}[t]{0.23\columnwidth}\raggedright
Monitor\strut
\end{minipage} & \begin{minipage}[t]{0.65\columnwidth}\raggedright
Displays and provides widget selection options related to the
\textbf{Monitor} menu, including \emph{Chamber Status Items},
\emph{Temperature}, \emph{Humidity}, \emph{Time Signals} and
\emph{General Items}.\strut
\end{minipage}\tabularnewline
\begin{minipage}[t]{0.04\columnwidth}\raggedright
3\strut
\end{minipage} & \begin{minipage}[t]{0.23\columnwidth}\raggedright
Start/Stop\strut
\end{minipage} & \begin{minipage}[t]{0.65\columnwidth}\raggedright
Displays and provides widget selection options related to the
\textbf{Operation Mode} menu, namely \emph{Constant Run} and
\emph{Program Run}.\strut
\end{minipage}\tabularnewline
\begin{minipage}[t]{0.04\columnwidth}\raggedright
4\strut
\end{minipage} & \begin{minipage}[t]{0.23\columnwidth}\raggedright
Color Options\strut
\end{minipage} & \begin{minipage}[t]{0.65\columnwidth}\raggedright
Displays a list of color options for personalization and
customization.\strut
\end{minipage}\tabularnewline
\begin{minipage}[t]{0.04\columnwidth}\raggedright
5\strut
\end{minipage} & \begin{minipage}[t]{0.23\columnwidth}\raggedright
Active Widget\strut
\end{minipage} & \begin{minipage}[t]{0.65\columnwidth}\raggedright
Displays the active widget of the menu buttons (\emph{Operation Mode},
\emph{Monitor}, \emph{Set Mode}, \emph{Setup})--the widget on which the
edit button was activated; this widget will be replaced with the
selected one.\strut
\end{minipage}\tabularnewline
\begin{minipage}[t]{0.04\columnwidth}\raggedright
6\strut
\end{minipage} & \begin{minipage}[t]{0.23\columnwidth}\raggedright
Delete\strut
\end{minipage} & \begin{minipage}[t]{0.65\columnwidth}\raggedright
Delete button; this button deletes the current widget.\strut
\end{minipage}\tabularnewline
\begin{minipage}[t]{0.04\columnwidth}\raggedright
7\strut
\end{minipage} & \begin{minipage}[t]{0.23\columnwidth}\raggedright
Close\strut
\end{minipage} & \begin{minipage}[t]{0.65\columnwidth}\raggedright
This button closes the current page (the widget page).\strut
\end{minipage}\tabularnewline
\bottomrule
\end{longtable}

\textbf{General Links:} The general links page consists of widget
selection options from four different categories as depicted in the
figure.

\begin{figure}
\centering

\includegraphics[width=0.95\textwidth]{figures/1849484488-S6_12_1-links-001.PNG}
\caption{General Link page}
\end{figure}\FloatBarrier

\textbf{Monitor Links:} The monitor links page consists of four separate
display categories for monitoring options, as depicted in the figure.

\begin{figure}
\centering

\includegraphics[width=0.95\textwidth]{figures/313774186-S6_12_1-monitor-001.PNG}
\caption{Monitor Link page}
\end{figure}\FloatBarrier

\textbf{Operation Links:} The operation links page consists of two
separate categories related to \emph{Constant Run} and \emph{Program
Run} options.

\begin{figure}
\centering

\includegraphics[width=0.95\textwidth]{figures/1969713226-S6_12_1-startstop-001.PNG}
\caption{Operation Link page}
\end{figure}\FloatBarrier

\textbf{Color Options:} The color option page consists of a list of
different color selection used for decoration or personalization, as
depicted in the figure.

\begin{figure}
\centering

\includegraphics[width=0.95\textwidth]{figures/2953613789-S6_12_1-colors-001.PNG}
\caption{Color selection options}
\end{figure}\FloatBarrier


\subsection{General Widget Editing
Procedure}

The following procedure outlines the general steps to edit the links on
the Home page. The procedure can be applied to Submonitor Page 2 or Page
3.

\textbf{\emph{Procedure:}}

\begin{enumerate}
\def\labelenumi{\arabic{enumi}.}
\item
  Decide on what shortcut or widget to edit (e.g., menu or operation
  button, or display panel).
\item
  Access the Home menu.
\item
  Press the edit button in the secondary menu (see arrow).

  \begin{figure}
  \centering
  
\includegraphics[width=0.95\textwidth]{figures/1006079423-S6_12_2-step3-001a.PNG}
  \caption{Set the Home page in edit mode}
  \end{figure}\FloatBarrier
\item
  Press the edit button of the desired widget to be edited (see arrow).
  Example: Edit the \emph{Start Prog. 1} widget.

  \begin{figure}
  \centering
  
\includegraphics[width=0.95\textwidth]{figures/4269976903-S6_12_2-step4-001a.PNG}
  \caption{Select a widget or link for editing}
  \end{figure}\FloatBarrier
\item
  Select a new widget to replace the current one. Example: Select
  \emph{Set Sampling} (see arrow). \textbf{\emph{Note:}} Program 2 or 3,
  etc., may be selected. This example illustrates that any widget can be
  used as a replacement. To remove the selected (current) widget, press
  \textbf{Delete}; to close the widget selection window, press
  \textbf{Close}.

  \begin{figure}
  \centering
  
\includegraphics[width=0.95\textwidth]{figures/4259717268-S6_12_2-step5-001aa.PNG}
  \caption{Select a new widget or delete the current one}
  \end{figure}\FloatBarrier
\item
  Press the check button, followed by Yes (in dialog box), to save the
  setting. Or, press the back button followed by Yes to cancel the
  editing task.

  \begin{figure}
  \centering
  
\includegraphics[width=0.95\textwidth]{figures/4111005989-S6_12_2-step6-001a.PNG}
  \caption{Save new settings or cancel entire operration}
  \end{figure}\FloatBarrier
\end{enumerate}


\subsection{General Widget Creating
Procedure}

The following procedure outlines the general steps to create a new link
on the Home page. The procedure can be applied to Submonitor Page 2 or
Page 3.

\textbf{\emph{Procedure:}}

\begin{enumerate}
\def\labelenumi{\arabic{enumi}.}
\item
  Decide on what shortcut or widget to create (e.g., menu or operation
  button, or display panel).
\item
  Access the Home menu.
\item
  Press the edit button in the secondary menu (see arrow).

  \begin{figure}
  \centering
  
\includegraphics[width=0.95\textwidth]{figures/1006079423-S6_12_2-step3-001a.PNG}
  \caption{Set the Home page in edit mode}
  \end{figure}\FloatBarrier
\item
  Press the plus (+) button to create a new widget (see arrow).

  \begin{figure}
  \centering
  
\includegraphics[width=0.95\textwidth]{figures/1671121483-S6_12_3-step4-001a.PNG}
  \caption{Creating a new widget}
  \end{figure}\FloatBarrier
\item
  The new widget is placed at the bottom of the display area (scroll
  down if necessary to access it).
\item
  Press (or Click) and drag the widget to a desired location.

  \begin{figure}
  \centering
  
\includegraphics[width=0.95\textwidth]{figures/2751405702-S6_12_3-step6-001a.PNG}
  \caption{Positioning the widget in the layout}
  \end{figure}\FloatBarrier
\item
  Press the edit button of the widget to create a new link (see arrow).

  \begin{figure}
  \centering
  
\includegraphics[width=0.95\textwidth]{figures/2095987255-S6_12_3-step7-001a.PNG}
  \caption{Editing the new widget}
  \end{figure}\FloatBarrier
\item
  Repeat steps 4 through 6 in Section 6.12.2.
\end{enumerate}


\chapter{Configuration Menu}

The GL Configuration menu, being part of the \textbf{Setup} menu
outlined in Chapter 6, is the administration page where different
settings and configuration options can be applied to the GL controller
and the chamber. For this reason, this page should be off limit to
regular users. User privileges can be enforced via the \textbf{Users
(Register User Password)} setup page. Refer to Section 7.4 for details.

Under this menu, the following can be configured: Set Communication,
Operation Process, Control Attainment Range, Time Signal Names, Display
Setup, Set Language, Set Sound, Set Date/Time, Users (register users and
passwords), Sensor Offset, Set Chamber Details, Set Option, Email,
Service and Macros.


\section{Set Communication}

The GL controller system supports four different communication methods:
RS485, RS232, GPIB and LAN. The LAN is the standard setting.

\begin{figure}
\centering

\includegraphics[width=0.95\textwidth]{figures/1213643087-S6_1-set-comm-003a.PNG}
\caption{The Set Communication page}
\end{figure}\FloatBarrier

\begin{longtable}[]{@{}lll@{}}
\toprule
\begin{minipage}[b]{0.04\columnwidth}\raggedright
No.\strut
\end{minipage} & \begin{minipage}[b]{0.23\columnwidth}\raggedright
Name\strut
\end{minipage} & \begin{minipage}[b]{0.65\columnwidth}\raggedright
Description\strut
\end{minipage}\tabularnewline
\midrule
\endhead
\begin{minipage}[t]{0.04\columnwidth}\raggedright
1\strut
\end{minipage} & \begin{minipage}[t]{0.23\columnwidth}\raggedright
Path indicator\strut
\end{minipage} & \begin{minipage}[t]{0.65\columnwidth}\raggedright
Path to Set Communication; it indicates the path to the Set
Communication via the Menu bar or the Hamburger button\strut
\end{minipage}\tabularnewline
\begin{minipage}[t]{0.04\columnwidth}\raggedright
2\strut
\end{minipage} & \begin{minipage}[t]{0.23\columnwidth}\raggedright
Communication Type\strut
\end{minipage} & \begin{minipage}[t]{0.65\columnwidth}\raggedright
Lists of types of supported communication: RS485, RS232, GPIB,
LAN.\strut
\end{minipage}\tabularnewline
\begin{minipage}[t]{0.04\columnwidth}\raggedright
3\strut
\end{minipage} & \begin{minipage}[t]{0.23\columnwidth}\raggedright
LAN Communication\strut
\end{minipage} & \begin{minipage}[t]{0.65\columnwidth}\raggedright
Configuration page for LAN communication; it is the standard
setting.\strut
\end{minipage}\tabularnewline
\begin{minipage}[t]{0.04\columnwidth}\raggedright
4\strut
\end{minipage} & \begin{minipage}[t]{0.23\columnwidth}\raggedright
Network Status\strut
\end{minipage} & \begin{minipage}[t]{0.65\columnwidth}\raggedright
Display network status, interface and configuration protocol\strut
\end{minipage}\tabularnewline
\bottomrule
\end{longtable}


\subsection{RS-485 Interface}

The following figure depicts the RS-485 setup screen. When attempting to
communicate via the RS-485 protocol, make sure the configuration on the
client side has a matching set of parameters illustrated in this figure
(described in the accompanying table). This setting overrides that on
the client device. RS-485 (Recommended Standard-485) is a serial
interface standard for multipoint communication lines that was adopted
by the Electronic Industries Association (EIA).

\textbf{\emph{Note:}} The RS-485 setting used here is the standard
protocol. Therefore, the \textbf{protocol} selection and \textbf{echo
back} options are not available.

\begin{figure}
\centering

\includegraphics[width=0.95\textwidth]{figures/3989606209-S6_1-RS485-cfg-002a.PNG}
\caption{RS-485 protocol and setting}
\end{figure}\FloatBarrier

\begin{longtable}[]{@{}lll@{}}
\toprule
\begin{minipage}[b]{0.04\columnwidth}\raggedright
No.\strut
\end{minipage} & \begin{minipage}[b]{0.23\columnwidth}\raggedright
Name\strut
\end{minipage} & \begin{minipage}[b]{0.65\columnwidth}\raggedright
Description / Operation\strut
\end{minipage}\tabularnewline
\midrule
\endhead
\begin{minipage}[t]{0.04\columnwidth}\raggedright
1\strut
\end{minipage} & \begin{minipage}[t]{0.23\columnwidth}\raggedright
Address\strut
\end{minipage} & \begin{minipage}[t]{0.65\columnwidth}\raggedright
Specify the address number from 1 to 16. The address number is assigned
to identify the chamber. Make sure that it is not the same as the
address of other chambers connected.\strut
\end{minipage}\tabularnewline
\begin{minipage}[t]{0.04\columnwidth}\raggedright
2\strut
\end{minipage} & \begin{minipage}[t]{0.23\columnwidth}\raggedright
Delimiter\strut
\end{minipage} & \begin{minipage}[t]{0.65\columnwidth}\raggedright
End packet signal with delimiter with option CR, LF or CR+LF. Make sure
to select CR+LF (default setting).\strut
\end{minipage}\tabularnewline
\begin{minipage}[t]{0.04\columnwidth}\raggedright
3\strut
\end{minipage} & \begin{minipage}[t]{0.23\columnwidth}\raggedright
Baud Rate\strut
\end{minipage} & \begin{minipage}[t]{0.65\columnwidth}\raggedright
Select from 4800, 9600, and 19200. Default setting is 9600.\strut
\end{minipage}\tabularnewline
\begin{minipage}[t]{0.04\columnwidth}\raggedright
4\strut
\end{minipage} & \begin{minipage}[t]{0.23\columnwidth}\raggedright
Parity\strut
\end{minipage} & \begin{minipage}[t]{0.65\columnwidth}\raggedright
Select from None, Odd, and Eve. Default setting is None.\strut
\end{minipage}\tabularnewline
\begin{minipage}[t]{0.04\columnwidth}\raggedright
5\strut
\end{minipage} & \begin{minipage}[t]{0.23\columnwidth}\raggedright
Data Bits\strut
\end{minipage} & \begin{minipage}[t]{0.65\columnwidth}\raggedright
Select from 8 and 7. Default setting is 8bits.\strut
\end{minipage}\tabularnewline
\begin{minipage}[t]{0.04\columnwidth}\raggedright
6\strut
\end{minipage} & \begin{minipage}[t]{0.23\columnwidth}\raggedright
Stop Bits\strut
\end{minipage} & \begin{minipage}[t]{0.65\columnwidth}\raggedright
Select from 1 and 2. Default setting is 1.\strut
\end{minipage}\tabularnewline
\bottomrule
\end{longtable}


\subsection{RS-232C Interface}

RS-232C (Recommended Standard-232) is a serial interface widely adopted
for transmission between computers and peripheral devices. It is based
on a communication standard of the EIA.

The RS-232C setup screen follows the similar configuration protocol of
the RS-485 with the exception that it does not have an address. When
attempting to communicate via the RS-232 protocol, make sure the
configuration on the client side has a matching parameter as illustrated
in this screen; it overrides that on the client device.


\subsection{GPIB Interface}

The General Purpose Interface Bus or GPIB is a standard parallel
interface used for attaching sensors and programmable instruments to a
computer. It is officially known as IEEE-488 (Standard No.~488 of the
Institute of Electrical and Electronic Engineers {[}USA{]}) which was
based on the HP-IB (Hewlett-Packard Interface Bus) standard of
Hewlett-Packard Company, often referred to as GI-IB (IEEE-488/HP-IB).

The GPIB setup screen and its configuration option is currently
unavailable/unsupported.


\subsection{Set LAN: Network Interface \&
Setup}

With the GL controller instrumentation connected to LAN, you can access
the chamber to control and operate it remotely, just as you would on the
HMI. The Set LAN setup screen displays the GL controller network
protocol and its configuration. By default, the GL controller uses DHCP
to join a LAN network. Its IP address will be provided by the DHCP
server of that network.

\textbf{\emph{Note:}} Details of network configuration and Set LAN
editing are provided in Chapter 8.

\begin{figure}
\centering

\includegraphics[width=0.95\textwidth]{figures/3124032764-S7_1_4N-LAN-network-001a.PNG}
\caption{LAN network configuration and protocol}
\end{figure}\FloatBarrier

\begin{longtable}[]{@{}lll@{}}
\toprule
\begin{minipage}[b]{0.04\columnwidth}\raggedright
No.\strut
\end{minipage} & \begin{minipage}[b]{0.23\columnwidth}\raggedright
Name\strut
\end{minipage} & \begin{minipage}[b]{0.65\columnwidth}\raggedright
Description / Operation\strut
\end{minipage}\tabularnewline
\midrule
\endhead
\begin{minipage}[t]{0.04\columnwidth}\raggedright
1\strut
\end{minipage} & \begin{minipage}[t]{0.23\columnwidth}\raggedright
Hostname\strut
\end{minipage} & \begin{minipage}[t]{0.65\columnwidth}\raggedright
Displays the hostname of the GL controller system; it can be used to
access the GL controller UI via a Web browser.\strut
\end{minipage}\tabularnewline
\begin{minipage}[t]{0.04\columnwidth}\raggedright
2\strut
\end{minipage} & \begin{minipage}[t]{0.23\columnwidth}\raggedright
DHCP\strut
\end{minipage} & \begin{minipage}[t]{0.65\columnwidth}\raggedright
Protocol option; when DHCP box is checked, IP address provided by DHCP
server is used; when unchecked, the GL controller uses a static IP
configuration, as shown. Refer to Chapter 8 for details.\strut
\end{minipage}\tabularnewline
\begin{minipage}[t]{0.04\columnwidth}\raggedright
3\strut
\end{minipage} & \begin{minipage}[t]{0.23\columnwidth}\raggedright
IP Address\strut
\end{minipage} & \begin{minipage}[t]{0.65\columnwidth}\raggedright
Displays the current IP address of the GL controller system\strut
\end{minipage}\tabularnewline
\begin{minipage}[t]{0.04\columnwidth}\raggedright
4\strut
\end{minipage} & \begin{minipage}[t]{0.23\columnwidth}\raggedright
Subnet Mask\strut
\end{minipage} & \begin{minipage}[t]{0.65\columnwidth}\raggedright
Displays the network subnet mask of the system\strut
\end{minipage}\tabularnewline
\begin{minipage}[t]{0.04\columnwidth}\raggedright
5\strut
\end{minipage} & \begin{minipage}[t]{0.23\columnwidth}\raggedright
Gateway\strut
\end{minipage} & \begin{minipage}[t]{0.65\columnwidth}\raggedright
Displays the gateway address for network connection points\strut
\end{minipage}\tabularnewline
\begin{minipage}[t]{0.04\columnwidth}\raggedright
6\strut
\end{minipage} & \begin{minipage}[t]{0.23\columnwidth}\raggedright
DNS1\strut
\end{minipage} & \begin{minipage}[t]{0.65\columnwidth}\raggedright
Displays the primary Domain Name System, DNS 1.\strut
\end{minipage}\tabularnewline
\begin{minipage}[t]{0.04\columnwidth}\raggedright
7\strut
\end{minipage} & \begin{minipage}[t]{0.23\columnwidth}\raggedright
DNS2\strut
\end{minipage} & \begin{minipage}[t]{0.65\columnwidth}\raggedright
Displays the secondary Domain Name System, DNS 2.\strut
\end{minipage}\tabularnewline
\begin{minipage}[t]{0.04\columnwidth}\raggedright
8\strut
\end{minipage} & \begin{minipage}[t]{0.23\columnwidth}\raggedright
Communication Mode\strut
\end{minipage} & \begin{minipage}[t]{0.65\columnwidth}\raggedright
Setting options to control communication mode.\strut
\end{minipage}\tabularnewline
\begin{minipage}[t]{0.04\columnwidth}\raggedright
9\strut
\end{minipage} & \begin{minipage}[t]{0.23\columnwidth}\raggedright
Delimiter\strut
\end{minipage} & \begin{minipage}[t]{0.65\columnwidth}\raggedright
Delimiter type to control communication and transmission.\strut
\end{minipage}\tabularnewline
\begin{minipage}[t]{0.04\columnwidth}\raggedright
10\strut
\end{minipage} & \begin{minipage}[t]{0.23\columnwidth}\raggedright
TCP Socket\strut
\end{minipage} & \begin{minipage}[t]{0.65\columnwidth}\raggedright
Communication port selection options.\strut
\end{minipage}\tabularnewline
\bottomrule
\end{longtable}


\subsection{LAN Communication Test: HMI or Web
Display}

With LAN connection, direct communication with the GL controller via its
native text commands can be accomplished remotely that can bypass
control on the HMI or the Web browser interface. This type of
communication can take place via this LAN Communication Test page
locally on the HMI or remotely on the Web display.

\begin{figure}
\centering

\includegraphics[width=0.95\textwidth]{figures/3641070169-S6_1_5-cmd-prompt-001a.PNG}
\caption{Command prompt and response buffer of LAN Communication Test
page}
\end{figure}\FloatBarrier

\begin{longtable}[]{@{}lll@{}}
\toprule
\begin{minipage}[b]{0.04\columnwidth}\raggedright
No.\strut
\end{minipage} & \begin{minipage}[b]{0.23\columnwidth}\raggedright
Name\strut
\end{minipage} & \begin{minipage}[b]{0.65\columnwidth}\raggedright
Description \& Operation\strut
\end{minipage}\tabularnewline
\midrule
\endhead
\begin{minipage}[t]{0.04\columnwidth}\raggedright
1\strut
\end{minipage} & \begin{minipage}[t]{0.23\columnwidth}\raggedright
Command Prompt\strut
\end{minipage} & \begin{minipage}[t]{0.65\columnwidth}\raggedright
Issue raw commands here; commands must contain command and delimiter,
using the correct syntax of command structure. The prompt supports
history commands; previous commands can be recalled with the up-arrow
key for editing and reinvocation.\strut
\end{minipage}\tabularnewline
\begin{minipage}[t]{0.04\columnwidth}\raggedright
2\strut
\end{minipage} & \begin{minipage}[t]{0.23\columnwidth}\raggedright
Response Data\strut
\end{minipage} & \begin{minipage}[t]{0.65\columnwidth}\raggedright
Displays responses returned by the GL controller; previously issued
command may be re-executed (re-invoked) by clicking on it.\strut
\end{minipage}\tabularnewline
\bottomrule
\end{longtable}

The following procedure outlines the steps to open the command prompt to
issues raw text commands to the GL controller, locally or remotely from
within the GL controller UI.

\textbf{\emph{Procedure:}} LAN Communication Test on HMI or Web Display

\begin{enumerate}
\def\labelenumi{\arabic{enumi}.}
\item
  Press/click \textbf{Setup} (in the menu bar or main home page);
  press/click Configuration; press/click Set Communication; press/click
  LAN Communication Test (see arrow).

  \begin{figure}
  \centering
  
\includegraphics[width=0.95\textwidth]{figures/911988068-S6_1_5-cmd-hmi-001a.PNG}
  \caption{Command prompt and response buffer}
  \end{figure}\FloatBarrier
\item
  In the command prompt (see arrow), enter: \textbf{date?} and press
  Enter; enter: \textbf{rom?} and press Enter; enter: \textbf{temp?} and
  press Enter. Results are shown as follows.

  \begin{figure}
  \centering
  
\includegraphics[width=0.95\textwidth]{figures/609353641-S6_1_5-step2-001a.PNG}
  \caption{Command data and response data}
  \end{figure}\FloatBarrier
\item
  To exit the LAN Communication Test page, access any menu outside of
  this page.
\end{enumerate}


\subsection{LAN Communication Test:
PuTTY}

The most practical approach is via a terminal emulator, such as Termite
or PuTTY, where raw text commands can be issued directly to the GL
controller. Refer to GL Controller Communications Manual for details on
GL controller native text commands and its instruction set.

Examples of GL controller text commands issued on the LAN Communication
Test page and via PuTTY terminal emulator are provided as follows.

\textbf{\emph{Procedure:}} Remote control via PuTTY

\begin{enumerate}
\def\labelenumi{\arabic{enumi}.}
\item
  Download the executable or installable package for the desired
  platform from the PuTTY official web site: https://www.putty.org.
\item
  Look up the IP address of the GL controller on the Set Communication
  page, indicated by the IPv4 Address field, under the \textbf{eth0}
  column (see arrow).

  \begin{figure}
  \centering
  
\includegraphics[width=0.95\textwidth]{figures/3188314483-S6_1_5-network-status-001a.PNG}
  \caption{How to locate IP address of GL controller system}
  \end{figure}\FloatBarrier
\item
  Launch PuTTY on your PC and follow the procedure in the diagram to
  connect to the GL controller.

  \begin{itemize}
  \tightlist
  \item
    Enter IP address found in (2). Example: 10.30.200.247
  \item
    Enter port number 10001 in the ``Port'' field (in place of default
    22)
  \item
    Select \textbf{Raw} for connection type
  \item
    Click Open to start the communication session
  \end{itemize}

  \begin{figure}
  \centering
  
\includegraphics[width=0.95\textwidth]{figures/2258019861-S6_1_5-putty-comm-001.PNG}
  \caption{Network configuration on PuTTY}
  \end{figure}\FloatBarrier
\item
  Enter ``date?'' and press Enter; enter ``rom?'' and press Enter; enter
  ``temp?'' and press Enter. Results are shown below.

  \begin{figure}
  \centering
  
\includegraphics[width=0.95\textwidth]{figures/1264660407-S6_1_5-putty-cmds-001.PNG}
  \caption{PuTTY input commands and output responses}
  \end{figure}\FloatBarrier
\item
  To end communication, click X and confirm the exit.
\end{enumerate}


\section{Set Date/Time}

To keep an accurate date/time of the data log, it is important to have
the correct date/time on the GL controller. Data log will bear a
date/time stamp used by the GL controller. By default, the GL controller
date/time is synchronized using the Network Time Protocol (NTP) server
provided by the Debian network time pool.

Based on this setting, the synchronization can only occur if the GL
controller has access to the Internet. Without Internet connection, the
NTP server of the Debian network time pool does not work, and the
date/time of the GL controller will be out of sync as it may drift from
the real time. The purpose of NTP is to apply skewness to the date/time
such that it always remains in sync. If your network has its own NTP
server, the GL controller can use it. Without a local NTP server, you
can still adjust the date/time manually or synchronize it using a local
device such as your laptop or PC.


\subsection{Set Date/Time on HMI}

On the HMI, date/time information of the GL controller is displayed on
the status bar, next to the login status icon. This information is also
found on the Set Date/Time page, as depicted in the following figure.

\begin{figure}
\centering

\includegraphics[width=0.95\textwidth]{figures/138463195-S6_2-datetime-hmi-002aa.PNG}
\caption{Date/Time setting on HMI}
\end{figure}\FloatBarrier

\begin{longtable}[]{@{}lll@{}}
\toprule
\begin{minipage}[b]{0.04\columnwidth}\raggedright
No.\strut
\end{minipage} & \begin{minipage}[b]{0.23\columnwidth}\raggedright
Name\strut
\end{minipage} & \begin{minipage}[b]{0.65\columnwidth}\raggedright
Description\strut
\end{minipage}\tabularnewline
\midrule
\endhead
\begin{minipage}[t]{0.04\columnwidth}\raggedright
1\strut
\end{minipage} & \begin{minipage}[t]{0.23\columnwidth}\raggedright
Set Date/Time\strut
\end{minipage} & \begin{minipage}[t]{0.65\columnwidth}\raggedright
Displays date/time setting options\strut
\end{minipage}\tabularnewline
\begin{minipage}[t]{0.04\columnwidth}\raggedright
2\strut
\end{minipage} & \begin{minipage}[t]{0.23\columnwidth}\raggedright
Time Zone\strut
\end{minipage} & \begin{minipage}[t]{0.65\columnwidth}\raggedright
Displays selected time zone currently used by the GL controller; apply
radio button to select a desired time zone from the list.\strut
\end{minipage}\tabularnewline
\begin{minipage}[t]{0.04\columnwidth}\raggedright
3\strut
\end{minipage} & \begin{minipage}[t]{0.23\columnwidth}\raggedright
NTP Setup\strut
\end{minipage} & \begin{minipage}[t]{0.65\columnwidth}\raggedright
Lists current NTP server\strut
\end{minipage}\tabularnewline
\begin{minipage}[t]{0.04\columnwidth}\raggedright
4\strut
\end{minipage} & \begin{minipage}[t]{0.23\columnwidth}\raggedright
Add\strut
\end{minipage} & \begin{minipage}[t]{0.65\columnwidth}\raggedright
Additional NTP can be added via this button\strut
\end{minipage}\tabularnewline
\bottomrule
\end{longtable}

The following procedure outlines the steps to adjust date/time and local
time zone.

\textbf{\emph{Procedure:}} Date/Time and local time zone adjustment on
HMI

\begin{enumerate}
\def\labelenumi{\arabic{enumi}.}
\tightlist
\item
  To adjust the date, press the calendar icon (see item 1) and select a
  new date.\\
\item
  To adjust the time, press the clock icon (see item 1) and select a new
  time.
\item
  To adjust local time zone, press the radio button (see item 2) and
  select a desired time zone on the list.
\item
  Press the Save button in the secondary menu to apply the new setting.
\end{enumerate}


\subsection{Set Date/Time on Remote
Display}

On the Web display, this information can be found on the Set Date/Time
page itself, as depicted in the figure.

\begin{figure}
\centering

\includegraphics[width=0.95\textwidth]{figures/4068765571-S6_2-datetime-web-002a.PNG}
\caption{Date/Time setting on Web display}
\end{figure}\FloatBarrier

\begin{longtable}[]{@{}lll@{}}
\toprule
\begin{minipage}[b]{0.04\columnwidth}\raggedright
No.\strut
\end{minipage} & \begin{minipage}[b]{0.23\columnwidth}\raggedright
Name\strut
\end{minipage} & \begin{minipage}[b]{0.65\columnwidth}\raggedright
Description\strut
\end{minipage}\tabularnewline
\midrule
\endhead
\begin{minipage}[t]{0.04\columnwidth}\raggedright
1\strut
\end{minipage} & \begin{minipage}[t]{0.23\columnwidth}\raggedright
Set Date/Time\strut
\end{minipage} & \begin{minipage}[t]{0.65\columnwidth}\raggedright
Displays date/time setting options: Automatic or Manual.\strut
\end{minipage}\tabularnewline
\begin{minipage}[t]{0.04\columnwidth}\raggedright
2\strut
\end{minipage} & \begin{minipage}[t]{0.23\columnwidth}\raggedright
NTP Setup\strut
\end{minipage} & \begin{minipage}[t]{0.65\columnwidth}\raggedright
Displays list of current NTP servers with priority queue from top to
bottom; each can be removed via its trash bin icon.\strut
\end{minipage}\tabularnewline
\begin{minipage}[t]{0.04\columnwidth}\raggedright
3\strut
\end{minipage} & \begin{minipage}[t]{0.23\columnwidth}\raggedright
Add\strut
\end{minipage} & \begin{minipage}[t]{0.65\columnwidth}\raggedright
Additional NTP server can be added via this button; it is added to the
bottom of the list, acting as the last queue on the list.\strut
\end{minipage}\tabularnewline
\bottomrule
\end{longtable}

The following procedure outlines the steps to adjust date and time
manually.

\begin{figure}
\centering

\includegraphics[width=0.95\textwidth]{figures/2307331095-S6_2_2-manual-datetime-setting-001bb.PNG}
\caption{Manually adjusting date/time}
\end{figure}\FloatBarrier

\textbf{\emph{Procedure:}} Manual Date/Time Setting

\textbf{\emph{Note:}} The calendar and clock inputs are Web browser
dependent. Each Web browser provides calendar and clock drop-down
pickers differently. For example, Firefox does not support the clock
picker. The screenshot in the above figure was based on Google Chrome.

\begin{enumerate}
\def\labelenumi{\arabic{enumi}.}
\tightlist
\item
  Click the \textbf{Manual} button (see arrow 1) to turn on calendar and
  clock options.
\item
  Enter new month/date/year values in the field (see arrow 2) or click
  the calendar icon (see arrow 3) to adjust the month (year) and date
  using the drop-down calendar.\\
\item
  Enter new hour:min:sec values in the time field (see arrow 4) or click
  the clock icon (see arrow 5) to adjust the time using the drop-down
  clock values.\\
\item
  Apply the setting using the Save button in the secondary buttons.
\end{enumerate}


\subsection{NTP Server Setup: Local Time
Server}

If the GL controller does not have access to the Internet, its clock
will drift and eventually be out of sync with the real clock. In order
to apply autonomous date/time and synchronization, the GL controller
must have access to a local time server. This local time server can be
added to the top of the NTP list to take priority.

\begin{figure}
\centering

\includegraphics[width=0.95\textwidth]{figures/3188116534-S6_2_3-ntp-001a.PNG}
\caption{NTP setting and adjustment}
\end{figure}\FloatBarrier

\begin{longtable}[]{@{}lll@{}}
\toprule
No. & Name & Description\tabularnewline
\midrule
\endhead
1 & NTP & Displays NTP hostname or IP address\tabularnewline
2 & Delete & Use the trash bin icon to remove NTP from the
list\tabularnewline
3 & Add & Add additional NTP to the bottom of the list\tabularnewline
4 & Save & Apply and store the current setting with the save
button\tabularnewline
\bottomrule
\end{longtable}

The following procedure outlines the steps how to use a local time
server.

\textbf{\emph{Procedure:}}

\begin{enumerate}
\def\labelenumi{\arabic{enumi}.}
\tightlist
\item
  Edit the top NTP entry with your local timer server using either its
  IP address or its hostname.
\item
  Remove the rest of the Debian NTP server using the trash bin. This
  procedure ensures that the system will not attempt to access Debian
  NTP server to apply date/time synchronization.
\item
  Save the setting via the Save button.
\end{enumerate}


\subsection{Restore Debian NTP
Server}

To restore the NTP server using Debian date/time protocol, apply the
following steps. To ensure successful date/time synchronization, four
different Debian NTP servers are used, numbered from 0 to 3. For
instance, if the first NTP server does not respond, the second one will
be used, and the process continues to the last one on the list.

\begin{longtable}[]{@{}ll@{}}
\toprule
NTP Server \# & NTP Hostname or IP Address\tabularnewline
\midrule
\endhead
1 & 0.debian.pool.ntp.org\tabularnewline
2 & 1.debian.pool.ntp.org\tabularnewline
3 & 2.debian.pool.ntp.org\tabularnewline
4 & 3.debian.pool.ntp.org\tabularnewline
\bottomrule
\end{longtable}


\section{Operation Process}

The \textbf{Operation Process} page provides options to manage and
control the chamber operation when it is interrupted abruptly. In
default setting, the chamber will stop any test after resuming from a
power outage (that occurred during testing); test operation continues
uninterrupted when door is open; door open warning is set to Off. These
three default settings are depicted in the following configure. However,
all of them are customizable.

\begin{figure}
\centering

\includegraphics[width=0.95\textwidth]{figures/834206548-S6_3-op-process-001a.PNG}
\caption{Operation Process options}
\end{figure}\FloatBarrier

\begin{longtable}[]{@{}lll@{}}
\toprule
\begin{minipage}[b]{0.04\columnwidth}\raggedright
No.\strut
\end{minipage} & \begin{minipage}[b]{0.23\columnwidth}\raggedright
Name\strut
\end{minipage} & \begin{minipage}[b]{0.65\columnwidth}\raggedright
Description / Operation\strut
\end{minipage}\tabularnewline
\midrule
\endhead
\begin{minipage}[t]{0.04\columnwidth}\raggedright
1\strut
\end{minipage} & \begin{minipage}[t]{0.23\columnwidth}\raggedright
Power Fail Recovery\strut
\end{minipage} & \begin{minipage}[t]{0.65\columnwidth}\raggedright
Chamber operation after power failure recovery. Two options are
available: (1) Continue Test, (2) Stop Test. Default setting is Stop
Test.\strut
\end{minipage}\tabularnewline
\begin{minipage}[t]{0.04\columnwidth}\raggedright
2\strut
\end{minipage} & \begin{minipage}[t]{0.23\columnwidth}\raggedright
Door Open Action\strut
\end{minipage} & \begin{minipage}[t]{0.65\columnwidth}\raggedright
Displays available options with Continue, Pause and Stop Test. The Pause
and Stop Test options can be adjusted with delay time in x
minutes.\strut
\end{minipage}\tabularnewline
\begin{minipage}[t]{0.04\columnwidth}\raggedright
3\strut
\end{minipage} & \begin{minipage}[t]{0.23\columnwidth}\raggedright
Door Open Warning\strut
\end{minipage} & \begin{minipage}[t]{0.65\columnwidth}\raggedright
Displays available options with On and Off. The On option can be
adjusted with delay time in x minutes.\strut
\end{minipage}\tabularnewline
\begin{minipage}[t]{0.04\columnwidth}\raggedright
4\strut
\end{minipage} & \begin{minipage}[t]{0.23\columnwidth}\raggedright
Save\strut
\end{minipage} & \begin{minipage}[t]{0.65\columnwidth}\raggedright
Use the \textbf{Save} button to apply and store new settings on alarm or
warning beeper.\strut
\end{minipage}\tabularnewline
\begin{minipage}[t]{0.04\columnwidth}\raggedright
5\strut
\end{minipage} & \begin{minipage}[t]{0.23\columnwidth}\raggedright
Back\strut
\end{minipage} & \begin{minipage}[t]{0.65\columnwidth}\raggedright
Use the Back button to exit the current page.\strut
\end{minipage}\tabularnewline
\bottomrule
\end{longtable}

The following procedure outlines the steps to set the chamber to stop
test after x minutes when the door is open.

\textbf{\emph{Procedure:}}

\begin{enumerate}
\def\labelenumi{\arabic{enumi}.}
\tightlist
\item
  Press/click the \textbf{Stop Test} button
\item
  Enter x minutes in the minute field. Refer to Sections 2.4.4 and 2.4.5
  on how to turn on the keyboard/keypad when operating on the HMI.\\
\item
  Press/click the \textbf{Save} icon/button to save the setting.
\item
  Press/click the \textbf{Back} icon/button to exit.
\end{enumerate}


\section{Users (Register User
Password)}

The GL controller supports and allows multiple operators with separate
user accounts to control and operate the instrumentation. Each user
account has specific roles with read/write privileges to access and
operate the system. By default, three user accounts exist on the system,
as depicted in the figure. The layout of the \textbf{Users} page
displays specific roles associated with these users; each may be
assigned with specific roles according to the selection in the drop-down
list.

\begin{figure}
\centering

\includegraphics[width=0.95\textwidth]{figures/795481158-S6_4-user-accnt-001a.PNG}
\caption{Users settings page}
\end{figure}\FloatBarrier

For HMI operation, the localhost user has specific roles to operate and
control the chamber, including settings in the configuration menu. Refer
to the following table.

For remote access and operation via a network, each user with specific
roles on read/write privileges can log in to the GL controller system to
operate the chamber. Refer to the following table.

\begin{figure}
\centering

\includegraphics[width=0.95\textwidth]{figures/1045510844-S6_4-user-accnt-002a.PNG}
\caption{The nomenclature of the User Settings submenu}
\end{figure}\FloatBarrier

\begin{longtable}[]{@{}lll@{}}
\toprule
\begin{minipage}[b]{0.04\columnwidth}\raggedright
No.\strut
\end{minipage} & \begin{minipage}[b]{0.23\columnwidth}\raggedright
Name\strut
\end{minipage} & \begin{minipage}[b]{0.65\columnwidth}\raggedright
Description\strut
\end{minipage}\tabularnewline
\midrule
\endhead
\begin{minipage}[t]{0.04\columnwidth}\raggedright
1\strut
\end{minipage} & \begin{minipage}[t]{0.23\columnwidth}\raggedright
Account Property\strut
\end{minipage} & \begin{minipage}[t]{0.65\columnwidth}\raggedright
Displays account username and password\strut
\end{minipage}\tabularnewline
\begin{minipage}[t]{0.04\columnwidth}\raggedright
2\strut
\end{minipage} & \begin{minipage}[t]{0.23\columnwidth}\raggedright
Access Properties\strut
\end{minipage} & \begin{minipage}[t]{0.65\columnwidth}\raggedright
Displays read/write privileges associated with the five menus: Home,
Operation Mode, Monitor, Set Mode and Setup. Each menu has three access
privileges in the drop-down list: RO (Read-Only), RW (Read-Write), NA
(No-Access). Different privileges can be assigned to each account using
options from the drop-down list.\strut
\end{minipage}\tabularnewline
\begin{minipage}[t]{0.04\columnwidth}\raggedright
3\strut
\end{minipage} & \begin{minipage}[t]{0.23\columnwidth}\raggedright
admin\strut
\end{minipage} & \begin{minipage}[t]{0.65\columnwidth}\raggedright
Displays administrator (admin) account. Default credentials for this
user account are username: admin, password: admin. This account is
generally assigned with complete powers to the GL controller
system.\strut
\end{minipage}\tabularnewline
\begin{minipage}[t]{0.04\columnwidth}\raggedright
4\strut
\end{minipage} & \begin{minipage}[t]{0.23\columnwidth}\raggedright
localhost\strut
\end{minipage} & \begin{minipage}[t]{0.65\columnwidth}\raggedright
User account granted with specific roles with read/write privileges to
access and control the instrumentation\strut
\end{minipage}\tabularnewline
\begin{minipage}[t]{0.04\columnwidth}\raggedright
5\strut
\end{minipage} & \begin{minipage}[t]{0.23\columnwidth}\raggedright
Add\strut
\end{minipage} & \begin{minipage}[t]{0.65\columnwidth}\raggedright
Create new user account using this button\strut
\end{minipage}\tabularnewline
\begin{minipage}[t]{0.04\columnwidth}\raggedright
6\strut
\end{minipage} & \begin{minipage}[t]{0.23\columnwidth}\raggedright
Delete\strut
\end{minipage} & \begin{minipage}[t]{0.65\columnwidth}\raggedright
Remove user account using this trash bin\strut
\end{minipage}\tabularnewline
\bottomrule
\end{longtable}

\textbf{\emph{Note:}} Refer to Section 2.6.2 for details on the
localhost and how to log in or out of the account.


\subsection{Add User Account}

The following procedure outlines the steps to create a user account
(called technician) and set certain access privileges to the menu
properties. The user can access all operation menus except
\textbf{Setup}.

\begin{figure}
\centering

\includegraphics[width=0.95\textwidth]{figures/260381379-S6_4_1-add-user-002a.PNG}
\caption{Creating or Adding a new user account}
\end{figure}\FloatBarrier

\textbf{\emph{Procedure:}}

\begin{enumerate}
\def\labelenumi{\arabic{enumi}.}
\tightlist
\item
  \textbf{Add}: Click the \textbf{add} button (callout 1). Refer to the
  callout in the above figure.\\
\item
  \textbf{Username}: Enter username under the \textbf{Username}
  column.\\
\item
  \textbf{Password}: Enter the password for this user in the
  \textbf{Password} fields (twice).
\item
  \textbf{Access Privileges}: Click and select access privilege (No
  Access, Read Only, Read Write) for each menu. Repeat the process for
  the rest of the menus.
\item
  \textbf{Save}: Click the \textbf{Save} button at the upper-right
  corner to save the settings.
\end{enumerate}

This new account will be available for use immediately after the
\textbf{Save} button is applied.


\subsection{Delete or Remove User
Account}

The following procedure outlines the steps to remove a user account from
the GL controller system.

\textbf{\emph{Procedure:}}

\begin{enumerate}
\def\labelenumi{\arabic{enumi}.}
\tightlist
\item
  Ensure that the user has logged out.
\item
  Locate the user account to be deleted.
\item
  Click the trash bin next to the account to be deleted.
\item
  \textbf{Save}: Click the \textbf{Save} button at the upper-right
  corner to save the settings.
\end{enumerate}


\section{Control Attainment
Range}

Conditions of the attainment range and the attainment time can be used
to determine whether the chamber has reached the exposure temperature
(i.e., set points). The status where the process value falls within the
attainment range of the set point for a certain period of time or longer
is considered as a set point attainment; that the control is being
attained.

\begin{figure}
\centering

\includegraphics[width=0.95\textwidth]{figures/2858873299-S6_5-attainment-range-001.PNG}
\caption{Attainment range settings}
\end{figure}\FloatBarrier

The processing of controlling attainment range is active in the
following conditions.

\begin{itemize}
\tightlist
\item
  Soak control of a program step
\item
  Temperature attainment output (optional)
\item
  Humidity attainment output (optional)
\end{itemize}

The attainment range must be smaller than the upper and lower
deviations, as depicted in the following graph.

\begin{figure}
\centering

\includegraphics[width=0.95\textwidth]{figures/1645091637-S6_5-attainment-setpoint-graph-001.PNG}
\caption{Temp/Humi attainment range setup}
\end{figure}\FloatBarrier

\textbf{\emph{Procedure:}}

\begin{enumerate}
\def\labelenumi{\arabic{enumi}.}
\tightlist
\item
  Set the attainment range (temperature, humidity). Input range:
  Temperature: 0.1 °C to 100.0 °C; humidity: 1 \%RH to 100 \%RH
\item
  Set the attainment time. Set the attainment time for each of
  temperature and humidity. Input range: 0 to 9999 seconds.
\item
  Set the attainment range and attainment time of the product
  temperature (only when option is installed). Input range: 0.1 °C to
  100.0 °C; 0 to 9999 seconds.
\end{enumerate}


\section{Sensor Offset
(calibration)}

Set the correction values of the install temperature input sensor. By
setting the correction values for the sensors input value, you can
compensate for changes over time and adjust the differences between the
chamber and other measuring instruments. Therefore, the goal of sensor
calibration is to match the values with those of a calibrated device
(the reference device) at all times, not to correct for control errors.
Also, these correction values are enabled for all measured values such
as the monitor screen and the trend graph.

Perform correction for the input value of the temperature (dry-bulb)
sensor, humidity (web-bulb) sensor, specimen temperature control sensor
(optional) for measuring the temperature and humidity in the test area.

The purpose of this offset adjustment is to adjust the process values
for the already calibrated instruments and not to correct control system
errors. In addition, the correction is reflected in all process values,
such as the monitor screen, trend graph, and monitor output (optional
output terminal for temperature recorder).

\textbf{\emph{Procedure:}}

\begin{enumerate}
\def\labelenumi{\arabic{enumi}.}
\tightlist
\item
  Set the correction value of the offset adjustment.
\item
  Set the correction value to be added to or subtracted from the sensor
  in the range of +/- 5 °C.
\item
  After entering the value, the value corrected to the sensor input
  value is displayed at the Process Value.
\end{enumerate}

\textbf{\emph{Procedure:}}

\begin{enumerate}
\def\labelenumi{\arabic{enumi}.}
\tightlist
\item
  Press Set/Calibrate sensor on the Configuration screen.
\item
  Enter the correction values. Press an entered value to display a
  numeric keypad, and then enter the correction value. After you entered
  the value, the corrected sensor input value will be displayed under
  the
\end{enumerate}

\begin{longtable}[]{@{}l@{}}
\toprule
\textbf{Reference}\tabularnewline
\midrule
\endhead
If precise calibration is required, contact ESPEC.\tabularnewline
\bottomrule
\end{longtable}


\section{Time Signal Names}

The names of time signals may be defined by the user. As depicted in the
following figure, both the full and short name can defined, along with
the On/Off labels.

\textbf{\emph{Note:}} The number of available time signals depend on the
options and/or specifications of the chamber.

\begin{figure}
\centering

\includegraphics[width=0.95\textwidth]{figures/3557207195-S6_7-TS-005a.PNG}
\caption{Names and labels of time signals}
\end{figure}\FloatBarrier

\begin{longtable}[]{@{}lll@{}}
\toprule
No. & Name & Description\tabularnewline
\midrule
\endhead
1 & Full & Displays full name of time signal.\tabularnewline
2 & Short & Displays short name of time signal\tabularnewline
3 & On & Displays On label; it can be defined with descriptive
name/label\tabularnewline
4 & Off & Displays Off label; it can be defined with descriptive
name/label\tabularnewline
\bottomrule
\end{longtable}

Once defined, time signal names and labels will appear in the display
monitor, Constant and Program setup pages.

\begin{figure}
\centering

\includegraphics[width=0.95\textwidth]{figures/4168980682-S6_7-TS-006-eg01-eg02a.PNG}
\caption{Effects of time signal names and labels}
\end{figure}\FloatBarrier

The following procedure outlines the steps to customize the time signal
names and their On/Off properties (labels), using Time Signal 1 as an
example.

\textbf{\emph{Procedure:}}

\begin{enumerate}
\def\labelenumi{\arabic{enumi}.}
\item
  Access the \textbf{Time Signal Names} page based on the menu sequence:
  Setup \textgreater{} Configuration \textgreater{} Time Signal Names.
  If this performance is done via the HMI, refer to Section 2.4.4 and
  2.4.5 on how to enable the keyboard/keypad for input.\\
\item
  Enter ``Dry Air Purge'' in the field indicated by callout 1; enter
  ``DAP'' indicated by callout 2, ``Start'' by callout 3 and ``Stop'' by
  callout 4.

  \begin{figure}
  \centering
  
\includegraphics[width=0.95\textwidth]{figures/3022334899-S6_7-step2-002a.PNG}
  \caption{Example: Customizing Time Signal 1 with descriptive names}
  \end{figure}\FloatBarrier
\item
  Save the setting using the \textbf{Save} icon (callout 5) and confirm
  the action in the pop-up dialog.
\end{enumerate}


\section{Set Chamber Details}

Chamber setting details is not available via the GL controller software
system.


\section{Display Setup}

The \textbf{Display Setup} page provides a number of customization
options for user interface theme, startup display tab/menu, HMI
screensaver timeout, user interface zooming and sub-tab (submenu)
position memory.

\textbf{\emph{Note:}} Customization options differ slightly between the
HMI display and the Web display. Only three options--startup display
tab, user interface theme and sub-tab position memorization--are
available for customization via the remote operation on the Web display.


\subsection{Display Setup on HMI}

The label associated with each button indicates its specific setting and
behavior. To customize the setting, press the desired button.

\begin{figure}
\centering

\includegraphics[width=0.95\textwidth]{figures/3796118120-S6_9-hmi-display-001.png}
\caption{Display options on HMI setting}
\end{figure}\FloatBarrier


\subsection{Display Setup on Web
Display}

The label associated with each button indicates its specific setting and
behavior. To customize the setting, press the desired button.

\begin{figure}
\centering

\includegraphics[width=0.95\textwidth]{figures/1505897642-S6_3-web-001.PNG}
\caption{Display options on Web display setting}
\end{figure}\FloatBarrier


\section{Set Option}

The \textbf{Set Option} has a few setting options on proportional band,
integral and derivative.

\begin{figure}
\centering

\includegraphics[width=0.95\textwidth]{figures/3355872501-S6_10-set-option-001.PNG}
\caption{GL controller set option}
\end{figure}\FloatBarrier


\section{Set Language}

The default language is English. Set your preferred language by
selecting one of the buttons. For a remote access via a Web browser, the
browser determines the selected language.

\begin{figure}
\centering

\includegraphics[width=0.95\textwidth]{figures/3814309202-S7_11-language-support-002.PNG}
\caption{Select your preferred language}
\end{figure}\FloatBarrier


\section{Set Email}

Two most important features in the \textbf{Set Email} page, among other
things, are the \textbf{admin} password recovery (i.e., revert to
factory setting) and automatic alert emails to authorized users. Both
features require that the GL controller system has access to the
Internet.

\begin{figure}
\centering

\includegraphics[width=0.95\textwidth]{figures/1708017426-S6_12-set-email-001ab.PNG}
\caption{Email configuration page}
\end{figure}\FloatBarrier

\begin{longtable}[]{@{}lll@{}}
\toprule
\begin{minipage}[b]{0.04\columnwidth}\raggedright
No.\strut
\end{minipage} & \begin{minipage}[b]{0.23\columnwidth}\raggedright
Name\strut
\end{minipage} & \begin{minipage}[b]{0.65\columnwidth}\raggedright
Description\strut
\end{minipage}\tabularnewline
\midrule
\endhead
\begin{minipage}[t]{0.04\columnwidth}\raggedright
1\strut
\end{minipage} & \begin{minipage}[t]{0.23\columnwidth}\raggedright
Host\strut
\end{minipage} & \begin{minipage}[t]{0.65\columnwidth}\raggedright
Displays the hostname or IP address of the email server\strut
\end{minipage}\tabularnewline
\begin{minipage}[t]{0.04\columnwidth}\raggedright
2\strut
\end{minipage} & \begin{minipage}[t]{0.23\columnwidth}\raggedright
Port\strut
\end{minipage} & \begin{minipage}[t]{0.65\columnwidth}\raggedright
Displays the port number used by the email server\strut
\end{minipage}\tabularnewline
\begin{minipage}[t]{0.04\columnwidth}\raggedright
3\strut
\end{minipage} & \begin{minipage}[t]{0.23\columnwidth}\raggedright
Send As\strut
\end{minipage} & \begin{minipage}[t]{0.65\columnwidth}\raggedright
Displays default sender name as \textbf{chamber\_controller}\strut
\end{minipage}\tabularnewline
\begin{minipage}[t]{0.04\columnwidth}\raggedright
4\strut
\end{minipage} & \begin{minipage}[t]{0.23\columnwidth}\raggedright
Authentication\strut
\end{minipage} & \begin{minipage}[t]{0.65\columnwidth}\raggedright
Allows options to authenticate outbound email; authentication is set by
default.\strut
\end{minipage}\tabularnewline
\begin{minipage}[t]{0.04\columnwidth}\raggedright
5\strut
\end{minipage} & \begin{minipage}[t]{0.23\columnwidth}\raggedright
SSL/TLS\strut
\end{minipage} & \begin{minipage}[t]{0.65\columnwidth}\raggedright
Allows options to use encryption for outbound email\strut
\end{minipage}\tabularnewline
\begin{minipage}[t]{0.04\columnwidth}\raggedright
6\strut
\end{minipage} & \begin{minipage}[t]{0.23\columnwidth}\raggedright
User\strut
\end{minipage} & \begin{minipage}[t]{0.65\columnwidth}\raggedright
Displays user for mail server authentication\strut
\end{minipage}\tabularnewline
\begin{minipage}[t]{0.04\columnwidth}\raggedright
7\strut
\end{minipage} & \begin{minipage}[t]{0.23\columnwidth}\raggedright
Password\strut
\end{minipage} & \begin{minipage}[t]{0.65\columnwidth}\raggedright
Password field for outbound email authentication\strut
\end{minipage}\tabularnewline
\begin{minipage}[t]{0.04\columnwidth}\raggedright
8\strut
\end{minipage} & \begin{minipage}[t]{0.23\columnwidth}\raggedright
Recovery email\strut
\end{minipage} & \begin{minipage}[t]{0.65\columnwidth}\raggedright
Email address used for recovering administration password of the GL
controller system\strut
\end{minipage}\tabularnewline
\begin{minipage}[t]{0.04\columnwidth}\raggedright
9\strut
\end{minipage} & \begin{minipage}[t]{0.23\columnwidth}\raggedright
Alarm recipient\strut
\end{minipage} & \begin{minipage}[t]{0.65\columnwidth}\raggedright
Enter email for each recipient to receive email alert from the GL
controller\strut
\end{minipage}\tabularnewline
\begin{minipage}[t]{0.04\columnwidth}\raggedright
10\strut
\end{minipage} & \begin{minipage}[t]{0.23\columnwidth}\raggedright
Save\strut
\end{minipage} & \begin{minipage}[t]{0.65\columnwidth}\raggedright
Save and apply new settings using the \textbf{Save} icon in the
secondary buttons\strut
\end{minipage}\tabularnewline
\begin{minipage}[t]{0.04\columnwidth}\raggedright
11\strut
\end{minipage} & \begin{minipage}[t]{0.23\columnwidth}\raggedright
Test Email\strut
\end{minipage} & \begin{minipage}[t]{0.65\columnwidth}\raggedright
Press/Click the button to test sending email alert; it requires access
to SMTP Office365 on the Internet\strut
\end{minipage}\tabularnewline
\begin{minipage}[t]{0.04\columnwidth}\raggedright
12\strut
\end{minipage} & \begin{minipage}[t]{0.23\columnwidth}\raggedright
Back\strut
\end{minipage} & \begin{minipage}[t]{0.65\columnwidth}\raggedright
Returns to the previous page, thus exiting the current email page.\strut
\end{minipage}\tabularnewline
\bottomrule
\end{longtable}


\subsection{Set Email: Password
Recovery}

To protect the GL controller system, the Account Recovery E-mail field
must contain an email address of the administrator (or owner) who
manages the GL controller system. This setup is to ensure that only the
administrator who manages the GL controller can proceed with the
password reset/recovery for the admin account. A one-time password will
be sent to this e-mail address to recover the administrator's account.

\begin{figure}
\centering

\includegraphics[width=0.95\textwidth]{figures/512749482-S6_12_1-password-recovery-002a.PNG}
\caption{Enter the admin email for password recovery}
\end{figure}\FloatBarrier

\textbf{\emph{Procedure:}}

\begin{enumerate}
\def\labelenumi{\arabic{enumi}.}
\item
  Enter the admin's email address in the \textbf{Account Recovery
  E-Mail} field (see callout 1, figure above).
\item
  Click/Press the \textbf{Save} icon (callout 2) to save the settings.
\end{enumerate}


\subsection{Set Email: Alert Email
Settings}

The alert email feature can notify the operator in case the chamber is
tripped with an alarm or with other issues. This notification may help
to prevent damage to test product or the chamber itself with the
operator acting to intervene the problem immediately by stopping the
chamber remotely.

The following procedure outlines the steps to set up e-mail alert to
receive notification of the operating conditions of the chamber.

\textbf{\emph{Procedure:}}

\begin{enumerate}
\def\labelenumi{\arabic{enumi}.}
\item
  Click \textbf{Add} (in the field of \textbf{Alarm Recipient})
\item
  Enter the recipient's email address.
\item
  Additional email can be added by clicking the \textbf{Add} button.

  \begin{figure}
  \centering
  
\includegraphics[width=0.95\textwidth]{figures/40339246-S6_12_2-step4-email-setting-001a.PNG}
  \caption{Setting alert email options}
  \end{figure}\FloatBarrier
\item
  The \textbf{X} button can be used to remove unwanted email address
  from the list.
\item
  Click Save (see callout 1) and confirm the action in the Yes/No pop-up
  window.
\item
  Click the \textbf{Test Email} icon (see callout 2) to test the alert
  email. A test email message will be sent to all the recipients on the
  list. If the GL controller fails to reach the Office365 mail server,
  the test will fail, but the setting can still be saved.
\end{enumerate}


\section{Set Sound}

The \textbf{Set Sound} page offers two options to control the Alarm
volume and Warning volume. By default, both volumes are set to maximum
level.

\begin{figure}
\centering

\includegraphics[width=0.95\textwidth]{figures/2321439942-S6_13-sound-page-001.PNG}
\caption{Alarm and Warning beep level}
\end{figure}\FloatBarrier

The slider bar on each volume control can be used to adjust volume
setting. The speaker icon at each slider bar can be used to turn off the
volume or turn it on to maximum level.

\begin{figure}
\centering

\includegraphics[width=0.95\textwidth]{figures/421474405-S6_13-sound-setting-001aa.PNG}
\caption{Volume control and adjustment}
\end{figure}\FloatBarrier

\begin{longtable}[]{@{}lll@{}}
\toprule
\begin{minipage}[b]{0.04\columnwidth}\raggedright
No.\strut
\end{minipage} & \begin{minipage}[b]{0.23\columnwidth}\raggedright
Name\strut
\end{minipage} & \begin{minipage}[b]{0.65\columnwidth}\raggedright
Description\strut
\end{minipage}\tabularnewline
\midrule
\endhead
\begin{minipage}[t]{0.04\columnwidth}\raggedright
1\strut
\end{minipage} & \begin{minipage}[t]{0.23\columnwidth}\raggedright
Slider Bar\strut
\end{minipage} & \begin{minipage}[t]{0.65\columnwidth}\raggedright
Use the slider bar to adjust the volume level to a desired
setting.\strut
\end{minipage}\tabularnewline
\begin{minipage}[t]{0.04\columnwidth}\raggedright
2\strut
\end{minipage} & \begin{minipage}[t]{0.23\columnwidth}\raggedright
On/Off button\strut
\end{minipage} & \begin{minipage}[t]{0.65\columnwidth}\raggedright
Use the speaker On/Off button to turn off the volume or turn it on to
maximum level.\strut
\end{minipage}\tabularnewline
\begin{minipage}[t]{0.04\columnwidth}\raggedright
3\strut
\end{minipage} & \begin{minipage}[t]{0.23\columnwidth}\raggedright
Save\strut
\end{minipage} & \begin{minipage}[t]{0.65\columnwidth}\raggedright
Use the \textbf{Save} button to apply and store new settings on alarm or
warning beeper.\strut
\end{minipage}\tabularnewline
\begin{minipage}[t]{0.04\columnwidth}\raggedright
4\strut
\end{minipage} & \begin{minipage}[t]{0.23\columnwidth}\raggedright
Back\strut
\end{minipage} & \begin{minipage}[t]{0.65\columnwidth}\raggedright
Use the Back button to exit the current page.\strut
\end{minipage}\tabularnewline
\bottomrule
\end{longtable}

\begin{longtable}[]{@{}l@{}}
\toprule
\begin{minipage}[b]{0.97\columnwidth}\raggedright
CAUTION\strut
\end{minipage}\tabularnewline
\midrule
\endhead
\begin{minipage}[t]{0.97\columnwidth}\raggedright
\textbf{It is strongly advised to keep the beep on whenever possible to
prevent any delay in the detection of an alarm or warning.} If the beep
is turned off, alarm and warning are not notified with sound and are
indicated only with red blinking of the operation lamp and on warning
occurrence screen. \textbf{Set the beep volume in accordance with the
ambient environment.}\strut
\end{minipage}\tabularnewline
\bottomrule
\end{longtable}


\section{Service}

The \textbf{Service} page is password protected. It is intended for use
by authorized ESPEC North America service personnel.

\begin{longtable}[]{@{}l@{}}
\toprule
\begin{minipage}[b]{0.97\columnwidth}\raggedright
CAUTION\strut
\end{minipage}\tabularnewline
\midrule
\endhead
\begin{minipage}[t]{0.97\columnwidth}\raggedright
\textbf{Do not use attempt to change settings in the Service page.}
Changing parameters at yours discretion may cause an accident. This
function is used by our service representatives to make various
adjustments. ESPEC is not liable for any accident and injury that may be
caused by the change of parameters at your discretion.\strut
\end{minipage}\tabularnewline
\bottomrule
\end{longtable}


\section{Macros}

Frequently used tasks can be automated by creating and running as a set
of instructions, called scripted macros. Macros are a series of scripted
commands and instructions grouped together to accomplish a certain task.
These scripted commands can be triggered automatically by the state of
the chamber or by an authorized operator through a manual manipulation.
Automated tasks through macro-scripted actions can range from sending
e-mail notification about test completion to synchronization of
operation between multiple chambers.

\textbf{\emph{Note:}} Macros operation page is inaccessible when
operated locally; this menu is grayed out on the HMI. It is accessible
and available for operation only via a remote access on a Web browser,
as shown in the figure below.

\begin{figure}
\centering

\includegraphics[width=0.95\textwidth]{figures/1027942748-S6_15-macros-web-001a.PNG}
\caption{Macros menu and its accessibility}
\end{figure}\FloatBarrier

When activated (via the button in the above figure), the \textbf{Macros}
operation page (with its classic user interface) is displayed and
operated completely inside the GL controller main display area, as
depicted in the following figure.

\begin{figure}
\centering

\includegraphics[width=0.95\textwidth]{figures/3442052129-S6_15-macros-ui-002.PNG}
\caption{Macros operation page}
\end{figure}\FloatBarrier

The following figure depicts the \textbf{Macros} setup page with a
default scripted action called \textbf{Alarm Emails}. The lock symbol
indicates that the contents of the \textbf{Alarm Emails} script cannot
be modified, since it was generated by the GL controller.

\begin{figure}
\centering

\includegraphics[width=0.95\textwidth]{figures/3517169944-macros-main-page-001a.PNG}
\caption{Macros main display page}
\end{figure}\FloatBarrier

The submenu of \textbf{Macros} consists of four main operation buttons
for managing and manipulating the macro scripted profiles.

\textbf{\emph{Procedure:}}

\begin{enumerate}
\def\labelenumi{\arabic{enumi}.}
\item
  \textbf{Create New}: A new macro script can be created via this
  \textbf{Create New} button.
\item
  \textbf{List of Macro Actions}: A list of all macro scripted profiles
  on the system. Click on its name to display its contents in the macro
  editor page (item 3).
\item
  \textbf{Macro Scripts}: The first macro script on the list (item 2) is
  listed in the main display by default. Its contents can be viewed
  using the macro editor (main display).
\item
  \textbf{Import from local file}: A macro scripted profile can be
  imported from the local computer. Apply this button to import a macro
  profile from the local computer. The macro editor will be launched to
  display the contents of the profile. The \textbf{Save As} button needs
  to be applied to save the imported profile; its name will appear under
  the list of macro actions (item 2).
\end{enumerate}

\textbf{\emph{Note:}} Many operations associated with the macros require
that the GL controller has access to the Internet.


\subsection{Macro Editor and Trigger
Options}

A macro script can be created to contain various trigger options. The
\textbf{Create New} button, when clicked, launches the macro editor,
within which the operator can compose the macro scripts to set different
alert and trigger options. The components of the macro editor are listed
as follows:

\begin{figure}
\centering

\includegraphics[width=0.95\textwidth]{figures/3811481287-P300-macro-editor-trigger-type-001a.PNG}
\caption{Macro trigger modes/options}
\end{figure}\FloatBarrier

\begin{enumerate}
\def\labelenumi{\arabic{enumi}.}
\item
  \textbf{Name}: A macro has a unique name to identify its action or
  task.
\item
  \textbf{Enable}: The macro action can be enabled or disabled. When
  enabled, trigger will take effect based on the chamber condition
  specified in the macro script.
\item
  \textbf{Trigger}: A macro may be triggered by any of the following
  types:

  \begin{itemize}
  \item
    \textbf{Always}: The macro will run every time. This type of trigger
    is not recommended.

    \begin{figure}
    \centering
    
\includegraphics[width=0.95\textwidth]{figures/81948110-P300-trigger-always-001.PNG}
    \caption{Trigger mode with always option}
    \end{figure}\FloatBarrier
  \item
    \textbf{Never (Manual Only)}: The macro must be manually triggered
    by an authorized operator or an API request.

    \begin{figure}
    \centering
    
\includegraphics[width=0.95\textwidth]{figures/3306130436-P300-macro-manual-trigger-001.PNG}
    \caption{Trigger mode with manual}
    \end{figure}\FloatBarrier
  \item
    \textbf{Program State}: The macro script will run when an execution
    state in the selected program has changed based on the parameters
    listed in the following figure.

    \begin{figure}
    \centering
    
\includegraphics[width=0.95\textwidth]{figures/2155600691-P300-macro-prog-state-001a.PNG}
    \caption{P300-macro-prog-state-001a.PNG}
    \end{figure}\FloatBarrier

    \begin{itemize}
    \tightlist
    \item
      \textbf{Program State}: Under the Program State options (drop-down
      menu in the above figure), a trigger alert can happen when there
      is a change in the program execution (designated as Program
      Changed), when a program has started or stopped (Program Started,
      Program Stopped); or a step within the executed program has
      changed (Step Changed), or a step has started or stopped (Step
      Started, Step Stopped). All of these options can be incorporated
      into the Program State trigger type.
    \item
      \textbf{Program}: A specific program may be selected to trigger
      the effect if the condition is met. If \textbf{Any} (default
      setting) is selected, any program will cause the trigger effect if
      the trigger condition during execution is met.
    \item
      \textbf{Step Number}: A specific step in the selected program can
      be used to trip the trigger effect if the condition is met, such
      as when step 5 in the selected program has completed its
      execution.
    \end{itemize}
  \item
    \textbf{Alarm State}: The macro will run when the state of an alarm
    has changed. The parameters that specify the alarm state are listed
    in the drop-down menu shown in the following figure. The alarm list
    is chamber and PLC dependent.

    \begin{figure}
    \centering
    
\includegraphics[width=0.95\textwidth]{figures/2246722414-P300-trigger-alarm-state-001a.PNG}
    \caption{Trigger alarm options}
    \end{figure}\FloatBarrier
  \item
    \textbf{Date/Time Trigger}: The macro will run at a specified time
    or date and time with periodic operation. When the date/time matches
    the configured ``Month'', ``Day of the Month'', ``Year'', ``Day of
    the Week'', ``Hour'', ``Minute'', and ``Second'' the macro will
    fire. This operation can be configured for a one-time trigger or a
    periodic trigger, as depicted in the following figures.

    
\includegraphics[width=0.95\textwidth]{figures/4216476246-P300-macro-date-time-trigger-one-shot-002.PNG}
    
\includegraphics[width=0.95\textwidth]{figures/3675414021-P300-trigger-datetime-periodic-001.PNG}
  \item
    Additional trigger types can be selected from the trigger list that
    include time signals, loop temperature, custom expression and
    logical operations. All of which have the same programming or
    scripting paradigm based on specific parameters and trigger
    conditions. A single macro script can be created to monitor a list
    of programs, their status or conditions, using a set of complex
    logical operations selected from the list of trigger types (such as,
    \textbf{and}, \textbf{or}, \textbf{not}).
  \end{itemize}
\item
  \textbf{Multiple Trigger Types}: Additional trigger types can be added
  via the (+/-) button, with options to insert additional trigger type
  using \textbf{Insert Before} or \textbf{Insert After} buttons. Trigger
  type can be removed from the list via the \textbf{Delete} button.

  \begin{figure}
  \centering
  
\includegraphics[width=0.95\textwidth]{figures/3590729099-P300-macro-3-trigger-001.PNG}
  \caption{Three trigger types in macro script}
  \end{figure}\FloatBarrier
\item
  \textbf{APPEND}: The action or actions of the trigger (item 3, above)
  can be implemented in the body of the \textbf{Operations} template.
  The \textbf{APPEND} button can be applied to add and compose the
  trigger operations. The components of the trigger operations are
  outlined as follows:

  \begin{figure}
  \centering
  
\includegraphics[width=0.95\textwidth]{figures/4157574869-P300-macro-operations-002a.PNG}
  \caption{Components of trigger oeprations}
  \end{figure}\FloatBarrier

  \begin{longtable}[]{@{}lll@{}}
  \toprule
  \begin{minipage}[b]{0.04\columnwidth}\raggedright
  Number\strut
  \end{minipage} & \begin{minipage}[b]{0.23\columnwidth}\raggedright
  Name\strut
  \end{minipage} & \begin{minipage}[b]{0.65\columnwidth}\raggedright
  Description/Operation\strut
  \end{minipage}\tabularnewline
  \midrule
  \endhead
  \begin{minipage}[t]{0.04\columnwidth}\raggedright
  a\strut
  \end{minipage} & \begin{minipage}[t]{0.23\columnwidth}\raggedright
  Conditional Statement\strut
  \end{minipage} & \begin{minipage}[t]{0.65\columnwidth}\raggedright
  The conditional statements of trigger operations consist of
  \textbf{If}, \textbf{Else If}, \textbf{Else} and \textbf{Wait For}.
  These conditional statements can be used to check the type(s) of
  trigger operation.\strut
  \end{minipage}\tabularnewline
  \begin{minipage}[t]{0.04\columnwidth}\raggedright
  b\strut
  \end{minipage} & \begin{minipage}[t]{0.23\columnwidth}\raggedright
  Type of Condition\strut
  \end{minipage} & \begin{minipage}[t]{0.65\columnwidth}\raggedright
  The type of trigger used in the conditional statement.\strut
  \end{minipage}\tabularnewline
  \begin{minipage}[t]{0.04\columnwidth}\raggedright
  c\strut
  \end{minipage} & \begin{minipage}[t]{0.23\columnwidth}\raggedright
  Additional Conditions\strut
  \end{minipage} & \begin{minipage}[t]{0.65\columnwidth}\raggedright
  Additional conditions can be added via (+/-) button. The available
  options are: Insert Before (current type), Insert After (current
  type), Delete.\strut
  \end{minipage}\tabularnewline
  \begin{minipage}[t]{0.04\columnwidth}\raggedright
  d\strut
  \end{minipage} & \begin{minipage}[t]{0.23\columnwidth}\raggedright
  Alert Method\strut
  \end{minipage} & \begin{minipage}[t]{0.65\columnwidth}\raggedright
  Available alert methods for the trigger operation. Default option is:
  Mail: Send a custom e-mail.\strut
  \end{minipage}\tabularnewline
  \begin{minipage}[t]{0.04\columnwidth}\raggedright
  e\strut
  \end{minipage} & \begin{minipage}[t]{0.23\columnwidth}\raggedright
  E-mail Address\strut
  \end{minipage} & \begin{minipage}[t]{0.65\columnwidth}\raggedright
  The operator's e-mail address in this block will be used to send an
  alert e-mail. Multiple e-mail addresses can used, one e-mail address
  per line.\strut
  \end{minipage}\tabularnewline
  \begin{minipage}[t]{0.04\columnwidth}\raggedright
  f\strut
  \end{minipage} & \begin{minipage}[t]{0.23\columnwidth}\raggedright
  Subject title\strut
  \end{minipage} & \begin{minipage}[t]{0.65\columnwidth}\raggedright
  A descriptive subject title is important in an alert e-mail.\strut
  \end{minipage}\tabularnewline
  \begin{minipage}[t]{0.04\columnwidth}\raggedright
  g\strut
  \end{minipage} & \begin{minipage}[t]{0.23\columnwidth}\raggedright
  Message\strut
  \end{minipage} & \begin{minipage}[t]{0.65\columnwidth}\raggedright
  The message in the e-mail should be brief and descriptive.\strut
  \end{minipage}\tabularnewline
  \begin{minipage}[t]{0.04\columnwidth}\raggedright
  h\strut
  \end{minipage} & \begin{minipage}[t]{0.23\columnwidth}\raggedright
  Manage Trigger Operations Steps\strut
  \end{minipage} & \begin{minipage}[t]{0.65\columnwidth}\raggedright
  With this button, management of trigger operations steps is possible
  with the \textbf{Inset After}, \textbf{Insert Before} or
  \textbf{Delete} options.\strut
  \end{minipage}\tabularnewline
  \begin{minipage}[t]{0.04\columnwidth}\raggedright
  i\strut
  \end{minipage} & \begin{minipage}[t]{0.23\columnwidth}\raggedright
  Add Trigger Operations Step\strut
  \end{minipage} & \begin{minipage}[t]{0.65\columnwidth}\raggedright
  The \textbf{APPEND} button appends an additional step at the bottom of
  the step list.\strut
  \end{minipage}\tabularnewline
  \bottomrule
  \end{longtable}
\item
  \textbf{Operations}: Components of the trigger \textbf{Operations} are
  outlined in the previous and the following figure. The default
  operation type is \textbf{E-Mail: Send a custom email}. Different
  operation types can be selected from the list. Each selected type is
  used in the conditional statement to trip the operation. Multiple
  types can be implemented by applying the (+/-) button. The operation
  type thus consists of three fields: (1) who to send the alert, (2)
  subject of the alert, (3) message of the alert. Multiple operations
  can be added using the \textbf{APPEND} button to create additional or
  multiple operations.

  \begin{figure}
  \centering
  
\includegraphics[width=0.95\textwidth]{figures/2901991376-P300-macro-operation-001c.PNG}
  \caption{Options of trigger components}
  \end{figure}\FloatBarrier
\item
  \textbf{File Manipulation}: Three file manipulation options are
  available in the main macro editor when the \textbf{Create New} option
  is applied. These are \textbf{Export to local file}, \textbf{Import
  from local file} and \textbf{Save}.

  \begin{figure}
  \centering
  
\includegraphics[width=0.95\textwidth]{figures/4293313447-P300-macro-file-buttons-001.PNG}
  \caption{File manipulation buttons}
  \end{figure}\FloatBarrier

  After the macro script has been composed and saved, additional options
  are available as follows:

  \begin{figure}
  \centering
  
\includegraphics[width=0.95\textwidth]{figures/2811041385-P300-macro-file-buttons-002.PNG}
  \caption{File manipulation options for macro editing}
  \end{figure}\FloatBarrier

  The current macro script can be deleted from the macro editor with the
  \textbf{Delete} button (trash bin). A warning appears to reconfirm the
  action.

  \begin{figure}
  \centering
  
\includegraphics[width=0.95\textwidth]{figures/2949811339-macro-delete-001.PNG}
  \caption{Reconfirm the deletion of macro script}
  \end{figure}\FloatBarrier

  This script can be invoked to test its operation by applying the
  \textbf{person} icon. If the macro operation involves sending out an
  alert e-mail, the recipient on the e-mail list will immediately
  receive the e-mail alert. The GL controller can send out e-mail alerts
  only if it has access to the Internet because, by default, it uses
  SMTP Office 365 for the email protocol.
\end{enumerate}


\subsection{Example: A Macro Script with Alarm
Alert}

The following example illustrates a simple macro script to send out an
e-mail alert when the chamber trips an alarm (any alarm). In order for
the macro setup to work, the GL controller must have access to the
Internet. The procedure for this example is outlined as follows.

\textbf{\emph{Procedure:}}

\begin{enumerate}
\def\labelenumi{\arabic{enumi}.}
\item
  Click \textbf{Macros} in the \textbf{Settings} submenu.
\item
  Click \textbf{Create New}.
\item
  Click the Name field in the macro editor, enter \textbf{ALARM01} for
  the macro name. Confirm that the trigger is enabled (with its box
  checked).
\item
  Click the type field under the \textbf{Trigger} type and select
  \textbf{Alarm State} from the list (as shown).

  \begin{figure}
  \centering
  
\includegraphics[width=0.95\textwidth]{figures/3382272209-P300-alarm-state-type-001.PNG}
  \caption{Select Alarm State type from the list}
  \end{figure}\FloatBarrier
\item
  Click the \textbf{Alarm State} field and select \textbf{Tripped} from
  the drop-down list as its parameter (as shown).

  \begin{figure}
  \centering
  
\includegraphics[width=0.95\textwidth]{figures/1086519463-P300-alarm-state-type-002.PNG}
  \caption{Set trip parameter for the Alarm State}
  \end{figure}\FloatBarrier
\item
  Confirm that \textbf{Any} is selected under the Alarms option. The
  complete selection type, state and alarms options is depicted below.

  \begin{figure}
  \centering
  
\includegraphics[width=0.95\textwidth]{figures/3946241683-P300-alarm-state-type-003.PNG}
  \caption{Parameters of alarm trigger type}
  \end{figure}\FloatBarrier
\item
  Click \textbf{APPEND} to add the operations instructions.
\item
  Confirm that the logic \textbf{If} is selected by default to check the
  condition, and the condition type \textbf{Always} is select. This is
  to ensure that an alert will be sent out whenever an alarm is tripped.
\item
  Confirm, under the Operation Type, that ``E-Mail: Send a custom
  email'' has been selected.
\item
  In the \textbf{Recipients} block, enter the operator's e-mail address.
  Multiple e-mail addresses can be used, with one e-mail address per
  line.
\item
  Enter the subject title of the e-mail in the \textbf{Subject}.
\item
  Enter the message to be included in the e-mail in the \textbf{Message
  body} box.
\item
  Click the \textbf{Save} button. The macro list now has the macro
  script listed by its file name. If you attempt to exit the macro
  editor (by clicking on other submenus or menus), a warning message
  will appear (as shown below).

  \begin{figure}
  \centering
  
\includegraphics[width=0.95\textwidth]{figures/235526585-const-temp-setting-f4t-warning-00.PNG}
  \caption{Macro script must be save before exiting the pane}
  \end{figure}\FloatBarrier
\item
  To test the macro script, click the \textbf{person} icon. The operator
  should receive an e-mail alert from the GL controller.
\end{enumerate}

The complete macro script is depicted in the following figure.

\begin{figure}
\centering

\includegraphics[width=0.95\textwidth]{figures/2534645778-P300-macro-operation-002a.PNG}
\caption{Macro profile with Alarm alert}
\end{figure}\FloatBarrier


\subsection{Example: Macro Script with Program State and Alarm
Alert}

Here is an example of the use of two logical \textbf{Or} conditions of
the trigger type to monitor two specific programs (\textbf{PROGTEST001}
and \textbf{PROG3TEST}) and trigger the e-mail operation. Two logical
\textbf{If} conditions monitor the program state and alarms; an alert is
sent out accordingly. The first \textbf{If} condition is set up to
monitor the program activity; the second \textbf{If} condition monitors
any alarm occurrences. \textbf{\emph{Note:}} This sample macro script is
presented to illustrate the flexibility of the macro editor. There are
numerous ways to achieve the same the task noted here.

\textbf{\emph{Procedure:}} To construct the macro script, proceed as
follows.

\begin{enumerate}
\def\labelenumi{\arabic{enumi}.}
\item
  Click \textbf{Create New} on the \textbf{Macro} submenu.
\item
  Click the Name field and enter \textbf{ALARM02}.
\item
  Confirm that the trigger is enabled (with its box checked).
\item
  Under \textbf{Trigger}:

  \begin{itemize}
  \item
    Click the Type field and select \textbf{Program State} from the
    drop-down list.
  \item
    Click the Program State field and select \textbf{Program Started}
    from the drop-down list.
  \item
    Click the Program field and select \textbf{PROG2TEST} on slot 2.
    \textbf{\emph{Note:}}: \textbf{PROG2TEST} must be made available on
    the program list.
  \item
    The last option on step number is set as default for the trigger to
    apply at any step.
  \item
    Click the (+/-) and select \textbf{Insert After} to add a new
    trigger type, as shown in the figure.

    \begin{figure}
    \centering
    
\includegraphics[width=0.95\textwidth]{figures/1688259577-macro-example-002a.PNG}
    \caption{Inserting additional trigger type}
    \end{figure}\FloatBarrier
  \item
    Click the Type field and select \textbf{Or} from the drop-down list.
  \item
    Click the Type field and select \textbf{Program State} from the
    drop-down list.
  \item
    Click the Program State field and select \textbf{Program Started}
    from the drop-down list.
  \item
    Click the Program field and select \textbf{PROG3TEST} on slot 3.
    \textbf{\emph{Note:}} \textbf{PROG3TEST} must be created and is
    available on the program list.
  \item
    The last option on step number is set as default for the trigger to
    apply at any step.
  \item
    Apply the (+/-) button to add two more trigger type steps to contain
    the logical \textbf{Or} and \textbf{Alarm State} as depicted in the
    following figure which depicts the complete configuration of the
    conditional statements in the trigger logic:

    \begin{figure}
    \centering
    
\includegraphics[width=0.95\textwidth]{figures/1707337315-macro-example-003.PNG}
    \caption{Trigger setup logic}
    \end{figure}\FloatBarrier
  \end{itemize}
\item
  Click \textbf{APPEND} under the \textbf{Operations} block to create
  the first step of the trigger operation.

  \begin{itemize}
  \item
    Confirm that \textbf{If} condition is selected (by default).
  \item
    Click the Type field and select \textbf{Always} from the list. This
    is to ensure that an e-mail alert will be sent out if the state of a
    specified program occurs.
  \item
    Confirm that \textbf{E-Mail: Send a custom email} is selected (by
    default) for Operation Type.
  \item
    Enter the recipient's e-mail address in the Recipients block, one
    e-mail address per line.
  \item
    Enter subject title in the Subject block.
  \item
    Enter a brief message in the Message Body block. The complete
    configuration is illustrated as follows:

    \begin{figure}
    \centering
    
\includegraphics[width=0.95\textwidth]{figures/2572712087-macro-example-step2-002.PNG}
    \caption{Conditions for trigger actions}
    \end{figure}\FloatBarrier
  \end{itemize}
\item
  Click the operation step number (number 1 in the circle, as shown
  below) and select \textbf{Insert After} to add a new step to the
  trigger operation.

  \begin{figure}
  \centering
  
\includegraphics[width=0.95\textwidth]{figures/786399888-macro-example-step2-003a.PNG}
  \caption{Inserting additional step of trigger operation}
  \end{figure}\FloatBarrier

  \begin{itemize}
  \tightlist
  \item
    Confirm that \textbf{If} condition is selected (by default).
  \item
    Click the Type field and select \textbf{Alarm State} from the list.
  \item
    Click the Program State field and select \textbf{Tripped} from the
    list.
  \item
    Confirm that \textbf{Any} is selected by default under the Alarms
    field; it is to ensure any alarm will trigger this action.
  \item
    Confirm that \textbf{E-Mail: Send a custom email} is selected (by
    default) for Operation Type.
  \item
    Enter the recipient's e-mail address in the Recipients block, one
    e-mail address per line.
  \item
    Enter subject title in the Subject block.
  \item
    Enter a brief message in the Message Body block to indicate alarm
    has been tripped. The complete configuration is illustrated as
    follows:
  \end{itemize}

  \begin{figure}
  \centering
  
\includegraphics[width=0.95\textwidth]{figures/3090880437-macro-example-step2-003.PNG}
  \caption{Conditions for step trigger operation}
  \end{figure}\FloatBarrier
\item
  Click \textbf{Save} to save this macro scripted profile.
\item
  \textbf{Testing the Macro Script}: Click the \textbf{person} icon to
  test and run the macro script.

  \begin{itemize}
  \item
    If the run is successful (that is, no errors are found in the
    script), a \textbf{Successful} message appears and an e-mail alert
    is sent to the recipient(s) listed in the Recipients block.

    \begin{figure}
    \centering
    
\includegraphics[width=0.95\textwidth]{figures/4092544407-run-macro-003.PNG}
    \caption{A macro script ran successfully}
    \end{figure}\FloatBarrier
  \item
    If the run failed, an error message appears as shown below. The
    macro script does not work.

    \begin{figure}
    \centering
    
\includegraphics[width=0.95\textwidth]{figures/349251575-run-macro-error-003.PNG}
    \caption{A macro script failed, errors found in the script}
    \end{figure}\FloatBarrier
  \end{itemize}
\end{enumerate}


\chapter{Network Setup \& Remote
Access}

This chapter discusses the networking features of the GL controller
operation as a networked device. It explains how the GL controller can
join a network, how to configure a different type of network protocol
for the GL controller and how to access the GL controller via its
hostname or its IP address on a Web browser.


\section{GL Controller on a
Network}

As stated in Section 2.1, the GL controller is capable of operating as a
standalone system or as a networked system. When the GL controller is
connected to a network, its UI can be accessed via a Web browser from
any device (e.g., PC, laptop, handheld, smartphone) on the network using
either the IP address or the hostname of the GL controller. The entire
GL controller system and chamber can be controlled remotely via this Web
browser. The following figure depicts the GL controller UI on Google
Chrome Web browser, accessed via its IP address.

\begin{figure}
\centering

\includegraphics[width=0.95\textwidth]{figures/3796224613-S2_1_2N-gl-homepage-network-003a.PNG}
\caption{GL controller as a networked system}
\end{figure}\FloatBarrier

The GL controller is programmed to join a DHCP (Dynamic Host
Configuration Protocol) network using an IP address provided by the DHCP
server. On a network that does not support DHCP, the GL controller uses
its preconfigured IP address (called fallback IP) to join the network;
this fallback IP address is based on a Class C private network
configuration. The GL controller can also be configured using a specific
static protocol with static IP address to join any static network
domain.

\begin{longtable}[]{@{}ll@{}}
\toprule
\begin{minipage}[b]{0.22\columnwidth}\raggedright
Network Type\strut
\end{minipage} & \begin{minipage}[b]{0.72\columnwidth}\raggedright
Description\strut
\end{minipage}\tabularnewline
\midrule
\endhead
\begin{minipage}[t]{0.22\columnwidth}\raggedright
DHCP\strut
\end{minipage} & \begin{minipage}[t]{0.72\columnwidth}\raggedright
By default, GL controller joins a DHCP network by requesting (and using)
an IP address from the DHCP server. Refer to Section 8.3.\strut
\end{minipage}\tabularnewline
\begin{minipage}[t]{0.22\columnwidth}\raggedright
Default Static\strut
\end{minipage} & \begin{minipage}[t]{0.72\columnwidth}\raggedright
By default, GL controller uses its preconfigured IP address
(192.168.0.83) when DHCP server is not available. Refer to Section
8.4.\strut
\end{minipage}\tabularnewline
\begin{minipage}[t]{0.22\columnwidth}\raggedright
Custom Static\strut
\end{minipage} & \begin{minipage}[t]{0.72\columnwidth}\raggedright
GL controller can be configured to use any static IP address in any
network domain. Refer to Section 8.5.\strut
\end{minipage}\tabularnewline
\bottomrule
\end{longtable}

Remote access to the GL controller system as a networked device offers
many potential advantages. Here are few examples.

\textbf{\emph{Summary:}}

\begin{itemize}
\item
  GL controller can be controlled and operated remotely
\item
  Data logs can be downloaded and stored on the operator's PC for
  analysis
\item
  Multiple users may access and view chamber operation status from their
  PC at the same time
\item
  Only authorized users on the network can access and control the GL
  controller system
\item
  User interface and operation are made easier via a Web browser; no
  need to utilize the on-screen keyboard/keypad for control inputs.
\item
  Automatic Over-the-Air update via Mender IO cloud service is done
  without the user's intervention, thus keeping the GL controller system
  up to date.
\item
  GL controller system is kept secure and protected through the
  implementation of the user policy and separate user accounts.
\item
  Alert notifications via email can be sent out to all or specific
  operators about the status or conditions of the chamber, such as
  warning or alarm.
\end{itemize}


\section{Web Browser
Compatibility}

The GL controller supports the following Web browsers: Chromium,
Google-Chrome, Mozilla Firefox, Microsoft Edge, Apple Safari and Opera.
Microsoft Internet Explorer 11 can also be used to access and operate
the GL controller. However, due to its slow performance, the use of
Microsoft Internet Explorer 11 is strongly discouraged.


\section{DHCP Network}

The following procedure illustrates how to connect the GL controller
system to a DHCP network, look up its hostname and its IP address, and
access its UI via a Web browser from a laptop or PC on the network.

\textbf{\emph{Procedure:}}

\begin{enumerate}
\def\labelenumi{\arabic{enumi}.}
\item
  Plug a Cat5 or Cat6 Ethernet cable into the RJ-45 Ethernet port at the
  rear of the chamber (see arrow); this can be done while the GL
  controller system is still running. \textbf{\emph{Note:}} Ethernet
  connection on the GL controller, like any other network devices, is
  hot-pluggable.

  \begin{figure}
  \centering
  
\includegraphics[width=0.95\textwidth]{figures/4166753972-S7-step2-001a.PNG}
  \caption{Connecting Ethernet cable}
  \end{figure}\FloatBarrier
\item
  Plug the other end of the cable into an Ethernet port (or a router)
  that connects to the main network.

  \begin{figure}
  \centering
  
\includegraphics[width=0.95\textwidth]{figures/1766916705-S8_3N-step2-001a.PNG}
  \caption{DHCP network setup}
  \end{figure}\FloatBarrier
\item
  After about 30 seconds, press the hamburger menu in the status bar
  (1), followed by Configuration (2) and Set Communication (3) in the
  system menu, as depicted in the following figure.

  \begin{figure}
  \centering
  
\includegraphics[width=0.95\textwidth]{figures/2575102693-S7_3_1-step3-002a.PNG}
  \caption{Looking up GL controller hostname and IP address}
  \end{figure}\FloatBarrier
\item
  The hostname of this GL controller system is \textbf{GL-EPL-3-P300},
  indicated by callout (1). The hostname of your GL controller will be
  different and unique, which is associated with the registered serial
  number. The IP address of the GL controller is listed next to
  \textbf{IPv4 Address}, indicated by callout (2); it is shown as
  \textbf{10.30.100.190}. \textbf{\emph{Note:}} The IP address of the GL
  controller is assigned to \textbf{eth0}. Thus, the IP address under
  the \textbf{eth0} column is the one used by the GL controller to join
  the main network. The IP address (192.168.1.100) assigned to (and
  listed under) \textbf{eth1} is for internal communication between the
  GL controller and the PLC.

  \begin{figure}
  \centering
  
\includegraphics[width=0.95\textwidth]{figures/1833686550-S7_3_1-step4-002a.PNG}
  \caption{Hostname and IP address of GL controller}
  \end{figure}\FloatBarrier
\item
  Connect a PC or a laptop based on the diagram illustrated in step 2.
  Turn on the computer and confirm that it is connected to the same
  network that the GL controller is a part of.
\item
  Launch a Web browser on the computer.

  \begin{itemize}
  \item
    To access the GL controller via its IP address, enter:
    http://10.30.100.190/ in the URL of the Web browser and press Enter.
    The home page of the GL controller will appear.
  \item
    To access the GL controller via its hostname, enter:
    http://GL-EPL-3-P300.local/ in the URL of the Web browser and press
    Enter. The home page of the GL controller will appear.
  \end{itemize}
\item
  The general URL of the GL controller is:

  \begin{itemize}
  \tightlist
  \item
    http://IP-address/
  \item
    http://hostname.local/
  \end{itemize}
\end{enumerate}


\section{Static Network using Fallback
IP}

This section outlines the steps to set up a private static network
connection between the GL controller system and a PC or a group of PC's
via a network hub where no DHCP is available. The goal is that the GL
controller UI can be accessed from any PC on the network.


\subsection{Static Network Connection and
Configuration}

By default, if the GL controller detects a network signal on its
\textbf{eth0} Ethernet port, it first applies DHCP protocol to request
an IP address. If the request fails, the preconfigured IP address
192.168.0.83, called fallback IP, is used to join the network.

\begin{figure}
\centering

\includegraphics[width=0.95\textwidth]{figures/2985967273-S7_4_1-gl-fallback-cfg-001a.PNG}
\caption{Fallback IP address}
\end{figure}\FloatBarrier

\begin{longtable}[]{@{}lll@{}}
\toprule
\begin{minipage}[b]{0.04\columnwidth}\raggedright
No.\strut
\end{minipage} & \begin{minipage}[b]{0.23\columnwidth}\raggedright
Category\strut
\end{minipage} & \begin{minipage}[b]{0.65\columnwidth}\raggedright
Description\strut
\end{minipage}\tabularnewline
\midrule
\endhead
\begin{minipage}[t]{0.04\columnwidth}\raggedright
1\strut
\end{minipage} & \begin{minipage}[t]{0.23\columnwidth}\raggedright
DHCP\strut
\end{minipage} & \begin{minipage}[t]{0.65\columnwidth}\raggedright
DHCP is enabled by default (from factory). Fallback IP address is used
in DHCP mode. Refer to Section 8.5 for custom static setting.\strut
\end{minipage}\tabularnewline
\begin{minipage}[t]{0.04\columnwidth}\raggedright
2\strut
\end{minipage} & \begin{minipage}[t]{0.23\columnwidth}\raggedright
Network protocol\strut
\end{minipage} & \begin{minipage}[t]{0.65\columnwidth}\raggedright
Displays factory settings for fallback IP network configuration; they
must not be changed. Refer to Section 8.5 for custom static network
setting.\strut
\end{minipage}\tabularnewline
\begin{minipage}[t]{0.04\columnwidth}\raggedright
3\strut
\end{minipage} & \begin{minipage}[t]{0.23\columnwidth}\raggedright
Fallback setting\strut
\end{minipage} & \begin{minipage}[t]{0.65\columnwidth}\raggedright
Displays notification of the fallback IP address when DHCP request
failed.\strut
\end{minipage}\tabularnewline
\bottomrule
\end{longtable}

The client PC must be configured to use the same Class C network with
the same subnet mask and gateway as follows.

\begin{longtable}[]{@{}ll@{}}
\toprule
Category & Parameter\tabularnewline
\midrule
\endhead
IP Address & 192.168.0.84 (recommended)\tabularnewline
Subnet Mask & 255.255.255.0\tabularnewline
Gateway & 192.168.0.1\tabularnewline
Preferred DNS server & 8.8.8.8\tabularnewline
Alternate DNS server & 8.8.4.4\tabularnewline
\bottomrule
\end{longtable}

The following procedure outlines the steps to configuration a static
network protocol on a PC running MS Windows 8/10/11.

\textbf{\emph{Procedure:}}

\begin{enumerate}
\def\labelenumi{\arabic{enumi}.}
\item
  Repeat step 1 in the procedure of Section 8.3 (above).
\item
  Set up Ethernet connection as depicted in the diagram.

  \begin{figure}
  \centering
  
\includegraphics[width=0.95\textwidth]{figures/3958540732-S8_4_1N-gl-static-class-C-network-001a.PNG}
  \caption{Static network configuration}
  \end{figure}\FloatBarrier
\item
  With the PC turned on, hold down the \textbf{Windows} key and press
  \textbf{R} to launch the Run Command dialog box.
\item
  In the Run dialog box, enter \textbf{ncpa.cpl} into the Open box field
  and press \textbf{Enter}.
\item
  Point and Right-Click the ``Local Area Connection'' icon, then click
  Properties from the drop-down menu (follow the arrows in the figure).
  \textbf{\emph{Note:}} The Local Area Connection icon is the one
  connected to the GL controller (or a hub). It is important to access
  the correct icon in case the PC has multiple Ethernet ports.

  \begin{figure}
  \centering
  
\includegraphics[width=0.95\textwidth]{figures/316319689-static-cfg-1a.PNG}
  \caption{Selecting the right Local Area Connection}
  \end{figure}\FloatBarrier
\item
  In the ``Local Area Connection Properties'' window, confirm that the
  check mark is placed in front of ``Internet Protocol Version 4
  (TCP/IPv4)'', as depicted in the figure. If not, check it. Click to
  highlight ``Internet Protocol Version 4 (TCP/IPv4)'' and then click
  Properties in the lower-right corner.

  \begin{figure}
  \centering
  
\includegraphics[width=0.95\textwidth]{figures/2493799905-static-cfg-2a.PNG}
  \caption{Setting TCP/IPv4 properties}
  \end{figure}\FloatBarrier
\item
  In the ``Internet Protocol Version (TCP/IPv4) Properties'' window,
  turn on the radio button for ``Use the following IP address:'' and
  enter the IP address (192.168.0.84), subnet mask (255.255.255.0) and
  gateway (192.168.0.1), as shown in the following figure.
\item
  In the ``Use the following DNS server addresses:'' section, enter the
  following DNS server values (also shown in the following figure):

  \begin{itemize}
  \item
    \textbf{Preferred DNS server}: 8.8.8.8
  \item
    \textbf{Alternate DNS server}: 8.8.4.4
  \end{itemize}
\item
  Turn on ``Validate settings upon exit'' with a check mark and click
  OK, as illustrated in the following figure.

  \begin{figure}
  \centering
  
\includegraphics[width=0.95\textwidth]{figures/2958266862-static-cfg-3a.PNG}
  \caption{The complete static IP config on the TCP/IPv4 connection}
  \end{figure}\FloatBarrier
\item
  Click OK to close ``Local Area Connection Properties'' window. Close
  out the Network window.
\item
  Returning to the GL controller: On the GL controller HMI, press the
  hamburger menu in the status bar (1), followed by Configuration (2)
  and Set Communication (3) in the system menu, as depicted below.

  \begin{figure}
  \centering
  
\includegraphics[width=0.95\textwidth]{figures/3032867281-S7_4_1-step11-001a.PNG}
  \caption{Looking up GL controller hostname and IP address}
  \end{figure}\FloatBarrier
\item
  The fallback IP address used by the GL controller is listed next to
  \textbf{IPv4 Address} under \textbf{eth0} (see arrow), shown as
  \textbf{192.168.0.83}. \textbf{\emph{Note:}} The GL controller IP
  address is assigned to \textbf{eth0}. The IP address (192.168.1.100)
  assigned to (and listed under) \textbf{eth1} is for internal
  communication between the GL controller and the PLC.

  \begin{figure}
  \centering
  
\includegraphics[width=0.95\textwidth]{figures/993438973-S7_4_1-fallback-ip-001a.PNG}
  \caption{Fallback IP address used by GL}
  \end{figure}\FloatBarrier
\item
  To access the GL controller, open a Web browser on your PC and
  navigate to: http://192.168.0.83/
\end{enumerate}


\subsection{Static Network via Network
Hub}

If multiple PCs are required to access and control the chamber, a
network hub can be used to connect all PCs with the GL controller, as
depicted in the diagram. The GL controller system will use its fallback
IP 192.168.0.83. All PCs on the network must use IP address 192.168.0.x,
where x is unique for each PC.

\begin{figure}
\centering

\includegraphics[width=0.95\textwidth]{figures/414296829-S8_4_2N-gl-fallback-static-IP-001a.PNG}
\caption{GL controller static network with multiple PCs}
\end{figure}\FloatBarrier

\textbf{\emph{Procedure:}}

\begin{enumerate}
\def\labelenumi{\arabic{enumi}.}
\tightlist
\item
  Repeat step 1 in the procedure of Section 8.3 (above) and set up
  Ethernet connection as depicted in the above diagram.
\item
  Follow steps 3-10 in Section 8.4.1 to set up IP address for each PC on
  the network.
\item
  To access the GL controller on each PC, open a Web browser on the PC
  and navigate to: http://192.168.0.83/
\end{enumerate}


\section{Custom Static Network}

For security reasons and to limit access to the network domain, your IT
may control what device can join the network. This means no device can
simply join the network without IT's approval. Each device must have a
unique static IP address assigned by the IT team. The GL controller
system can be configured with a unique static IP address to join such a
network; it uses a custom static IP address with network settings given
by your IT, as depicted in the following table.

\begin{longtable}[]{@{}ll@{}}
\toprule
\begin{minipage}[b]{0.22\columnwidth}\raggedright
Name\strut
\end{minipage} & \begin{minipage}[b]{0.72\columnwidth}\raggedright
Description\strut
\end{minipage}\tabularnewline
\midrule
\endhead
\begin{minipage}[t]{0.22\columnwidth}\raggedright
IP Address\strut
\end{minipage} & \begin{minipage}[t]{0.72\columnwidth}\raggedright
Unique IP address for the GL Controller, assigned by IT. Example:
10.30.200.247\strut
\end{minipage}\tabularnewline
\begin{minipage}[t]{0.22\columnwidth}\raggedright
Subnet Mask\strut
\end{minipage} & \begin{minipage}[t]{0.72\columnwidth}\raggedright
Correct subnet to help identify the network and host device on the
network. Example: 255.255.0.0\strut
\end{minipage}\tabularnewline
\begin{minipage}[t]{0.22\columnwidth}\raggedright
Gateway\strut
\end{minipage} & \begin{minipage}[t]{0.72\columnwidth}\raggedright
Correct default gateway to allow network access. Example:
10.30.0.1\strut
\end{minipage}\tabularnewline
\begin{minipage}[t]{0.22\columnwidth}\raggedright
DNS1\strut
\end{minipage} & \begin{minipage}[t]{0.72\columnwidth}\raggedright
Default primary domain name server to resolve hostnames (optional).
Example: 10.30.30.31 or 1.1.1.1\strut
\end{minipage}\tabularnewline
\begin{minipage}[t]{0.22\columnwidth}\raggedright
DNS2\strut
\end{minipage} & \begin{minipage}[t]{0.72\columnwidth}\raggedright
Default secondary domain name server to resolve hostnames (optional).
Example: 10.30.30.23 or 8.8.8.8\strut
\end{minipage}\tabularnewline
\bottomrule
\end{longtable}

The fallback IP configuration in the Set LAN page will be replaced with
that given by your IT. The IP configuration in the table will be used as
an example in the following step-by-step procedure.

\textbf{\emph{Procedure:}}

\begin{enumerate}
\def\labelenumi{\arabic{enumi}.}
\item
  Repeat step 1 in the procedure of Section 8.3 (above).
\item
  Set up Ethernet connection as depicted in the following diagram.

  \begin{figure}
  \centering
  
\includegraphics[width=0.95\textwidth]{figures/2178670663-S8_5N-gl-custom-static-step2-001a.PNG}
  \caption{Static network configuration}
  \end{figure}\FloatBarrier
\item
  On the GL controller HMI, press the hamburger button (1) in the status
  bar. Check to confirm the floating keyboard is enabled under the
  system menu. A diagonal bar over the keyboard indicates it is
  disabled. Press the keyboard (2) to enable it; press Configuration
  (3), followed by Set Communication (4), as depicted in the figure.

  \begin{figure}
  \centering
  
\includegraphics[width=0.95\textwidth]{figures/956115099-S7_5-step3-002a.PNG}
  \caption{Looking up GL controller hostname and IP address}
  \end{figure}\FloatBarrier
\item
  Press Set LAN (Network Setup) page, then uncheck the DHCP box (see
  arrow 1); enter \textbf{10.30.200.247} in IP address field (see arrow
  2). Enter the rest of the parameters listed in the table above, as
  depicted in the figure. Press Save (arrow 3) and confirm the action.

  \begin{figure}
  \centering
  
\includegraphics[width=0.95\textwidth]{figures/3913041737-S7_5-step4-custom-static-001aa.PNG}
  \caption{Configure custom static IP}
  \end{figure}\FloatBarrier
\item
  Back on the Set Communication page, confirm that your custom IP
  setting is displayed under the \textbf{eth0} column, as depicted in
  the following figure.

  \begin{figure}
  \centering
  
\includegraphics[width=0.95\textwidth]{figures/1286559998-S7_5-step5-001a.PNG}
  \caption{Confirming static IP address}
  \end{figure}\FloatBarrier
\item
  Confirm that the computer on the network uses the same network
  protocol. Launch a Web browser on the computer.

  \begin{itemize}
  \item
    To access the GL controller via its IP address, enter:
    http://10.30.200.247/ in the URL of the Web browser and press Enter.
    The home page of the GL controller will appear.
  \item
    If your static network has a DNS to resolve hostnames, enter:
    http://hostname.local/ in the URL of the Web browser and press
    Enter. The home page of the GL controller will appear.
  \end{itemize}
\end{enumerate}


\chapter{Network Monitor}

GL controller can communicate with each other on the local network. The
hostname in the status bar itself is a clickable link; when clicked, a
drop-down menu displays a link to the Network Monitor page and a list of
clickable links to other GL controller systems detected on the local
network.

\begin{figure}
\centering

\includegraphics[width=0.95\textwidth]{figures/2210352599-S8-network-monitor-001aa.PNG}
\caption{List of GL Controller on the local network}
\end{figure}\FloatBarrier

\begin{longtable}[]{@{}lll@{}}
\toprule
\begin{minipage}[b]{0.04\columnwidth}\raggedright
No.\strut
\end{minipage} & \begin{minipage}[b]{0.23\columnwidth}\raggedright
Name\strut
\end{minipage} & \begin{minipage}[b]{0.65\columnwidth}\raggedright
Description\strut
\end{minipage}\tabularnewline
\midrule
\endhead
\begin{minipage}[t]{0.04\columnwidth}\raggedright
1\strut
\end{minipage} & \begin{minipage}[t]{0.23\columnwidth}\raggedright
Network Monitor\strut
\end{minipage} & \begin{minipage}[t]{0.65\columnwidth}\raggedright
A clickable link that opens in a new page to manage all GL controller
systems detected on the local network. Operation of the Network Monitor
is bound to user privileges; a user must therefore have specific
privileges to use it.\strut
\end{minipage}\tabularnewline
\begin{minipage}[t]{0.04\columnwidth}\raggedright
2\strut
\end{minipage} & \begin{minipage}[t]{0.23\columnwidth}\raggedright
List of GL Controllers\strut
\end{minipage} & \begin{minipage}[t]{0.65\columnwidth}\raggedright
Displays a list of all GL controller systems (hostname and URL via IP
address) detected on the local network. Any GL controller system on the
list can be accessed directly by clicking on its hostname/IP. The back
button of the Web browser can be used to return to the GL controller
that initiated the call.\strut
\end{minipage}\tabularnewline
\bottomrule
\end{longtable}

\textbf{\emph{Note:}} The \textbf{Network Monitor} feature is not
operational on the HMI. The GL controller system must be part of a
network and its UI must be accessed remotely via a control PC on the
network.


\section{Network Monitor Page and
EWC}

The Network Monitor page can perform as a central server monitor where
all the GL controller systems on the local network can be displayed on a
single monitor, as depicted in the following figure. A operator can
remotely monitor and manage all the GL controller systems from a
centralized location.

\begin{figure}
\centering

\includegraphics[width=0.95\textwidth]{figures/3981361552-network-monitor-001aa.PNG}
\caption{Network Monitor page displaying all EWC devices on the local
network}
\end{figure}\FloatBarrier

To access and manage the Network Monitor page,

\textbf{\emph{Procedure:}}

\begin{enumerate}
\def\labelenumi{\arabic{enumi}.}
\tightlist
\item
  Log in as admin or as a user with permission to control the Network
  Monitor page
\item
  Click the E logo or the hostname of EWC
\item
  Click the Network Monitor link on the drop-down menu
\end{enumerate}

The following figure illustrates the layout of the Network Monitor page
with seven (7) important components. Each is listed and described as
follows:

\begin{figure}
\centering

\includegraphics[width=0.95\textwidth]{figures/3312785380-network-monitor-page-001a1.PNG}
\caption{Components and menu options of the Network Monitor page}
\end{figure}\FloatBarrier

\begin{enumerate}
\def\labelenumi{\arabic{enumi}.}
\item
  \textbf{Back Button}: This button closes the Network Monitor page and
  returns to the \textbf{Overview} of EWC.
\item
  \textbf{Refresh Button}: The refresh button updates the display of the
  Network Monitor page to apply any new configuration.
\item
  \textbf{Monitor Tab}: By default, only one tab with default name
  \textbf{ALL CHAMBERS} is active and ready to populate the selected EWC
  devices (as depicted in the above figure). Multiple tabs with unique
  and descriptive names can be created to display different groups of
  EWC devices.
\item
  \textbf{EWC on the Network}: By default, each EWC is displayed with
  its hostname for identification, but a unique name may be created in
  place of the hostname. As shown in the above figure, each EWC is
  displayed based on its operating status, such as Standby, Constant or
  Program, using it's default color. Each EWC can be accessed directly
  and displayed inside the Monitor tab as shown in the following figure
  for F4T-EP-RD. To close F4T-EP-RD, click the X button on the top
  right. To open it entirely in a new Web browser tab, click the expand
  button next to its hostname.

  \begin{figure}
  \centering
  
\includegraphics[width=0.95\textwidth]{figures/2622322270-ewc-open-inside-monitor-page-001.PNG}
  \caption{EWC displayed inside the Network Monitor tab}
  \end{figure}\FloatBarrier

  The color coding scheme allows the operator to quickly identify a
  chamber with issues, such as the one depicted in the following figure.

  \begin{figure}
  \centering
  
\includegraphics[width=0.95\textwidth]{figures/1392474873-chamber-error-001.PNG}
  \caption{Chamber error identifiable by its color}
  \end{figure}\FloatBarrier
\item
  \textbf{Monitor Area}: The monitor area is based on the size of the
  monitor screen which defines the real estate of the display to
  accommodate EWC devices.\\
\item
  \textbf{Edit Layout}: The \textbf{Edit Layout} (pen icon) is a
  clickable link/button that provides the edit menu with options to
  create additional tabs to populate EWC devices, as shown below.

  \begin{figure}
  \centering
  
\includegraphics[width=0.95\textwidth]{figures/4269699095-edit-menu-002b.PNG}
  \caption{Edit Layout and its menu options}
  \end{figure}\FloatBarrier

  From left to right, the four buttons are: Revert (back button), Save,
  Add Chamber/Device, Edit/Create Tab. Detailed operations of these
  buttons will be discussed in the following section.
\item
  \textbf{User}: Identity of the current user is listed here for
  reference.
\end{enumerate}


\section{Managing Network
Monitor}

When Network Monitor page is open for the first time, it has an empty
tab with default tab name \textbf{ALL CHAMBERS}.

\begin{figure}
\centering

\includegraphics[width=0.95\textwidth]{figures/3569009129-network-monitor-tab-001.PNG}
\caption{Empty network monitor tab}
\end{figure}\FloatBarrier

The \textbf{Edit Layout} (pen icon) button can be used to edit or create
new tabs to populate and display EWC devices. The four buttons in the
menu (labeled in the figure) can be used to manage and edit the layout.

\begin{figure}
\centering

\includegraphics[width=0.95\textwidth]{figures/3437174497-edit-menu-002aa.PNG}
\caption{Edit Layout features}
\end{figure}\FloatBarrier

These four buttons are described as follows:

\begin{enumerate}
\def\labelenumi{\arabic{enumi}.}
\item
  \textbf{Edit Tab}: The default tab name \textbf{All Chambers} can be
  edited via this button. Additional tabs can be added via the
  \textbf{+} button. Any tab on the list can be removed with the trash
  bin button. The up/down arrows on the right can be used to select a
  tab to be edited.

  \begin{figure}
  \centering
  
\includegraphics[width=0.95\textwidth]{figures/2734111444-edit-menu-001a2.PNG}
  \caption{Editing network tabs}
  \end{figure}\FloatBarrier
\item
  \textbf{Add Chamber/Device}: A new EWC device can be selected to add
  into the display tab via (1) the \textbf{Create New} button or (2) the
  EWC list as shown in the following figure. The first option permits
  manual configuration, while the second option applies the default
  settings for the selected EWC device which includes hostname, IP
  address, type of EWC and its version, as depicted in the following
  figure (for WebDev25).

  \begin{figure}
  \centering
  
\includegraphics[width=0.95\textwidth]{figures/1416981892-edit-menu-001a1.PNG}
  \caption{Adding new EWC devices}
  \end{figure}\FloatBarrier
\item
  \textbf{Save}: To apply all settings made in the network tab, the
  \textbf{Save} button must be applied.
\item
  \textbf{Revert}: This is the back button to exit the \textbf{Edit
  Layout}.
\end{enumerate}

In the following section, we provide an example how to set up the
Network Monitor page.


\section{Example 1}

In this section, we provide a step-by-step example how to set up and
configure a network monitor page to display three EWC devices in the
default Monitor tab. This example can be completed only if you have at
least two EWC devices, one as a host and the other as a client on your
network. \textbf{\emph{Note:}} All EWC device hostnames and IP addresses
used in this example are for illustration purposes only. You must use
the hostnames and IP address of your EWC devices found on your network.

\begin{enumerate}
\def\labelenumi{\arabic{enumi}.}
\item
  Select one EWC on the network to host a Network Monitor page.
\item
  Log in as admin or as a user with Network Monitor privilege.
\item
  Click the E logo or the hostname of EWC.

  \begin{figure}
  \centering
  
\includegraphics[width=0.95\textwidth]{figures/3817190245-example1-001.PNG}
  \caption{Accessing the Network Monitor page}
  \end{figure}\FloatBarrier
\item
  Click the \textbf{Network Monitor} link on the drop-down menu.
\item
  Click \textbf{Edit Layout} (pen icon) in the Network Monitor page.

  \begin{figure}
  \centering
  
\includegraphics[width=0.95\textwidth]{figures/756686761-ex1-step5.PNG}
  \caption{Accessing the Edit Layout option}
  \end{figure}\FloatBarrier
\item
  Click the \textbf{Add Chamber/Device} button. Two options are
  available for adding the EWC device into the Monitor tab: (1) Create
  New and (2) Select from the drop-down list. The first option is for a
  manual configuration; the second option applies the default
  configuration. Option 2 will be used for this example. From the
  drop-down menu, click to select the desired EWC to populate the
  Monitor tab.

  \begin{figure}
  \centering
  
\includegraphics[width=0.95\textwidth]{figures/3877497672-ex-2-step6.PNG}
  \caption{Selecting EWC to populate the Monitor tab}
  \end{figure}\FloatBarrier

  Default configuration of the selected EWC has been applied as shown in
  the pop-up window, which includes hostname, IP address, type of EWC
  and version.

  \begin{figure}
  \centering
  
\includegraphics[width=0.95\textwidth]{figures/3454823858-ex-2a-step6.PNG}
  \caption{Accept default setting}
  \end{figure}\FloatBarrier

  Click \textbf{ACCEPT} to populate this device in the Monitor tab.

  \begin{figure}
  \centering
  
\includegraphics[width=0.95\textwidth]{figures/3453184212-ex-2b-step6.PNG}
  \caption{First EWC populated in the Monitor tab}
  \end{figure}\FloatBarrier
\item
  Apply step 6 again to add two additional EWC devices from the list.
  The selected EWC devices will be populated below the first one in a
  single-column profile. The check mark in each box indicates the device
  selection.

  \begin{figure}
  \centering
  
\includegraphics[width=0.95\textwidth]{figures/460455672-ex-2-step7.PNG}
  \caption{Adding a second EWC}
  \end{figure}\FloatBarrier
\item
  The pen icon appearing in each EWC object indicates that the Monitor
  tab is still in edit mode. Click and drag to relocate or arrange each
  EWC in the Monitor tab, as shown below.

  \begin{figure}
  \centering
  
\includegraphics[width=0.95\textwidth]{figures/3732380629-ex-2-step8.PNG}
  \caption{Rearranging EWC devices in the Monitor tab}
  \end{figure}\FloatBarrier
\item
  Click \textbf{Save} to apply the settings. The Network Monitor page
  configuration is complete with three EWC devices in the Monitor tab
  (called \textbf{ALL CHAMBERS}).

  \begin{figure}
  \centering
  
\includegraphics[width=0.95\textwidth]{figures/4243869309-ex-2-step9.PNG}
  \caption{Selected EWC devices in the Monitor tab}
  \end{figure}\FloatBarrier
\item
  Each EWC in the Monitor tab is a clickable link that, when clicked,
  opens the selected EWC inside the Network Monitor page, as depicted
  below.

  \begin{figure}
  \centering
  
\includegraphics[width=0.95\textwidth]{figures/2251353925-scp220-001.PNG}
  \caption{EWC opened inside the Monitor tab}
  \end{figure}\FloatBarrier

  The expand button (1) next to the hostname can be used to open
  WebDemoSCP220 in a new Web browser tab so that the Network Monitor
  page is free to manage other EWC devices. The X button (2) can be used
  to exit and close WebDemoSCP220 to get back to the Network Monitor
  page.
\item
  The back button (indicated by the arrow) can be used to return to the
  host EWC for its chamber control and operation.

  \begin{figure}
  \centering
  
\includegraphics[width=0.95\textwidth]{figures/3640026851-ex2-step11.PNG}
  \caption{Returning to the \textbf{Overview} page}
  \end{figure}\FloatBarrier

  \begin{figure}
  \centering
  
\includegraphics[width=0.95\textwidth]{figures/3504587795-original-ewc-001.PNG}
  \caption{The host EWC in \textbf{Overview} home page}
  \end{figure}\FloatBarrier

  The E logo (in the \textbf{Overview} page) and the back button (in the
  Network Monitor page) can be used to toggle back and forth between the
  EWC's Overview page and its Network Monitor page to monitor other EWC
  devices on the network. The toggle operation can be performed during
  any operating mode of the host EWC, not just in Standby as shown in
  the figure.
\end{enumerate}


\section{Example 2}

In this section, four more examples are given to illustrate the
practical application of the Network Monitor page. The EWC device in
Example 1 will be used as the host EWC in following examples.


\subsection{Adding a new Monitor
tab}

This short example illustrates how to create and add a new Monitor tab
(called \textbf{TESTS}). It is then populated with two EWC devices.

\begin{enumerate}
\def\labelenumi{\arabic{enumi}.}
\item
  Log into EWC as admin or as a user with Network Monitor privilege.
\item
  Click the E logo (or the hostname).
\item
  Click \textbf{Network Monitor} link on the drop-down menu.
\item
  Click \textbf{Edit Layout} (pen icon).

  \begin{figure}
  \centering
  
\includegraphics[width=0.95\textwidth]{figures/3613019116-ex2-step4.PNG}
  \caption{Access the Edit Layout option}
  \end{figure}\FloatBarrier
\item
  Click \textbf{Edit Tabs}; then, click \textbf{+} in the \textbf{Tab
  Editor} window and edit the Display Name field with \textbf{TESTS},
  and click \textbf{CLOSE}.

  \begin{figure}
  \centering
  
\includegraphics[width=0.95\textwidth]{figures/2058622900-ex2-step4b.PNG}
  \caption{Creating TESTS in Tab Editor}
  \end{figure}\FloatBarrier
\item
  Click the \textbf{TESTS} tab, then click \textbf{Add Chamber/Device}
  to select and populate two EWC devices.

  \begin{figure}
  \centering
  
\includegraphics[width=0.95\textwidth]{figures/3482745320-ex2-step6.PNG}
  \caption{Populating two EWC devices in TESTS}
  \end{figure}\FloatBarrier

  Rearrange the layout and click \textbf{Save} to complete the
  configuration.
\item
  \textbf{TESTS} tab and \textbf{ALL CHAMBERS} tab can be toggled to
  display two groups of EWC devices, as illustrated in the following
  figure. An active tab is highlighted and underlined in blue (as
  indicated by the arrows).

  \begin{figure}
  \centering
  
\includegraphics[width=0.95\textwidth]{figures/3398004567-ex2-step7abc.PNG}
  \caption{Toggling between two Monitor tabs}
  \end{figure}\FloatBarrier
\item
  Additional tabs may be created and added into the tab bar by repeating
  steps 5-7, as depicted in the following figure.

  \begin{figure}
  \centering
  
\includegraphics[width=0.95\textwidth]{figures/1525972342-ex2-step8a.PNG}
  \caption{Multiple Monitor tabs in the tab bar}
  \end{figure}\FloatBarrier
\end{enumerate}


\subsection{Removing an EWC
Device}

This short example illustrates how to remove an EWC device from a
Monitor tab. As a concrete example, WebDev25 will be removed from the
\textbf{ALL CHAMBERS} tab.

\begin{enumerate}
\def\labelenumi{\arabic{enumi}.}
\item
  Starting from the Network Monitor page, click the \textbf{ALL
  CHAMBERS} tab.
\item
  Click the \textbf{Edit Layout} (pen icon).
\item
  Click the pen icon of the WebDev25 object, then click \textbf{DELETE}
  in the Chamber/Device Editor window as shown in the following figure.

  \begin{figure}
  \centering
  
\includegraphics[width=0.95\textwidth]{figures/1239063280-ex3-step3.PNG}
  \caption{Deleting WebDev25}
  \end{figure}\FloatBarrier
\item
  Click \textbf{Save} to apply new settings to this tab. Only two EWC
  devices remain in this tab.
\end{enumerate}


\subsection{Adding a New EWC
Device}

This short example illustrates how to manually add a new EWC device
(using WebDev25 as an example) to the \textbf{TESTS} tab using the
\textbf{Create New} button.

\begin{enumerate}
\def\labelenumi{\arabic{enumi}.}
\item
  Starting from the Network Monitor page, click the \textbf{ALL
  CHAMBERS} tab. \textbf{\emph{Note:}} It is not necessary to select
  \textbf{ALL CHAMBERS}, but selecting it allows a complete manual
  configuration (as will be apparent at step 4).\\
\item
  Click the \textbf{Edit Layout} (pen icon).
\item
  Click the \textbf{Add Chamber/Device} button and click \textbf{Create
  New}.
\item
  Apply the following steps to the Chamber/Device Editor:

  \begin{figure}
  \centering
  
\includegraphics[width=0.95\textwidth]{figures/1146931544-ex4-step4.PNG}
  \caption{Editing the Monito Tab}
  \end{figure}\FloatBarrier

  \begin{itemize}
  \item
    Click the \textbf{Show on Tab} spin button to select \textbf{TESTS}
    from the list.
  \item
    Enter \textbf{MyWebDev25} in the Display Name field.
    \textbf{\emph{Note:}} To display the EWC device by its hostname,
    \textbf{WebDev25} may be entered here. This field cannot be left
    blank or an \textbf{INVALID NAME} will be displayed.
  \item
    Enter hostname as ``webdev25.local''.
  \item
    Enter its IP address (10.30.200.239).
  \item
    Click the Device Version field to select 2.?.? from the list. The
    selected EWC is known to have version 2.x.x. The correct version
    must be selected or an error will occur.
  \item
    Click \textbf{ACCEPT}.
  \end{itemize}
\item
  Click \textbf{Save}.
\item
  Click the \textbf{TESTS} tab to view its new display and notice that
  WebDev25 is now identified as \textbf{MyWebDev25}.

  \begin{figure}
  \centering
  
\includegraphics[width=0.95\textwidth]{figures/356898404-ex4-step11.PNG}
  \caption{Viewing new display of TESTS tab}
  \end{figure}\FloatBarrier
\end{enumerate}

Warning! At step 4 above, it is crucial to enter the correct IP address,
since the system tries to locate the device based on that IP address on
its network. A wrong IP address would produce an EWC error as depicted
in the following figure. Here, steps 1-5 were repeated to create a new
EWC device with a non-existing IP address.

\begin{figure}
\centering

\includegraphics[width=0.95\textwidth]{figures/1775694629-ex4-step12.PNG}
\caption{Unknown EWC device}
\end{figure}\FloatBarrier

To remove a misconfigured EWC device, such as \textbf{MYTEST2}, refer to
the previous example on \textbf{Removing an EWC Device}.


\subsection{Removing a Monitor
Tab}

This short example illustrates how to remove \textbf{TESTS} from the
Monitor tab.

\begin{enumerate}
\def\labelenumi{\arabic{enumi}.}
\item
  Starting from the Network Monitor page, click the \textbf{Edit Layout}
  button (from within any active tab).
\item
  Click \textbf{Edit Tabs}.
\item
  In the Tab Editor window, click to select \textbf{TESTS} in the tab
  list, then click the trash bin (to remove it) and click CLOSE.

  \begin{figure}
  \centering
  
\includegraphics[width=0.95\textwidth]{figures/402015291-ex3-step4-5.PNG}
  \caption{Removing TESTS tab from the list}
  \end{figure}\FloatBarrier
\item
  Click \textbf{Save} to apply the changes. \textbf{TESTS} tab is now
  removed from the Monitor tab and no longer appears in the tab bar.
\end{enumerate}


\section{APPENDIX}


\subsection{A.1 Program Creation
Sheet}

Sometimes, it is helpful to construct a table to outline the steps of
the program before entering them into the GL controller program editor
(see Section 5.3). The following table illustrates an example how to
prepare these steps. This program conducts \emph{test} on temperature
fluctuation.

\begin{longtable}[]{@{}lllllllllll@{}}
\toprule
\begin{minipage}[b]{0.09\columnwidth}\raggedright
Step No.\strut
\end{minipage} & \begin{minipage}[b]{0.05\columnwidth}\raggedright
Temp SP\strut
\end{minipage} & \begin{minipage}[b]{0.09\columnwidth}\raggedright
Temp Mode\strut
\end{minipage} & \begin{minipage}[b]{0.06\columnwidth}\raggedright
Humi SP\strut
\end{minipage} & \begin{minipage}[b]{0.06\columnwidth}\raggedright
Humid Mode\strut
\end{minipage} & \begin{minipage}[b]{0.12\columnwidth}\raggedright
Time (h:min)\strut
\end{minipage} & \begin{minipage}[b]{0.08\columnwidth}\raggedright
Refrig\strut
\end{minipage} & \begin{minipage}[b]{0.04\columnwidth}\raggedright
TS 1\strut
\end{minipage} & \begin{minipage}[b]{0.05\columnwidth}\raggedright
TS 2\strut
\end{minipage} & \begin{minipage}[b]{0.03\columnwidth}\raggedright
Counter Step\strut
\end{minipage} & \begin{minipage}[b]{0.03\columnwidth}\raggedright
Counter Count\strut
\end{minipage}\tabularnewline
\midrule
\endhead
\begin{minipage}[t]{0.09\columnwidth}\raggedright
1\strut
\end{minipage} & \begin{minipage}[t]{0.05\columnwidth}\raggedright
-30\strut
\end{minipage} & \begin{minipage}[t]{0.09\columnwidth}\raggedright
Soak Ctrl\strut
\end{minipage} & \begin{minipage}[t]{0.06\columnwidth}\raggedright
5\strut
\end{minipage} & \begin{minipage}[t]{0.06\columnwidth}\raggedright
Off\strut
\end{minipage} & \begin{minipage}[t]{0.12\columnwidth}\raggedright
1:00:00\strut
\end{minipage} & \begin{minipage}[t]{0.08\columnwidth}\raggedright
Auto\strut
\end{minipage} & \begin{minipage}[t]{0.04\columnwidth}\raggedright
\strut
\end{minipage} & \begin{minipage}[t]{0.05\columnwidth}\raggedright
\strut
\end{minipage} & \begin{minipage}[t]{0.03\columnwidth}\raggedright
\strut
\end{minipage} & \begin{minipage}[t]{0.03\columnwidth}\raggedright
0\strut
\end{minipage}\tabularnewline
\begin{minipage}[t]{0.09\columnwidth}\raggedright
2\strut
\end{minipage} & \begin{minipage}[t]{0.05\columnwidth}\raggedright
20\strut
\end{minipage} & \begin{minipage}[t]{0.09\columnwidth}\raggedright
Soak Ctrl\strut
\end{minipage} & \begin{minipage}[t]{0.06\columnwidth}\raggedright
5\strut
\end{minipage} & \begin{minipage}[t]{0.06\columnwidth}\raggedright
Off\strut
\end{minipage} & \begin{minipage}[t]{0.12\columnwidth}\raggedright
1:00:00\strut
\end{minipage} & \begin{minipage}[t]{0.08\columnwidth}\raggedright
Auto\strut
\end{minipage} & \begin{minipage}[t]{0.04\columnwidth}\raggedright
\strut
\end{minipage} & \begin{minipage}[t]{0.05\columnwidth}\raggedright
\strut
\end{minipage} & \begin{minipage}[t]{0.03\columnwidth}\raggedright
\strut
\end{minipage} & \begin{minipage}[t]{0.03\columnwidth}\raggedright
0\strut
\end{minipage}\tabularnewline
\begin{minipage}[t]{0.09\columnwidth}\raggedright
3\strut
\end{minipage} & \begin{minipage}[t]{0.05\columnwidth}\raggedright
55\strut
\end{minipage} & \begin{minipage}[t]{0.09\columnwidth}\raggedright
Soak Ctrl\strut
\end{minipage} & \begin{minipage}[t]{0.06\columnwidth}\raggedright
5\strut
\end{minipage} & \begin{minipage}[t]{0.06\columnwidth}\raggedright
Off\strut
\end{minipage} & \begin{minipage}[t]{0.12\columnwidth}\raggedright
1:00:00\strut
\end{minipage} & \begin{minipage}[t]{0.08\columnwidth}\raggedright
Auto\strut
\end{minipage} & \begin{minipage}[t]{0.04\columnwidth}\raggedright
\strut
\end{minipage} & \begin{minipage}[t]{0.05\columnwidth}\raggedright
\strut
\end{minipage} & \begin{minipage}[t]{0.03\columnwidth}\raggedright
\strut
\end{minipage} & \begin{minipage}[t]{0.03\columnwidth}\raggedright
0\strut
\end{minipage}\tabularnewline
\begin{minipage}[t]{0.09\columnwidth}\raggedright
4\strut
\end{minipage} & \begin{minipage}[t]{0.05\columnwidth}\raggedright
65\strut
\end{minipage} & \begin{minipage}[t]{0.09\columnwidth}\raggedright
Soak Ctrl\strut
\end{minipage} & \begin{minipage}[t]{0.06\columnwidth}\raggedright
5\strut
\end{minipage} & \begin{minipage}[t]{0.06\columnwidth}\raggedright
Off\strut
\end{minipage} & \begin{minipage}[t]{0.12\columnwidth}\raggedright
1:00:00\strut
\end{minipage} & \begin{minipage}[t]{0.08\columnwidth}\raggedright
Auto\strut
\end{minipage} & \begin{minipage}[t]{0.04\columnwidth}\raggedright
\strut
\end{minipage} & \begin{minipage}[t]{0.05\columnwidth}\raggedright
\strut
\end{minipage} & \begin{minipage}[t]{0.03\columnwidth}\raggedright
\strut
\end{minipage} & \begin{minipage}[t]{0.03\columnwidth}\raggedright
0\strut
\end{minipage}\tabularnewline
\begin{minipage}[t]{0.09\columnwidth}\raggedright
5\strut
\end{minipage} & \begin{minipage}[t]{0.05\columnwidth}\raggedright
180\strut
\end{minipage} & \begin{minipage}[t]{0.09\columnwidth}\raggedright
Soak Ctrl\strut
\end{minipage} & \begin{minipage}[t]{0.06\columnwidth}\raggedright
5\strut
\end{minipage} & \begin{minipage}[t]{0.06\columnwidth}\raggedright
Off\strut
\end{minipage} & \begin{minipage}[t]{0.12\columnwidth}\raggedright
1:00:00\strut
\end{minipage} & \begin{minipage}[t]{0.08\columnwidth}\raggedright
Auto\strut
\end{minipage} & \begin{minipage}[t]{0.04\columnwidth}\raggedright
\strut
\end{minipage} & \begin{minipage}[t]{0.05\columnwidth}\raggedright
\strut
\end{minipage} & \begin{minipage}[t]{0.03\columnwidth}\raggedright
\strut
\end{minipage} & \begin{minipage}[t]{0.03\columnwidth}\raggedright
0\strut
\end{minipage}\tabularnewline
\begin{minipage}[t]{0.09\columnwidth}\raggedright
6\strut
\end{minipage} & \begin{minipage}[t]{0.05\columnwidth}\raggedright
23\strut
\end{minipage} & \begin{minipage}[t]{0.09\columnwidth}\raggedright
On\strut
\end{minipage} & \begin{minipage}[t]{0.06\columnwidth}\raggedright
5\strut
\end{minipage} & \begin{minipage}[t]{0.06\columnwidth}\raggedright
Off\strut
\end{minipage} & \begin{minipage}[t]{0.12\columnwidth}\raggedright
1:00:00\strut
\end{minipage} & \begin{minipage}[t]{0.08\columnwidth}\raggedright
Auto\strut
\end{minipage} & \begin{minipage}[t]{0.04\columnwidth}\raggedright
\strut
\end{minipage} & \begin{minipage}[t]{0.05\columnwidth}\raggedright
\strut
\end{minipage} & \begin{minipage}[t]{0.03\columnwidth}\raggedright
\strut
\end{minipage} & \begin{minipage}[t]{0.03\columnwidth}\raggedright
0\strut
\end{minipage}\tabularnewline
\bottomrule
\end{longtable}

Here is another program that performs Low Humidity test. This program
requires a chamber with humidity option. As illustrated in the following
Table, deviation soak between plus and minus 10 have been utilized. Both
temperature and humidity have been controlled via \emph{soak}. Time
signal TS 9 and TS 10 have been utilized. \textbf{\emph{Note:}} TS 1 and
2 are generally the default time signals.

\begin{longtable}[]{@{}lllllllllllll@{}}
\toprule
\begin{minipage}[b]{0.06\columnwidth}\raggedright
Step No.\strut
\end{minipage} & \begin{minipage}[b]{0.06\columnwidth}\raggedright
Temp Mode\strut
\end{minipage} & \begin{minipage}[b]{0.05\columnwidth}\raggedright
Temp SP\strut
\end{minipage} & \begin{minipage}[b]{0.04\columnwidth}\raggedright
Dev(+)\strut
\end{minipage} & \begin{minipage}[b]{0.04\columnwidth}\raggedright
Dev(-)\strut
\end{minipage} & \begin{minipage}[b]{0.06\columnwidth}\raggedright
Humi Mode\strut
\end{minipage} & \begin{minipage}[b]{0.05\columnwidth}\raggedright
Humi SP\strut
\end{minipage} & \begin{minipage}[b]{0.05\columnwidth}\raggedright
Duration\strut
\end{minipage} & \begin{minipage}[b]{0.04\columnwidth}\raggedright
Refrig\strut
\end{minipage} & \begin{minipage}[b]{0.03\columnwidth}\raggedright
TS9\strut
\end{minipage} & \begin{minipage}[b]{0.03\columnwidth}\raggedright
Ts10\strut
\end{minipage} & \begin{minipage}[b]{0.07\columnwidth}\raggedright
Counter Step\strut
\end{minipage} & \begin{minipage}[b]{0.08\columnwidth}\raggedright
Counter Loop\strut
\end{minipage}\tabularnewline
\midrule
\endhead
\begin{minipage}[t]{0.06\columnwidth}\raggedright
1\strut
\end{minipage} & \begin{minipage}[t]{0.06\columnwidth}\raggedright
Soak\strut
\end{minipage} & \begin{minipage}[t]{0.05\columnwidth}\raggedright
60\strut
\end{minipage} & \begin{minipage}[t]{0.04\columnwidth}\raggedright
10\strut
\end{minipage} & \begin{minipage}[t]{0.04\columnwidth}\raggedright
-10\strut
\end{minipage} & \begin{minipage}[t]{0.06\columnwidth}\raggedright
Soak\strut
\end{minipage} & \begin{minipage}[t]{0.05\columnwidth}\raggedright
10\strut
\end{minipage} & \begin{minipage}[t]{0.05\columnwidth}\raggedright
2:00\strut
\end{minipage} & \begin{minipage}[t]{0.04\columnwidth}\raggedright
Auto\strut
\end{minipage} & \begin{minipage}[t]{0.03\columnwidth}\raggedright
On\strut
\end{minipage} & \begin{minipage}[t]{0.03\columnwidth}\raggedright
On\strut
\end{minipage} & \begin{minipage}[t]{0.07\columnwidth}\raggedright
\strut
\end{minipage} & \begin{minipage}[t]{0.08\columnwidth}\raggedright
\strut
\end{minipage}\tabularnewline
\begin{minipage}[t]{0.06\columnwidth}\raggedright
2\strut
\end{minipage} & \begin{minipage}[t]{0.06\columnwidth}\raggedright
Soak\strut
\end{minipage} & \begin{minipage}[t]{0.05\columnwidth}\raggedright
40\strut
\end{minipage} & \begin{minipage}[t]{0.04\columnwidth}\raggedright
10\strut
\end{minipage} & \begin{minipage}[t]{0.04\columnwidth}\raggedright
-10\strut
\end{minipage} & \begin{minipage}[t]{0.06\columnwidth}\raggedright
Soak\strut
\end{minipage} & \begin{minipage}[t]{0.05\columnwidth}\raggedright
15\strut
\end{minipage} & \begin{minipage}[t]{0.05\columnwidth}\raggedright
2:00\strut
\end{minipage} & \begin{minipage}[t]{0.04\columnwidth}\raggedright
Auto\strut
\end{minipage} & \begin{minipage}[t]{0.03\columnwidth}\raggedright
On\strut
\end{minipage} & \begin{minipage}[t]{0.03\columnwidth}\raggedright
On\strut
\end{minipage} & \begin{minipage}[t]{0.07\columnwidth}\raggedright
\strut
\end{minipage} & \begin{minipage}[t]{0.08\columnwidth}\raggedright
\strut
\end{minipage}\tabularnewline
\begin{minipage}[t]{0.06\columnwidth}\raggedright
3\strut
\end{minipage} & \begin{minipage}[t]{0.06\columnwidth}\raggedright
Soak\strut
\end{minipage} & \begin{minipage}[t]{0.05\columnwidth}\raggedright
20\strut
\end{minipage} & \begin{minipage}[t]{0.04\columnwidth}\raggedright
10\strut
\end{minipage} & \begin{minipage}[t]{0.04\columnwidth}\raggedright
-10\strut
\end{minipage} & \begin{minipage}[t]{0.06\columnwidth}\raggedright
Soak\strut
\end{minipage} & \begin{minipage}[t]{0.05\columnwidth}\raggedright
15\strut
\end{minipage} & \begin{minipage}[t]{0.05\columnwidth}\raggedright
2:00\strut
\end{minipage} & \begin{minipage}[t]{0.04\columnwidth}\raggedright
Auto\strut
\end{minipage} & \begin{minipage}[t]{0.03\columnwidth}\raggedright
On\strut
\end{minipage} & \begin{minipage}[t]{0.03\columnwidth}\raggedright
On\strut
\end{minipage} & \begin{minipage}[t]{0.07\columnwidth}\raggedright
\strut
\end{minipage} & \begin{minipage}[t]{0.08\columnwidth}\raggedright
\strut
\end{minipage}\tabularnewline
\begin{minipage}[t]{0.06\columnwidth}\raggedright
4\strut
\end{minipage} & \begin{minipage}[t]{0.06\columnwidth}\raggedright
Soak\strut
\end{minipage} & \begin{minipage}[t]{0.05\columnwidth}\raggedright
10\strut
\end{minipage} & \begin{minipage}[t]{0.04\columnwidth}\raggedright
10\strut
\end{minipage} & \begin{minipage}[t]{0.04\columnwidth}\raggedright
-10\strut
\end{minipage} & \begin{minipage}[t]{0.06\columnwidth}\raggedright
Soak\strut
\end{minipage} & \begin{minipage}[t]{0.05\columnwidth}\raggedright
15\strut
\end{minipage} & \begin{minipage}[t]{0.05\columnwidth}\raggedright
2:00\strut
\end{minipage} & \begin{minipage}[t]{0.04\columnwidth}\raggedright
Auto\strut
\end{minipage} & \begin{minipage}[t]{0.03\columnwidth}\raggedright
On\strut
\end{minipage} & \begin{minipage}[t]{0.03\columnwidth}\raggedright
On\strut
\end{minipage} & \begin{minipage}[t]{0.07\columnwidth}\raggedright
\strut
\end{minipage} & \begin{minipage}[t]{0.08\columnwidth}\raggedright
\strut
\end{minipage}\tabularnewline
\begin{minipage}[t]{0.06\columnwidth}\raggedright
5\strut
\end{minipage} & \begin{minipage}[t]{0.06\columnwidth}\raggedright
Soak\strut
\end{minipage} & \begin{minipage}[t]{0.05\columnwidth}\raggedright
40\strut
\end{minipage} & \begin{minipage}[t]{0.04\columnwidth}\raggedright
10\strut
\end{minipage} & \begin{minipage}[t]{0.04\columnwidth}\raggedright
-10\strut
\end{minipage} & \begin{minipage}[t]{0.06\columnwidth}\raggedright
Soak\strut
\end{minipage} & \begin{minipage}[t]{0.05\columnwidth}\raggedright
20\strut
\end{minipage} & \begin{minipage}[t]{0.05\columnwidth}\raggedright
2:00\strut
\end{minipage} & \begin{minipage}[t]{0.04\columnwidth}\raggedright
Auto\strut
\end{minipage} & \begin{minipage}[t]{0.03\columnwidth}\raggedright
On\strut
\end{minipage} & \begin{minipage}[t]{0.03\columnwidth}\raggedright
On\strut
\end{minipage} & \begin{minipage}[t]{0.07\columnwidth}\raggedright
\strut
\end{minipage} & \begin{minipage}[t]{0.08\columnwidth}\raggedright
\strut
\end{minipage}\tabularnewline
\begin{minipage}[t]{0.06\columnwidth}\raggedright
6\strut
\end{minipage} & \begin{minipage}[t]{0.06\columnwidth}\raggedright
Soak\strut
\end{minipage} & \begin{minipage}[t]{0.05\columnwidth}\raggedright
20\strut
\end{minipage} & \begin{minipage}[t]{0.04\columnwidth}\raggedright
10\strut
\end{minipage} & \begin{minipage}[t]{0.04\columnwidth}\raggedright
-10\strut
\end{minipage} & \begin{minipage}[t]{0.06\columnwidth}\raggedright
Soak\strut
\end{minipage} & \begin{minipage}[t]{0.05\columnwidth}\raggedright
30\strut
\end{minipage} & \begin{minipage}[t]{0.05\columnwidth}\raggedright
2:00\strut
\end{minipage} & \begin{minipage}[t]{0.04\columnwidth}\raggedright
Auto\strut
\end{minipage} & \begin{minipage}[t]{0.03\columnwidth}\raggedright
On\strut
\end{minipage} & \begin{minipage}[t]{0.03\columnwidth}\raggedright
On\strut
\end{minipage} & \begin{minipage}[t]{0.07\columnwidth}\raggedright
\strut
\end{minipage} & \begin{minipage}[t]{0.08\columnwidth}\raggedright
\strut
\end{minipage}\tabularnewline
\begin{minipage}[t]{0.06\columnwidth}\raggedright
7\strut
\end{minipage} & \begin{minipage}[t]{0.06\columnwidth}\raggedright
Soak\strut
\end{minipage} & \begin{minipage}[t]{0.05\columnwidth}\raggedright
40\strut
\end{minipage} & \begin{minipage}[t]{0.04\columnwidth}\raggedright
10\strut
\end{minipage} & \begin{minipage}[t]{0.04\columnwidth}\raggedright
-10\strut
\end{minipage} & \begin{minipage}[t]{0.06\columnwidth}\raggedright
Soak\strut
\end{minipage} & \begin{minipage}[t]{0.05\columnwidth}\raggedright
40\strut
\end{minipage} & \begin{minipage}[t]{0.05\columnwidth}\raggedright
2:00\strut
\end{minipage} & \begin{minipage}[t]{0.04\columnwidth}\raggedright
Auto\strut
\end{minipage} & \begin{minipage}[t]{0.03\columnwidth}\raggedright
On\strut
\end{minipage} & \begin{minipage}[t]{0.03\columnwidth}\raggedright
On\strut
\end{minipage} & \begin{minipage}[t]{0.07\columnwidth}\raggedright
\strut
\end{minipage} & \begin{minipage}[t]{0.08\columnwidth}\raggedright
\strut
\end{minipage}\tabularnewline
\begin{minipage}[t]{0.06\columnwidth}\raggedright
8\strut
\end{minipage} & \begin{minipage}[t]{0.06\columnwidth}\raggedright
Soak\strut
\end{minipage} & \begin{minipage}[t]{0.05\columnwidth}\raggedright
10\strut
\end{minipage} & \begin{minipage}[t]{0.04\columnwidth}\raggedright
10\strut
\end{minipage} & \begin{minipage}[t]{0.04\columnwidth}\raggedright
-10\strut
\end{minipage} & \begin{minipage}[t]{0.06\columnwidth}\raggedright
Soak\strut
\end{minipage} & \begin{minipage}[t]{0.05\columnwidth}\raggedright
20\strut
\end{minipage} & \begin{minipage}[t]{0.05\columnwidth}\raggedright
2:00\strut
\end{minipage} & \begin{minipage}[t]{0.04\columnwidth}\raggedright
Auto\strut
\end{minipage} & \begin{minipage}[t]{0.03\columnwidth}\raggedright
On\strut
\end{minipage} & \begin{minipage}[t]{0.03\columnwidth}\raggedright
\strut
\end{minipage} & \begin{minipage}[t]{0.07\columnwidth}\raggedright
\strut
\end{minipage} & \begin{minipage}[t]{0.08\columnwidth}\raggedright
\strut
\end{minipage}\tabularnewline
\begin{minipage}[t]{0.06\columnwidth}\raggedright
9\strut
\end{minipage} & \begin{minipage}[t]{0.06\columnwidth}\raggedright
Soak\strut
\end{minipage} & \begin{minipage}[t]{0.05\columnwidth}\raggedright
20\strut
\end{minipage} & \begin{minipage}[t]{0.04\columnwidth}\raggedright
10\strut
\end{minipage} & \begin{minipage}[t]{0.04\columnwidth}\raggedright
-10\strut
\end{minipage} & \begin{minipage}[t]{0.06\columnwidth}\raggedright
Soak\strut
\end{minipage} & \begin{minipage}[t]{0.05\columnwidth}\raggedright
30\strut
\end{minipage} & \begin{minipage}[t]{0.05\columnwidth}\raggedright
2:00\strut
\end{minipage} & \begin{minipage}[t]{0.04\columnwidth}\raggedright
Auto\strut
\end{minipage} & \begin{minipage}[t]{0.03\columnwidth}\raggedright
On\strut
\end{minipage} & \begin{minipage}[t]{0.03\columnwidth}\raggedright
\strut
\end{minipage} & \begin{minipage}[t]{0.07\columnwidth}\raggedright
\strut
\end{minipage} & \begin{minipage}[t]{0.08\columnwidth}\raggedright
\strut
\end{minipage}\tabularnewline
\begin{minipage}[t]{0.06\columnwidth}\raggedright
10\strut
\end{minipage} & \begin{minipage}[t]{0.06\columnwidth}\raggedright
Soak\strut
\end{minipage} & \begin{minipage}[t]{0.05\columnwidth}\raggedright
10\strut
\end{minipage} & \begin{minipage}[t]{0.04\columnwidth}\raggedright
10\strut
\end{minipage} & \begin{minipage}[t]{0.04\columnwidth}\raggedright
-10\strut
\end{minipage} & \begin{minipage}[t]{0.06\columnwidth}\raggedright
Soak\strut
\end{minipage} & \begin{minipage}[t]{0.05\columnwidth}\raggedright
40\strut
\end{minipage} & \begin{minipage}[t]{0.05\columnwidth}\raggedright
2:00\strut
\end{minipage} & \begin{minipage}[t]{0.04\columnwidth}\raggedright
Auto\strut
\end{minipage} & \begin{minipage}[t]{0.03\columnwidth}\raggedright
On\strut
\end{minipage} & \begin{minipage}[t]{0.03\columnwidth}\raggedright
\strut
\end{minipage} & \begin{minipage}[t]{0.07\columnwidth}\raggedright
\strut
\end{minipage} & \begin{minipage}[t]{0.08\columnwidth}\raggedright
\strut
\end{minipage}\tabularnewline
\begin{minipage}[t]{0.06\columnwidth}\raggedright
11\strut
\end{minipage} & \begin{minipage}[t]{0.06\columnwidth}\raggedright
Soak\strut
\end{minipage} & \begin{minipage}[t]{0.05\columnwidth}\raggedright
40\strut
\end{minipage} & \begin{minipage}[t]{0.04\columnwidth}\raggedright
10\strut
\end{minipage} & \begin{minipage}[t]{0.04\columnwidth}\raggedright
-10\strut
\end{minipage} & \begin{minipage}[t]{0.06\columnwidth}\raggedright
Soak\strut
\end{minipage} & \begin{minipage}[t]{0.05\columnwidth}\raggedright
20\strut
\end{minipage} & \begin{minipage}[t]{0.05\columnwidth}\raggedright
2:00\strut
\end{minipage} & \begin{minipage}[t]{0.04\columnwidth}\raggedright
Auto\strut
\end{minipage} & \begin{minipage}[t]{0.03\columnwidth}\raggedright
On\strut
\end{minipage} & \begin{minipage}[t]{0.03\columnwidth}\raggedright
\strut
\end{minipage} & \begin{minipage}[t]{0.07\columnwidth}\raggedright
\strut
\end{minipage} & \begin{minipage}[t]{0.08\columnwidth}\raggedright
\strut
\end{minipage}\tabularnewline
\begin{minipage}[t]{0.06\columnwidth}\raggedright
12\strut
\end{minipage} & \begin{minipage}[t]{0.06\columnwidth}\raggedright
Soak\strut
\end{minipage} & \begin{minipage}[t]{0.05\columnwidth}\raggedright
20\strut
\end{minipage} & \begin{minipage}[t]{0.04\columnwidth}\raggedright
10\strut
\end{minipage} & \begin{minipage}[t]{0.04\columnwidth}\raggedright
-10\strut
\end{minipage} & \begin{minipage}[t]{0.06\columnwidth}\raggedright
Soak\strut
\end{minipage} & \begin{minipage}[t]{0.05\columnwidth}\raggedright
45\strut
\end{minipage} & \begin{minipage}[t]{0.05\columnwidth}\raggedright
2:00\strut
\end{minipage} & \begin{minipage}[t]{0.04\columnwidth}\raggedright
Auto\strut
\end{minipage} & \begin{minipage}[t]{0.03\columnwidth}\raggedright
On\strut
\end{minipage} & \begin{minipage}[t]{0.03\columnwidth}\raggedright
\strut
\end{minipage} & \begin{minipage}[t]{0.07\columnwidth}\raggedright
\strut
\end{minipage} & \begin{minipage}[t]{0.08\columnwidth}\raggedright
\strut
\end{minipage}\tabularnewline
\begin{minipage}[t]{0.06\columnwidth}\raggedright
13\strut
\end{minipage} & \begin{minipage}[t]{0.06\columnwidth}\raggedright
Soak\strut
\end{minipage} & \begin{minipage}[t]{0.05\columnwidth}\raggedright
10\strut
\end{minipage} & \begin{minipage}[t]{0.04\columnwidth}\raggedright
10\strut
\end{minipage} & \begin{minipage}[t]{0.04\columnwidth}\raggedright
-10\strut
\end{minipage} & \begin{minipage}[t]{0.06\columnwidth}\raggedright
Soak\strut
\end{minipage} & \begin{minipage}[t]{0.05\columnwidth}\raggedright
80\strut
\end{minipage} & \begin{minipage}[t]{0.05\columnwidth}\raggedright
2:00\strut
\end{minipage} & \begin{minipage}[t]{0.04\columnwidth}\raggedright
Auto\strut
\end{minipage} & \begin{minipage}[t]{0.03\columnwidth}\raggedright
On\strut
\end{minipage} & \begin{minipage}[t]{0.03\columnwidth}\raggedright
\strut
\end{minipage} & \begin{minipage}[t]{0.07\columnwidth}\raggedright
\strut
\end{minipage} & \begin{minipage}[t]{0.08\columnwidth}\raggedright
\strut
\end{minipage}\tabularnewline
\bottomrule
\end{longtable}


\subsection{A.2 Program Creation
Sheet}

Copy and use the following programming table to prepare a program.

\begin{longtable}[]{@{}llllllllllllllllll@{}}
\toprule
Step & Temp & Humi & Time & Soak & Refrig & 1 & 2 & 3 & 4 & 5 & 6 & 7 &
8 & 9 & 10 & 11 & 12\tabularnewline
\midrule
\endhead
& & & & & & & & & & & & & & & & &\tabularnewline
& & & & & & & & & & & & & & & & &\tabularnewline
& & & & & & & & & & & & & & & & &\tabularnewline
& & & & & & & & & & & & & & & & &\tabularnewline
& & & & & & & & & & & & & & & & &\tabularnewline
& & & & & & & & & & & & & & & & &\tabularnewline
& & & & & & & & & & & & & & & & &\tabularnewline
& & & & & & & & & & & & & & & & &\tabularnewline
& & & & & & & & & & & & & & & & &\tabularnewline
& & & & & & & & & & & & & & & & &\tabularnewline
& & & & & & & & & & & & & & & & &\tabularnewline
\bottomrule
\end{longtable}

Programming optional features.

Processing at End of Operation

☐ Stop; ☐ Power OFF; ☐ Hold Last Step. ☐ Constant Operation
(No.~1)(No.~2)(No.~3) ☐ Operate Next Program No.~( ), Step ( ) Example:
☑ Stop

Counter A

Repeat Cycle: Repeating Start Step ( ) Step ; Repeating End Step ( )
Step

Counter B

Repeat Cycle: Repeating Start Step ( ) Step ; Repeating End Step ( )
Step

Start Setting (Ramp Operation)

☐ Process Value ☐ Set Point ( °C) ( \%RH) Example: ☑ Ste Point (35°C)
(25\%RH)

\textbf{\emph{Note:}} GL controller supports multiple counters which
therefore will not be limited to just Counter A and Counter B. Each
counter and its loop can be invoked by entering the loop number under
the \emph{Counter} column (refer to the first table in this Appendix)
and then set the start step. Multiple counters can be created by
manipulating the start/stop step. The following table illustrates how
counters are manipulated in a program.

\begin{longtable}[]{@{}lllllllllll@{}}
\toprule
\begin{minipage}[b]{0.08\columnwidth}\raggedright
Step No.\strut
\end{minipage} & \begin{minipage}[b]{0.07\columnwidth}\raggedright
Temp SP\strut
\end{minipage} & \begin{minipage}[b]{0.09\columnwidth}\raggedright
Temp Mode\strut
\end{minipage} & \begin{minipage}[b]{0.06\columnwidth}\raggedright
Humi SP\strut
\end{minipage} & \begin{minipage}[b]{0.06\columnwidth}\raggedright
Humid Mode\strut
\end{minipage} & \begin{minipage}[b]{0.12\columnwidth}\raggedright
Time (h:min)\strut
\end{minipage} & \begin{minipage}[b]{0.07\columnwidth}\raggedright
Refrig\strut
\end{minipage} & \begin{minipage}[b]{0.04\columnwidth}\raggedright
TS 1\strut
\end{minipage} & \begin{minipage}[b]{0.05\columnwidth}\raggedright
TS 2\strut
\end{minipage} & \begin{minipage}[b]{0.03\columnwidth}\raggedright
Counter Step\strut
\end{minipage} & \begin{minipage}[b]{0.03\columnwidth}\raggedright
Counter Count\strut
\end{minipage}\tabularnewline
\midrule
\endhead
\begin{minipage}[t]{0.08\columnwidth}\raggedright
1\strut
\end{minipage} & \begin{minipage}[t]{0.07\columnwidth}\raggedright
-30\strut
\end{minipage} & \begin{minipage}[t]{0.09\columnwidth}\raggedright
Soak Ctrl\strut
\end{minipage} & \begin{minipage}[t]{0.06\columnwidth}\raggedright
5\strut
\end{minipage} & \begin{minipage}[t]{0.06\columnwidth}\raggedright
Off\strut
\end{minipage} & \begin{minipage}[t]{0.12\columnwidth}\raggedright
1:00:00\strut
\end{minipage} & \begin{minipage}[t]{0.07\columnwidth}\raggedright
Auto\strut
\end{minipage} & \begin{minipage}[t]{0.04\columnwidth}\raggedright
\strut
\end{minipage} & \begin{minipage}[t]{0.05\columnwidth}\raggedright
\strut
\end{minipage} & \begin{minipage}[t]{0.03\columnwidth}\raggedright
1\strut
\end{minipage} & \begin{minipage}[t]{0.03\columnwidth}\raggedright
0\strut
\end{minipage}\tabularnewline
\begin{minipage}[t]{0.08\columnwidth}\raggedright
2\strut
\end{minipage} & \begin{minipage}[t]{0.07\columnwidth}\raggedright
20\strut
\end{minipage} & \begin{minipage}[t]{0.09\columnwidth}\raggedright
Soak Ctrl\strut
\end{minipage} & \begin{minipage}[t]{0.06\columnwidth}\raggedright
5\strut
\end{minipage} & \begin{minipage}[t]{0.06\columnwidth}\raggedright
Off\strut
\end{minipage} & \begin{minipage}[t]{0.12\columnwidth}\raggedright
1:00:00\strut
\end{minipage} & \begin{minipage}[t]{0.07\columnwidth}\raggedright
Auto\strut
\end{minipage} & \begin{minipage}[t]{0.04\columnwidth}\raggedright
\strut
\end{minipage} & \begin{minipage}[t]{0.05\columnwidth}\raggedright
\strut
\end{minipage} & \begin{minipage}[t]{0.03\columnwidth}\raggedright
1\strut
\end{minipage} & \begin{minipage}[t]{0.03\columnwidth}\raggedright
0\strut
\end{minipage}\tabularnewline
\begin{minipage}[t]{0.08\columnwidth}\raggedright
3\strut
\end{minipage} & \begin{minipage}[t]{0.07\columnwidth}\raggedright
55\strut
\end{minipage} & \begin{minipage}[t]{0.09\columnwidth}\raggedright
Soak Ctrl\strut
\end{minipage} & \begin{minipage}[t]{0.06\columnwidth}\raggedright
5\strut
\end{minipage} & \begin{minipage}[t]{0.06\columnwidth}\raggedright
Off\strut
\end{minipage} & \begin{minipage}[t]{0.12\columnwidth}\raggedright
1:00:00\strut
\end{minipage} & \begin{minipage}[t]{0.07\columnwidth}\raggedright
Auto\strut
\end{minipage} & \begin{minipage}[t]{0.04\columnwidth}\raggedright
\strut
\end{minipage} & \begin{minipage}[t]{0.05\columnwidth}\raggedright
\strut
\end{minipage} & \begin{minipage}[t]{0.03\columnwidth}\raggedright
2\strut
\end{minipage} & \begin{minipage}[t]{0.03\columnwidth}\raggedright
3\strut
\end{minipage}\tabularnewline
\begin{minipage}[t]{0.08\columnwidth}\raggedright
4\strut
\end{minipage} & \begin{minipage}[t]{0.07\columnwidth}\raggedright
65\strut
\end{minipage} & \begin{minipage}[t]{0.09\columnwidth}\raggedright
Soak Ctrl\strut
\end{minipage} & \begin{minipage}[t]{0.06\columnwidth}\raggedright
5\strut
\end{minipage} & \begin{minipage}[t]{0.06\columnwidth}\raggedright
Off\strut
\end{minipage} & \begin{minipage}[t]{0.12\columnwidth}\raggedright
1:00:00\strut
\end{minipage} & \begin{minipage}[t]{0.07\columnwidth}\raggedright
Auto\strut
\end{minipage} & \begin{minipage}[t]{0.04\columnwidth}\raggedright
\strut
\end{minipage} & \begin{minipage}[t]{0.05\columnwidth}\raggedright
\strut
\end{minipage} & \begin{minipage}[t]{0.03\columnwidth}\raggedright
3\strut
\end{minipage} & \begin{minipage}[t]{0.03\columnwidth}\raggedright
2\strut
\end{minipage}\tabularnewline
\begin{minipage}[t]{0.08\columnwidth}\raggedright
5\strut
\end{minipage} & \begin{minipage}[t]{0.07\columnwidth}\raggedright
180\strut
\end{minipage} & \begin{minipage}[t]{0.09\columnwidth}\raggedright
Soak Ctrl\strut
\end{minipage} & \begin{minipage}[t]{0.06\columnwidth}\raggedright
5\strut
\end{minipage} & \begin{minipage}[t]{0.06\columnwidth}\raggedright
Off\strut
\end{minipage} & \begin{minipage}[t]{0.12\columnwidth}\raggedright
1:00:00\strut
\end{minipage} & \begin{minipage}[t]{0.07\columnwidth}\raggedright
Auto\strut
\end{minipage} & \begin{minipage}[t]{0.04\columnwidth}\raggedright
\strut
\end{minipage} & \begin{minipage}[t]{0.05\columnwidth}\raggedright
\strut
\end{minipage} & \begin{minipage}[t]{0.03\columnwidth}\raggedright
1\strut
\end{minipage} & \begin{minipage}[t]{0.03\columnwidth}\raggedright
4\strut
\end{minipage}\tabularnewline
\begin{minipage}[t]{0.08\columnwidth}\raggedright
6\strut
\end{minipage} & \begin{minipage}[t]{0.07\columnwidth}\raggedright
23\strut
\end{minipage} & \begin{minipage}[t]{0.09\columnwidth}\raggedright
On\strut
\end{minipage} & \begin{minipage}[t]{0.06\columnwidth}\raggedright
5\strut
\end{minipage} & \begin{minipage}[t]{0.06\columnwidth}\raggedright
Off\strut
\end{minipage} & \begin{minipage}[t]{0.12\columnwidth}\raggedright
1:00:00\strut
\end{minipage} & \begin{minipage}[t]{0.07\columnwidth}\raggedright
Auto\strut
\end{minipage} & \begin{minipage}[t]{0.04\columnwidth}\raggedright
\strut
\end{minipage} & \begin{minipage}[t]{0.05\columnwidth}\raggedright
\strut
\end{minipage} & \begin{minipage}[t]{0.03\columnwidth}\raggedright
4\strut
\end{minipage} & \begin{minipage}[t]{0.03\columnwidth}\raggedright
4\strut
\end{minipage}\tabularnewline
\bottomrule
\end{longtable}

The following figure depicts the loops created by the program in the
above table.

\begin{figure}
\centering

\includegraphics[width=0.95\textwidth]{figures/3818980958-appdxa-2-3.PNG}
\caption{Program with multiple looped counters}
\end{figure}\FloatBarrier


\subsection{B.1 Specification of HMI (Touch Screen
Display)}

\begin{longtable}[]{@{}lll@{}}
\toprule
\begin{minipage}[b]{0.17\columnwidth}\raggedright
Item\strut
\end{minipage} & \begin{minipage}[b]{0.41\columnwidth}\raggedright
Specification\strut
\end{minipage} & \begin{minipage}[b]{0.33\columnwidth}\raggedright
Application\strut
\end{minipage}\tabularnewline
\midrule
\endhead
\begin{minipage}[t]{0.17\columnwidth}\raggedright
HMI / Display\strut
\end{minipage} & \begin{minipage}[t]{0.41\columnwidth}\raggedright
10.1-inch TFT LCD touch screen display 1280x800px resolution High
brightness: 850 cd/m2 Half life period: 50,000h Video Interface: HDMI
Touch Panel Interface: USB-C Power Supply: Power Jack (DC 7.0V -
14.0V)\strut
\end{minipage} & \begin{minipage}[t]{0.33\columnwidth}\raggedright
User Interface for GL Controller and Chamber Operation\strut
\end{minipage}\tabularnewline
\bottomrule
\end{longtable}


\subsection{C.1 Specification of AAEON
Up2Board}

The GL controller software is an embedded computer powered by Debian
GNU/Linux distribution running as the operating system (OS). The
embedded computer hardware is an x86 64-bit architecture system board,
called AAEON UP2Board, designed by
\href{https://www.aaeon.com/en/product/detail/iot-gateway-maker-boards-up-squared/specification}{AEEON}.
The following is a list of the specifications of this hardware.

\begin{longtable}[]{@{}ll@{}}
\toprule
\begin{minipage}[b]{0.29\columnwidth}\raggedright
Component\strut
\end{minipage} & \begin{minipage}[b]{0.65\columnwidth}\raggedright
Details\strut
\end{minipage}\tabularnewline
\midrule
\endhead
\begin{minipage}[t]{0.29\columnwidth}\raggedright
CPU\strut
\end{minipage} & \begin{minipage}[t]{0.65\columnwidth}\raggedright
Intel® Celeron® Processor N3350\strut
\end{minipage}\tabularnewline
\begin{minipage}[t]{0.29\columnwidth}\raggedright
Graphics\strut
\end{minipage} & \begin{minipage}[t]{0.65\columnwidth}\raggedright
Intel® HD Graphics\strut
\end{minipage}\tabularnewline
\begin{minipage}[t]{0.29\columnwidth}\raggedright
Memory\strut
\end{minipage} & \begin{minipage}[t]{0.65\columnwidth}\raggedright
Up to 8GB onboard LPDDR4\strut
\end{minipage}\tabularnewline
\begin{minipage}[t]{0.29\columnwidth}\raggedright
Storage\strut
\end{minipage} & \begin{minipage}[t]{0.65\columnwidth}\raggedright
Up to 128GB onboard eMMC\strut
\end{minipage}\tabularnewline
\begin{minipage}[t]{0.29\columnwidth}\raggedright
I/O\strut
\end{minipage} & \begin{minipage}[t]{0.65\columnwidth}\raggedright
HDMI 1.4b x 1; DP 1.2 x 1\strut
\end{minipage}\tabularnewline
\begin{minipage}[t]{0.29\columnwidth}\raggedright
Camera\strut
\end{minipage} & \begin{minipage}[t]{0.65\columnwidth}\raggedright
MIPI-CSI via 21 Pin FPC connector x 1; MIPI-CSI via 31 Pin FPC connector
x 1\strut
\end{minipage}\tabularnewline
\begin{minipage}[t]{0.29\columnwidth}\raggedright
USB\strut
\end{minipage} & \begin{minipage}[t]{0.65\columnwidth}\raggedright
USB 2.0 x 2 (via 10 Pin Header x 1); USB 3.2 Gen 1 x 3 (Type-A); USB 3.2
Gen 1 OTG x 1 (Micro B)\strut
\end{minipage}\tabularnewline
\begin{minipage}[t]{0.29\columnwidth}\raggedright
Expansion\strut
\end{minipage} & \begin{minipage}[t]{0.65\columnwidth}\raggedright
40-pin GPIO x 1; Full-sized mPCIe x 1 (PCIe \& SATA {[}x1{]}, USB2.0
{[}x1{]}); M.2 2230 E-Key x 1 (PCIe {[}x1{]}, USB2.0 {[}x1{]}); SATA III
x 1; 60 Pin EXHAT x 1\strut
\end{minipage}\tabularnewline
\begin{minipage}[t]{0.29\columnwidth}\raggedright
Display Interface\strut
\end{minipage} & \begin{minipage}[t]{0.65\columnwidth}\raggedright
HDMI 1.4b x 1; DP 1.2 x 1; eDP 1.3 x 1\strut
\end{minipage}\tabularnewline
\begin{minipage}[t]{0.29\columnwidth}\raggedright
Ethernet\strut
\end{minipage} & \begin{minipage}[t]{0.65\columnwidth}\raggedright
GbE x 2 (Realtek 8111H CG)\strut
\end{minipage}\tabularnewline
\begin{minipage}[t]{0.29\columnwidth}\raggedright
Security\strut
\end{minipage} & \begin{minipage}[t]{0.65\columnwidth}\raggedright
fTPM\strut
\end{minipage}\tabularnewline
\begin{minipage}[t]{0.29\columnwidth}\raggedright
RTC\strut
\end{minipage} & \begin{minipage}[t]{0.65\columnwidth}\raggedright
Yes\strut
\end{minipage}\tabularnewline
\begin{minipage}[t]{0.29\columnwidth}\raggedright
OS support\strut
\end{minipage} & \begin{minipage}[t]{0.65\columnwidth}\raggedright
Windows® 10 IoT Enterprise LTSC 2021; Windows IoT Core Linux Ubuntu
20.04 LTS; Linux Yocto 3.1\strut
\end{minipage}\tabularnewline
\begin{minipage}[t]{0.29\columnwidth}\raggedright
Power\strut
\end{minipage} & \begin{minipage}[t]{0.65\columnwidth}\raggedright
DC 5V / 4-6A\strut
\end{minipage}\tabularnewline
\begin{minipage}[t]{0.29\columnwidth}\raggedright
Dimension\strut
\end{minipage} & \begin{minipage}[t]{0.65\columnwidth}\raggedright
856 mm x 900 mm\strut
\end{minipage}\tabularnewline
\bottomrule
\end{longtable}


\subsection{D.1 FAQ}

Answers to frequently asked questions (FAQ) are compiled, summarized and
displayed in the following table. \textbf{\emph{Note:}} As more FAQ and
answers become available, they will be appended in the table for future
release.

\begin{longtable}[]{@{}llll@{}}
\toprule
\begin{minipage}[b]{0.08\columnwidth}\raggedright
No.\strut
\end{minipage} & \begin{minipage}[b]{0.30\columnwidth}\raggedright
Question\strut
\end{minipage} & \begin{minipage}[b]{0.33\columnwidth}\raggedright
Answer\strut
\end{minipage} & \begin{minipage}[b]{0.17\columnwidth}\raggedright
Reference\strut
\end{minipage}\tabularnewline
\midrule
\endhead
\begin{minipage}[t]{0.08\columnwidth}\raggedright
1\strut
\end{minipage} & \begin{minipage}[t]{0.30\columnwidth}\raggedright
I want to change the buzzer operation in the event of the occurrence of
an ``alarm or warning.''\strut
\end{minipage} & \begin{minipage}[t]{0.33\columnwidth}\raggedright
ON/OFF switch can be controlled via {[}Setup{]} → {[}Configuration{]} →
{[}Set Sound{]}. If the beep is turned off, alarm and warning are
indicated only with red notifications on the HMI.\strut
\end{minipage} & \begin{minipage}[t]{0.17\columnwidth}\raggedright
2.8 Setting the Alarm and Warning Buzzer\strut
\end{minipage}\tabularnewline
\begin{minipage}[t]{0.08\columnwidth}\raggedright
2\strut
\end{minipage} & \begin{minipage}[t]{0.30\columnwidth}\raggedright
What should I do when a system abnormality has occurred?\strut
\end{minipage} & \begin{minipage}[t]{0.33\columnwidth}\raggedright
Notify the sales representative or customer support for
assistance.\strut
\end{minipage} & \begin{minipage}[t]{0.17\columnwidth}\raggedright
\strut
\end{minipage}\tabularnewline
\begin{minipage}[t]{0.08\columnwidth}\raggedright
3\strut
\end{minipage} & \begin{minipage}[t]{0.30\columnwidth}\raggedright
I want to turn off the HMI immediately.\strut
\end{minipage} & \begin{minipage}[t]{0.33\columnwidth}\raggedright
To turn off the HMI immediately, press {[}Setup{]} → {[}Configuration{]}
→ {[}Display Setup{]}, then press \textbf{Turn Off Display}. To turn it
back on, tab your finger on the screen.\strut
\end{minipage} & \begin{minipage}[t]{0.17\columnwidth}\raggedright
7.9.1 Display Setup on HMI\strut
\end{minipage}\tabularnewline
\begin{minipage}[t]{0.08\columnwidth}\raggedright
4\strut
\end{minipage} & \begin{minipage}[t]{0.30\columnwidth}\raggedright
I want to set the HMI in sleep mode.\strut
\end{minipage} & \begin{minipage}[t]{0.33\columnwidth}\raggedright
Several options are available for setting the HMI in sleep mode. Press
{[}Setup{]} → {[}Configuration{]} → {[}Display Setup{]}, then select
\textbf{HMI Screensaver timeout}. To turn it back on, tab your finger on
the screen.\strut
\end{minipage} & \begin{minipage}[t]{0.17\columnwidth}\raggedright
7.9.1 Display Setup on HMI\strut
\end{minipage}\tabularnewline
\begin{minipage}[t]{0.08\columnwidth}\raggedright
5\strut
\end{minipage} & \begin{minipage}[t]{0.30\columnwidth}\raggedright
``OFF'' is displayed for a setup value.\strut
\end{minipage} & \begin{minipage}[t]{0.33\columnwidth}\raggedright
When the equipment is stopped or the humidity control is set to OFF,
``OFF'' is displayed.\strut
\end{minipage} & \begin{minipage}[t]{0.17\columnwidth}\raggedright
2.3.3 Monitor\strut
\end{minipage}\tabularnewline
\begin{minipage}[t]{0.08\columnwidth}\raggedright
6\strut
\end{minipage} & \begin{minipage}[t]{0.30\columnwidth}\raggedright
``---'' is displayed as the temperature (humidity) process value.\strut
\end{minipage} & \begin{minipage}[t]{0.33\columnwidth}\raggedright
If the temperature (humidity) process value is abnormal or the humidity
control is set to OFF, ``---'' is displayed.\strut
\end{minipage} & \begin{minipage}[t]{0.17\columnwidth}\raggedright
2.3.3 Monitor\strut
\end{minipage}\tabularnewline
\begin{minipage}[t]{0.08\columnwidth}\raggedright
7\strut
\end{minipage} & \begin{minipage}[t]{0.30\columnwidth}\raggedright
I want to change the refrigeration setup from Auto to Manual (arbitrary
state).\strut
\end{minipage} & \begin{minipage}[t]{0.33\columnwidth}\raggedright
Access {[}Set Mode{]}, select a Constant setup, then apply refrigeration
manual setting as needed. When you want to change the mode from Auto to
Manual during operation, make the manual setting first.\strut
\end{minipage} & \begin{minipage}[t]{0.17\columnwidth}\raggedright
4.3 Refrigeration Setting\strut
\end{minipage}\tabularnewline
\begin{minipage}[t]{0.08\columnwidth}\raggedright
8\strut
\end{minipage} & \begin{minipage}[t]{0.30\columnwidth}\raggedright
When is the setting of time enabled?\strut
\end{minipage} & \begin{minipage}[t]{0.33\columnwidth}\raggedright
At the moment when the time signal is set, the contact output in
question changes Contact rating: 250VAC, 5A\strut
\end{minipage} & \begin{minipage}[t]{0.17\columnwidth}\raggedright
\strut
\end{minipage}\tabularnewline
\begin{minipage}[t]{0.08\columnwidth}\raggedright
9\strut
\end{minipage} & \begin{minipage}[t]{0.30\columnwidth}\raggedright
Is the trend graph data erased when the circuit breaker is turned off or
when power failure occurs?\strut
\end{minipage} & \begin{minipage}[t]{0.33\columnwidth}\raggedright
Data are stored in the non-volatile memory of the GL controller system;
they cannot not be erased by power failure or shutdown.\strut
\end{minipage} & \begin{minipage}[t]{0.17\columnwidth}\raggedright
4.9.7 Monitor: Trend; 4.9.8 Trend Graph Manipulation Buttons\strut
\end{minipage}\tabularnewline
\begin{minipage}[t]{0.08\columnwidth}\raggedright
10\strut
\end{minipage} & \begin{minipage}[t]{0.30\columnwidth}\raggedright
What is the ramp operation?\strut
\end{minipage} & \begin{minipage}[t]{0.33\columnwidth}\raggedright
The ramp operation is to control the temperatures from the set point of
the previous step to that of the present step at a certain ramp.\strut
\end{minipage} & \begin{minipage}[t]{0.17\columnwidth}\raggedright
5.3.1 Step-Wise Programming\strut
\end{minipage}\tabularnewline
\begin{minipage}[t]{0.08\columnwidth}\raggedright
11\strut
\end{minipage} & \begin{minipage}[t]{0.30\columnwidth}\raggedright
What is the soak time control?\strut
\end{minipage} & \begin{minipage}[t]{0.33\columnwidth}\raggedright
After the operation of the step is started, the equipment waits until
the process value attains the attainment range of the set point, and
then starts the counting of time. Thus, the soak time in the temperature
set point is the same as the setup time.\strut
\end{minipage} & \begin{minipage}[t]{0.17\columnwidth}\raggedright
5.3.1 Step-Wise Programming\strut
\end{minipage}\tabularnewline
\begin{minipage}[t]{0.08\columnwidth}\raggedright
12\strut
\end{minipage} & \begin{minipage}[t]{0.30\columnwidth}\raggedright
Can the soak time and ramp control be set simultaneously?\strut
\end{minipage} & \begin{minipage}[t]{0.33\columnwidth}\raggedright
The ramp control and soak time cannot be set simultaneously.\strut
\end{minipage} & \begin{minipage}[t]{0.17\columnwidth}\raggedright
5.3.1 Step-Wise Programming\strut
\end{minipage}\tabularnewline
\begin{minipage}[t]{0.08\columnwidth}\raggedright
13\strut
\end{minipage} & \begin{minipage}[t]{0.30\columnwidth}\raggedright
Can the program in operation be edited or saved?\strut
\end{minipage} & \begin{minipage}[t]{0.33\columnwidth}\raggedright
During program execution, that program can be loaded into the program
editor for editing, but the edited version cannot be saved back in the
same slot that was loaded from for execution; it can however be saved in
a different slot (and under a different name).\strut
\end{minipage} & \begin{minipage}[t]{0.17\columnwidth}\raggedright
5.6.4 Manipulating a Program during its Execution\strut
\end{minipage}\tabularnewline
\begin{minipage}[t]{0.08\columnwidth}\raggedright
14\strut
\end{minipage} & \begin{minipage}[t]{0.30\columnwidth}\raggedright
Can the program in operation be deleted?\strut
\end{minipage} & \begin{minipage}[t]{0.33\columnwidth}\raggedright
During program execution, that program cannot be deleted.\strut
\end{minipage} & \begin{minipage}[t]{0.17\columnwidth}\raggedright
5.6.4 Manipulating a Program during its Execution\strut
\end{minipage}\tabularnewline
\begin{minipage}[t]{0.08\columnwidth}\raggedright
15\strut
\end{minipage} & \begin{minipage}[t]{0.30\columnwidth}\raggedright
Can the program contain loops (counters) to reduce the number of
steps?\strut
\end{minipage} & \begin{minipage}[t]{0.33\columnwidth}\raggedright
Program may contain multiple loops (i.e., counters) to repeat steps or
groups of steps.\strut
\end{minipage} & \begin{minipage}[t]{0.17\columnwidth}\raggedright
5.3.1 Step-Wise programming; 5.3.6 Example 6: Setting Multiple Counters
in a Program\strut
\end{minipage}\tabularnewline
\begin{minipage}[t]{0.08\columnwidth}\raggedright
16\strut
\end{minipage} & \begin{minipage}[t]{0.30\columnwidth}\raggedright
Can I connect my GL chamber to my main network and access it
remotely?\strut
\end{minipage} & \begin{minipage}[t]{0.33\columnwidth}\raggedright
Your GL chamber can join the main network using DHCP or static. The UI
can be accessed via a Web browser.\strut
\end{minipage} & \begin{minipage}[t]{0.17\columnwidth}\raggedright
8 Network Setup \& Remote Access\strut
\end{minipage}\tabularnewline
\begin{minipage}[t]{0.08\columnwidth}\raggedright
17\strut
\end{minipage} & \begin{minipage}[t]{0.30\columnwidth}\raggedright
Can I set my GL chamber to use a static IP address on the network?\strut
\end{minipage} & \begin{minipage}[t]{0.33\columnwidth}\raggedright
Your GL chamber can use a static IP address on a static network\strut
\end{minipage} & \begin{minipage}[t]{0.17\columnwidth}\raggedright
8.5 Custom Static Network\strut
\end{minipage}\tabularnewline
\begin{minipage}[t]{0.08\columnwidth}\raggedright
18\strut
\end{minipage} & \begin{minipage}[t]{0.30\columnwidth}\raggedright
Can I access other GL chambers on the network?\strut
\end{minipage} & \begin{minipage}[t]{0.33\columnwidth}\raggedright
Working with multiple GL chambers on the network, you can choose one GL
controller (on the network) to monitor other GL chambers on the
network.\strut
\end{minipage} & \begin{minipage}[t]{0.17\columnwidth}\raggedright
2.1 GL Controller Operation Options; 9 Network Monitor\strut
\end{minipage}\tabularnewline
\bottomrule
\end{longtable}


